% ترجمة أنتجها الذكاء الاصطناعي. التصحيحات والتنضيد من إنتاج OpenAI Codex — GPT-6 Astra بمستوى الاستدلال Ultra. تبقى أسماء إيمي نوتر والمؤلفين المشاركين كما وردت في الأعمال. لا يُنسب إلى الناشر عمل ترجمة بشرية أو مراجعة علمية بشرية مستقلة.
% تبقى الحقوق لأصحابها. يستند التوزيع المجاني لأغراض غير تجارية وتعليمية بحتة إلى افتراض الناشر عدم اعتراض أصحاب الحقوق، لا إلى إذن ثابت. الإتاحة المجانية لا تمنح بذاتها حق نشر عمل مشتق. يمكن لأصحاب الحقوق الاعتراض عبر الرابط أدناه؛ وسيسحب الناشر المادة المعنية بناء على طلبهم. لا يمنح هذا البيان ترخيصًا جديدًا ولا ينقل أعمال الآخرين إلى الملك العام.
% https://github.com/KokunoYumeto/modern-latex-manuscripts/issues/new
\documentclass[11pt]{article}
\pdfvariable suppressoptionalinfo 767
\errorcontextlines=200
\usepackage{fontspec}
\usepackage{amsmath,amssymb,mathtools,array,booktabs,graphicx,adjustbox,pdflscape}
\usepackage[a4paper,margin=16mm,includefoot]{geometry}
\usepackage[english,bidi=basic,layout=sectioning lists contents footnotes captions tabular graphics]{babel}
\makeatletter
\@ifundefined{babelsubrl}{\DeclareRobustCommand\babelsubrl[1]{\leavevmode{\bbl@textdir\@ne#1}}}{}
\makeatother
\babelprovide[import=ar,main,onchar=ids fonts]{arabic}
\babelprovide[import=fa]{persian}
\babelfont{rm}{TeX Gyre Pagella}
\babelfont{sf}{TeX Gyre Heros}
\babelfont{tt}{Latin Modern Mono}
\babelfont[arabic]{rm}[Path=fonts/,UprightFont=Scheherazade-Regular.ttf,BoldFont=Scheherazade-Bold.ttf,ItalicFont=Scheherazade-Regular.ttf,ItalicFeatures={FakeSlant=0.15},BoldItalicFont=Scheherazade-Bold.ttf,BoldItalicFeatures={FakeSlant=0.15},Renderer=HarfBuzz,Scale=1.08]{Scheherazade-Regular.ttf}
\babelfont[persian]{rm}[Path=fonts/,UprightFont=Scheherazade-Regular.ttf,BoldFont=Scheherazade-Bold.ttf,ItalicFont=Scheherazade-Regular.ttf,ItalicFeatures={FakeSlant=0.15},BoldItalicFont=Scheherazade-Bold.ttf,BoldItalicFeatures={FakeSlant=0.15},Renderer=HarfBuzz,Scale=1.08]{Scheherazade-Regular.ttf}
\usepackage[hyperfootnotes=false]{hyperref}
\hypersetup{unicode=true,hidelinks,hypertexnames=false,pdflang=ar,pdfauthor={},pdftitle={أعمال إيمي نوتر — النص العربي المتاح},pdfcreator={},pdfproducer={}}
\makeatletter
\pdfstringdefDisableCommands{\let\babelsubrl\@firstofone}
\makeatother
\setlength{\parindent}{1.25em}
\setlength{\parskip}{0pt}
\setlength{\emergencystretch}{10em}
\sloppy\raggedbottom\allowdisplaybreaks
\providecommand{\A}{\Delta}
\providecommand{\D}{\Delta}
\providecommand{\R}{\varrho}
\providecommand{\N}{\operatorname{N}}
\providecommand{\modu}[1]{\;\operatorname{mod} #1}
\providecommand{\mcell}[1]{\ensuremath{\begin{gathered}#1\end{gathered}}}
\providecommand{\glqq}{\textquotedblbase}
\providecommand{\grqq}{\textquotedblleft}
\providecommand{\srcspaced}[1]{{\bfseries #1}}
\providecommand{\srcnumdisplay}[3][3.4em]{%
  \par\smallskip\noindent
  \babelsublr{\begin{minipage}[t]{#1}\vspace*{0pt}\srcLTR{#2}\end{minipage}%
  \begin{minipage}[t]{\dimexpr\linewidth-#1\relax}\vspace*{0pt}\centering
  \(\displaystyle #3\)\end{minipage}}\par\smallskip}
\newcommand{\srcRTL}[1]{\relax\ifmmode\text{\mbox{\babelsubrl{\foreignlanguage{arabic}{#1}}}}\else\babelsubrl{\foreignlanguage{arabic}{#1}}\fi}
\newcommand{\srcLTR}[1]{\relax\ifmmode\text{\mbox{\babelsublr{\foreignlanguage{english}{#1}}}}\else\babelsublr{\foreignlanguage{english}{#1}}\fi}
\newcommand{\srcfn}[2]{\begingroup\renewcommand{\thefootnote}{\srcLTR{#1}}\footnote{#2}\addtocounter{footnote}{-1}\endgroup}
\newcommand{\srcfnmark}[1]{\begingroup\renewcommand{\thefootnote}{\srcLTR{#1}}\footnotemark\addtocounter{footnote}{-1}\endgroup}
\newcommand{\srcfntext}[2]{\begingroup\renewcommand{\thefootnote}{\srcLTR{#1}}\footnotetext{#2}\endgroup}
\providecommand{\Lsys}{\mathfrak{L}}
\providecommand{\Jdom}{\mathfrak{J}}
\providecommand{\Gdom}{\mathfrak{G}}
\providecommand{\rcdots}{\,\cdots\,}
\providecommand{\pair}[2]{(#1\,|\,#2)}
\providecommand{\eps}{\varepsilon}
\RequirePackage{tikz}
\newcommand{\NoetherRosetteGlyph}{%
  \begin{tikzpicture}[x=0.1363584905660377pt,y=-0.1363584905660377pt,baseline=(current bounding box.south)]
    \path[use as bounding box] (0,0) rectangle (108,108);
    \fill[even odd rule]
    (33.5,84) -- (32.5,83) -- (30.5,83) -- (29,81.5) -- (29,76.5) -- (30.5,75) -- (31.5,75) -- (32.5,74) -- (34.5,74) -- (36,72.5) -- (36,71.5) -- (35.5,71) -- (33.5,71) -- (33,70.5) -- (33,69.5) -- (30.5,67) -- (28.5,69) -- (26.5,69) -- (24.5,71) -- (23.5,71) -- (22.5,72) -- (21.5,72) -- (20.5,73) -- (18.5,73) -- (17.5,72) -- (16.5,73) -- (16,72.5) -- (16,68.5) -- (18.5,66) -- (19.5,66) -- (20,65.5) -- (20,64.5) -- (26,58.5) -- (26,56.5) -- (24,54.5) -- (24,53.5) -- (17,46.5) -- (17,45.5) -- (16.5,45) -- (15.5,45) -- (15,44.5) -- (15,42.5) -- (15.5,42) -- (17.5,42) -- (19.5,44) -- (20.5,44) -- (21,44.5) -- (20,45.5) -- (21,46.5) -- (21,47.5) -- (22,48.5) -- (22,49.5) -- (29,56.5) -- (29,57.5) -- (27,59.5) -- (27,60.5) -- (26,61.5) -- (26,62.5) -- (24,64.5) -- (24,65.5) -- (22,67.5) -- (22.5,68) -- (24.5,68) -- (25.5,67) -- (28.5,67) -- (29.5,66) -- (31.5,66) -- (32.5,65) -- (33.5,65) -- (34.5,64) -- (36.5,64) -- (37.5,63) -- (38.5,63) -- (39.5,62) -- (41.5,62) -- (42.5,61) -- (44.5,61) -- (45.5,60) -- (46.5,60) -- (47.5,59) -- (48.5,59) -- (49,58.5) -- (49,57.5) -- (49.5,57) -- (51.5,57) -- (52,56.5) -- (50.5,55) -- (49.5,55) -- (48.5,54) -- (47.5,54) -- (46.5,53) -- (45.5,53) -- (43,50.5) -- (43,49.5) -- (42,48.5) -- (42,47.5) -- (41.5,47) -- (40.5,47) -- (39.5,48) -- (33.5,48) -- (32,46.5) -- (32,44.5) -- (34,42.5) -- (34,40.5) -- (33.5,40) -- (32.5,40) -- (31.5,39) -- (26.5,39) -- (26,38.5) -- (26,35.5) -- (27.5,34) -- (31.5,34) -- (33,35.5) -- (33,36.5) -- (34.5,38) -- (37.5,38) -- (38,37.5) -- (38,36.5) -- (41,33.5) -- (41,31.5) -- (40,30.5) -- (40,27.5) -- (39,26.5) -- (39,25.5) -- (38,24.5) -- (38,23.5) -- (37,22.5) -- (37,19.5) -- (37.5,19) -- (38.5,20) -- (39.5,20) -- (41,21.5) -- (41,24.5) -- (42,25.5) -- (42,27.5) -- (43,28.5) -- (43,30.5) -- (44,31.5) -- (44,35.5) -- (45,36.5) -- (45,37.5) -- (45.5,38) -- (46.5,38) -- (47,38.5) -- (47,40.5) -- (48,41.5) -- (48,42.5) -- (49,43.5) -- (49,44.5) -- (50,45.5) -- (50,48.5) -- (53,51.5) -- (52,52.5) -- (52,55.5) -- (53.5,57) -- (56.5,57) -- (58.5,59) -- (59.5,59) -- (60.5,60) -- (61.5,60) -- (62,60.5) -- (62,61.5) -- (63,62.5) -- (63,64.5) -- (63.5,65) -- (64.5,65) -- (65.5,64) -- (68.5,64) -- (69.5,65) -- (73.5,65) -- (74,65.5) -- (74,66.5) -- (75,67.5) -- (75,68.5) -- (75.5,69) -- (76,68.5) -- (76,66.5) -- (76.5,66) -- (78.5,68) -- (80.5,68) -- (81.5,69) -- (84.5,69) -- (85,68.5) -- (84,67.5) -- (84,65.5) -- (83.5,65) -- (82.5,65) -- (82,64.5) -- (82.5,64) -- (84.5,64) -- (86,65.5) -- (86,66.5) -- (88.5,69) -- (89.5,69) -- (90,69.5) -- (90,70.5) -- (91,71.5) -- (91,72.5) -- (93,74.5) -- (91.5,76) -- (90.5,76) -- (88.5,74) -- (87.5,74) -- (86.5,73) -- (85.5,73) -- (83.5,71) -- (82.5,71) -- (81.5,70) -- (80.5,70) -- (79.5,69) -- (77.5,69) -- (74.5,72) -- (73.5,72) -- (73,72.5) -- (73,74.5) -- (73.5,75) -- (78.5,75) -- (79.5,74) -- (80.5,74) -- (81,74.5) -- (81,78.5) -- (80,79.5) -- (80,80.5) -- (78.5,82) -- (73.5,82) -- (73,81.5) -- (73,79.5) -- (72,78.5) -- (72,77.5) -- (71.5,77) -- (70.5,77) -- (69.5,76) -- (68.5,76) -- (68,76.5) -- (68,77.5) -- (67,78.5) -- (67,79.5) -- (66.5,80) -- (65.5,80) -- (65,80.5) -- (66,81.5) -- (64.5,83) -- (63,81.5) -- (63,79.5) -- (62,78.5) -- (62,77.5) -- (61,76.5) -- (61,74.5) -- (60,73.5) -- (60,72.5) -- (59,71.5) -- (59,68.5) -- (58,67.5) -- (58,64.5) -- (56,62.5) -- (56,61.5) -- (55.5,61) -- (54.5,62) -- (53.5,62) -- (52,63.5) -- (52,64.5) -- (51,65.5) -- (51,66.5) -- (50,67.5) -- (50,68.5) -- (49,69.5) -- (49,70.5) -- (48,71.5) -- (48,75.5) -- (46,77.5) -- (46,79.5) -- (44.5,81) -- (42.5,81) -- (42,80.5) -- (44,78.5) -- (44,77.5) -- (43.5,77) -- (41.5,77) -- (39.5,75) -- (38.5,76) -- (37.5,76) -- (37,76.5) -- (37,78.5) -- (36,79.5) -- (38,81.5) -- (38,82.5) -- (37.5,83) -- (34.5,83) -- cycle
    (65.5,22) -- (63.5,22) -- (63,21.5) -- (63.5,21) -- (65.5,21) -- (66,21.5) -- cycle
    (55.5,28) -- (53.5,28) -- (53,27.5) -- (53.5,27) -- (54.5,27) -- (57.5,24) -- (59.5,24) -- (60,24.5) -- (57.5,27) -- (56.5,27) -- cycle
    (64.5,26) -- (64,25.5) -- (64.5,25) -- (65,25.5) -- cycle
    (63.5,29) -- (63,28.5) -- (63,27.5) -- (63.5,27) -- (64.5,27) -- (65,27.5) -- cycle
    (58.5,55) -- (57,53.5) -- (57,52.5) -- (56,51.5) -- (56,49.5) -- (57,48.5) -- (57,46.5) -- (60,43.5) -- (60,42.5) -- (59.5,42) -- (58.5,42) -- (58,41.5) -- (59,40.5) -- (59,39.5) -- (61.5,37) -- (62.5,37) -- (63.5,36) -- (65.5,36) -- (66.5,37) -- (68.5,37) -- (70,35.5) -- (69,34.5) -- (69,32.5) -- (68,31.5) -- (68.5,31) -- (74.5,31) -- (75,31.5) -- (75,32.5) -- (76,33.5) -- (76,35.5) -- (75,36.5) -- (75,37.5) -- (72,40.5) -- (72,41.5) -- (71,42.5) -- (72,43.5) -- (72,44.5) -- (72.5,45) -- (73.5,45) -- (74.5,46) -- (75.5,46) -- (76.5,45) -- (78.5,45) -- (79.5,44) -- (80,44.5) -- (80,47.5) -- (79.5,48) -- (77.5,48) -- (76.5,49) -- (74.5,49) -- (73.5,50) -- (71.5,50) -- (70.5,51) -- (68.5,51) -- (67.5,52) -- (65.5,52) -- (64.5,53) -- (63.5,53) -- (62.5,54) -- (59.5,54) -- cycle
    (86.5,50) -- (85.5,50) -- (85,49.5) -- (86,48.5) -- (86,46.5) -- (85.5,46) -- (82.5,46) -- (81.5,47) -- (81,46.5) -- (82,45.5) -- (82,44.5) -- (82.5,44) -- (84.5,44) -- (86.5,42) -- (88.5,42) -- (89.5,41) -- (92.5,41) -- (93,41.5) -- (93,42.5) -- (94,43.5) -- (94,44.5) -- (93.5,45) -- (92.5,45) -- (88.5,49) -- (87.5,49) -- cycle
    (44,42.5) -- (43.5,42) -- (43,42.5) -- (43.5,43) -- cycle
    (45,43.5) -- (44.5,43) -- (44,43.5) -- (44.5,44) -- cycle
    (80.5,60) -- (79.5,60) -- (79,59.5) -- (79,56.5) -- (81,54.5) -- (81,53.5) -- (82.5,52) -- (83,52.5) -- (83,53.5) -- (82,54.5) -- (82,55.5) -- (81.5,56) -- (80.5,56) -- (80,56.5) -- (80,57.5) -- (81,58.5) -- (81,59.5) -- cycle
    (55.5,54) -- (55,53.5) -- (55.5,53) -- (56,53.5) -- cycle
    (43,65.5) -- (42.5,65) -- (42,65.5) -- (42.5,66) -- cycle
    (61,67.5) -- (60.5,67) -- (60,67.5) -- (60.5,68) -- cycle
    (47,69.5) -- (47,68.5) -- (46.5,68) -- (46,68.5) -- (46,69.5) -- (46.5,70) -- cycle
    (62,68.5) -- (61.5,68) -- (61,68.5) -- (61.5,69) -- cycle
    (67.5,97) -- (65.5,97) -- (64,95.5) -- (64,94.5) -- (61.5,92) -- (60.5,92) -- (54.5,86) -- (53.5,86) -- (49.5,90) -- (48.5,90) -- (47.5,91) -- (46.5,91) -- (45.5,92) -- (45,91.5) -- (45,90.5) -- (45.5,90) -- (46.5,90) -- (48.5,88) -- (49.5,88) -- (54.5,83) -- (55.5,83) -- (60.5,88) -- (61.5,88) -- (62.5,89) -- (63.5,89) -- (64,89.5) -- (64,90.5) -- (64.5,91) -- (66.5,91) -- (67,90.5) -- (66,89.5) -- (67,88.5) -- (65,86.5) -- (65,85.5) -- (65.5,85) -- (66.5,85) -- (67,85.5) -- (67,86.5) -- (68,87.5) -- (68,90.5) -- (69,91.5) -- (69,94.5) -- (70,95.5) -- (69.5,96) -- (68.5,96) -- cycle
    (40.5,96) -- (36.5,96) -- (36,95.5) -- (37,94.5) -- (37,91.5) -- (40.5,88) -- (41.5,88) -- (42,88.5) -- (41,89.5) -- (41,92.5) -- (42,93.5) -- (42,94.5) -- cycle;
  \end{tikzpicture}%
}
\providecommand{\noethpIIrosette}{\makebox[9.44444pt][c]{\raisebox{-3.85pt}[0pt][0pt]{\NoetherRosetteGlyph}}}
\providecommand{\NoetherRosetteInline}{\makebox[17.72671698pt][c]{\raisebox{-3.85pt}[10.87671698pt][3.85pt]{\NoetherRosetteGlyph}}}


\newenvironment{NoetherWideTable}[2]{%
  \newgeometry{layoutwidth=#1,layoutheight=#2,margin=16mm,includefoot}%
  \paperwidth=#1\paperheight=#2\pagewidth=#1\pageheight=#2
  \begin{otherlanguage}{english}\begin{landscape}%
}{%
  \end{landscape}\end{otherlanguage}
}
\AddToHook{env/NoetherWideTable/after}{%
  \restoregeometry
  \paperwidth=210mm\paperheight=297mm\pagewidth=210mm\pageheight=297mm
}
\newtagform{NoetherLTR}[\srcLTR]{(}{)}
\usetagform{NoetherLTR}
\begin{document}
\selectlanguage{arabic}
\begin{center}
{\Large\bfseries 1. في تكوين منظومة الملازمات لحدودية متجانسة ثلاثية من الدرجة الرابعة.}\par
\vspace{1.2em}
بقلم \srcspaced{إيمي نوتر}.\par
\vspace{0.8em}
(مقتطف من أطروحة المؤلفة.)\par
\vspace{1.2em}
\srcLTR{Sitz. Ber. d. Physikal.-mediz. Sozietät in Erlangen 39 (1907), S. 176--179}
\end{center}

تتناول أعمال \srcspaced{غوردان} و\srcspaced{مايزانو} و\srcspaced{باسكال} منظومة الملازمات لحدودية متجانسة ثلاثية من الدرجة الرابعة\srcfn{1)}{\srcLTR{\srcspaced{P. Gordan}, Über das volle Formensystem der ternären biquadratischen Form $f=x_1^3x_2+x_2^3x_3+x_3^3x_1$ (Math. Annalen Bd. XVII, S. 217 bis 233. 1880); \srcspaced{G. Maisano}, 1. Sistemi completi dei primi cinque gradi della forma ternaria biquadratica e degl' invarianti, covarianti e contravarianti di sesto grado (Giorn. di Battaglini XIX). 2. Sui covarianti indipendenti di $6^\circ$ grado nei coefficienti della forma biquadratica ternaria (Rend. Circ. Mat. di Palermo I. 1887); \srcspaced{E. Pascal}, Contributo alla teoria della forma ternaria biquadratica e delle sue varie decomposizioni in fattori (Memoria premiata dalla R. Accademia delle scienze fisiche e matematiche di Napoli. 1905).}}. ويقيم السيد \srcspaced{غوردان} المنظومة الكاملة، المؤلفة من 54 تكوينا، للحدودية الخاصة ذاتية التشاكل (\srcspaced{automorphic form}) $f=x_1^3x_2+x_2^3x_3+x_3^3x_1$، مستندا إلى مبادئ مماثلة لما وضعه لمنظومات الملازمات في المجال الثنائي.

يعرض السيد \srcspaced{مايزانو}، للحدودية العامة من الدرجة الرابعة، الحدوديات حتى الرتبة الخامسة\srcfn{2)}{يُقصد بـ\glqq{}الرتبة\grqq{} الدرجة بالنسبة إلى المعاملات، وبـ\glqq{}الدرجة\grqq{} الدرجة بالنسبة إلى المتغيرات.} بما فيها، وكذلك بعض اللامتغيرات والملازمات الموافقة للتغير (\srcspaced{covariants}) والملازمات المخالفة للتغير (\srcspaced{contravariants}) من رتب أعلى، وذلك وفق الطريقة التي طبقها السيد \srcspaced{غوردان} في المجلد الأول من Math. Annalen على الحدودية المتجانسة التكعيبية الثلاثية. أما السيد \srcspaced{باسكال} فيعنى، مستفيدا من نتائج مايزانو، أساسا بمسألة تحلل الحدودية من الدرجة الرابعة إلى عوامل\srcfn{1)}{يحاول السيد \srcspaced{مايزانو} في مذكرته القصيرة الثانية إثبات اعتماد خطي للملازمات الثلاثة الموافقة للتغير ذات الرتبة 6 والدرجة 6. وعلى خلاف هذا البرهان غير المكتمل تماما، يعتقد السيد \srcspaced{باسكال} أنه أثبت الاستقلال الخطي للملازمات الثلاثة الموافقة للتغير ذات الرتبة 6، وكذلك استقلال اللامتغيرات الثلاثة ذات الرتبة 9. غير أن حساباتي بينت أن ثمة بالفعل علاقة خطية بين الملازمات الثلاثة، وعلاقة أخرى بين اللامتغيرات الثلاثة؛ والصيغ الصريحة، بترميز \srcspaced{Pascal}، هي:
\[
0=20\Omega_1+6\Omega_2-3\Omega_3+4f\cdot C_1-12f\cdot C_2-4A\cdot\Delta.
\]
\[
0=4A^3-15A\cdot B+30C-90D+9E.
\]}. 

\srcspaced{غاية أبحاثي هي إقامة منظومة الملازمات للحدودية العامة المتجانسة الثلاثية من الدرجة الرابعة؛ ولن أقدم أولا إلا الأسس الرئيسة، أي سأقيم ما يسمى \glqq{}منظومة كاملة نسبيا\grqq{}}\srcfn{2)}{في شأن التسمية، انظر \srcspaced{Gordan-Kerschensteiner}, Vorlesungen über Invariantentheorie, Bd. II, S. 227.}\srcspaced{.}

الفكرة الأساسية في تكوين منظومات الحدوديات المتجانسة الثلاثية هي نفسها في المجال الثنائي. فمن منظومة أولى كاملة نسبيا ــ هي منظومة الحدوديات المنقولة من المجال الثنائي ــ ننتقل، وفق قانون معين، إلى منظومات ذات مقاييس أعلى فأعلى، إلى أن تصبح منظومة أحد المقاييس ــ مع اتخاذ المقياس حدودية أساسية ــ منتهية ومعلومة، أو إلى أن يمكن اختزال مقياس إلى حدوديات تكون اللامتغيرات عوامل لها. وفي الحالة الأولى تنشأ المنظومة الكاملة مطلقا من إجراء عملية الاشتقاق والتقليص (\srcspaced{transvection; Überschiebung}) على المنظومة الكاملة نسبيا بواسطة منظومة المقياس؛ أما في الحالة الثانية فتتطابق المنظومة الكاملة نسبيا والمنظومة الكاملة مطلقا. وبسبب تناهي منظومات الملازمات الذي برهنه السيد \srcspaced{هيلبرت} بوجه عام، لا بد لهذه العملية أن تنتهي.

في حالتنا يمكن اختزال المقياس $(abc)$ للمنظومة الكاملة نسبيا الأولى إلى المقياسين
\[
\Delta=(abc)^2a_x^2b_x^2c_x^2\quad\text{و}\quad \nu=(abu)^4;
\]
ومن ذلك يسهل الحصول على المنظومة الكاملة نسبيا بالنسبة إلى المقياس $\nu$. ونختار الآن الحدوديات الآتية سلسلة للمقاييس:
\[
\nu;\qquad \nu^{(\nu)}=(\nu\nu_1x)^4=s_x^4,
\]
\[
\nu(s)=(ss'u)^4\ldots
\]
ونضم، بوصفهما مقياسين إضافيين، الحدوديتين من الدرجة الثانية اللتين تظهران عند تكوين المنظومة الكاملة نسبيا بالنسبة إلى المقياس $s$، وهما $u_\varrho^2$ و$t_x^2$. ويمكن حينئذ البرهنة على أن المقياس التالي لـ$s$، أي $(ss'u)^4$، قابل للاختزال إلى المقياس $(\varrho,t)$. وبما أن المنظومة المشتركة لحدوديتين من الدرجة الثانية منتهية ومعلومة، فقد تحقق بذلك الإتمام المبين أعلاه. وتتألف المنظومة الكاملة نسبيا بالنسبة إلى المقياس $(\varrho,t)$ من 331 تكوينا. ــ

وأضيف هنا بعض البيانات المتعلقة بالطريقة المستعملة.

يمكن تعريف العملية الأساسية لتوليد الحدوديات، وهي \srcspaced{عملية الطي}، في المجال الثلاثي على النحو الآتي:

إذا استبدلنا، في جداء رمزي
\[
s_x^m t_x^n u_\sigma^\mu u_\tau^\nu
\]
أزواج العوامل
\[
\begin{array}{@{}c@{\qquad}c@{\qquad}c@{\qquad}c@{\qquad}c@{}}
 & s_xt_x & u_\sigma u_\tau & s_xu_\sigma\ \text{أو}\ s_xu_\tau & t_xu_\tau\ \text{أو}\ t_xu_\sigma\\[0.45em]
\text{على الترتيب بـ} & (stu) & (\sigma\tau x) & s_\sigma\ \text{أو}\ s_\tau & t_\tau\ \text{أو}\ t_\sigma\\[0.45em]
\text{الطي} & \mathrm{I} & \mathrm{II} & \mathrm{III} & \mathrm{IV}
\end{array}
\]
فإن الحدوديات الناتجة تكون قد نشأت من الحدودية الأصلية بالطي\srcfn{1)}{\srcspaced{Gordan}. Math. Annalen, Bd. XVII, S. 219.}.

ويمكن دائما وضع $s_\sigma=0;\ t_\tau=0$. وفي العلاقة بين عمليات الطي المفردة تصح المبرهنة الآتية:

\srcspaced{عمليتا الطي I وII هما عمليتا الطي الأساسيتان، ويمكن منهما تركيب عمليتي الطي III وIV بصرف النظر عن ترتيب التركيب. وبعبارة أخرى: لتكوين جميع الحدوديات الناتجة بالطي من عبارة معطاة، لا يلزم إلا تطبيق عمليتي الطي I وII}\srcfn{2)}{تستثنى من المبرهنة الحالة التي ينعدم فيها $n$ و$\mu$ معا، أو $m$ و$\nu$ معا؛ وعندئذ تختزل \glqq{}سلسلة الملازمات\grqq{} إلى العنصر القطري.}.

نعرّف \glqq{}سلسلة الملازمات\grqq{} بأنها حدودية ابتدائية مع جميع الحدوديات الناتجة منها بطيها في نفسها. وبحسب المبرهنة يمكن ترتيب سلسلة الملازمات $s_x^m t_x^n u_\sigma^\mu u_\tau^\nu$ في مخطط مستطيل. وعند التقدم فيه نحصل على:
\begin{enumerate}
\item عند الانتقال عمودا واحدا إلى اليمين، جميع الحدوديات الناشئة من الحدوديات المجاورة بعملية الطي I،
\item عند الانتقال صفا واحدا إلى الأسفل، جميع الحدوديات الناشئة من الحدوديات الواقعة أعلاه بعملية الطي II، ومن ثم جميع الحدوديات الناشئة بالطي.
\end{enumerate}

في ظل هذه الفروض يمكن صياغة المبرهنات المتعلقة بالمختزلات في أعم صورها، ومع التعريف الآتي
\begin{center}
\glqq{}المختزل سلسلة ملازمات قابلة للاختزال\grqq{}
\end{center}
يمكن صياغتها على النحو الآتي:

\begin{center}
\srcspaced{إذا كانت الحدودية الابتدائية في سلسلة ملازمات قابلة للاختزال من حيث إن أحد عناصرها قد نشأ بالطي مع مختزل، وكانت الحدودية النهائية في السلسلة قد نشأت من ذلك العنصر نفسه بالطي، فإن سلسلة الملازمات كلها قابلة للاختزال.}
\end{center}

وتجدر تسمية طريقتي الاختزال الإضافيتين الآتيتين:
\begin{enumerate}
\item ما يسمى \glqq{}الاختزال المزدوج\grqq{}؛ أي اختزال حدودية بطريقتين للحصول على \srcspaced{علاقة بين} الحدوديات الأعلى رتبة،
\item \glqq{}الطي مع الحدوديات المتحللة إلى عوامل\grqq{}، وهو يتسم جزئيا \srcspaced{بطابع} الاختزال بالمختزلات، وجزئيا \srcspaced{بطابع \glqq{}الاختزال المزدوج\grqq{}}.
\end{enumerate}
\clearpage
\setcounter{footnote}{0}

\providecommand{\Afield}{\mathfrak A}
\providecommand{\Bfield}{\mathfrak B}
\providecommand{\Cfield}{\mathfrak C}
\providecommand{\Sdomain}{\mathfrak S}
\providecommand{\Kfield}{\mathfrak K}
\providecommand{\Rfield}{\mathfrak R}
\providecommand{\Xreal}{\mathfrak X}
\providecommand{\Yreal}{\mathfrak Y}
\providecommand{\qtext}[1]{\glqq #1\grqq{}}
\begin{center}
{\Large\bfseries 2. في تكوين منظومة الملازمات للحدودية المتجانسة الثلاثية من الدرجة الرابعة\srcfn{*)}{نشر مستخرج قصير من المقدمة والفصل الأول في \glqq{}Sitzungsber. der phys.-med. Sozietät Erlangen 1907\grqq{}، ص. 176 إلى 179.}.}\par
\vspace{0.8em}
بقلم الآنسة \emph{\srcLTR{Emmy Noether}} في إرلنغن.\par
\vspace{1.2em}
\srcLTR{Journal f. d. reine u. angew. Math. 134 (1908)}، ص. 23--90، وجدولان
\end{center}

\vspace{3em}

\begin{center}
\textbf{المقدمة.}
\end{center}

تتناول أعمال \emph{غوردان} و\emph{مايزانو} و\emph{باسكال}\srcfn{**)}{\srcLTR{\emph{P. Gordan}, Über das volle Formensystem der ternären biquadratischen Form
\[
f=x_1^3x_2+x_2^3x_3+x_3^3x_1.
\]
(Math. Annalen, Bd. XVII (1880), S. 217--233.)
\emph{G. Maisano}, 1) \emph{Sistemi completi dei primi cinque gradi della forma ternaria biquadratica e degl' invarianti, covarianti e contravarianti di sesto grado}. (Giorn. di Battaglini XIX.)
2) \emph{Sui covarianti indipendenti di $6^\circ$ grado nei coefficienti della forma biquadratica ternaria}. (Rend. Circ. Mat. di Palermo I, 1887.)
\emph{E. Pascal}, \emph{Contributo alla teoria della forma ternaria biquadratica e delle sue varie decomposizioni in fattori}. (Memoria premiata dalla R. Accademia delle scienze fisiche e matematiche di Napoli. 1905.)}} منظومة الملازمات للحدودية المتجانسة الثلاثية من الدرجة الرابعة. ويقيم السيد \emph{غوردان} المنظومة الكاملة، المؤلفة من 54 تكوينا، للحدودية الخاصة ذاتية التشاكل: \(f=x_1^3x_2+x_2^3x_3+x_3^3x_1\)، مستندا إلى مبادئ مماثلة لما وضعه لمنظومات الملازمات في المجال الثنائي.
يقيم السيد \emph{مايزانو} للحدودية العامة المتجانسة الثلاثية من الدرجة الرابعة الملازمات حتى الرتبة الخامسة بما فيها الرتبة الخامسة نفسها\srcfn{***)}{يقصد بـ\glqq{}الرتبة\grqq{} الدرجة بالنسبة إلى المعاملات، وبـ\glqq{}الدرجة\grqq{} الدرجة بالنسبة إلى المتغيرات.}، كما يقيم بعض اللامتغيرات والملازمات الموافقة للتغير والملازمات المخالفة للتغير من رتب أعلى، وفق الطريقة التي طبقها السيد \emph{غوردان} على الحدودية المتجانسة الثلاثية من الدرجة الثالثة في المجلد الأول من Math. Annalen. أما السيد \emph{باسكال} فيتناول، مستفيدا من نتائج \emph{مايزانو}، بوجه رئيس مسألة تحلل الحدودية من الدرجة الرابعة إلى عوامل\srcfn{*)}{يحاول السيد \emph{مايزانو} في المذكرة القصيرة الثانية إثبات اعتماد خطي بين الملازمات الثلاثة الموافقة للتغير من الرتبة السادسة والدرجة السادسة. وعلى خلاف هذا البرهان غير المكتمل تماما، يعتقد السيد \emph{باسكال} أنه أثبت الاستقلال الخطي للملازمات الثلاثة الموافقة للتغير من الرتبة السادسة، وكذلك استقلال اللامتغيرات الثلاثة من الرتبة التاسعة. غير أن وجود علاقة خطية بالفعل بين الملازمات الثلاثة من جهة، وعلاقة أخرى بين اللامتغيرات الثلاثة من جهة أخرى، تبينه الصيغتان (13)، § 11 و(22)، § 17، ملاحظة، من هذا العمل، حيث ترد العلاقتان صراحة. ووفقا لإفادة السيد \emph{باسكال}، يرجع التناقض إلى خطأ عددي في تعبير ملازمه المخالف للتغير $p$ (المسمى هنا $\varrho$) في ص. 46 من بحثه.}.
\emph{هدف الأبحاث الآتية هو إقامة منظومة الملازمات للحدودية العامة المتجانسة الثلاثية من الدرجة الرابعة؛ ولا يقدم هذا العمل إلا الأسس الرئيسة، إذ يقيم ما يسمى \glqq{}منظومة كاملة نسبيا\grqq{}}\srcfn{**)}{للتسمية، راجع § 3a.}\emph{.}
يرتبط هذا العمل ارتباطا وثيقا بعمل \emph{غوردان}؛ غير أن المبادئ الواردة هناك بصيغة عامة جدا كان لا بد من تفصيلها. فبمساعدة مبرهنة في ترابط عمليات الطي، وبإدخال \glqq{}سلسلة الملازمات\grqq{} (§ 1)، تصاغ مبرهنات الاختزال صياغة دقيقة وتستكمل (§ 3)، في حين أن الإنشاء التكراري لتوسعات خاصة (§ 2) يتيح الوسيلة الحسابية لإجراء الاختزالات فعليا.
الفكرة الأساسية في تكوين منظومات الملازمات للحدوديات الثلاثية هي نفسها في المجال الثنائي. فمن منظومة أولى كاملة نسبيا ــ هي منظومة الحدوديات المنقولة من المجال الثنائي ــ ننتقل وفق قانون معين إلى منظومات ذات مقياس أعلى فأعلى، إلى أن تصبح منظومة أحد المقاييس ــ مع اتخاذ المقياس حدودية أساسية ــ منتهية ومعروفة، أو إلى أن يمكن اختزال مقياس إلى حدوديات تتخذ اللامتغيرات عوامل لها. ففي الحالة الأولى تنشأ المنظومة الكاملة مطلقا بإجراء عملية الاشتقاق والتقليص للمنظومة الكاملة نسبيا على منظومة المقياس؛ أما في الحالة الثانية فتتطابق المنظومتان الكاملة نسبيا والكاملة مطلقا. وبسبب تناهي منظومات الملازمات الذي أثبته السيد \emph{هيلبرت} بوجه عام، لا بد لهذا الإجراء من أن ينتهي.
في حالتنا يختزل المقياس $(abc)$ للمنظومة الأولى الكاملة نسبيا (§ 4) إلى المقياسين
\[
\Delta=(abc)^2a_x^2b_x^2c_x^2\quad\text{و}\quad \nu=(abu)^4
\]
(§ 5)، ثم توجد المنظومة الكاملة نسبيا بترديد $\nu$ (§ 6). وقد وردت هذه النتائج من قبل في عمل \emph{غوردان} للحدودية العامة من الدرجة الرابعة؛ إلا أن طريقتنا في الاشتقاق أيسر نقلا إلى حدوديات من درجات أعلى.
ونختار سلسلة المقاييس الآن على أنها الحدوديات الآتية (§ 7):
\[
\nu,\qquad \nu(\nu)=(\nu\nu_1x)^4=s_x^4,\qquad \nu(s)=(ss'u)^4,\ldots .
\]

عند تكوين المنظومة الكاملة نسبيا بترديد $s$، يضم إلى المقاييس ملازمان تربيعيان واردان في المنظومة، هما $u_\varrho^2$ و$t_x^2$ (§ 9 و§ 10)، ومن ثم تكون المنظومة الكاملة نسبيا بترديد $(s,\varrho,t)$؛ ثم يبين (§ 17) أن المقياس $(ss'u)^4$ قابل للاختزال إلى المقياسين $\varrho$ و$t$. وبذلك تتحول المنظومة الأعلى التالية إلى منظومة كاملة نسبيا بترديد $(\varrho,t)$. ولما كانت المنظومة المشتركة لحدوديتين من الدرجة الثانية منتهية ومعروفة، وتتألف في الحالة العامة من 20 تكوينا\srcfn{*)}{راجع \emph{Clebsch-Lindemann}, Vorlesungen über Geometrie. (Bd. I, S. 288 ff.)، منظومة حدوديتين متوافقتين في التدرج. \emph{Gordan}, Über Büschel von Kegelschnitten. (Math. Ann. Bd. XIX. S. 530)، منظومة حدوديتين متخالفتين في التدرج.}، فإن إقامة المنظومة الكاملة نسبيا بترديد $(\varrho,t)$ تحقق النهاية الموصوفة أعلاه. أما إجراء عملية الاشتقاق والتقليص لهذه المنظومة على منظومة $(\varrho,t)$ لتكوين المنظومة الكاملة مطلقا فيبقى مؤجلا.
توجد 331 تكوينا بوصفها منظومة كاملة نسبيا بترديد $(\varrho,t)$، وقد رتبت في الجدول المرفق بحسب درجتها بالنسبة إلى المتغيرين $x$ و$u$.
وفي الختام نقدم نظرة عامة على الرموز المستعملة:
\begin{align*}
f&=a_x^4,\\
\theta&=\theta_x^4u_\vartheta^2=(abu)^2a_x^2b_x^2,\\
K&=K_x^6u_k^3=(a\theta u)u_\vartheta^2a_x^3\theta_x^3,\\
j&=u_j^6=(a\theta u)^4u_\vartheta^2,\\
\Delta&=\Delta_x^6=a_\vartheta^2a_x^2\theta_x^4=(abc)^2a_x^2b_x^2c_x^2,\\
N&=N_x^8u_\eta=(a\Delta u)a_x^3\Delta_x^5,\\
\nu&=u_\nu^4=(abu)^4,\\
H&=H_x^2u_\eta^4=(\nu\nu_1x)^2u_\nu^2u_{\nu_1}^2,\\
L&=L_x^3u_l^6=(\nu\eta x)H_x^2u_\nu^3u_\eta^3,\\
g&=g_x^6=(\nu\eta x)^4H_x^2,\\
\sigma&=u_\sigma^6=H_\nu^2u_\nu^2u_\eta^4,\\
s&=s_x^4=(\nu\nu_1x)^4,\\
Z&=Z_x^4u_\zeta^2=(ss'u)^2s_x^2s_x'^2,\\
\varrho&=u_\varrho^2=\theta_\nu^4u_\vartheta^2,\\
t&=t_x^2=a_\eta^4H_x^2,\\
i&=a_\nu^4,\\
J&=s_\nu^4.
\end{align*}
\begin{center}
{\Large\bfseries الفصل الأول. \quad مبرهنات عامة في الحدوديات المتجانسة الثلاثية.}
\end{center}

\begin{center}
\textbf{§ 1.}\quad \emph{عملية الطي. سلاسل الملازمات.}
\end{center}

العملية الأساسية لتكوين الحدوديات هي \emph{عملية الطي}، ويمكن تعريفها في المجال الثلاثي على النحو الآتي. ليكن لدينا حاصل ضرب رمزي:
\[
s_x^m t_x^n u_\sigma^\mu u_\tau^\nu .
\]
فإذا استبدلنا أزواج العوامل
\[
s_xt_x\qquad u_\sigma u_\tau\qquad s_xu_\sigma\ \text{أو}\ s_xu_\tau
\qquad t_xu_\tau\ \text{أو}\ t_xu_\sigma
\]
على الترتيب بـ
\[
(stu)\qquad (\sigma\tau x)\qquad s_\sigma\ \text{أو}\ s_\tau
\qquad t_\tau\ \text{أو}\ t_\sigma,
\]
الطي
\[
\text{I}\qquad\text{II}\qquad\text{III}\qquad\text{IV},
\]
فتكون الحدوديات الناتجة قد نشأت من الحدودية الأصلية بالطي\srcfn{*)}{\emph{Gordan}, Math. Annalen Bd. XVII S. 219.}.
ويمكن النظر إلى العبارة $s_x^m t_x^n u_\sigma^\mu u_\tau^\nu$ إما على أنها حاصل ضرب فعلي لحدوديتين $S\cdot T$، أو على أنها حدودية واحدة ذات عدة صفوف من الرموز:
\[
s_x^m t_x^n u_\sigma^\mu u_\tau^\nu=A_x^{m+n}u_\lambda^{\mu+\nu}.
\]
ففي الحالة الأولى نتحدث عن طي الحدودية $S$ مع $T$، وفي الحالة الثانية عن \glqq{}طي الحدودية $A$ في نفسها\grqq{}. ومن باب الإيجاز نفترض دائما فيما يأتي ــ وهو ما يتحقق فعلا في جميع الحدوديات التي سننظر فيها لاحقا ــ الشرط الآتي:
\srcnumdisplay{(a)}{\begin{gathered}
s_\sigma^\lambda=0,\\
t_\tau^\varkappa=0
\end{gathered}}
لكل $\lambda$ و$\varkappa$ غير مساويين للصفر، وكذلك عند الجمع مع أي عمليات طي أخرى. ويعني ذلك أن الحدودية $s_x^m t_y^n u_\sigma^\mu u_\tau^\nu$ هي ما يسمى \glqq{}صيغة نظامية بالنسبة إلى عملية أوميغا\grqq{}؛ أي إنها تنعدم عند تطبيق عملية $\Omega$ بالنسبة إلى المتغيرين $x$ و$u$، وكذلك بالنسبة إلى المتغيرين $y$ و$v$.
ومن ثم لا يمثل الشرط (a) أي تقييد، بل يمكن تحقيقه دائما\srcfn{*)}{راجع \emph{Gordan}, Math. Annalen Bd. V S. 104.}.
وفي ترابط عمليات الطي المفردة تصح المبرهنة الآتية:
\emph{المبرهنة I. عمليتا الطي I وII هما عمليتا الطي الأساسيتان، ويمكن منهما تركيب عمليتي الطي III وIV بصرف النظر عن ترتيب التركيب. وبعبارة أخرى: لتكوين جميع الحدوديات الناتجة بالطي من عبارة معطاة، لا يلزم إلا تطبيق عمليتي الطي I وII.}
\textbf{البرهان:} ينتج من مبرهنة المتطابقات، أو من مبرهنة حاصل ضرب المصفوفات، مع مراعاة (a):
\[
(\widehat{st}\sigma x)=-t_\sigma,\qquad
(\widehat{\sigma\tau}su)=-s_\tau,
\]
\[
(\widehat{st}\tau x)=s_\tau,\qquad
(\widehat{\sigma\tau}tu)=t_\sigma,
\]
\[
(stu)(\sigma\tau x)=s_\tau+t_\sigma-u_x\cdot s_\tau t_\sigma .
\]
ومن ثم تنشأ عمليتا الطي III وIV بتركيب عمليتي الطي I وII، سواء عند تطبيق عملية الطي II على عملية الطي I، أي على العوامل $(stu)u_\sigma u_\tau$، أم عند تطبيق عملية الطي I على عملية الطي II، أي على العوامل $(\sigma\tau x)s_xt_x$.
نعرف \emph{سلسلة الملازمات}\srcfn{**)}{راجع \emph{Clebsch}, Über eine Fundamentalaufgabe der Invariantentheorie (Abh. der Gött. Ges. d. Wiss. Bd. XVII): § 17 و§ 18، الخاتمة. وتختلف \glqq{}سلسلة الملازمات\grqq{} عن \glqq{}المنظومة المكافئة المختزلة بالمعنى الدقيق\grqq{} المدخلة هناك بمبدأ الترتيب بحسب الملازمات الأعلى.}\setcounter{footnote}{9} بأنها ملازم ابتدائي مع جميع الملازمات التي نشأت منه بطيه في نفسه، ونرمز إلى سلسلة الملازمات بواسطة الملازم الابتدائي. ويقصد بـ\glqq{}الملازم الأعلى\grqq{} هنا كل ملازم يشتمل على عدد أكبر من عمليات الطي في نفسه؛ أما من عمليتي الطي المتكافئتين $s_\tau$ و$t_\sigma$ فلا بد من تعيين إحداهما اصطلاحا على أنها الأعلى. وينتج من قانون تكوين سلسلة الملازمات ما يأتي:
1) كل ملازم في سلسلة الملازمات خطي في رموز الملازم الابتدائي.
2) تتعين سلسلة الملازمات تعيينا وحيدا في حالة الملازمات الابتدائية ذات صفين من الرموز المتخالفة التدرج؛ أما في الملازمات الابتدائية ذات صفوف رموز أكثر، فلا بد لجعلها وحيدة من تمييز عمليات طي مخصوصة من بين العمليات المرتبطة بمبرهنة المتطابقات.
ووفقا للمبرهنة I يمكن ترتيب سلسلة الملازمات $s_x^m t_x^n u_\sigma^\mu u_\tau^\nu$ ($m>n;\ \mu>\nu$) بحسب المخطط المستطيل الآتي:
\begingroup\small
\[
\begin{array}{c|c|c|c|c}
s_x^m t_x^n u_\sigma^\mu u_\tau^\nu
& (stu)
& (stu)^2
& \cdots
& \noethpIIrosette\\[0.3em]
(\sigma\tau x)
& t_\sigma,\ s_\tau
& t_\sigma(stu),\ s_\tau(stu)
& \cdots
& \vdots\\[0.3em]
(\sigma\tau x)^2
& t_\sigma(\sigma\tau x),\ s_\tau(\sigma\tau x)
& t_\sigma^2,\ t_\sigma s_\tau,\ s_\tau^2
& \cdots
& \vdots\\[0.3em]
\vdots&\vdots&\vdots&\ddots&\vdots\\[0.3em]
(\sigma\tau x)^\nu
& t_\sigma(\sigma\tau x)^{\nu-1},\ s_\tau(\sigma\tau x)^{\nu-1}
& t_\sigma^2(\sigma\tau x)^{\nu-2},\ t_\sigma s_\tau(\sigma\tau x)^{\nu-2},\ s_\tau^2(\sigma\tau x)^{\nu-2}
& \cdots
& \vdots\\[0.3em]
\vdots
& t_\sigma(\sigma\tau x)^\nu
& t_\sigma^2(\sigma\tau x)^{\nu-1},\ t_\sigma s_\tau(\sigma\tau x)^{\nu-1}
& \cdots
& \vdots\\[0.3em]
\vdots
& \vdots
& t_\sigma^2(\sigma\tau x)^\nu
& \cdots
& \vdots
\end{array}
\]
\endgroup
\begingroup\small
\[
\begin{array}{r@{\quad}c}
\noethpIIrosette\srcfnmark{*)}
&
\begin{array}{c|c|c}
(stu)^n & \cdots & \cdots\\[0.3em]
t_\sigma(stu)^{n-1},\ s_\tau(stu)^{n-1}
& s_\tau(stu)^n & \cdots\\[0.3em]
t_\sigma^2(stu)^{n-2},\ t_\sigma s_\tau(stu)^{n-2},\ s_\tau^2(stu)^{n-2}
& t_\sigma s_\tau(stu)^{n-1},\ s_\tau^2(stu)^{n-1}
& s_\tau^2(stu)^n .
\end{array}
\end{array}
\]
\endgroup
\srcfntext{*)}{هنا، وكذلك في المواضع المماثلة الآتية، يلحق الجزء السفلي الموسوم بالعلامة \NoetherRosetteInline{} بالموضع الواقع أعلى اليمين والموسوم بالعلامة نفسها في المخطط.}
ونحصل هنا، إذا
1) تقدمنا عمودا واحدا نحو اليمين، على جميع الحدوديات الناشئة من الحدوديات المجاورة بعملية الطي I،
2) تقدمنا صفا واحدا نحو الأسفل، على جميع الحدوديات الناشئة من الحدوديات الواقعة فوقها بعملية الطي II،
ومن ثم على جميع الحدوديات الناشئة بالطي.
كل ملازم اختياري من سلسلة الملازمات الكلية، من الصورة $t_\sigma^\kappa s_\tau^\lambda(stu)^\mu$ أو $t_\sigma^\kappa s_\tau^\lambda(\sigma\tau x)^\nu$، يكون الملازم الابتدائي لسلسلة ملازمات جديدة. وتتألف هذه السلسلة من جميع ملازمات المخطط المستطيل التي يكون فيها الملازم $t_\sigma^\kappa s_\tau^\lambda(stu)^\mu$ أو $t_\sigma^\kappa s_\tau^\lambda(\sigma\tau x)^\nu$ رأس الزاوية العليا اليسرى، والتي تشترك في العامل $t_\sigma^\kappa s_\tau^\lambda$.
ملاحظة ختامية. ترد على المبرهنة I حالة استثنائية عندما يكون العددان $n$ و$\mu$ (أو، على الترتيب، $m$ و$\nu$) مساويين للصفر؛ ففي الحالة الأولى لا تعود عمليات الطي I وII وIV موجودة، بينما تبقى عملية الطي III، وفي الحالة المقابلة لا تعود عمليات الطي I وII وIII موجودة، بينما تبقى عملية الطي IV. وفي هذه الحالة نسمي العنصر القطري من المخطط الآتي سلسلة الملازمات:
\[
\begin{array}{c@{\qquad}c@{\qquad}c}
\begin{array}{ccccc}
s_x^m u_\tau^\nu &&&&\\
&s_\tau&&&\\
&&s_\tau^2&&\\
&&&\ddots&\\
&&&&s_\tau^\nu
\end{array}
&\text{أو، في الحالة المقابلة،}&
\begin{array}{ccccc}
t_x^n u_\sigma^\mu &&&&\\
&t_\sigma&&&\\
&&t_\sigma^2&&\\
&&&\ddots&\\
&&&&t_\sigma^\mu
\end{array}
\end{array}
\]
وباتساق مع غياب عمليتي الطي I وII، تختزل توسعة سلسلة الملازمات بحسب الأشكال القطبية في هذه الحالة دائما إلى العنصر الابتدائي؛ غير أن المبرهنات المتعلقة بالمختزلات تظل صحيحة.
\begin{center}
\textbf{§ 2.}\quad \emph{توسعات بحسب الأشكال القطبية لسلسلة الملازمات}\srcfn{*)}{انظر Study, Methoden zur Theorie der ternären Formen (Teubner 1889) II. § 7. وخلافا للاشتقاق الوارد هناك، يتيح اشتقاقنا طريقة التعيين الحسابي السريع للمعاملات العددية $C_{pr\varrho}$. ولو أجري تخصيص الحالتين a)، a') في توسعة كليبش--شتودي لكان نصها:
\[
(stu)^\mu u_\tau^{\nu-\varkappa}v_\tau^\varkappa s_x^m t_x^{\,n-\lambda}(tuv)^\lambda
=\sum_{p,r,\varrho} C'_{pr\varrho}
\bigl[s_\tau^\varrho(stu)^{\mu+r-\varrho}\bigr]_{\widehat{uv}^{\lambda-r+\varrho}v^{\varkappa-\varrho},\,w^p,\,(w=\widehat{xw})},
\]
ومن ثم لا تسمح بالتبسيط الذي يظهر في توسعتنا في أثناء الحساب نفسه، متى علمنا أن جميع الحدود التي تحتوي العامل $u_x$ يجب إغفالها.}.
\end{center}

للحدوديات المتجانسة ذات $n$ من المتغيرات، أقيم صراحة نوعان من التوسعات المتسلسلة\srcfn{**)}{انظر Gordan, Über Kombinanten. § 2 und 5 (Math. Ann. Bd. V).}:
1) للحدوديات المتجانسة التي لكل منها صف واحد من المتغيرات متخالفة التدرج: تكون $s_x^m u_\tau^\nu$ متسلسلة تتقدم بحسب قوى $u_x$، ومعاملاتها \glqq{}صيغ نظامية\grqq{}؛
2) للحدوديات المتجانسة ذات صفين من المتغيرات متوافقة التدرج $s_x^m t_x^n$، توسعة بحسب الأشكال القطبية للملازمات الأولية الموافقة للتغير (أي الأشكال القطبية لسلسلة الملازمات $s_x^m t_x^n$)، وتكتب في الحالة الثلاثية كما يأتي:
\[
\tag{I}
s_x^m t_x^{n-\lambda}t_y^\lambda
=\sum_{\varrho}
\frac{\binom{m}{\varrho}\binom{\lambda}{\varrho}}
 {\binom{m+n-\varrho+1}{\varrho}}\cdot
\bigl[(st\widehat{xy})^\varrho s_x^{m-\varrho}t_x^{n-\varrho}\bigr]_{y^{\lambda-\varrho}}.
\]

نجمع التوسعتين، بأن نقيم للحدوديات المتجانسة التي لكل منها صفان من المتغيرات متخالفة التدرج:
\[
s_x^m t_x^{n-\lambda}t_y^\lambda u_\sigma^\mu u_\tau^{\nu-\varkappa}v_\tau^\varkappa
\]
توسعات بحسب قوى $u_x$، تكون معاملاتها أشكالا قطبية مركبة لسلسلة الملازمات $s_x^m t_x^n u_\sigma^\mu u_\tau^\nu$، مأخوذة بالنسبة إلى المتغيرين $x$ و$u$.
يكفي أن نقيم المتسلسلة لكل صفين من صفوف الرموز (أو المتغيرات، على الترتيب) المتخالفة التدرج، إذ يمكن رد الحدوديات المتجانسة ذات عدد أكبر من صفوف الرموز (أو المتغيرات) إلى هذه الحالة، بجمع كل صفين من صفوف الرموز (أو المتغيرات) في صف جديد واحد.
ولكي نجعل جميع ملازمات السلسلة لا تحتوي إلا على صفين من الرموز أو المتغيرات في كل حالة، نجري التخصيصات الآتية:
\[
\begin{array}{ll}
\text{a) } \sigma=\widehat{st}, & \text{a') } y=\widehat{uv},\\[0.3em]
\text{b) } s=\widehat{\sigma\tau}, & \text{b') } v=\widehat{xy},
\end{array}
\]
وننفذ التوسعة للتخصيصين a)، a').
بالنقل المباشر للتوسعة I نحصل على أشكال قطبية مركبة للحدوديات المتجانسة $(stu)^\mu s_x^m t_x^{n-\lambda}t_y^\lambda$؛ أما للحدوديات المتجانسة $(stu)^\mu u_\tau^{\nu-\varkappa}v_\tau^\varkappa s_x^m t_x^{n-\lambda}t_y^\lambda$ فتنشأ، بدلا من الأشكال القطبية المركبة، حدود يستلزم تقييمها إدخال متغيرات مساعدة جديدة باستمرار، لا يساوى بعضها ببعض إلا في النهاية، مما يعقد الحساب بلا ضرورة.
ولذلك نسلك طريقا غير مباشر، فنعبر عن الأشكال القطبية المركبة لجميع ملازمات السلسلة، أي العبارات
\[
\bigl[s_\tau^\varrho(stu)^{\mu-\varrho+r}\bigr]_{y^\lambda}^{v^\varkappa}
\]
بدلالة الحدود الابتدائية لهذه الأشكال القطبية، ونحصل، بعكس منظومة المعادلات الارتدادية الناشئة، على التوسعة المنشودة بحسب الأشكال القطبية.
وفقا لمبدأ النقل، تصح للشكل القطبي ذي الرتبة $\lambda$ من الحدودية المتجانسة $s_x^m t_x^n:[s_x^m t_x^n]_{y^\lambda}$ ($\lambda\le n$) التوسعة الآتية\srcfn{*)}{Gordan, Die Resultante binärer Formen (Kap. I, § 5). Rend. Circ. Mat. di Palermo XXII. 1906.} (أما الحالة $\lambda>n$ فترد، بتوسعة $[s_y^m t_y^n]_{x^{m+n-\lambda}}$، إلى الحالة الأولى):
\[
\bigl[s_x^m t_x^n\bigr]_{y^\lambda}
=\sum_r (-1)^r\,
\frac{\binom{m}{r}\binom{\lambda}{r}}{\binom{m+n}{r}}\cdot
(st\widehat{xy})^r s_x^{m-r}t_x^{n-\lambda}t_y^{\lambda-r},
\]
وبالمثل للشكل القطبي ذي الرتبة $\varkappa$ من $u_\sigma^\mu u_\tau^\nu:[u_\sigma^\mu u_\tau^\nu]^{v^\varkappa}$
\[
\bigl[u_\sigma^\mu u_\tau^\nu\bigr]^{v^\varkappa}
=\sum_\varrho (-1)^\varrho\,
\frac{\binom{\mu}{\varrho}\binom{\varkappa}{\varrho}}{\binom{\mu+\nu}{\varrho}}\,
(\sigma\tau\widehat{uv})^\varrho u_\sigma^{\mu-\varrho}u_\tau^{\nu-\varkappa}v_\tau^{\varkappa-\varrho}.
\]

وبالضرب نحصل، مع إدخال التخصيصين a) وa')، على
\[
\begin{aligned}
\bigl[s_x^m t_x^n(stu)^\mu u_\tau^\nu\bigr]_{\widehat{uv}^{\lambda}}^{v^\varkappa}
={}&
\sum_{r,\varrho}c_{r\varrho}\cdot
(stx\widehat{uv})^r(\widehat{st}\tau\widehat{uv})^\varrho
s_x^{m-r}t_x^{n-\lambda}(tuv)^{\lambda-r}(stu)^{\mu-\varrho}
u_\tau^{\nu-\varkappa}v_\tau^{\varkappa-\varrho},
\end{aligned}
\]
ومنها، بالتوسعة بحسب قوى $u_x$ مع تمييز الحد الأول (وتعني $\sum'_{r,\varrho}$ إغفال الحد $r=0$، $\varrho=0$) ومراعاة البند (a) في ص. 26، نحصل على
\[
\tag{II}
\begin{aligned}
&s_x^m t_x^{n-\lambda}(tuv)^\lambda(stu)^\mu u_\tau^{\nu-\varkappa}v_\tau^\varkappa\\
&\quad =
\bigl[s_x^m t_x^n(stu)^\mu u_\tau^\nu\bigr]_{\widehat{uv}^{\lambda}}^{v^\varkappa}\\
&\qquad
-\sum_{r,\varrho}' c_{r\varrho}\cdot s_\tau^\varrho(stu)^{\mu-\varrho+r}
u_\tau^{\nu-\varkappa}v_\tau^{\varkappa-\varrho}
s_x^{m-r}t_x^{n-\lambda}(tuv)^{\lambda-r+\varrho}\cdot v_x^r\\
&\qquad
+u_x\cdot\sum_{\varrho}\cdot\sum_{r=1}^{\lambda} c_{r\varrho}\cdot
s_\tau^\varrho(stu)^{\mu-\varrho+r-1}(stv)
u_\tau^{\nu-\varkappa}v_\tau^{\varkappa-\varrho}
s_x^{m-r}t_x^{n-\lambda}(tuv)^{\lambda-r+\varrho}\cdot v_x^{r-1}\\
&\qquad
\vdots\\
&\qquad
-(-1)^\lambda u_x^\lambda\cdot\sum_\varrho c_{\lambda\varrho}\cdot
s_\tau^\varrho(stu)^{\mu-\varrho}(stv)^\lambda
u_\tau^{\nu-\varkappa}v_\tau^{\varkappa-\varrho}
s_x^{m-\lambda}t_x^{n-\lambda}(tuv)^\varrho,
\end{aligned}
\]
حيث
\[
c_{r\varrho}=(-1)^{r+\varrho}\cdot
\frac{\binom{m}{r}\binom{\lambda}{r}}{\binom{m+n}{r}}\cdot
\frac{\binom{\mu}{\varrho}\binom{\varkappa}{\varrho}}{\binom{\mu+\nu}{\varrho}}
\qquad
\begin{array}{l}
\text{إذا كان }\lambda\le n,\\
\varkappa\le\nu.
\end{array}
\]
وبالمثل لسائر تجميعات الأعداد
\[\lambda,n;\varkappa,\nu.\]
والعبارة
\[s_x^m t_x^{n-\lambda}(tuv)^\lambda(stu)^\mu u_\tau^{\nu-\varkappa}v_\tau^\varkappa\]
تكون، إذن، قد ردت إلى شكل قطبي مركب، ومع استعمال المتطابقة
\[
(stv)u_\tau=(stu)v_\tau-s_\tau(tuv),
\]
إلى مجموع من عبارات منشأة على نحو مماثل، تقابل ملازمات أعلى في سلسلة الملازمات.
ومن خلال تكرار العملية نصل إلى التوسعة:
\[
\tag{III}
s_x^m t_x^{n-\lambda}(tuv)^\lambda(stu)^\mu u_\tau^{\nu-\varkappa}v_\tau^\varkappa
=
\sum_{p,r,\varrho}u_x^p\cdot C_{pr\varrho}\cdot
\bigl[s_\tau^\varrho(stu)^{\mu-\varrho+r}\bigr]_{\widehat{uv}^{\lambda-r+\varrho}v^{\varkappa-\varrho+p},\,v_x^{r-p}}.
\]
وتحسب المعاملات العددية $C_{pr\varrho}$ بصورة وحيدة وعلى نحو ارتدادي انطلاقا من المعاملات $c_{r\varrho}$، $c'_{r\varrho}$ لمنظومة المعادلات الناتجة من تكرار (II.) (انظر المثال ص. 42). وتنشأ توسعات مماثلة عند سائر التخصيصات.
\begin{center}
\textbf{§ 3.}\quad \emph{مبرهنات الاختزال.}
\end{center}

سنربط الآن المبرهنات المعروفة في اختزال الملازمات ومنظومات الملازمات بالفقرتين 1 و2؛ ولهذا تلزم بعض التعريفات المنقولة من نظرية الملازمات الثنائية.
أ) نسمي \emph{منظومة كاملة نسبيا} بترديد سلسلة معطاة من الملازمات منظومة ملازمات تتميز بأن جميع الإنشاءات الناتجة من طي حاصل ضرب اعتباطي لهذه الملازمات يمكن التعبير عنها بدوال نسبية صحيحة في ملازمات المنظومة، باستثناء حد جمعي مؤلف من تعبيرات نشأت عن طي ملازمات المنظومة بمنظومة المقياس\srcfn{*)}{انظر Gordan-Kerschensteiner, Vorlesungen über Invariantentheorie. Bd. II. S. 227.}.
ب) نسمي ملازما \emph{قابلا للاختزال} إذا أمكن التعبير عنه بملازمات تكون اللامتغيرات عوامل فيها، أو بـ\glqq{}ملازمات أعلى\grqq{}؛ أي بملازمات أعلى من سلسلة الملازمات الكلية التي ينتمي إليها الملازم، أو بملازمات تحتوي رموز المقياس برتبة أعلى.
يمكن الآن صياغة المبرهنات المتعلقة بالمختزلات\srcfn{**)}{انظر Gordan, Math. Annalen Bd. XVII. S. 222.} في أعم صورها على النحو الآتي:
\textbf{تعريف:} \emph{المختزل سلسلة ملازمات تكون قابلة للاختزال.}
\textbf{المبرهنة II.} \emph{إذا كان الملازم الابتدائي لسلسلة ملازمات قابلا للاختزال لأن أحد عناصره قد نشأ بطيه مع مختزل (أي له مختزل عامل)، وكان الملازم النهائي لسلسلة الملازمات قد نشأ بطي العنصر نفسه، فإن سلسلة الملازمات الكلية قابلة للاختزال}\srcfn{*)}{يمكن إيضاح التبسيط الذي تحققه المبرهنة II بالمثال الآتي: للملازم الخاص $f=x_1^3x_2+x_2^3x_3+x_3^3x_1$ يكون $a_y^2a_x^2u_y^2$ مختزلا. ومن ثم تستنتج المبرهنة II اختزال سلسلة الملازمات
\[
\theta_\nu(\vartheta\nu x)\theta_x^3u_\vartheta u_\nu^2
=
a_\nu^2(abu)-\underbrace{a_\nu b_\nu(aba)}_{0},
\]
لأن $\theta_\nu^4=a_\nu^2b_\nu^2(abu)^2$ نشأ من العنصر $a_\nu^2(abu)$ بالطي. أما عند Gordan, a. a. O. S. 231، فيجرى اختزال الملازمات
\[
\theta_\nu(\vartheta\nu x),\quad \theta_\nu(\vartheta\nu x)^2,\quad
\theta_\nu^2,\quad \theta_\nu^2(\vartheta\nu x),\quad
\theta_\nu^2(\vartheta\nu x)^2
\]
كل على حدة.}.
\textbf{البرهان:} تنشأ سلسلة الملازمات من العنصر الابتدائي بطيه في نفسه؛ أي بطي أعلى مع المختزل من جهة، وبطي المختزل في نفسه من جهة أخرى. وفي الحالتين، بحسب § 2، يمكن تمثيل جميع ملازمات السلسلة بأشكال قطبية (عمليات اشتقاق وتقليص) على سلسلة ملازمات المختزل.
لا يعود الاستدلال من الملازم الابتدائي إلى سلسلة الملازمات الكلية صالحا في طرائق الاختزال الأخرى. وهي:
1) ما يسمى \emph{\glqq{}الاختزال المزدوج\grqq{}}.
ونعني بذلك اختزال تعبير له مختزلان مستقلان أحدهما عن الآخر بوصفهما عاملين، بطريقتين، فتنتج علاقات بين الملازمات الأعلى التي انتهى إليها الاختزال.
بالتطبيق المنهجي لـ\glqq{}الاختزال المزدوج\grqq{} يجب، من جهة، الحصول على جميع العلاقات\srcfn{**)}{فمثلا، وجد بـ\glqq{}الاختزال المزدوج\grqq{} اختزال الملازم $a_\eta^3$ (§ 10)، وهو الملازم الوحيد الذي أدخله Maisano في المنظومة دون حاجة؛ ثم العلاقات المذكورة في حاشية المقدمة بين الملازمات الثلاثة الموافقة للتغير واللامتغيرات الثلاثة.}؛ ومن جهة أخرى، إذا لم ينتج \glqq{}الاختزال المزدوج\grqq{}، عند تطبيقه على تعبيرات مختلفة، علاقات جديدة، كان لدينا معيار لاستقلال الملازمات الأعلى وكذلك تحقق من الحساب.
2) \emph{الطي مع ملازمات قابلة للتحلل.}
يقصد بـ\glqq{}الملازمات القابلة للتحلل\grqq{} --- مع تعميم طفيف للمفهوم المألوف --- الملازمات التي يمكن التعبير عنها بحواصل ضرب ملازمات أدنى درجة وبـ\glqq{}ملازمات أعلى\grqq{}. وتتحول حواصل ضرب الملازمات عند طي واحد إلى حواصل ضرب؛ وكذلك حواصل ضرب الملازمات \emph{غير الخطية} عند طي مرتين في التخصيصين a') وb') (§ 2)، وفقا لمبرهنة حاصل الضرب:
\[
a_yb_y a_xb_x=\frac12a_y^2\cdot b_x^2+\frac12b_y^2\cdot a_x^2-\frac12(ab\widehat{yx})^2.
\]
ومن ثم فإن الأشكال القطبية الأولى (عمليات الاشتقاق والتقليص) وبعض الأشكال القطبية الثانية للملازمات القابلة للتحلل هي أيضا ملازمات قابلة للتحلل. وينتج:
إن \emph{عنصرا} ناشئا عن طي واحد (أو طيين، على الترتيب) مع ملازم قابل للتحلل، ضمن عملية اشتقاق وتقليص واحدة (أو مزدوجة، على الترتيب) على ملازم قابل للتحلل، يصبح بدوره ملازما قابلا للتحلل:
أ) إذا تحللت الملازمات الأعلى التالية في سلسلة ملازمات الملازم القابل للتحلل، أو صارت قابلة للاختزال؛ وكذلك إذا حدث، في حالة الطي الواحد، أحد التخصيصين $y=\widehat{uv}$ أو $v=\widehat{xy}$؛
ب) إذا أدت سائر الطيات الأحادية إلى ملازمات أعلى (انظر § 12، B. $H_k$ و$H_k(k\eta x)$).
وقد يصبح مثل هذا العنصر من عملية اشتقاق وتقليص ملازما قابلا للتحلل أيضا:
ج) إذا أدت سائر الطيات الأحادية إلى ملازمات أدنى تتحلل بصورة مستقلة عن الطي مع الملازم القابل للتحلل، أي يمكن التعبير عنها بحواصل ضرب وبالملازم المراد اختزاله؛ غير أنه يجب في كل مرة التحقق حسابيا أولا من أن المعامل العددي للعنصر المعني لا ينعدم بصورة مطابقة. وليس هذا إلا \glqq{}الاختزال المزدوج\grqq{} للملازم الأدنى (انظر § 23، C. $H_q(\theta H u)(\theta s u)$).
تنشأ ملازمات قابلة للتحلل من جميع الطيات الأحادية مع المحددات اليعقوبية\srcfn{*)}{انظر Gordan, a. a. O. § 3.}، وكذلك من بعض الطيات الأعلى. ونحصل على
\[
(st\widehat{yx})s_y=\frac12\,t\cdot s_y^2-\frac12\,s\cdot t_y^2+\frac12(st\widehat{yx})^2,
\]
\[
(st\widehat{yx})^2s_y^2=\frac13\,t\cdot s_y^3-\frac13\,s\cdot t_y^3+s_y(st\widehat{yx})^2-\frac13(st\widehat{yx})^3.
\]
وتصح صيغ مماثلة للملازمات الثنوية
\[
(\sigma\tau\widehat{vu})v_\sigma
\quad\text{و}\quad
(\sigma\tau\widehat{vu})v_\sigma^2.
\]
وفقا لمبرهنات هذه الفقرة، تترتب القاعدة الآتية لإنشاء جميع الملازمات غير القابلة للاختزال الناشئة من طي $S$ مع $T$:
نبدأ بأدنى الطيات، ونتقدم في تكوين سلسلة الملازمات $ST$ على مسار اعتباطي محدد مسبقا (مثلا سطرا فسطر أو متناظرا بالنسبة إلى العنصر القطري)، حتى نصل إلى ملازم له مختزل عامل. وتحذف سلسلة الملازمات التي يحددها هذا الملازم بمجرد أن يحقق الملازم النهائي الشرط المبين في المبرهنة II. وبعد استبعاد جميع سلاسل الملازمات القابلة للاختزال، تزال، بتطبيق الاختزال المزدوج، جميع الملازمات الأخرى القابلة للاختزال، وكذلك جميع الملازمات القابلة للتحلل. والمواضع في سلسلة الملازمات الكلية التي تغطيها مرتين سلسلتا ملازمات قابلة للاختزال مستقلتان إحداهما عن الأخرى تتيح اختزال ملازمات أعلى (انظر § 11، D).
\begin{center}
\textbf{الفصل الثاني. \quad المنظومات الأساسية الكاملة نسبيا.}
\textbf{§ 4.}\quad\emph{المنظومة الأولى الكاملة نسبيا.}
\end{center}

سنقتصر فيما يلي على الحدودية المتجانسة الثلاثية من الدرجة الرابعة، باستثناء المبرهنة الأولى. غير أن نتائج الفقرتين 5 و6 يمكن، عند اختيار المقاييس اختيارا مناسبا، نقلها بسهولة إلى حدوديات متجانسة من درجات أعلى، كما أنها تصح للحدودية المتجانسة الثلاثية من الدرجة الثالثة\srcfn{*)}{بالنسبة إلى الحدودية المتجانسة الثلاثية من الدرجة الثالثة، تكون الصيغ الموافقة، المشتقة بطريقة مماثلة، كما يأتي:
\[
(abc)a_x^2b_y^2c_z^2=\frac13(sxy)(sxz)(syz)+\frac23\Delta_x\Delta_y\Delta_z\cdot(xyz),\qquad (\mathrm{IV}.)
\]
\[
\begin{gathered}
s=(abc)(abu)(acu)(bcu),\qquad \Delta=(abc)^2a_xb_xc_x,\\
t=\Delta_\vartheta(\Delta\theta u)^2u_\vartheta=\theta_s^2u_su_\vartheta^2,
\qquad T=a_t^3,
\end{gathered}
\]
\[
\left.
\begin{aligned}
(a\theta u)^2&=\frac13 u_xs,\\
a_\vartheta&=\frac13u_x\cdot\Delta,
& a_\vartheta(a\theta u)&=0,
& a_\vartheta(a\theta u)^2&=s,\\
a_\vartheta^2&=\Delta,
& a_\vartheta^2(a\theta u)&=0.
\end{aligned}
\right\}\tag{1.}
\]
}.
يمكن الحصول على منظومة أولى كاملة نسبيا للحدوديات المتجانسة من درجة اعتباطية بموجب المبرهنة المعروفة:
\clearpage
\noindent\textbf{المبرهنة III.}
\emph{إن الحدوديات المنقولة من المجال الثنائي، أي الحدوديات التي تنشأ من حدوديات المنظومة الخاصة بالحدودية المتجانسة الثنائية الموافقة، وفقا لمبدأ كليبش للنقل، عن طريق متاخمة المحددات، تؤلف منظومة كاملة نسبيا بترديد $(abc)$.}
\noindent\textbf{البرهان:}
إن علاقة في مجال الحدوديات المتجانسة الثنائية تنص على أن الحدوديات $Q'_1,\ldots,Q'_n$ تؤلف منظومة كاملة لحدودية أساسية ثنائية:
\[
Q'-F(Q')=0=\sum_\lambda\bigl\{(ab)c_x+(bc)a_x+(ca)b_x\bigr\}\varphi_\lambda\srcfnmark{*)}
\]
\srcfntext{*)}{Gordan--Kerschensteiner, a. a. O. S. 134.}
وتقابل هذه العلاقة، بمتاخمة المحددات، العلاقة الآتية في مجال الحدوديات المتجانسة الثلاثية:
\[
Q-F(Q_i)=\sum_\lambda u_x\cdot(abc)\varphi_\lambda .
\]
وهذه هي المعادلة التعريفية لمنظومة كاملة نسبيا بترديد $(abc)$. وفي حالة الحدودية المتجانسة الثلاثية من الدرجة الرابعة، تتألف منظومة الحدوديات المنقولة من الحدوديات الآتية:
\[
\begin{gathered}
f=a_x^4,\qquad
\theta=\theta_x^4u_\vartheta^2=(abu)^2a_x^2b_x^2,
\qquad
K=K_x^6u_k^3=(a\theta u)u_\vartheta^2a_x^3\theta_x^3,\\
j=u_j^6=(a\theta u)^4u_\vartheta^2,\qquad
\nu=u_\nu^4=(abu)^4,
\end{gathered}
\]
وتؤلف الحدوديات الأربع الأولى منها منظومة كاملة نسبيا بترديد $((abc),\nu)$.
\begin{center}
\textbf{§ 5.}\quad\emph{اختزال المقياس $(abc)$ إلى المقياسين $\Delta$ و$\nu$.}
\end{center}

المقياس $(abc)$، المعطى بعامل واحد بين قوسين فقط، يُختزل ـ وفقًا للتعريف (§ 3a) ـ إلى حدوديات المنظومة. وبعبارة أخرى: ينبغي وضع سلسلة الملازمات $(abc)$ في صيغة نظامية بصورة وحيدة (انظر § 1).
\[
a_s^2a_xu_s=\frac13a_s^3\cdot u_x=\frac13S\cdot u_x. \tag{2.}
\]
\[
\left.
\begin{aligned}
(a\Delta u)^2&=a_s-\frac{S}{6}u_x^2,\\
(a\Delta u)^3&=0.
\end{aligned}
\right\} \tag{3.}
\]
\[
\left.
\begin{aligned}
\theta(s)=(ss_1x)^2&=2a_t+\frac13S\cdot\theta-\frac13Tu_x^2,\\
\theta(t)=(tt_1x)^2&=-\frac13S\cdot a_t+\frac23T\cdot a_s+\frac1{18}S^2\cdot\theta-\frac1{18}u_x^2\cdot S\cdot T.
\end{aligned}
\right\} \tag{4.}
\]

سلسلة المقاييس: $(abc)$؛ $\Delta$؛ $s$؛ $t$ (تتألف منظومة المقياس $t$ من الحدودية الوحيدة $t$).
سيتضح أن الحدوديتين:
\[
\Delta=(abc)^2a_x^2b_x^2c_x^2,\qquad \nu=(abu)^4
\]
تكفيان بوصفهما مقياسين، أي إن الحدوديات
\[
\begin{gathered}
\Delta=(abc)^2a_x^2b_x^2c_x^2,
\qquad a_\nu=(abc)(bcu)^3a_x^3,\\
a_\nu^2=(abc)^2(bcu)^2a_x^2,
\qquad i=a_\nu^4=(abc)^4
\end{gathered}
\]
تعرّف سلسلة الملازمات $(abc)$. ولا تنتمي الحدودية $a_\nu^3$ إلى سلسلة الملازمات $a_\nu$، وفقًا للمتطابقة:
\[
a_\nu^3=(abc)^3a_x(bcu)=\frac13u_x\cdot(abc)^4=\frac13i\cdot u_x.
\]

لاختزال مقياس إلى مقياس أعلى، علينا، وفقًا للتعريف، أن نختزل عمليات الطي مع المقياس ــ أعمها، أو جميع عمليات الطي الخاصة الداخلة في الاعتبار ــ إلى عمليات طي مع المقياس الأعلى.
في حالتنا، تنشأ أعم عمليات الطي من العبارات الآتية:
\[
(abc)a_x^3b_y^3c_z^3,
\qquad
(abc)a_x^2b_y^2c_z^2a_xb_xc_x
\]
(مأخوذتان من الحدوديات المتجانسة من الدرجة الثالثة)،
\[
(abc)a_xb_yc_za_x^2b_x^2c_x^2
\]
(مأخوذة من الحدوديات المتجانسة من الدرجة الثانية)، وبأخذ الأشكال القطبية لهذه العبارات.
لحساب هذه العبارات بواسطة الأشكال القطبية لـ$\Delta$ ولسلسلة الملازمات $a_\nu$، أو بالأحرى حدودها الابتدائية، نسلك، كما في § 2، الطريق غير المباشر. ونطوّر الحدود القطبية

\[
(abc)^2a_x^2b_y^2c_z^2,
\qquad
 a_\nu(\nu xy)(\nu yz)(\nu zx)a_xa_ya_z=(abc)(bc\widehat{xy})(bc\widehat{yz})(bc\widehat{zx})a_xa_ya_z\quad\text{إلخ.}
\]
بوصفها تراكيب خطية لعبارات ذات العامل الرمزي $(abc)$، فنحصل بعكس منظومة المعادلات على العلاقات المطلوبة، وكذلك على علاقة بين الأشكال القطبية المأخوذة بحسب توليفات مختلفة للمتغيرات. وللتقييم تُستخدم مبرهنة المتطابقات ومبرهنة حاصل ضرب المحددات.

تكون منظومة المعادلات، بالنسبة إلى $(abc)a_x^3b_y^3c_z^3$، كما يأتي:
\begingroup\small
\[
\begin{aligned}
1)\quad
\sum_3a_\nu(\nu yz)^3a_x^3
&=6(abc)a_x^3b_y^3c_z^3-6\sum_3(abc)a_x^3b_y^2b_zc_z^2c_y\srcfnmark{*)},\\[0.4em]
2)\quad
3(abc)^2a_x^2b_y^2c_z^2\cdot(xyz)-\frac12\sum_3a_\nu^2(\nu yz)^2a_x^2\cdot(xyz)
&=2\sum_3(abc)a_x^3b_y^2b_zc_z^2c_y\\
&\quad-4\sum_3(abc)a_xa_ya_zb_y^2b_xc_z^2c_x,\\[0.4em]
3)\quad
a_\nu(\nu xy)(\nu yz)(\nu zx)a_xa_ya_z
&=2\sum_3(abc)a_xa_ya_zb_y^2b_xc_z^2c_x,\\[0.4em]
4)\quad
\sum_3a_\nu(\nu yz)^3a_x^3-\frac32\sum_3a_\nu^2(\nu yz)^2a_x^2\cdot(xyz)+\frac16i\cdot(xyz)^3
&=6\sum_3(abc)a_xa_ya_zb_y^2b_xc_z^2c_x.
\end{aligned}
\]
\endgroup
\srcfntext{*)}{$\sum_3$ يشير إلى مجموع من ثلاثة حدود ناشئة من الحد الابتدائي بالتبديل الدوري للمتغيرات. وتصح المعادلات أيضًا لكل حد من حدود المجموع على حدة.}

ومن ذلك ينتج، بالنسبة إلى $(abc)a_x^3b_y^3c_z^3$، وبحساب مماثل بالنسبة إلى $(abc)a_x^2b_y^2c_z^2a_xb_xc_x$ و$(abc)a_xb_yc_za_x^2b_x^2c_x^2$:
\[
\begin{aligned}
(abc)a_x^3b_y^3c_z^3
&=\frac32(abc)^2a_x^2b_y^2c_z^2\cdot(xyz)
+\frac32a_\nu(\nu xy)(\nu yz)(\nu zx)a_xa_ya_z
-\frac1{12}i\cdot(xyz)^3,\\[0.4em]
(\mathrm{IV}.)\quad
(abc)a_x^2b_y^2c_z^2a_xb_xc_x
&=\frac23(abc)^2a_x^2b_yb_xc_zc_x\cdot(xyz)
+\frac13a_\nu(\nu xy)(\nu yz)(\nu zx)a_x^3\\
&\quad-\frac16a_\nu^2(\nu xy)(\nu zx)a_x^2\cdot(xyz),\\[0.4em]
(abc)a_xb_yc_za_x^2b_x^2c_x^2
&=\frac16(abc)^2a_x^2b_x^2c_x^2\cdot(xyz).
\end{aligned}
\]

وهكذا حصلنا على منظومة كاملة نسبيا بترديد $(\Delta,\nu)$، تتألف من الملازمات الأربع: $f,\theta,K,j$.

وينتج من ذلك: أن جميع عمليات طي $f$ على $\theta$ غير المنقولة من المجال الثنائي، وكذلك عملية الطي $(a\theta u)^2$ القابلة في الحالة الثنائية للاختزال إلى $(ab)^4$، يمكن التعبير عنها بالرمزين $\Delta$ و$\nu$.

ونحصل على صيغ الاختزال الأساسية لما يأتي، بتخصيص التوسعة IV، أو على نحو أقصر بحساب مباشر (تبديل الرموز)، للتعبيرات $a_\vartheta(a\theta u)^2$ و$a_\vartheta^2(a\theta u)^2$ و$(a\theta u)^2$ وفق الفرض الآتي:
\[
\begin{aligned}
a_\vartheta(a\theta u)^2
&=(abc)(bcu)(abu)^2-\frac13(abc)(bcu)\{(bcu)a_x-u_x\cdot(abc)\}^2,\\
a_\vartheta^2(a\theta u)^2
&=(abc)^2(abu)^2-\frac13(abc)^2\{(bcu)a_x-u_x\cdot(abc)\}^2.
\end{aligned}
\]
بالنسبة إلى $(a\theta u)^2$\srcfn{*)}{Gordan--Kerschensteiner، المرجع المذكور، ص. 92.}:
\[
\begin{aligned}
(abu)a_y^3b_x^3&=\frac32u_\vartheta(\vartheta yx)\theta_y^2+\frac14(\nu yx)^3,\\
(abu)(acu)^3b_x^3&=\frac32u_\vartheta(\vartheta\widehat{au}x)(\theta au)^2+\frac14(\nu\widehat{au}x)^3.
\end{aligned}
\]
\begingroup\small
\[
\begin{array}{r@{\quad}|l@{\qquad}c|}
& \underline{(a\theta u)^2}=\frac16\nu\cdot f-\frac23u_x\cdot a_\nu+\frac23u_x^2\cdot a_\nu^2-\frac{i}{18}u_x^4
& \bigstar\\[0.55em]
a_\vartheta=\frac13u_x\cdot\Delta
& \underline{a_\vartheta(a\theta u)}=0
& \\[0.55em]
& \underline{a_\vartheta^2}=\Delta
&
\end{array}
\]
\vspace{-0.5em}
\[
\begin{array}{r@{\quad}c@{\quad}|l@{\qquad}|l}
(1.) & \bigstar
& \underline{(a\theta u)^3}=0
& \underline{(a\theta u)^4}=j\\[0.55em]
&
& \underline{a_\vartheta(a\theta u)^2}=-\frac56a_\nu+\frac76u_x\cdot a_\nu^2-\frac{i}{9}u_x^3
& \underline{a_\vartheta(a\theta u)^3}=0\\[0.55em]
&
& \underline{a_\vartheta^2(a\theta u)}=0
& \underline{a_\vartheta^2(a\theta u)^2}=\frac23a_\nu^2-\frac{i}{9}u_x^2
\end{array}
\]
\endgroup

ونأخذ، إلى جانب ذلك، الصيغة الآتية، الناشئة هي أيضًا من اختزال المقياس $(abc)$:
\[
\begin{array}{r@{\qquad}l}
(2.) & a_\nu^3a_xu_\nu=\frac13i\cdot u_x.
\end{array}
\]

\emph{من الصيغتين (1.) و(2.)، ومن الصيغ المناظرة لهما والمكوّنة للحدودية الأساسية $\nu$، تُشتق جميع صيغ الاختزال اللاحقة بأخذ الأشكال القطبية}، باستثناء العلاقات الخاصة بالحدوديات القابلة للتحلل. ونرى من الصيغ (1.):

تؤدي سلسلة الملازمات كما يأتي:
\[
\begin{array}{rcl}
a_\vartheta(a\theta u)^2 &\text{إلى}& \text{الرمز }\nu\text{ وحده},\\[0.2em]
a_\vartheta &\text{إلى}& \text{الرمزين }\Delta\text{ و}\nu,\\[0.2em]
(a\theta u)^2 &\text{إلى}& \text{الرموز }j,\Delta\text{ و}\nu,\\[0.2em]
(a\theta u)^2\text{ عند الطي I} &\text{إلى}& \text{الرمزين }j\text{ و}\nu,\\[0.2em]
(a\theta u)^2\text{ عند الطي II} &\text{إلى}& \text{الرمزين }\Delta\text{ و}\nu.
\end{array}
\]

يمكننا إذن أن نستبدل:
\begingroup
\renewcommand{\labelenumi}{\arabic{enumi})}
\begin{enumerate}
\item بترديد $(\Delta,\nu)$: نستبدل عمليات الطي مع $j$ بعمليات مناظرة مع $(a\theta u)^2$ أو $(a\theta u)^3$؛
\item بترديد $(\nu)$: نستبدل عمليات الطي مع $\Delta$ بعمليات مناظرة مع $a_\vartheta$ أو $a_\vartheta(a\theta u)$، أو أيضًا بعمليات خاصة مع $(a\theta u)^2$ (لا تفضي إلا إلى طي واحد من النوع I).
\end{enumerate}
\endgroup
وبوجه خاص نحصل على:
\[
\begin{aligned}
(a\theta u)^2a_y^2\theta_x^2&=\frac13(jyx)^2\modu{\nu},&
(a\theta u)^2a_y\theta_y&=-\frac16(jyx)^2\modu{\nu},\\[0.4em]
(a\theta u)^2v_\vartheta^2&=\frac13(\Delta uv)^2\modu{\nu},&
(a\theta u)(a\theta v)u_\vartheta v_\vartheta&=-\frac16(\Delta uv)^2\modu{\nu},\\[0.4em]
(a\theta u)^2a_yv_\vartheta^2&=\frac13(\Delta uv)^2\Delta_y\modu{\nu}.
\end{aligned}
\]

\begin{center}
\textbf{§ 6.}\quad\emph{اختزال المقياس $(\Delta,\nu)$ إلى المقياس $(\nu)$.}
\end{center}

تتيح لنا صيغ الاختزال في الفقرة السابقة أن نقيم مباشرة المنظومة الكاملة نسبيا بترديد $\nu$؛ وسنرى أنه لا يلزم أن نضيف إلى منظومة الملازمات المنقولة إلا الملازمين $\Delta$ و$(a\Delta u)=N$ (كما في حالة الحدودية من الدرجة الثالثة، وهذا صحيح بوجه عام).

ولاختزال المقياس $\Delta$، ننظر في جميع عمليات الطي الخاصة الداخلة في الاعتبار، أي عمليات طي المنظومة الكاملة نسبيا بترديد $(\Delta,\nu)$ مع منظومة $\Delta$، ونبدأ بالملازمات الأقل درجة في المعاملات، أي عمليات طي $f$ مع $\Delta$.

ولاختزال الملازمات $(a\Delta u)^2$ و$(a\Delta u)^3$ و$(a\Delta u)^4$، نستبدل، وفق § 5، رموز $\Delta$ بملازمات أدنى من سلسلة الملازمات؛ أي إننا نوسّع التعبيرات
\[
a_\vartheta b_\vartheta(b\theta u)^2,
\qquad a_\vartheta(a\theta u)b_\vartheta(b\theta u)^2,
\qquad a_\vartheta(a\theta u)b_\vartheta(b\theta u)^3
\]
1) بحسب الأشكال القطبية لسلسلة الملازمات $a_\vartheta$ (الرمزان $\Delta$ و$\nu$)،

2) بحسب الأشكال القطبية لسلسلة الملازمات $b_\vartheta(b\theta u)^2$ (الرمز $\nu$)

ونحصل على صيغ الاختزال (انظر الحساب أدناه):
\[
\begin{aligned}
\text{a)}\quad \underline{(a\Delta u)^2}
&=\frac12(\vartheta \nu x)^2-\frac25f\cdot a_\nu^2
-\frac45u_x\cdot\theta_\nu(\vartheta \nu x)^2
+u_x^2\left\{\frac16i f-\frac1{40}s\right\},\\
(3.)\quad\text{b)}\quad \underline{(a\Delta u)^3}
&=-\frac3{10}\theta_\nu(\vartheta \nu x),\\
\text{c)}\quad \underline{(a\Delta u)^4}
&=\frac{21}{10}\theta_\nu^2-\frac3{10}H
-\frac65u_x\cdot\theta_\nu^3+\frac1{10}u_x^2\cdot\varrho.
\end{aligned}
\]

(وتؤدي أيضًا إلى النتائج نفسها عملية الاختزال المزدوج للتعبيرات $a_\vartheta^2(b\theta u)^2$ و$a_\vartheta^2(b\theta u)^3$ و$a_\vartheta^2(a\theta u)^2(b\theta u)^2$ و$a_\vartheta^2(a\theta u)(b\theta u)^3$ بعد حذف $a_\vartheta^2$.)
ومن الصيغ (3.) ينتج أن $(a\Delta u)^2$ مختزل بالنسبة إلى المقياس $\nu$. ومن ثم:

1) اختزال منظومة $\Delta$: تتخذ سلسلة الملازمات $(\Delta\Delta' u)^2\equiv a_\vartheta^2(a\Delta u)^2$ المختزل $(a\Delta u)^2$ عاملًا لها.

2) اختزال المنظومة الناشئة عن الطي مع المقياس $\Delta$:

تتخذ سلسلة الملازمات $(\theta\Delta u)^2$ المختزل $(a\Delta u)^2$ عاملًا لها.

بالنسبة إلى الملازمين $(\theta\Delta u)$ و$\Delta_\vartheta$ نحصل على:
\[
(a\Delta u)^2(abu)=(\theta\Delta u)-u_x\cdot\Delta_\vartheta(\theta\Delta u),
\]
\[
(a\Delta u)(a\Delta b)=-\frac12\Delta_\vartheta+\frac12u_x\cdot\Delta_\vartheta^2.
\]

لكن من المختزل $(\theta\Delta u)$ ينتج اختزال المنظومة الناشئة عن الطي مع المقياس $\Delta$.

\bigskip

وعليه تتألف المنظومة الكاملة نسبيا بترديد $\nu$ من الملازمات الستة:
\[
f,\ \theta,\ K,\ j,\ \A,\ N.
\]

وبوصفه مثالًا على التوسعة المتسلسلة في § 2، نورد اشتقاق الصيغة (3.)a بالتفصيل؛ أما في الحسابات اللاحقة فستقام مباشرة الصيغة النهائية III. (وقد حذفت الأشكال القطبية للملازمات المنعدمة أثناء الحساب؛ ويعطي السطر الأول القيم المدرجة بعد التوسعة II، ويعطي السطر الثاني القيم النهائية المستخرجة عوديا بدءًا من المعادلة الأخيرة في المنظومة.)

ووفقًا للتوسعة II ينتج:
\[
\begin{aligned}
\text{1)}\quad&m=3.\quad n=4.\quad \lambda=2.\quad \mu=0.\quad \nu=\varkappa=1,\\[0.25em]
\underline{a_\vartheta b_\vartheta(b\theta u)^2}
&=[a_\vartheta]_{\widehat{ub^2}}b+\frac67\{a_\vartheta b_\vartheta(a\theta u)(b\theta u)-u_x\cdot a_\vartheta b_\vartheta(a\theta b)(b\theta u)\}\\
&\quad-\frac17\{a_\vartheta(a\theta u)^2b_\vartheta-2u_x\cdot a_\vartheta(a\theta u)(a\theta b)b_\vartheta+u_x^2\cdot a_\vartheta(a\theta b)^2b_\vartheta\}\\
&=\frac23(a\A u)^2+\frac15[a_\vartheta(a\theta u)^2]b+\frac1{30}f\cdot[a_\vartheta^2(a\theta u)^2]-\frac25u_x\cdot[a_\vartheta(a\theta u)^2]b^2\\
&\quad-\frac1{15}u_x\cdot[a_\vartheta^2(a\theta u)^2]b+\frac15u_x^2\cdot[a_\vartheta(a\theta u)^2]b^3+\frac1{30}u_x^2\cdot[a_\vartheta^2(a\theta u)^2]b^2;\\[0.7em]
\text{2)}\quad&m=2.\quad n=3.\quad \lambda=1.\quad \mu=1.\quad \nu=\varkappa=1,\\[0.25em]
\underline{a_\vartheta(a\theta u)b_\vartheta(b\theta u)}
&=\frac25\{a_\vartheta(a\theta u)^2b_\vartheta-u_x\cdot a_\vartheta(a\theta u)(a\theta b)b_\vartheta\}+\frac12a_\vartheta^2(b\theta u)^2\\
&\quad-\frac15\{a_\vartheta^2(b\theta u)(a\theta u)-u_x\cdot a_\vartheta^2(a\theta b)(b\theta u)\}\\
&=\frac12(a\A u)^2+\frac25[a_\vartheta(a\theta u)^2]b+\frac1{15}[a_\vartheta^2(a\theta u)^2]\cdot f\\
&\quad-\frac25u_x\cdot[a_\vartheta(a\theta u)^2]b^2-\frac1{10}u_x\cdot[a_\vartheta^2(a\theta u)^2]b+\frac1{30}u_x^2\cdot[a_\vartheta^2(a\theta u)^2]b^2;\\[0.7em]
\text{3)}\quad&m=2.\quad n=3.\quad \lambda=1.\quad \mu=1.\quad \nu=1.\quad \varkappa=2,
\end{aligned}
\]
\[
\begin{aligned}
\underline{a_\vartheta(a\theta b)b_\vartheta(b\theta u)}
&=\frac25\{a_\vartheta(a\theta b)(a\theta u)b_\vartheta-u_x\cdot a_\vartheta(a\theta b)^2b_\vartheta\}\\
&=\frac25[a_\vartheta(a\theta u)^2]b^2+\frac1{30}[a_\vartheta^2(a\theta u)^2]b-\frac25u_x\cdot[a_\vartheta(a\theta u)^2]b^3\\
&\quad-\frac1{30}u_x\cdot[a_\vartheta^2(a\theta u)^2]b^2;\\[0.7em]
\text{4)}\quad&m=1.\quad n=2.\quad \lambda=0.\quad \mu=2.\quad \nu=\varkappa=1,\\[0.25em]
\underline{a_\vartheta(a\theta u)^2b_\vartheta}
&=[a_\vartheta(a\theta u)^2]b+\frac23a_\vartheta^2(a\theta u)(b\theta u)\\
&=[a_\vartheta(a\theta u)^2]b+\frac16[a_\vartheta^2(a\theta u)^2]\cdot f-\frac16u_x\cdot[a_\vartheta^2(a\theta u)^2]b;\\[0.7em]
\text{5)}\quad&m=1.\quad n=2.\quad \lambda=0.\quad \mu=2.\quad \nu=1.\quad \varkappa=2\quad(\varkappa>\nu),\\[0.25em]
\underline{a_\vartheta(a\theta u)(a\theta b)b_\vartheta}
&=[a_\vartheta(a\theta u)^2]b^2+\frac13a_\vartheta^2(a\theta b)(b\theta u)\\
&=[a_\vartheta(a\theta u)^2]b^2+\frac1{12}[a_\vartheta^2(a\theta u)^2]b-\frac1{12}u_x\cdot[a_\vartheta^2(a\theta u)^2]b^2;\\[0.7em]
\text{6)}\quad&m=1.\quad n=2.\quad \lambda=0.\quad \mu=2.\quad \nu=1.\quad \varkappa=3,\\[0.25em]
\underline{a_\vartheta(a\theta b)^2b_\vartheta}&=[a_\vartheta(a\theta u)^2]b^3;\\[0.7em]
\text{7)}\quad&m=2.\quad n=4.\quad \lambda=2.\quad \mu=\nu=\varkappa=0.\quad(\text{بعد التوسعة I}),\\[0.25em]
\underline{a_\vartheta^2(b\theta u)^2}&=[a_\vartheta^2]_{\widehat{ub^2}}+\frac1{10}\{[a_\vartheta^2(a\theta u)^2]\cdot f-2u_x\cdot[a_\vartheta^2(a\theta u)^2]b+u_x^2\cdot[a_\vartheta^2(a\theta u)^2]b^2\};\\[0.7em]
\text{8)}\quad&m=1.\quad n=3.\quad \lambda=1.\quad \mu=1.\quad \nu=\varkappa=0.\quad(\text{التوسعة I}),\\[0.25em]
\underline{a_\vartheta^2(a\theta u)(b\theta u)}&=\frac14[a_\vartheta^2(a\theta u)^2]\cdot f-\frac14u_x\cdot[a_\vartheta^2(a\theta u)^2]b;\\[0.7em]
\text{9)}\quad&m=1.\quad n=3.\quad \lambda=1.\quad \mu=\varkappa=1.\quad \nu=0.\quad(\text{التوسعة I}),\\[0.25em]
\underline{a_\vartheta^2(a\theta b)(b\theta u)}&=\frac14[a_\vartheta^2(a\theta u)^2]b-\frac14u_x\cdot[a_\vartheta^2(a\theta u)^2]b^2.
\end{aligned}
\]

ومن 1) و4) ينتج:
\[
\begin{aligned}
(a\A u)^2
&=\frac65[a_\vartheta(a\theta u)^2]b+\frac15f\cdot[a_\vartheta^2(a\theta u)^2]+\frac35u_x\cdot[a_\vartheta(a\theta u)^2]b^2-\frac3{20}u_x\cdot[a_\vartheta^2(a\theta u)^2]b\\
&\quad-\frac3{10}u_x^2\cdot[a_\vartheta(a\theta u)^2]b^3-\frac1{20}u_x^2\cdot[a_\vartheta^2(a\theta u)^2]b^2,\\[0.4em]
(a\A u)^2&=-a_\nu b_\nu+\frac35f\cdot a_\nu^2+\frac45u_x\cdot a_\nu^2b_\nu-\frac3{20}u_x^2\cdot a_\nu^2b_\nu^2-\frac1{12}u_x^2\cdot i\cdot f.
\end{aligned}
\]

(للتحويل إلى رموز $\theta$، قارن § 8 و§ 9.)
ويمكن حساب الصيغة (3.)a على نحو أقصر بالنقل المباشر للتوسعة I مع إدخال المتغير المساعد $w=x\widehat{u}b$، في حين أن الطريق الذي سلكناه يصير أبسط ابتداءً من الصيغتين (3.)b و(3.)c؛ وفي (3.)b نعلم مسبقًا، بموجب § 9، أن جميع الحدود ذات العامل $u_x$ يجب حذفها. ولحساب (3.)b و(3.)c نحصل على:
\[
\begin{aligned}
(3.)\mathrm{b)}\quad 1)&\quad a_\vartheta(a\theta u)b_\vartheta(b\theta u)^2=\frac12(a\A u)^3+\frac45[a_\vartheta(a\theta u)^2]_{\widehat{ub}}b+\frac{77}{360}[a_\vartheta^2(a\theta u)^2]_{\widehat{ub}},\\
2)&\quad a_\vartheta(a\theta u)b_\vartheta(b\theta u)^2=[a_\vartheta(a\theta u)^2]_{\widehat{ub}}b+\frac{13}{36}[a_\vartheta^2(a\theta u)^2]_{\widehat{ub}};\\[0.5em]
(3.)\mathrm{c)}\quad 1)&\quad a_\vartheta(a\theta u)b_\vartheta(b\theta u)^3=\frac12(a\A u)^4+\frac65[a_\vartheta(a\theta u)^2]_{\widehat{ub^2}}b+\frac{53}{120}[a_\vartheta^2(a\theta u)^2]_{\widehat{ub^2}}\\
&\qquad-\frac65u_x\cdot[a_\vartheta(a\theta u)^2]_{\widehat{ub^2}}b^2-\frac35u_x\cdot[a_\vartheta^2(a\theta u)^2]_{\widehat{ub^2}}b+\frac{19}{120}u_x^2\cdot[a_\vartheta^2(a\theta u)^2]_{\widehat{ub^3}}b^2,\\
2)&\quad a_\vartheta(a\theta u)b_\vartheta(b\theta u)^3=\frac34[a_\vartheta^2(a\theta u)^2]_{\widehat{ub^2}}.
\end{aligned}
\]

ملاحظة لاحقة: تحسب على نحو مماثل عمليات طي $f$ على $j$ القابلة للاختزال وفق § 5. فنحصل على
\[
\tag{4.}
\begin{aligned}
a_j&=-\theta_\nu+u_x\cdot\left\{\frac92\theta_\nu^2-\frac12H\right\}-3u_x^2\cdot\theta_\nu^3+\frac12u_x^3\cdot\varrho,\\
a_j^2&=\frac{21}{10}\theta_\nu^2-\frac3{10}H-2u_x\cdot\theta_\nu^3+\frac12u_x^2\cdot\varrho,\\
a_j^3&=-\frac7{10}\theta_\nu^3+\frac12u_x\cdot\varrho,\\
a_j^4&=\frac35\varrho.
\end{aligned}
\]

ومن (3.) و(4.) ينتج، بالنقل من المجال الثنائي،
\[
(\theta\theta'u)^2=\frac13 j\cdot f\modu{\nu}\srcfn{*)}{Gordan--Kerschensteiner، المرجع المذكور، ص. 181.}.
\]
وبقية ملازمات السلسلة $(\theta\theta'u)^2$، والملازم $\theta'_\vartheta$، قابلة للاختزال إلى رموز $\nu$. ومن ثم يمكننا أن نستبدل:
\[
\modu{\nu}:\quad \text{عمليات الطي مع } f\cdot j \text{ بعمليات مناظرة مع }(\theta\theta'u)^2.
\]
ونحصل كذلك على:
\[
(\vartheta\vartheta'x)^2=\frac43 f\cdot\A,\qquad
\theta_{\vartheta'}(\vartheta\vartheta'x)=-N.
\]

\begin{center}
{\Large\bfseries الفصل الثالث.\quad المنظومة الكاملة نسبيا بترديد $(s,\varrho,t)$.}
\end{center}

\begin{center}
\textbf{\S\ 7.}\quad\emph{نظرة عامة إلى تكوين المنظومة.}
\end{center}

للانتقال من المنظومة الكاملة نسبيا بترديد $\nu$ إلى المنظومة الكاملة نسبيا ذات المقياس الأعلى التالي، يجب، بحسب تعريف \S\ 3، تكوين منظومة الملازمات الخاصة بـ$\nu$ بالنسبة إلى هذا المقياس الأعلى وطيّها مع ملازمات المنظومة الكاملة نسبيا بترديد $\nu$. غير أن $\nu=u_\nu^4$ يقابل الحدودية $f$، على نحو ثنائي، بوصفه ملازمًا مخالفًا للتغير من الدرجة الرابعة. فإذا اعتبرنا $\nu$ حدودية أساسية، حصلنا على منظومة كاملة نسبيا مؤلفة من ستة ملازمات بترديد $\nu(\nu)=(\nu\nu_1x)^4=s_x^4$.

وعليه، لتكوين المنظومة الكاملة نسبيا بترديد $s$ (المنظومة III)، ينبغي أن نركّب
\[
\begin{array}{ll}
\text{المنظومة I:} & f,\ \theta,\ K,\ j,\ \A,\ N\quad \text{على}\vphantom{\dfrac11}\\[0.3em]
\text{المنظومة II:} & \nu,\ H=(\nu\nu_1x)^2u_\nu^2u_{\nu_1}^2,\ L=(\nu\eta x)H_x^2u_\nu^3u_\eta^3,\ g=(\nu\eta x)^4H_x^2,\\
& \sigma=H_\eta^2u_\nu^2u_\eta^4,\ (\nu\sigma x).
\end{array}
\]
وسيتبين (الصيغة (13.)) أن الملازم $g$ قابل للاختزال، أي إن المنظومة II لا تتألف إلا من خمسة ملازمات.

وفي أثناء الحساب تُلحق الحدوديتان المتجانستان من الدرجة الثانية
\[
\varrho=u_\varrho^2=\theta_\nu^4u_\vartheta^2,\qquad t=t_x^2=a_\eta^4H_x^2
\]
بوصفهما مقياسين، فنحصل على منظومة كاملة نسبيا بترديد $(s,\varrho,t)$؛ ويُعد المقياس $(\varrho,t)$ أعلى من المقياس $s$.

تُرتّب الملازمات المفردة بحسب رتبتها بالنسبة إلى المعاملات.

ونجعل عملية الطي
\[
\theta_\eta(K_\eta,\theta_\iota,\ldots)
\]
أعلى رتبة من عملية الطي
\[
H_\vartheta(H_\kappa,L_\vartheta,\ldots).
\]
وبذلك، بحسب \S\ 1 و\S\ 3b، يتحدد ترتيب الملازمات المفردة ذات الرتبة الواحدة تحديدًا وحيدًا.

يُجرى تكوين الملازمات ذات الرتبة الواحدة في جدول:

I. بيان أيّ ملازمات من المنظومتين I وII ينبغي طيّ بعضها مع بعض للحصول على الرتبة المعينة بالنسبة إلى المعاملات.

II. تدوين الملازمات غير القابلة للاختزال وفق مخطط سلسلة الملازمات.

III. اختزال سلاسل الملازمات أو الملازمات القابلة للاختزال أو القابلة للتحلل، أي المواضع التي تنقطع فيها سلسلة الملازمات الكاملة، وبذلك إثبات أن الإنشاءات غير القابلة للاختزال تنحصر في المذكورة.

IV. النتائج، وهي: 1) بيان المختزلات المستحدثة، 2) بيان الملازمات التي لا تدخل، في صورة جداءات مع ملازمات من المنظومة نفسها، في عمليات طي مع ملازمات المنظومة الأخرى.

وسيتبين أنه لا يظهر إلا ما يلي:

1) عمليات طي ملازم واحد من كل من المنظومتين I وII.

2) عمليات طي قوى $\theta$ أو $H$، على الترتيب، أو ملازمين من إحدى المنظومتين، مع ملازم واحد من المنظومة الثانية.

وسنحصل على المخطط الآتي للطي:
\[
\begin{array}{r@{\ }c@{\ }l@{\qquad}c@{\qquad}r@{\ }c@{\ }l}
\multicolumn{3}{c}{\text{المنظومة I}} & \text{مطوية مع} & \multicolumn{3}{c}{\text{المنظومة II}}\\[0.25em]
1. & \text{الرتبة} & f & & 2. & \text{الرتبة} & \nu\\
2. & ,, & \theta & & 4. & ,, & H\\
3. & ,, & K,j,\A & & 6. & ,, & L,\sigma\\
4. & ,, & N,\theta^2,f\cdot j & & 8. & ,, & (\nu\sigma x),H^2\\
5. & ,, & \theta\cdot K & & 10. & ,, & H\cdot L\\
6. & ,, & \theta^3,\A\cdot j & & 12. & ,, & H^3.
\end{array}
\]

وللتحقق من الحساب، يُستعان، إلى جانب «الاختزال المزدوج»، بالملازمين الخاصين الآتيين:
\[
\begin{array}{ll}
\text{I.}&\begin{array}{rlrlrl}
f&=\displaystyle\sum_3x_1^3x_2,& a_\nu^2&=\dfrac{i}{6}u_x^2,& s&=\dfrac{i}{3}f,\\
a_\eta^2&=-\dfrac19i\theta,& \A_\nu^2&=\dfrac{i}{6}\theta,& \sigma&=-\dfrac{i}{6}j,\\
g&=-\dfrac23i\A,& \varrho&=0,& t&=0,
\end{array}\\[1.2em]
\text{II.}&\begin{array}{rlrlrl}
f&=\displaystyle\sum_3x_1^4,& s&=\dfrac43if,& a_\eta^2&=\dfrac29i\theta,\\
\A_\nu^2&=\dfrac1{15}i\theta,& \sigma&=\dfrac43ij,& g&=\dfrac43i\A,\\
\varrho&=0,& t&=0.
\end{array}
\end{array}
\]

ملاحظة لاحقة: تقابل الصيغ \srcLTR{(1.)} و\srcLTR{(2.)} و\srcLTR{(3.)} و\srcLTR{(4.)}، بالنسبة إلى الحدودية الأساسية $\nu$، الصيغ الآتية:
\[
\tag{5.}
\begin{array}{@{}c|c|c@{}}
\cdot & \underline{H_\nu}=\dfrac13u_x\cdot\sigma & \cdot\\[0.45em]
\underline{(\nu\eta x)^2}=A & \underline{H_\nu(\nu\eta x)}=0 & \underline{H_\nu^2}=\sigma\\[0.45em]
\underline{(\nu\eta x)^3}=0 & \underline{H_\nu(\nu\eta x)^2}=B & \underline{H_\nu^2(\nu\eta x)}=0\\[0.45em]
\underline{(\nu\eta x)^4}=g & \underline{H_\nu(\nu\eta x)^3}=0 & \underline{H_\nu^2(\nu\eta x)^2}=C,
\end{array}
\]
\[
\begin{aligned}
A&=\frac16s\cdot\nu-\frac23u_x\cdot s_\nu+\frac23u_x^2\cdot s_\nu^2-\frac{J}{18}u_x^4,\\
B&=-\frac56s_\nu+\frac76u_x\cdot s_\nu^2-\frac{J}{9}u_x^3,\\
C&=\frac23s_\nu^2-\frac{J}{9}u_x^2.
\end{aligned}
\]
\[
\tag{6.}
s_\nu^3u_\nu s_x=\frac13u_x\cdot s_\nu^4=\frac13J\cdot u_x.
\]

\[
\tag{7.}
\begin{aligned}
\text{a)}\quad (\nu\sigma x)^2
&=\frac12(Hsu)^2-\frac25\nu\cdot s_\nu^2-\frac45u_x\cdot s_\eta(Hsu)^2+u_x^2\left\{\frac16J\cdot\nu-\frac1{40}(ss'u)^4\right\},\\
\text{b)}\quad (\nu\sigma x)^3&=-\frac3{10}s_\eta(Hsu),\\
\text{c)}\quad (\nu\sigma x)^4&=\frac{21}{10}s_\eta^2-\frac3{10}Z-\frac65u_x\cdot s_\eta^3+\frac1{10}u_x^2\cdot s_\eta^4.
\end{aligned}
\]
\[
\tag{8.}
\begin{aligned}
\text{a)}\quad g'_\nu&=-s_\eta+u_x\left\{\frac92s_\eta^2-\frac12Z\right\}-3u_x^2\cdot s_\eta^3+\frac12u_x^3\cdot s_\eta^4,\\
\text{b)}\quad g_\nu^2&=\frac{21}{10}s_\eta^2-\frac3{10}Z-2u_x\cdot s_\eta^3+\frac12u_x^2\cdot s_\eta^4,\\
\text{c)}\quad g_\nu^3&=-\frac7{10}s_\eta^3+\frac12u_x\cdot s_\eta^4,\\
\text{d)}\quad g_\nu^4&=\frac35s_\eta^4.
\end{aligned}
\]

\begin{center}
\textbf{\S\ 8.}\quad\emph{الملازمات من الرتبة الثالثة (المنظومة III).}
\end{center}

\noindent\textbf{طيّ $f$ مع $\nu$.}

\noindent\emph{الملازمات غير القابلة للاختزال:}
\[
\begin{array}{cccc}
a_\nu&\cdot&\cdot&\cdot\\
\cdot&a_\nu^2&\cdot&\cdot\\
\cdot&\cdot&\cdot&\cdot\\
\cdot&\cdot&\cdot&a_\nu^4=i.
\end{array}
\]

\noindent\emph{الاختزال:}
\[
a_\nu^3=\frac13iu_x\qquad\text{(الصيغة (2.)).}
\]
\noindent\emph{النتائج:} 1) $a_\nu^3$ هو مختزل. 2) لا يدخل $f$ إلا في جداءات مع الملازمات المخالفة للتغير $(j)$ عند طيّه مع $\nu$. وينطبق الأمر نفسه على $\nu$ بالنسبة إلى $f$.

\begin{center}
\textbf{\S\ 9.}\quad\emph{الملازمات من الرتبة الرابعة (المنظومة III).}
\end{center}

\noindent\textbf{طيّ $\theta$ مع $\nu$.}

\noindent\emph{الملازمات غير القابلة للاختزال:}
\[
\begin{array}{cccc}
(\vartheta\nu x)&\theta_\nu&\cdot&\cdot\\
(\vartheta\nu x)^2&\theta_\nu(\vartheta\nu x)&\theta_\nu^2&\cdot\\
\cdot&\theta_\nu(\vartheta\nu x)^2&\cdot&\theta_\nu^3.
\end{array}
\]

\noindent\emph{الاختزالات:} انطلاقًا من المختزل $a_\nu^3$، وبالاختزال المزدوج للتعبيرات
$a_\nu^3(abu)$، $a_\nu^3b_\nu$، $a_\nu^3b_\nu(abu)$، وفق الفرض:
\[
\begin{aligned}
\theta_\nu^2(\vartheta\nu x)
&=a_\nu^2(\widehat{ab}\nu x)(abu)-\frac13(\nu\nu_1x)^3
=a_\nu b_\nu(\widehat{ab}\nu x)(abu)+\frac16(\nu\nu_1x)^3\\
&=a_\nu^3(abu)-a_\nu^2b_\nu(abu)=2a_\nu^2b_\nu(abu)=\frac23a_\nu^3(abu),\\[.4em]
\theta_\nu^2(\vartheta\nu x)^2
&=a_\nu^2(\widehat{ab}\nu x)^2-\frac13(\nu\nu_1x)^4
=a_\nu b_\nu(\widehat{ab}\nu x)^2+\frac16(\nu\nu_1x)^4\\
&=i f-2a_\nu^3b_\nu+a_\nu^2b_\nu^2-\frac13s
=2a_\nu^3b_\nu-2a_\nu^2b_\nu^2+\frac16s
=\frac23i f-\frac23a_\nu^3b_\nu-\frac16s,\\[.4em]
\theta_\nu^3(\vartheta\nu x)&=a_\nu^2b_\nu(\widehat{ab}\nu x)(abu)=a_\nu^3b_\nu(abu).
\end{aligned}
\]

تنتج الصيغ:
\[
\tag{9}
\begin{array}{c|c|c}
\underline{\theta_\nu^2(\vartheta\nu x)}=0
&\underline{\theta_\nu^3(\vartheta\nu x)}=0
&\underline{\theta_\nu^4u_\vartheta^2}=u_\varrho^2=\varrho\quad(\text{المقياس الملحق، \S\ 7}),\\[0.6em]
\underline{\theta_\nu^2(\vartheta\nu x)^2}=\dfrac49 i f-\dfrac16s
&&
\end{array}
\]

\noindent\emph{النتائج:} 1) $\theta_\nu^2(\vartheta\nu x)^2$ هو مختزل.
2) لا يدخل $\nu$ إلا في جداءات مع الملازمات المخالفة للتغير $(g)$ عند طيّها مع $\theta$.

\begin{center}
\textbf{\S\ 10.}\quad\emph{الملازمات من الرتبة الخامسة (المنظومة III).}
\end{center}

\[
\text{طيّ }\left\{\begin{array}{l}
K,j,\A \text{ مع }\nu,\\
f\text{ مع }H.
\end{array}\right.
\]

\noindent\emph{الملازمات غير القابلة للاختزال:}
\[
\begin{array}{c@{\qquad}c@{\qquad}c@{\qquad}c}
\text{A)} & \begin{array}{cccc}
\cdot&K_\nu&\cdot&\cdot\\
(k\nu x)^2&\cdot&K_\nu^2&\cdot\\
(k\nu x)^3&K_\nu(k\nu x)^2&\cdot&\cdot\\
\cdot&K_\nu(k\nu x)^3&\cdot&\cdot
\end{array}
&
\text{B)}\ \begin{array}{c}(j\nu x)\\(j\nu x)^2\\(j\nu x)^3\\ \cdot\end{array}
&
\text{C)}\ \begin{array}{c}\A_\nu\\\cdot\\\cdot\\\cdot\end{array}
\end{array}
\]
\[
\text{D)}\qquad
\begin{array}{cccc}
(aHu)&(aHu)^2&\cdot&\cdot\\
a_\eta&a_\eta(aHu)&a_\eta(aHu)^2&\cdot\\
\cdot&a_\eta^2&\cdot&\cdot
\end{array}
\]

\noindent\emph{الاختزالات:}

\noindent A) سلسلة الملازمات $K_\nu^3$: المختزل $a_\nu^3$.

الملازم $K_\nu^2(k\nu x)^2$: المختزل $\theta_\nu^2(\vartheta\nu x)^2$.
\[
(k\nu x),\quad K_\nu(k\nu x),\quad K_\nu^2(k\nu x)
\]
هي ملازمات قابلة للتحلل بحسب \S\ 3.

\noindent B) الملازم
\[
\begin{aligned}
(j\nu x)^4&=(\widehat{a\theta}\nu x)^4
=a_\nu^4\theta+\theta_\nu^4 f-4a_\nu^3\theta_\nu
-4a_\nu\{a_\nu\theta_x-(\widehat{a\theta}\nu x)\}^3\\
&\quad+6a_\nu^2\{a_\nu\theta_x-(\widehat{a\theta}\nu x)\}^2
=\R f+\frac53 i\theta-6a_\nu^2(\widehat{a\theta}\nu x)^2
+4a_\nu(\widehat{a\theta}\nu x)^3,
\end{aligned}
\]
ومن ثم
\[
(j\nu x)^4=\R f+\frac53 i\theta\modu{(\A,\nu)},\qquad (\text{انظر خاتمة \S\ 5}).
\]

\noindent C) الملازم $\A_\nu^4$: المختزل $\theta_\nu^4$.

الملازمان $\A_\nu^2$ و$\A_\nu^3$: المختزل $a_\nu^3$ أو $\theta_\nu^4$، وذلك باستبدال الملازم $\A$ بملازمات أدنى من سلسلة الملازمات (خاتمة \S\ 5)، أي بالاختزال المزدوج للتعبيرات
\[
a_\vartheta a_\nu^3,\qquad
a_\vartheta a_\nu^3\theta_\nu,\qquad
a_\vartheta\theta_\nu^4.
\]
يؤدي الحساب المزدوج لـ$\A_\nu^3$ إلى اختزال الملازم الأعلى $a_\eta^3$.

\noindent\emph{صيغ الاختزال:}
\[
\tag{10}
\begin{aligned}
\text{a)}\quad \A_\nu^2
&=\frac35a_\eta^2-\frac3{10}(asu)^2+\frac13 i\theta
+\frac15u_x^2(2a_\R^2-t),\\
\text{b)}\quad \A_\nu^3
&=-\frac3{10}a_\R+\frac35u_x\left(\frac{13}{10}a_\R^2-\frac25t\right),\\
\text{c)}\quad \A_\nu^4&=a_\R^2-\frac25t.
\end{aligned}
\]

\noindent\emph{اشتقاق الصيغ:}
\[
\begin{aligned}
\text{a)}\quad
 a_\vartheta a_\nu^3
&=\frac13 i\theta
=[a_\vartheta]_{\nu^3}+\frac67[a_\vartheta^2]_{\nu^2}
+\frac65[a_\vartheta(a\theta u)^2]_{\nu^2}\,\widehat{\nu}x^2
+\frac{11}{30}[a_\vartheta^2(a\theta u)^2]_{\nu}\,\widehat{\nu}x^2\\
&\quad-\frac67u_x[a_\vartheta^2]_{\nu^3}
-\frac16u_x[a_\vartheta^2(a\theta u)^2]_{\nu^2}\,\widehat{\nu}x^2,\\
\frac13 i\theta
&=\A_\nu^2-\frac23u_x\A_\nu^3-a_\nu a_{\nu_1}(\nu\nu_1x)^2
+\frac25a_\nu^2(\nu\nu_1x)^2
+\frac15u_xa_\nu^2a_{\nu_1}(\nu\nu_1x)^2;
\end{aligned}
\]
\[
\begin{aligned}
\text{b)}\quad
 a_\vartheta a_\nu^3\theta_\nu
&=0=[a_\vartheta]_{\nu^4}+\frac9{14}[a_\vartheta^2]_{\nu^3}
+\frac35[a_\vartheta(a\theta u)^2]_{\nu^2}\,\widehat{\nu}x^2
+\frac{17}{120}[a_\vartheta^2(a\theta u)^2]_{\nu}\,\widehat{\nu}x^2\\
&\quad-\frac9{14}u_x[a_\vartheta^2]_{\nu^4}
-\frac1{24}u_x[a_\vartheta^2(a\theta u)^2]_{\nu^2}\,\widehat{\nu}x^2,\\
0&=\frac56\A_\nu^3-\frac12u_x\A_\nu^4-\frac14a_\eta^3+\frac1{20}u_xa_\eta^4;
\end{aligned}
\]
\[
\begin{aligned}
\text{b')}\quad
 a_\vartheta\theta_\nu^4
&=a_\R=a_\vartheta\theta_\nu\{a_\nu-(a\theta\widehat{\nu}x)\}^3\\
&=-\frac32[a_\vartheta^2]_{\nu^3}
+\frac35[a_\vartheta(a\theta u)^2]_{\nu^2}\,\widehat{\nu}x^2
-\frac{37}{120}[a_\vartheta^2(a\theta u)^2]_{\nu}\,\widehat{\nu}x^2\\
&\quad+\frac32u_x[a_\vartheta^2]_{\nu^4}
+\frac{49}{120}u_x[a_\vartheta^2(a\theta u)^2]_{\nu^2}\,\widehat{\nu}x^2,\\
a_\R&=-\frac32\A_\nu^3+\frac32u_x\A_\nu^4-\frac{11}{20}a_\eta^3+\frac7{20}u_xa_\eta^4;
\end{aligned}
\]
\[
\begin{aligned}
\text{c)}\quad
 a_\vartheta^2\theta_\nu^4
&=a_\R^2=[a_\vartheta^2]_{\nu^4}+\frac35[a_\vartheta^2(a\theta u)^2]_{\nu^2}\,\widehat{\nu}x^2,\\
a_\R^2&=\A_\nu^4+\frac25a_\eta^4.
\end{aligned}
\]

\noindent D) اختزال $a_\eta^3$ عن طريق الاختزال المزدوج لـ$\A_\nu^3$ (C).

أما سائر الاختزالات فتُستخرج بحساب يقابل، على نحو ثنائي، الحساب الوارد في \S\ 9.

\noindent\emph{صيغ الاختزال:}
\[
\tag{11}
\begin{array}{rl|rl}
\underline{a_\eta^2(aHu)}&=-\dfrac16(asu)^3,&
\underline{a_\eta^2(aHu)^2}&=\dfrac49 i \nu-\dfrac16(asu)^4,\\[0.7em]
\underline{a_\eta^3}&=-a_\R+\dfrac35u_xa_\R^2+\dfrac15u_xt,
&\underline{a_\eta^3(aHu)}&=0,\\[0.7em]
&&\underline{a_\eta^4}&=t\quad(\text{ملحق بوصفه مقياسًا}).
\end{array}
\]

\noindent\emph{النتائج:}
1) تُعَدّ $(j\nu x)^4$ و$\A_\nu^2$ و$a_\eta^2(aHu)$ مختزلات.

2) لا يدخل $\nu$ إلا في جداءات مع الملازمات المخالفة للتغير عند طيّها مع $K,j,\A$؛ وينطبق الأمر نفسه على $f$ بالنسبة إلى $H$، وعلى $j$ بالنسبة إلى $\nu$، لأن $\theta_\nu(\vartheta\cdot j,\nu,x)^3$ يصبح قابلًا للاختزال بحسب \S\ 11 B.

\begin{center}
\textbf{\S\ 11.\quad الملازمات من الرتبة السادسة (المنظومة III).}
\end{center}
\[
\text{طيّ }\left\{\begin{array}{l}
\theta^2,\ f\cdot j,\ N \text{ مع }\nu,\\
\theta\text{ مع }H.
\end{array}\right.
\]


\noindent\emph{الملازمات غير القابلة للاختزال:}
\[
\renewcommand{\arraystretch}{1.25}
\begin{array}{@{}r@{\quad}c@{\qquad\qquad}r@{\quad}c@{\qquad\qquad}r@{\quad}c@{}}
\text{A)}&(\vartheta^2\nu x)^3&\text{B)}&a_\nu(j\nu x)&\text{C)}&N_\nu,\\[0.45em]
&\cdot&&\cdot&&\\[-0.1em]
&\theta_\nu(\vartheta^2\nu x)^3&&a_\nu^2(j\nu x)&&
\end{array}
\]

\[
\renewcommand{\arraystretch}{1.25}
\begin{array}{@{}r@{\quad}c|c|c|c@{}}
&\cdot&(\theta Hu)&(\theta Hu)^2&\cdot\\
\text{D)}&(\vartheta\eta x)&\theta_\eta&H_\vartheta(\theta Hu),\ \theta_\eta(\theta Hu)&\theta_\eta(\theta Hu)^2\\
&(\vartheta\eta x)^2&H_\vartheta(\vartheta\eta x),\ \theta_\eta(\vartheta\eta x)&\theta_\eta^2&\theta_\eta^2(\theta Hu)\\
&\cdot&\theta_\eta(\vartheta\eta x)^2&\cdot&\cdot
\end{array}
\]
\noindent\emph{الاختزالات:}


\noindent A) يتحلل $(\vartheta^2\nu x)^4$ بحسب الصيغة (3.)a:
\[
(a\A\widehat{\vartheta}x)^2=\frac12(\vartheta^2\nu x)^4-\frac35 f\cdot a_\nu^2(\nu\vartheta x)^2
=\frac13\A^2+f\cdot\A_\vartheta^2.
\]

\noindent B) يتحلل $a_\nu(j\nu x)^2$ بحسب الصيغة (9.)a، إذ نستبدل، وفق خاتمة \S\ 6، $f\cdot j$ بـ$(\theta\theta'u)^2$، ونراعي قابلية $\theta_\vartheta'$ للاختزال:
\[
\frac13a_\nu(j\nu x)^2=(\widehat{\theta\theta'}\nu x)^2\theta_\nu
=\theta_\nu^3\cdot\theta-\theta_\nu^2\theta_\nu'
=\theta_\nu^3\cdot\theta+\theta_\nu^2(\widehat{\vartheta}\nu\theta'u)
=\theta_\nu^3\cdot\theta-\frac23\theta_\nu^3\cdot\theta.
\]
الملازمان $a_\nu(j\nu x)^3$ و$a_\nu^2(j\nu x)^2$؛ والمختزلان $a_j$ و$(j\nu x)^4$:
\[
a_\nu(j\nu x)^3=a_j(j\nu x)^3-(j\nu x)^3(j\nu\widehat{a}u),
\]
\[
a_\nu^2(j\nu x)^2=a_j^2(j\nu x)^2-2a_j(j\nu x)^2(j\nu\widehat{a}u)+(j\nu x)^2(j\nu\widehat{a}u)^2.
\]

\noindent C) سلسلة الملازمات $N_\nu^2$؛ المختزل $\A_\nu^2$.
يتحلل $(n\nu x)$ و$N_\nu(n\nu x)$، كلٌّ منهما بوصفه عملية الاشتقاق والتقليص على المحددات اليعقوبية، بحسب \S\ 3.

\noindent D) \emph{نظرة عامة إلى الاختزالات:}

a) سلسلة الملازمات $\theta_\eta^2H_\vartheta$: المختزل $a_\eta^2(aHu)$. (تفضي إلى ملازمات ذات رموز أعلى.)

b) سلسلة الملازمات $\theta_\eta H_\vartheta$: المختزل $a_\eta^2(aHu)$ بالاختزال المزدوج.
(تفضي إلى ملازمات ذات رموز أعلى، وإلى ملازمات أعلى في سلسلة الملازمات الكلية، أي إلى سلسلة الملازمات $\theta_\eta^2$.) وبدلًا من $\theta_\eta H_\vartheta$ نعتبر سلسلة الملازمات
\[
\theta_\eta(\vartheta\eta x)(\theta Hu)=\theta_\eta H_\vartheta+\theta_\eta^2,
\]
ونستبدل فيها الرموز $\theta$ بالملازم الأدنى $(abu)$، فنصل بذلك إلى الاختزال المزدوج لسلسلة الملازمات
\[
a_\eta^2(aHu)(abu)=\theta_\eta H_\vartheta+\frac32\theta_\eta^2\modu{s}
\]
حتى الملازم النهائي
\[
a_\eta^3b_\eta(bHu)(abu)=\theta_\eta^3H_\vartheta+\frac12\theta_\eta^4.
\]

c) سلسلة الملازمات $\theta_\eta^3$: بالاختزال المزدوج للملازمات التي تنتمي في آن واحد إلى سلسلتي الملازمات $\theta_\eta H_\vartheta$ و$\theta_\eta^2H_\vartheta$ ـ أي ملازمات سلسلة الملازمات $\theta_\eta^2H_\vartheta$ ـ إلى ملازمات ذات رموز أعلى من جهة، وإلى ملازمات ذات رموز أعلى وإلى سلسلة الملازمات الأعلى $\theta_\eta^3$ من جهة أخرى.

d) الملازمات $\theta_\eta^2(\vartheta\eta x)$ و$\theta_\eta^2(\vartheta\eta x)^2$ و$\theta_\eta^3(\vartheta\eta x)$: بالاختزال المزدوج باستخدام المختزل $a$، على نحو مماثل لحساب \S\ 9.

e) الملازمان $\theta_\eta^2(\theta Hu)^2$ و$\theta_\eta^2(\vartheta\eta x)^2$: بالاختزال المزدوج للتعبيرين $\A_\nu^2(a\A u)^4$ و$(a\A\widehat{\nu}x)^4$ بواسطة المختزلين $\A_\nu^2$ و$(a\A u)^4$.

f) الملازمان $H_\vartheta$ و$H_\vartheta^2$: ملازمان قابلان للتحلل. ويتبين ذلك لـ$H_\vartheta^2$ بالاختزال المزدوج للتعبير $\A_\nu^2(a\A u)^2$، أو أيضًا بالحساب المباشر للجداءين $(a_\nu^2)^2$ و$a_\nu^2\cdot b_{\nu_1}$ بدلالة ملازمات المنظومة. (ويعطي الحساب المماثل لـ$(a_\nu)^2$ علاقة بين ملازمات قابلة للتحلل.)

g) اختزال الملازمات الأعلى من كون ملازمات سلسلة الملازمات $\theta\cdot H$ تقبل طريقتين للاختزال (أي تنتمي إلى سلسلتي ملازمات قابلتين للاختزال)، وذلك كما يأتي:
\[
\begin{array}{rl}
g&\text{بالاختزال المزدوج لـ }\theta_\eta^2(\vartheta\eta x)^2:\text{ يقبل الاختزالين d) وe)؛}\\[0.3em]
s_\vartheta^2(\theta s u)&\text{بالاختزال المزدوج لـ }\theta_\eta^3(\vartheta\eta x):\text{ يقبل الاختزالين c) وd)؛}\\[0.3em]
s_\vartheta^2(\theta s u)^2&\text{بالاختزال المزدوج لـ }\theta_\eta^2H_\vartheta^2:\text{ يقبل الاختزال a) ويكون}\\
&\text{المختزل }\theta_\nu^2(\vartheta\nu x)^2\text{ عاملًا فيه؛}\\[0.3em]
s_\nu^2&\text{بالاختزال المزدوج لـ }\theta_\eta^3H_\vartheta:\text{ يقبل الاختزالين a) وc).}
\end{array}
\]
ويُثبت استقلال الملازمات الباقية بالاختزال المزدوج لسلسلة الملازمات $\A_\nu^2(a\A u)^2\equiv\theta_\eta^2$.
\noindent\emph{تنفيذ الاختزالات:}

إن وجود الاختزالات a) وb) وc) وd) واضح قبليًا، لأن العلاقات الناتجة عن الاختزال المزدوج مستقلة بعضها عن بعض بحسب قانون التكوين المبين. أما في الاختزالات e) وf) وg)، فيجب أن يثبت الحساب عدم نشوء أي متطابقات.

نعطي أيضًا الصيغ الناتجة من الاختزالات a) وb) وc) وd)، متقدمين على نحو متناظر مع الحد القطري في سلسلة الملازمات $\theta\cdot H$ حتى الملازم $\theta_\eta$، ومع الإشارة أحيانًا إلى الحساب؛ لأن بعض هذه الصيغ يُستعمل في الاختزالات e) وf) وg)، وبعضها يُستعمل لاحقًا.

\noindent 1) $H_\vartheta$ و$H_\vartheta^2$ بحسب f). انطلاقًا من الفرض\srcfn{*)}{تعني العلامة $\sim$ بدلًا من علامة المساواة أن جداءات $u_x$ في ملازمات أعلى من الملازمات الناشئة عن الطيّين III وIV قد أُغفلت.}
\[
\begin{aligned}
a_\nu\cdot b_{\nu_1}^2
&=\nu\cdot a_\nu b_\nu^2-(aHu)(bHu)b_\eta
\sim \nu\cdot a_\nu b_\nu^2-f\cdot a_\eta(aHu)^2-(abu)(aHu)b_\eta+(abu)^2b_\eta,\\
a_\nu^2\cdot b_{\nu_1}^2
&=b_{\nu_1}^2\{a_{\nu_1}u_\nu-(a_\nu\widehat{\nu}_1u)\}^2
\sim \nu\cdot a_\nu^2b_\nu^2-2a_\eta b_\eta(aHu)(bHu)-\frac16(asu)^2(bsu)^2\\
&\sim \nu\cdot a_\nu^2b_\nu^2-2a_\nu^2(aHu)(abu)-\frac16(asu)^2(bsu)^2+(\vartheta\eta x)^2(\theta Hu)^2.
\end{aligned}
\]
وبإدخال قيمة $\theta_\eta H_\vartheta$ تنتج الصيغ:
\[
\tag{12}
\begin{aligned}
\text{a)}\quad H_\vartheta&\sim 2a_\nu\cdot a_\nu^2+2\nu\cdot\theta_\nu(\vartheta\nu x)^2+2f\cdot a_\eta(aHu)^2-2\theta_\eta,\\
\text{b)}\quad H_\vartheta^2&\sim (a_\nu^2)^2+2\theta_\eta^2+\frac23(\theta su)^2+\frac1{12}s\cdot\nu-\frac19 i f\cdot\nu-\frac16f\cdot(asu)^4.
\end{aligned}
\]

\noindent 2) $\theta_\eta H_\vartheta$ بحسب b):
\[
\begin{aligned}
a_\eta^2(aHu)(abu)
&\sim [a_\eta^2(aHu)]_{\widehat{bu}}+\frac12 f\cdot[a_\eta^2(aHu)^2]
=a_\eta b_\eta(aHu)(abu)\\
&\quad+a_\eta(\widehat{ab}\eta x)(abu)(aHu)-\theta_\eta(\vartheta\eta x)(\theta Hu)+\frac12\theta_\eta^2+\frac14(\nu\eta x)^2.
\end{aligned}
\]
ومن الصيغتين (5.) و(11.) ينتج:
\[
\theta_\eta H_\vartheta\sim-\frac32\theta_\eta^2-\frac14(\theta su)^2-\frac1{12}s\cdot\nu+\frac29 i f\cdot\nu.
\]

\noindent 3) يُحسب $\theta_\eta H_\vartheta(\theta Hu)$ بحسب b) على نحو مماثل لـ$\theta_\eta H_\vartheta$:
\[
\theta_\eta H_\vartheta(\theta Hu)\sim-\theta_\eta^2(\theta Hu)-\frac16(\theta su)^3.
\]

\noindent 4) $\theta_\eta H_\vartheta(\vartheta\eta x)$ بحسب b)؛ و$\theta_\eta^2(\vartheta\eta x)$ بحسب d) (قارن \S\ 9):
\[
\theta_\eta H_\vartheta(\vartheta\eta x)\sim-\theta_\eta^2(\vartheta\eta x)-\frac16s_\vartheta(\theta su)
\sim\frac13(\vartheta\R x)-\frac16s_\vartheta(\theta su),
\]
\[
\theta_\eta^2(\vartheta\eta x)\sim\frac23a_\eta^3(abu)\sim-\frac13(\vartheta\R x).
\]


\noindent 5)
\[
\begin{array}{ccc}
\theta_\eta H_\vartheta^2\text{ بحسب b)},&
\theta_\eta^2H_\vartheta\text{ بحسب a)},&
\theta_\eta^2H_\vartheta\text{ بحسب b)، ومنه }\theta_\eta^3\text{ بحسب c)},\\
\text{من }a_\eta^2(aHb)(bHu);&
\text{من }a_\eta^2(Hab)(abu);&
\text{من }a_\eta^3(bHu)(abu):
\end{array}
\]
\[
\theta_\eta H_\vartheta^2\sim-\frac32\theta_\eta^2H_\vartheta-\frac14s_\vartheta(\theta su)^2+\frac16s_\nu-\frac49ia_\nu
\sim-\theta_\varrho-\frac16s_\nu-\frac19ia_\nu,
\]
\[
\theta_\eta^2H_\vartheta\sim\frac23\theta_\varrho-\frac16s_\vartheta(\theta su)^2+\frac29s_\nu-\frac29ia_\nu,
\]
\[
\theta_\eta^3H_\vartheta\sim\frac23\theta_\varrho-\frac23\theta_\eta^3+\frac5{36}s_\nu\quad\text{ومنها}
\]
\[
\theta_\eta^3\sim\frac14s_\vartheta(\theta su)^2-\frac18s_\nu+\frac13ia_\nu.
\]

\noindent 6) $\theta_\eta^2(\theta Hu)^2$ بحسب e):
\[
\A_\nu^2(a\A u)^4=[(a\A u)^4]_{\nu^2}=-\frac{12}{5}\theta_\eta^2(\theta Hu)^2-\frac3{10}\sigma-\frac1{20}(\theta su)^4+\nu\cdot\varrho
\quad(\text{بحسب الصيغة (3.) c}),
\]
\[
\A_\nu^2(a\A u)^4=[\A_\nu^2]^4_{\widehat{u a}}=\frac35\theta_\eta^2(\theta Hu)^2+\frac15\sigma-\frac3{10}(\theta su)^4+\frac13ij
\quad(\text{بحسب الصيغة (10.) a}),
\]
\[
\theta_\eta^2(\theta Hu)^2=-\frac16\sigma+\frac1{12}(\theta su)^4-\frac19ij+\frac13\nu\cdot\varrho.
\]

\noindent 7) $\theta_\eta^2(\vartheta\eta x)^2$ بحسب d) وe). ومن ذلك اختزال الملازم الأعلى $g$.

بحساب مماثل لحساب \S\ 9 ينتج:
\[
\begin{aligned}
\theta_\eta^2(\vartheta\eta x)^2
&=\frac12t\cdot f-\frac12a_\eta^2b_\eta-\frac1{12}g=\frac23t\cdot f-\frac23a_\eta^3b_\eta-\frac16g\\
&=-\frac16g-\frac13(\vartheta\varrho x)^2+\frac4{15}f\cdot a_\varrho^2+\frac8{15}f\cdot t\quad(\text{بحسب الصيغة (11.)}).
\end{aligned}
\]
وبحساب مماثل لما سبق ينتج:
\[
(a\A\widehat{\nu}x)^4=[(a\A u)^4]_{\widehat{\nu}x^4}=\frac{21}{10}\theta_\eta^2(\vartheta\eta x)^2+\frac7{10}s_\vartheta^2-\frac3{10}g
\quad(\text{بحسب الصيغة (3.) c}),
\]
\[
\begin{aligned}
(a\A\widehat{\nu}x)^4&=-\frac13i\A+6[\A_\nu^2]a^2-4[\A_\nu^3]a+\A_\nu^4\cdot f\\
&=-\frac{36}{5}\theta_\eta^2(\vartheta\eta x)^2-\frac95s_\vartheta^2-\frac35g-\frac35(\vartheta\varrho x)^2\\
&\quad+\frac53i\A+\frac{37}{25}f\cdot a_\varrho^2+\frac{74}{25}f\cdot t\quad(\text{بحسب الصيغة (10.)}).
\end{aligned}
\]
ومن ثم
\[
\begin{aligned}
\frac{93}{10}\theta_\eta^2(\vartheta\eta x)^2
&=-\frac3{10}g-\frac52s_\vartheta^2-\frac35(\vartheta\varrho x)^2\\
&\quad+\frac53i\A+\frac{37}{25}f\cdot a_\varrho^2+\frac{74}{25}f\cdot t,
\end{aligned}
\]
وبمقارنة القيمتين لـ$\theta_\eta^2(\vartheta\eta x)^2$ بحسب g):
\[
\tag{13.)\srcfnmark{*)}}
0=g-2s_\vartheta^2+2(\vartheta\varrho x)^2+\frac43i\A-\frac45f\cdot a_\varrho^2-\frac85f\cdot t.
\]
\srcfntext{*)}{هذه هي العلاقة المذكورة في حاشية المقدمة بين الملازمات الثلاثة الموافقة للتغير ذات الرتبة 6 والدرجة 6.}

\noindent 8) $\theta_\eta^2H_\vartheta(\theta Hu)$ بحسب a) وb)، ومنه $\theta_\eta^3(\theta Hu)$ بحسب c) (تزول الملازمات التي تحتوي العامل $u_x$):
\[
\theta_\eta^2H_\vartheta(\theta Hu)=-\frac16s_\vartheta(\theta su)^3-\frac13(\nu\varrho x),
\]
\[
\theta_\eta^2H_\vartheta(\theta Hu)=-\frac13\theta_\eta^3(\theta Hu)-\frac13(\nu\varrho x),
\]
\[
\theta_\eta^3(\theta Hu)=\frac12s_\vartheta(\theta su)^3.
\]

\noindent 9) $\theta_\eta^2H_\vartheta(\vartheta\eta x)$ بحسب a) وb)، ومنه $\theta_\eta^3(\vartheta\eta x)$ بحسب c). و$\theta_\eta^3(\vartheta\eta x)$ بحسب d)، ومنه اختزال الملازم الأعلى $s_\vartheta^2(\theta su)$ بحسب g) (تزول الملازمات التي تحتوي العامل $u_x$):
\[
\theta_\eta^2H_\vartheta(\vartheta\eta x)=\frac13(atu),
\]
\[
\theta_\eta^2H_\vartheta(\vartheta\eta x)=\frac13(atu)+\frac13\theta_\eta^3(\vartheta\eta x)-\frac16s_\vartheta^2(\theta su),
\]
\[
\theta_\eta^3(\vartheta\eta x)=\frac12s_\vartheta^2(\theta su),
\]
\[
\theta_\eta^3(\vartheta\eta x)=\frac25\theta_\varrho(\vartheta\varrho x)+\frac15(atu):
\]
\srcnumdisplay{(14.)}{%
s_\vartheta^2(\theta su)=\frac45\theta_\varrho(\vartheta\varrho x)+\frac25(atu).
}

\noindent 10) $\theta_\eta^2H_\vartheta^2$ بحسب a) وباستخدام المختزل $\theta_\nu^2(\vartheta\nu x)^2$، و$\theta_\eta^3H_\vartheta$ بحسب a) وc)، و$\theta_\eta^4$ بحسب c)؛ ومن ذلك اختزال الملازمين الأعلى $s_\vartheta^2(\theta su)^2$ و$s_\nu^2$ بحسب g):
\[
\begin{aligned}
\text{بحسب a)}\quad
\theta_\eta^2H_\vartheta^2&=[a_\eta^2(aHu)^2]b^2+\frac13[a_\eta^4]_{bu^2}-\frac13H_\nu^2(\nu\eta x)^2\\
&=\frac49i\cdot a_\nu^2-\frac16s_\vartheta^2(\theta su)^2-\frac5{18}s_\nu^2+\frac13(atu)^2+\frac1{54}J\cdot u_x^2.
\end{aligned}
\]
\[
\begin{aligned}
\text{باستخدام }\theta_\nu^2(\vartheta\nu x)^2\quad
\theta_\eta^2H_\vartheta^2
&=[\theta_\nu^2(\vartheta\nu x)^2]_{\nu_1}+\frac13[\theta_\nu^4]_{\nu_1}\widehat{\nu}_1x^2-\frac13s_\vartheta^2(\theta su)^2\\
&=\frac49i\cdot a_\nu^2-\frac16s_\nu^2-\frac13s_\vartheta^2(\theta su)^2+\frac13(\nu\varrho x)^2.
\end{aligned}
\]
\[
\begin{aligned}
\text{بحسب a)}\quad
\theta_\eta^3H_\vartheta&=-[a_\eta^2(aHu)]_{bu}b^2-\frac12[a_\eta^2(aHu)^2]b^2-\frac13[a_\eta^3]b+\frac1{12}[a_\eta^4]_{bu^2}\\
&\quad+\frac12u_x[a_\eta^2(aHu)^2]b^3\\
&=-\frac29ia_\nu^2+\frac1{12}s_\nu^2+\frac4{15}\theta_\varrho^2-\frac7{90}(\nu\varrho x)^2+\frac1{30}(atu)^2+\frac2{27}i^2\cdot u_x^2-\frac1{36}J\cdot u_x^2.
\end{aligned}
\]
\[
\begin{aligned}
\text{بحسب c)}\quad
\theta_\eta^3H_\vartheta:\quad a_\eta^3(Hab)(bHu)&=\frac12(atu)^2
=\theta_\eta^2H_\vartheta(\theta Hu)(\vartheta\eta x)+\frac14H_\nu^2(\nu\eta x)^2\\
&\quad+\frac12\theta_\eta^2H_\vartheta^2
=\frac32\theta_\eta^2H_\vartheta^2+\theta_\eta^3H_\vartheta-u_x[a_\eta^2(aHu)^2]b^3\\
&\quad+\frac14H_\nu^2(\nu\eta x)^2,
\end{aligned}
\]
ومن ثم
\[
\begin{aligned}
\theta_\eta^3H_\vartheta&=-\frac23i\cdot a_\nu^2+\frac5{12}s_\nu^2+\frac12(\nu\varrho x)^2-\frac12(atu)^2\\
&\quad+\frac4{27}i^2\cdot u_x^2-\frac1{12}J\cdot u_x^2.
\end{aligned}
\]
بحسب c)
\[
\theta_\eta^4=\frac49ia_\nu^2-\frac16s_\nu^2+\frac{16}{15}\theta_\varrho^2-\frac{14}{45}(\nu\varrho x)^2+\frac2{15}(atu)^2.
\]
بحسب g)، من القيمتين لـ$\theta_\eta^2H_\vartheta^2$:
\par\smallskip\noindent
\begin{minipage}[c]{3.6em}
\centering (15.)
\end{minipage}%
\begin{minipage}[c]{\dimexpr\linewidth-3.6em\relax}
\centering
\(\displaystyle
\begin{aligned}
s_\vartheta^2(\theta su)^2
&=\frac23s_\nu^2-2(atu)^2+2(\nu\varrho x)^2-\frac19J\cdot u_x^2\\
&=\frac89i\cdot a_\nu^2+\frac8{15}\theta_\varrho^2-\frac{14}{15}(atu)^2+\frac{38}{45}(\nu\varrho x)^2-\frac4{27}i^2\cdot u_x^2.
\end{aligned}
\)
\end{minipage}
\par\smallskip
بحسب g)، من القيمتين لـ$\theta_\eta^3H_\vartheta$:
\par\smallskip\noindent
\begin{minipage}[c]{3.6em}
\centering (16.)
\end{minipage}%
\begin{minipage}[c]{\dimexpr\linewidth-3.6em\relax}
\centering
\(\displaystyle
s_\nu^2=\frac43i\cdot a_\nu^2+\frac45\theta_\varrho^2+\frac85(atu)^2-\frac{26}{15}(\nu\varrho x)^2+\frac16J\cdot u_x^2-\frac29i^2\cdot u_x^2.
\)
\end{minipage}
\par\smallskip
تعطي الصيغة (16.) الأساس لاختزال المقياس $(ss'u)^4$، والصيغة (13.) لاختزال منظومة $s$. (\S\ 17.)

\noindent\emph{النتائج:}

1) تُعَدّ $a_\nu(j\nu x)^3$ و$\theta_\eta H_\vartheta$ و$\theta_\eta^2(\vartheta\eta x)$ و$\theta_\eta^2(\theta Hu)^2$ مختزلات، وتُعَدّ $a_\nu(j\nu x)^2$ و$H_\vartheta$ و$H_\vartheta^2$ ملازمات قابلة للتحلل.

2) ملازم المنظومة II: $j(\nu)=g=(\nu\eta x)^4H_x^2$ قابل للاختزال.

3) لا يدخل $\nu$ إطلاقًا في قوى ولا في جداءات مع ملازمات المنظومة نفسها عند الطيّ مع المنظومة I، لأن الملازم المخالف للتغير الوحيد $g$ يصبح قابلًا للاختزال.

لا يدخل $\theta$ في الطيّ في الجداءات أو القوى إلا عند اعتباره ملازمًا مخالفًا للتغير، وكذلك $H$ عند اعتباره ملازمًا موافقًا للتغير (قارن \S\ 13 B).
\begin{center}
\textbf{\S\ 12.}\quad\emph{الملازمات من الرتبة السابعة.}
\end{center}
\[
\text{طيّ }\left\{\begin{array}{l}
\theta\cdot K\text{ مع }\nu,\\
K,j,\A\text{ مع }H,\\
f\text{ مع }L,\sigma.
\end{array}\right.
\]


\noindent\emph{الملازمات غير القابلة للاختزال:}
\[
\renewcommand{\arraystretch}{1.35}
\begin{array}{@{}r@{\quad}c|c|c|c@{}}
&\cdot&\cdot&(KHu)^2&\cdot\\
&\cdot&K_\eta&\cdot&K_\eta(KHu)^2\\
\text{B)}&(k\eta x)^2&\cdot&K_\eta^2&\cdot\\
&(k\eta x)^3&H_k(k\eta x)^2,\ K_\eta(k\eta x)^2&\cdot&\cdot\\
&\cdot&K_\eta(k\eta x)^3&\cdot&\cdot
\end{array}\, ,
\]
\[
\renewcommand{\arraystretch}{1.35}
\begin{array}{@{}r@{\quad}c@{\qquad}c@{\qquad}c@{\qquad\qquad}r@{\quad}c@{}}
\text{C)}&(j\eta x)&H_j&\cdot&\text{D)}&(\A Hu)\\
&\cdot&H_j(j\eta x)&H_j^2,&&\A_\eta,
\end{array}
\]
\[
\renewcommand{\arraystretch}{1.35}
\begin{array}{@{}r@{\quad}c@{\qquad}c@{\qquad}c@{\qquad}c@{}}
&\cdot&(aLu)^2&(aLu)^3&\cdot\\
\text{E)}&a_l&\cdot&a_l(aLu)^2&a_l(aLu)^3\\
&\cdot&a_l^2&\cdot&\cdot
\end{array}
\]
\noindent\emph{الاختزالات:}

\noindent A) يتحلل $(\vartheta\cdot k,\nu,x)^4$ لأنه ناتج من إجراء عملية الاشتقاق والتقليص مرة واحدة على الملازم القابل للتحلل $(\vartheta^2\nu x)^4$.

\noindent B) سلسلة الملازمات $H_kK_\eta$: المختزل $H_\vartheta\theta_\eta$.

سلسلة الملازمات $K_\eta^2(KHu)$: المختزل $a_\eta^2(aHu)$.

سلسلة الملازمات $K_\eta^2(k\eta x)$: المختزل $\theta_\eta^2(\vartheta\eta x)$.

يؤدي الاختزال المزدوج الناتج من ذلك لسلسلة الملازمات $K_\eta^3$ إلى اختزال سلسلة الملازمات الأعلى $(j\eta x)^3$.

الملازمات $(KHu)$ و$(k\eta x)$ و$K_\eta(k\eta x)$ قابلة للتحلل بحسب \S\ 3.

يتحلل $H_k(KHu)$ و$K_\eta(KHu)$ لأنهما ناتجان من طيّ واحد مع الملازم القابل للتحلل $K_\nu(k\nu x)$ ومع المحدد اليعقوبي
\[
K_\nu^2\equiv\theta_\nu^2(a\theta u)\equiv a_\nu^2(a\theta u)\modu{\nu^2}.
\]
يقابل التمثيل المزدوج لـ$K_\nu^2$ علاقةٌ بين ملازمات قابلة للتحلل. فنحصل على:
\[
H_k(KHu)=2K_{\nu_1}(\nu\nu_1k)=2\theta_{\nu_1}(\nu\nu_1\widehat{a\theta})
=2\theta_\nu^2\cdot a_\nu-a_\nu^2\cdot\theta_\nu+f\cdot\theta_\eta(\theta Hu)^2-f\cdot \nu\cdot\theta_\nu^3\quad(\text{مبرهنة الجداء})\modu{(\nu^3)},
\]
\[
\begin{aligned}
K_\eta(KHu)&=-K_\nu^2(\nu\nu_1x)=-a_\nu^2(a\theta u)(\nu\nu_1x)=-\theta_\nu^2(a\theta u)(\nu\nu_1x)\\
&=a_\eta(aHu)^2\cdot\theta+a_\nu^2\cdot\theta_\nu
=-\theta_\eta(\theta Hu)^2\cdot f-\theta_\nu^2\cdot a_\nu.
\end{aligned}
\]
\par\smallskip\noindent
\begin{minipage}[c]{4.8em}
\centering (17.)\srcfnmark{*)}
\end{minipage}%
\begin{minipage}[c]{\dimexpr\linewidth-4.8em\relax}
\centering
\(\displaystyle
0=a_\nu^2\cdot\theta_{\nu_1}+\theta_\nu^2\cdot a_{\nu_1}+f\cdot\theta_\eta(\theta Hu)^2+\theta\cdot a_\eta(aHu)^2\modu{(\nu^3)}.
\)
\end{minipage}
\par\smallskip
\srcfntext{*)}{حُذفت في جميع العلاقات بين الملازمات القابلة للتحلل الجداءاتُ التي يكون أحد عواملها ملازمًا مخالفًا للتغير، مثل $f\cdot\theta_\nu^3\cdot \nu$ في هذه الحالة؛ إذ لا شأن لها عند تطبيق الصيغ على المنظومة $s_{IV}$.}
تتحلل الملازمات $H_k$ و$H_k(k\eta x)$ و$H_k^2$ و$H_k^2(k\eta x)$ لأنها ناتجة من طيّ واحد مع الملازمين القابلين للتحلل $H_\vartheta$ و$H_\vartheta^2$ (قارن \S\ 3. 2)، a) وb)).

لدينا
\[
H_k=H_\vartheta(a\theta u)\sim[H_\vartheta]_{\widehat{ua}}-f\cdot H_\vartheta(\theta Hu)\qquad(\text{تخصيص }y=\widehat{uv}\text{ من 2) a}).
\]
\[
H_k(k\eta x)=H_\vartheta a_\eta-H_\vartheta\theta_\eta\cdot f,
\qquad
H_\vartheta a_\eta=\frac54[H_\vartheta]a-\frac14H_\vartheta a_\vartheta
\quad(\S\ 3. 2)b),
\]
\[
H_\vartheta^2(a\theta u)=[H_\vartheta^2]_{\widehat{ua}},
\qquad
H_\vartheta^2a_\eta=[H_\vartheta^2]a=H_k^2(k\eta x)+H_\vartheta^2\theta_\eta\cdot f.
\]

\noindent C) سلسلة الملازمات $(j\eta x)^3$: قابلة للاختزال بالاختزال المزدوج لسلسلة الملازمات $K_\eta^3$ بحسب B)، انطلاقًا من التعبير:
\[
\begin{aligned}
0&=a_\eta^3(a\theta u)=\theta_\eta^3(a\theta u)
=\{\theta_\eta+(\widehat{a\theta}\eta x)\}^3(a\theta u)-\theta_\eta^3(a\theta u)\\
&=\frac12(j\eta x)^3\modu{(\nu^3,s,\R,t,i,J)}.\qquad(\text{قارن خاتمة \S\ 5.})
\end{aligned}
\]
يصح التعبير المناظر حتى الملازم النهائي في سلسلة الملازمات:
\[
0=H_\vartheta^2(\widehat{a\theta}\eta x)a_\eta^3
=H_\vartheta^2(\widehat{a\theta}\eta x)\theta_\eta^3=\frac12H_j^2(j\eta x)^4\modu{(\nu^3,s,\R,t,i,J)}.
\]
الملازم $(j\eta x)^2$: يتحلل 1) باستقطاب العلاقة الخاصة بالملازم القابل للتحلل $H_\vartheta^2$ مرتين ((12.) b)؛

2) بالحساب المباشر للجداء $a_\nu\cdot\theta_\nu^3$؛ ومن ذلك ينتج تحلل الملازم الأعلى $(\A Hu)^2$. فنحصل على:
\par\smallskip\noindent
\begin{minipage}[c]{4.8em}
\centering (18.)
\end{minipage}%
\begin{minipage}[c]{\dimexpr\linewidth-4.8em\relax}
\centering
\(\displaystyle
\left.\begin{array}{ll}
\text{a)}&\dfrac13(\A Hu)^2+\dfrac13(j\eta x)^2=\theta_\nu^2\cdot a_{\nu_1}^2,\\[0.6em]
\text{b)}&(j\eta x)^2=\dfrac43\theta_\nu^3\cdot a_{\nu_1},
\end{array}\right\}\modu{(\nu^3,s,\R,t,i,J)}
\)
\end{minipage}
\par\smallskip
انطلاقًا من التعبير:
\[
\begin{aligned}
H_\vartheta^2(a\theta u)^2=\frac13(\A Hu)^2
&=2\theta_\eta^2(a\theta u)(aHu)+\theta_\nu^2\cdot a_{\nu_1}^2\quad(\text{مبرهنة الجداء})\\
&=(a\theta u)(aHu)\{a_\eta-(\widehat{a\theta}\eta x)\}^2+\theta_\nu^2\cdot a_{\nu_1}^2,
\end{aligned}
\]
\[
\begin{aligned}
a_{\nu_1}\cdot\theta_\nu^3
&=-(a\widehat{\nu\nu_1}u)\theta_\nu^2(\theta\widehat{\nu\nu_1}u)-\frac12(a\widehat{\nu\nu_1}u)\theta_\nu\theta_{\nu_1}(\theta\widehat{\nu\nu_1}u)\\
&=-\frac32\theta_\eta^2(\theta Hu)(aHu)=-\frac32\theta_\eta^2(\theta Hu)(a\theta u)\\
&=-\frac32\{a_\eta-(\widehat{a\theta}\eta x)\}^2\{(aHu)-(a\theta u)\}(a\theta u).
\end{aligned}
\]

\noindent D) سلسلة الملازمات $\A_\eta(\A Hu)$: المختزل $\A_\nu^2$.

الملازم $(\A Hu)^2$ قابل للتحلل بحسب C).

\noindent E) سلسلة الملازمات $a_l^2(aLu)$: المختزل $a_\eta^2(aHu)$.
الملازمان $(aLu)$ و$a_l(aLu)$: يتحللان لأنهما ناتجان من عملية الاشتقاق والتقليص على محددات يعقوبية بحسب \S\ 3.

\noindent F) سلسلة الملازمات $a_\sigma^3$: المختزل $a_\eta^3$.

الملازمان $a_\sigma^2$ و$a_\sigma^3$: المختزلان $a_\nu^3$ و$a_\eta^3$، بالاختزال المزدوج للتعابير المناظرة
\[
H_\nu a_\nu^3\quad\text{و}\quad H_\nu a_\eta^3,
\qquad
H_\nu a_\nu^3a_\eta\quad\text{و}\quad H_\nu a_\eta^4
\]
(قارن \S\ 10. C)، اختزال $\A_\nu^2$ و$\A_\eta^3$).

ومن ذلك، مع مراعاة (16.)، نحصل على اختزال الملازمين الأعلى $a_\nu s_\nu(asu)^2$ و$a_\nu^2s_\nu(asu)^2$ والملازمات المكتوبة بالرموز $\R,t(a_\nu,a_\eta^2)$:
\par\smallskip\noindent
\begin{minipage}[c]{4.8em}
\centering (19.)
\end{minipage}%
\begin{minipage}[c]{\dimexpr\linewidth-4.8em\relax}
\centering
\(\displaystyle
\left.\begin{array}{rcl}
a_\nu s_\nu(asu)^2&=&\dfrac13 i\theta_\nu^2-\dfrac1{18}iH,\\[0.5em]
a_\nu^2s_\nu(asu)^2&=&0,
\end{array}\right\}\modu{(\R,t)}
\)
\end{minipage}
\par\smallskip
الملازم $a_\sigma$: يتحلل باستقطاب (12.)b، أو على نحو أقصر بالتوسيع المباشر للجداء $a_\nu^2\cdot\theta_\nu^3$، انطلاقًا من التعبير:
\[
a_\nu^2\cdot\theta_{\nu_1}^3=a_\nu^2\theta_{\nu_1}\{a_{\nu_1}-(\widehat{a\theta}\nu_1x)\}^2
=-2a_\eta(aHu)(a\theta H)(\widehat{a\theta}\eta x)+(\widehat{a\theta}\nu_1x)^2a_\nu^2\theta_{\nu_1}
\]
ومع مراعاة اختزال $H_j(j\eta x)^2$ (C):
\par\smallskip\noindent
\begin{minipage}[c]{4.8em}
\centering (20.)
\end{minipage}%
\begin{minipage}[c]{\dimexpr\linewidth-4.8em\relax}
\centering
\(\displaystyle a_\sigma=2a_\nu^2\cdot\theta_{\nu_1}^3\modu{(s,\R,t,i)}.\)
\end{minipage}
\par\smallskip

\noindent\emph{النتائج:}

1) الملازمات $(j\eta x)^3$ و$\A_\eta(\A Hu)$ و$a_\sigma^2$ مختزلات، والملازمات $(j\eta x)^2$ و$(\A Hu)^2$ و$a_\sigma$ قابلة للتحلل.

2) لا يدخل $f$ في الطيّ مع $L$ إلا في جداءات مع ملازمات متخالفة التدرج؛ ولا يدخل $f$، ومن ثم $K$ و$N$ أيضًا، مطلقًا في الطيّ مع $\sigma$، وبالتالي مع $(\nu\sigma x)$ أيضًا. أما $j$ فلا يدخل إلا في جداءات مع ملازمات متخالفة التدرج، وأما $\A$ فلا يدخل في جداءات مطلقًا، وذلك في الطيّ مع $H$ و$L$ (وهذا ينتج من $\A_\eta(\A Hu)$ و$(\A Hu)^2$). وكذلك لا يدخل $\sigma$، وبالتالي $(\nu\sigma x)$، في الطيّ في جداءات مع أي ملازم من المنظومة I. ومع مراعاة نتائج \S\ 11، لا يبقى من الجداءات في المنظومة II إلا قوى $H$ وجداءات $H$ مع $L$.

\begin{center}
\textbf{\S\ 13.}\quad\emph{الملازمات من الرتبة الثامنة (المنظومة III).}
\end{center}
\[
\text{طيّ }\left\{\begin{array}{l}
\A\cdot j\text{ مع }\nu,\\
\theta^2,
 f\cdot j,
 N\text{ مع }H,\\
\theta\text{ مع }L,\sigma.
\end{array}\right.
\]


\noindent\emph{الملازمات غير القابلة للاختزال:}
\[
\renewcommand{\arraystretch}{1.28}
\begin{array}{@{}r@{\quad}c@{\qquad\qquad}r@{\quad}c@{\qquad}c@{}}
\text{A)}&\A_\nu(j\nu x),&\text{B)}&(\vartheta^2\eta x)^3&\cdot\\
&&&(\vartheta^2\eta x)^4&H_\vartheta(\vartheta^2\eta x)^3,\ \theta_\eta(\vartheta^2\eta x)^3,
\end{array}
\]
\[
\renewcommand{\arraystretch}{1.28}
\begin{array}{@{}r@{\quad}c@{\qquad}c@{\qquad\qquad}r@{\quad}c@{}}
\text{C)}&(aHu)H_j&\cdot&\text{D)}&H_\eta;\ N_\eta,\\
&\cdot&a_\eta(aHu)H_j&&
\end{array}
\]
\[
\renewcommand{\arraystretch}{1.30}
\begin{array}{@{}r@{\quad}c@{\qquad}c@{\qquad}c@{\qquad}c@{\qquad}c@{}}
&\cdot&\cdot&(\theta Lu)^2&(\theta Lu)^3&\cdot\\
&\cdot&\theta_l&\cdot&L_\vartheta(\theta Lu)^2,\ \theta_l(\theta Lu)^2&\theta_l(\theta Lu)^3\\
\text{E)}&(\vartheta lx)^2&\cdot&\theta_l^2&\cdot&\cdot\\
&\cdot&\theta_l(\vartheta lx)^2&\cdot&\cdot&\cdot
\end{array}
\]
\noindent\emph{الاختزالات:}

\noindent A) $\A_\nu(j\nu x)^3$: المختزل $a_\nu(j\nu x)^3$.

الملازم $\A_\nu(j\nu x)^2$ قابل للاختزال بواسطة الملازم القابل للتحلل $a_\nu(j\nu x)^2$ بحسب \S\ 3. 2)b، مع مراعاة المختزل $\theta_{\vartheta j}$. إذ نحصل بحسب \S\ 11 B) على:
\[
\frac3{10}\A_\nu(j\nu x)^2+\frac35a_\nu(j\nu\vartheta)a_\vartheta(j\nu x)
+\frac1{10}a_\nu(j\nu\vartheta)^2
=\frac35\theta_\nu^3\cdot\theta_{\vartheta\nu}^{\prime2}+\frac25\theta_\nu^3\theta_{\vartheta\nu}\theta_{\vartheta\nu}'.
\]
(المختزلان $a_j$ و$\theta_{\vartheta j}$).

\noindent B) تتحلل الملازمات $H_\vartheta(\vartheta'\eta x)^2$ و$H_\vartheta^2(\vartheta'\eta x)$ و$H_\vartheta^2(\vartheta'\eta x)^2$ لأنها ناتجة من طيّ مرة واحدة في كل خطوة متعاقبة، مع الملازمين القابلين للتحلل $H_\vartheta$ و$H_\vartheta^2$، أو $H_\vartheta a_\eta$ و$H_\vartheta^2a_\eta$ على الترتيب، بحسب \S\ 3. 2)b) وبحسب العلاقة:
\[
H_\vartheta(\vartheta'\eta x)^2=2H_\vartheta a_\eta^2\cdot f-2H_\vartheta a_\eta b_\eta.
\]

\noindent C) سلسلة الملازمات $a_\eta(j\eta x)^2$: المختزلان $a_j$ و$(j\eta x)^3$.

يتحلل الملازم $a_\eta(j\eta x)$: المختزل $a_j$ وطيّ واحد مع الملازم القابل للتحلل $(j\eta x)^2$ بحسب \S\ 3. 2)a (التخصيص $y=uv$).

يتحلل الملازم $a_\eta H_j$: 1) بطيّ واحد مع الملازم القابل للتحلل $a_\nu(j\nu x)^2$؛

2) بطيّ مزدوج خاص مع الملازم القابل للتحلل $(j\eta x)^2$؛ ومن ذلك تنتج علاقة بين ملازمات قابلة للتحلل.

من $a_\nu(j\nu x)^2$:
\[
0=\theta_\nu^3\cdot\theta_\nu+2a_\eta H_j\modu{(s,\R,t,i,J)}.
\]
من $(j\eta x)^2$:
\[
0=\theta_\nu^3\cdot\theta_\nu-\frac32a_\eta H_j\modu{(s,\R,t,i,J)}
\]
(حُذفت الجداءات التي تتضمن ملازمات مخالفة للتغير)، انطلاقًا من التعبير:
\[
0=\frac12a_\nu(abu)^2\theta_\nu^3
+\frac12a_\nu(abu)\theta_{\nu_1}^3(\theta bu)
-\frac34(\widehat{j\eta}bu)^2+\frac34H_j(\widehat{j\eta}bu)
-\frac18f\cdot H_j^2
\]
وبتبادل الرموز في $a_\nu(abu)(\theta bu)\theta_{\nu_1}^3$.

\noindent D) سلسلة الملازمات $N_\eta(NHu)$: المختزل $\A_\nu^2$.

الملازمات $(NHu)$ و$(n\eta x)$ و$N_\eta(n\eta x)$ و$H_n(NHu)$ قابلة للتحلل (قارن \S\ 12 B)، ويتحلل $(NHu)^2$ بطيّ واحد مع الملازم $(\A Hu)^2$.

\noindent E) سلاسل الملازمات $\theta_lL_\vartheta$, $\theta_l^2(\theta Lu)^2$, $\theta_l^2(\vartheta lx)$:

المختزلات: $\theta_\eta H_\vartheta$, $\theta_\eta^2(\theta Hu)^2$, $\theta_\eta^2(\vartheta\eta x)$.

تتحلل الملازمات $(\theta Lu)$ و$(\vartheta lx)$ و$L_\vartheta$ و$L_\vartheta(\theta Lu)$ و$\theta_l(\theta Lu)$ و$L_\vartheta(\vartheta lx)$ و$\theta_l(\vartheta lx)$ و$L_\vartheta^2$ و$L_\vartheta^2(\theta Lu)$ و$\theta_l^2(\theta Lu)$. (وهي تقابل الملازمات القابلة للتحلل في \S\ 12 B) بالثنوية.)

\noindent F) سلسلة الملازمات $\theta_\sigma(\vartheta\sigma x)$: المختزل $a_\sigma^2$.

يتحلل الملازمان $(\vartheta\sigma x)$ و$(\vartheta\sigma x)^2$ لأنهما ناتجان من طيّ واحد مع الملازم القابل للتحلل $a_\sigma$؛ ويتحلل $\theta_\sigma$، الناتج من طيّ مزدوج، لأن سلسلة الملازمات تحقق
\[
a_\nu^2(a\theta u)\theta_{\nu_1}^3\equiv H_j^2(j\eta x)\equiv0\;\text{بترديد }(s,\R,t,i,J):
\]
\[
\theta_\sigma=a_\sigma(abu)^2=a_\nu^2(abu)^2\cdot\theta_\nu^3.
\]

\noindent\emph{النتائج:}

لا يدخل $\theta^3$ أو $\theta^2K$ في الطيّ مع $H$؛ ولا يدخل $\theta$ في الطيّ مع $L$ إلا بوصفه ملازمًا مخالفًا للتغير، وفي قوى أو جداءات، ولا يدخل مطلقًا في الطيّ مع $\sigma$ و$(\nu\sigma x)$.

لا يدخل $N$ في جداءات في الطيّ مع أي ملازم من المنظومة II، ولا مع قوى أو جداءات ملازمات من المنظومة II أيضًا. ومع مراعاة نتائج الفقرات الأخيرة، لا يبقى من الجداءات في المنظومة I إلا الجداءات $f\cdot j$ و$\A\cdot j$، وقوى $\theta$، وجَداءات $\theta\cdot K$.

\begin{center}
\textbf{\S\ 14.}\quad\emph{الملازمات من الرتبة التاسعة (المنظومة III).}
\end{center}
\[
\text{طيّ }\left\{\begin{array}{l}
\theta\cdot K\text{ مع }H,\\
K,j,\A\text{ مع }L,\\
j,\A\text{ مع }\sigma,\\
f\text{ مع }H^2.
\end{array}\right.
\]


\noindent\emph{الملازمات غير القابلة للاختزال:}
\[
\renewcommand{\arraystretch}{1.35}
\begin{array}{@{}r@{\quad}c@{\qquad}c@{\qquad\qquad}r@{\quad}c@{\qquad}c@{\qquad}c@{\qquad}c@{}}
&\cdot&\cdot&&\cdot&\cdot&(KLu)^3&\cdot\\
\text{A)}&(\vartheta\cdot k,\eta,x)^4&\cdot&\text{B)}&\cdot&\cdot&\cdot&K_l(KLu)^3\\
&\cdot&H_k(\vartheta\cdot k,\eta,x)^4&&\cdot&K_l^2&\cdot&\cdot\\
&&&& (klx)^3&\cdot&\cdot&\cdot\\
&&&&\cdot&K_l(Klx)^3&\cdot&\cdot
\end{array}
\]
\[
\renewcommand{\arraystretch}{1.35}
\begin{array}{@{}r@{\quad}c@{\qquad}c@{\qquad}r@{\quad}c@{\qquad}r@{\quad}c@{\qquad}r@{\quad}c@{\qquad}c@{}}
\text{C)}&L_j&\cdot&\text{D)}&\A_l&\text{E)}&(j\sigma x),&\text{G)}&(aH^2u)^3&(aH^2u)^4\\
&\cdot&L_j^2&&&&&&\cdot&a_\eta(aH^2u)^3
\end{array}
\]
\noindent\emph{الاختزالات:}

A) تتحلل الملازمات $H_\vartheta(k\eta x)^3$ و$H_\vartheta^2(k\eta x)^2$ و$H_\vartheta^2(k\eta x)^3$ لأنها ناتجة من طيّ مرة واحدة في كل خطوة متعاقبة، مع ملازمات قابلة للتحلل.

B) سلاسل الملازمات $K_lL_k$, $K_l^2(KLu)$, $K_l^2(klx)$:

المختزلات: $\theta_\eta H_\vartheta$, $a_\eta^2(aHu)$, $\theta_\eta^2(\vartheta\eta x)$.

تتحلل الملازمات $K_l$ و$K_l(KLu)$ و$K_l(klx)$ و$(KLu)^2$ و$(klx)^2$ لأنها نواتج عمليات الاشتقاق والتقليص على محددات يعقوبية (\S\ 3).

الملازمات $L_k$, $L_k(KLu)$, $L_k(KLu)^2$, $L_k(klx)$, $L_k(klx)^2$, $L_k^2$, $L_k^2(KLu)$, $L_k^2(klx)$, $L_k^3$, $K_l(KLu)^2$, $K_l(klx)^2$ تتحلل لأنها ناتجة من طيّ واحد مع الملازمات القابلة للتحلل:
\[
L_\vartheta,
\ H_k(KHu),
\ L_\vartheta(\vartheta lx),
\ H_k^2,
\ L_\vartheta^2,
\ L_\vartheta^2(\theta Lu),
\ K_\eta(KHu),
\ \theta_l(\vartheta lx)
\]
بحسب \S\ 3. 2), a) وb).

C) سلسلة الملازمات $(jlx)^3$: المختزل $(j\eta x)^3$.

الملازمان $(jlx)^2$ و$L_j(jlx)$: ناتجان من طيّ واحد مع الملازم القابل للتحلل $(j\eta x)^2$.

D) سلسلة الملازمات $\A_l(\A Lu)$: المختزل $\A_\nu^2$.

الملازمان $(\A Lu)^2$ و$(\A Lu)^3$: طيّ واحد مع الملازم القابل للتحلل $(\A Hu)^2$.

E) سلسلة الملازمات $(j\sigma x)^2$: المختزل $a_\sigma^2$ بإحلال ملازمات أدنى من السلسلة محل $j$ (\S\ 5):
\[
a_\sigma^2(a\theta u)^2=\frac13(j\sigma x)^2,
\qquad
 a_\sigma^2(a\theta\sigma x)^2=\frac13(j\sigma x)^4.
\]
F) سلسلة الملازمات $\A_\sigma^2$: المختزل $a_\sigma^2$.

الملازم $\A_\sigma$: المختزل $a_\sigma^2$ بإحلال ملازمات أدنى من السلسلة محل $\A$:
\[
a_\vartheta a_\sigma^2=\frac23\A_\sigma\modu{\nu}.
\]

\noindent\emph{النتائج:} الملازم $j$ وحده يدخل في الطيّ مع $\sigma$ و$(\nu\sigma x)$.

لا يدخل $N$ في الطيّ مع $L$، لأن الملازم المناظر لـ$\A_l$، وهو $N_l$، يتحلل.

\begin{center}
\textbf{\S\ 15.}\quad\emph{الملازمات من الرتبة العاشرة (المنظومة III).}
\end{center}
\[
\text{طيّ }\left\{\begin{array}{l}
\theta^2,
\ f\cdot j\text{ مع }L,\\
\theta\text{ مع }H^2.
\end{array}\right.
\]
\noindent\emph{الملازمات غير القابلة للاختزال:}
\[
\begin{array}{c@{\quad}l@{\qquad}c@{\quad}l}
\text{A)}&(\vartheta^2lx)^4,
\ \theta_l(\vartheta^2lx)^4&
\text{B)}&(aLu)^2L_j,
\ a_l(aLu)^2L_j,\\[0.8em]
\text{C)}&(\theta H^2u)^3,
\ (\theta H^2u)^4,
\ H_\vartheta(\theta H^2u)^3,
\ \theta_\eta(\theta H^2u)^3.
\end{array}
\]

\noindent\emph{الاختزالات:}

A) وB): تحلل الملازمات الأخرى بطيّ واحد مع ملازمات قابلة للتحلل.

C) تحلل الملازمات الأخرى بحسب \S\ 13 B) بالنظر من الجهة المقابلة وفق الثنوية.

\noindent\emph{النتائج:}

لا يدخل $\theta^2\cdot K$ في الطيّ مع $L$؛ وبالمثل لا يدخل $H^2L$ في الطيّ مع $K$. ولا يدخل $H^3$ و$H^2L$ في الطيّ مع $\theta$. وبذلك يثبت مخطط الطيّ في \S\ 7.

\begin{center}
\textbf{\S\ 16.}\quad\emph{الملازمات من الرتب 11 و12 و13 و15 (المنظومة III).}
\end{center}

\noindent\emph{الرتبة 11.}
\[
\text{طيّ }\left\{\begin{array}{l}
\theta\cdot K\text{ مع }L,\\
K,j\text{ مع }H^2,\\
j\text{ مع }(\nu\sigma x),\\
f\text{ مع }H\cdot L.
\end{array}\right.
\]
\noindent\emph{الملازمات غير القابلة للاختزال:}
\[
\begin{array}{c@{\quad}l@{\qquad}c@{\quad}l}
\text{A)}&(\vartheta\cdot k,l,x)^5,
\ \theta_l(\vartheta\cdot k,l,x)^5&
\text{B)}&(KH^2u)^4,
\ K_\eta(KH^2u)^4,\\[0.8em]
\text{C)}&(\nu\sigma j),&
\text{D)}&(a,H\cdot L,u)^4.
\end{array}
\]

\noindent\emph{الرتبة 12.}
\[
\text{طيّ }\left\{\begin{array}{l}
\theta^3\text{ مع }L,\\
\theta\text{ مع }H\cdot L.
\end{array}\right.
\]
\noindent\emph{الملازمات غير القابلة للاختزال:}
\[
\text{E)}\ (\vartheta^3lx)^6,
\qquad
\text{F)}\ (\theta,H\cdot L,u)^4,
\qquad L_\vartheta(\theta,H\cdot L,u^4).
\]

\noindent\emph{الرتبة 13.}
\[
\text{طيّ }K\text{ مع }H\cdot L.
\]
\noindent\emph{الملازمات غير القابلة للاختزال:}
\[
\text{G)}\ (K,H\cdot L,u)^5,
\qquad K_\eta(K,H\cdot L,u)^5.
\]

\noindent\emph{الرتبة 15.}
\[
\text{طيّ }K\text{ مع }H_3.
\]
\noindent\emph{الملازم غير القابل للاختزال:}
\[
\text{H)}\ (KH^3u)^6.
\]

\noindent\emph{الاختزالات:}

الملازمان $H_{j^3}^2$, $H_{j^4}^2$. المختزل $L_j^3$ من العلاقة:
\[
(\nu lx)L_x^3=-\frac12H^2\modu{(\sigma,s)}.
\]
تحلل الملازمات الأخرى أو قابليتها للاختزال بحسب اعتبارات الفقرات السابقة.

\noindent\emph{النتائج:}

بحسب مخطط الطيّ في \S\ 7 ونتائج \S\S\ 8 إلى 15، تكون الملازمات غير القابلة للاختزال للمنظومة III (117 ملازما) قد استوفيت بذلك.



\clearpage

\begin{center}
{\Large\bfseries الفصل IV.\quad المنظومة الكاملة نسبيا بالنسبة إلى المقياس $(\varrho,t)$.}
\end{center}

\begin{center}
\textbf{\S\ 17.}\quad\emph{اختزال المقياس $(ss'u)^4$ ومنظومة $s$.}
\end{center}

لتكوين المنظومة الكاملة نسبيا الأعلى التالية، ينبغي، قياسا على \S\ 7، تكوين منظومة $s$ بالنسبة إلى المقياس التالي
\[
\nu(s)=(ss'u)^4
\]
وطيّها مع ملازمات المنظومة الكاملة نسبيا بالنسبة إلى المقياس $s$. وسنبيّن أنه عند إدخال المقياس $(\varrho,t)$ بوصفه مقياسا أعلى
\[
\begin{array}{ll}
\text{A)}&\text{المقياس }(ss'u)^4\text{ يصبح قابلا للاختزال،}\vphantom{\dfrac11}\\
\text{B)}&\text{منظومة }s\text{ تتكون من الملازم الوحيد }s\text{.}
\end{array}
\]

A) ينتج اختزال المقياس $(ss'u)^4$ مباشرة من المختزل $s_\nu^2$ الذي وجد بالصيغة (16.) في \S\ 11. فمن
\[
 s_\nu^2=\frac43 i a_\nu^2+\frac45\theta_\varrho^2+\frac85(atu)^2-
 \frac{26}{15}(\nu\varrho x)^2+\frac16J\cdot u_x^2-
 \frac29i^2\cdot u_x^2
\]
ينتج، بأخذ الأشكال القطبية من $x$ إلى $\nu_1$:
\srcnumdisplay{(21.)\srcfnmark{*)}}{%
\begin{aligned}
(ss'u)^4&=-\frac23J\cdot \nu+6s_\nu^2s_{\nu_1}^2\\
&=\frac{24}{5}\theta_\nu^2\theta_\varrho^2+
\frac{48}{5}a_\nu^2(atu)^2-
\frac{52}{5}H_\varrho^2+
\frac43i\cdot(asu)^4+\frac13J\cdot \nu-
\frac49i^2\cdot \nu,
\end{aligned}
}
وبذلك ينتج اختزال المقياس $(ss'u)^4$.

\srcfntext{*)}{باستعمال الصيغة (10.)c، تنتج من الصيغة (21.) العلاقة بين الثوابت الثلاثة من الرتبة 9، المذكورة في الحاشية على المقدمة:
\srcnumdisplay{(22.)}{%
\begin{aligned}
(ass')^4&=\frac{24}{5}\theta_\nu^2\theta_\varrho^2a_\nu^2a_\vartheta^2+
\frac{48}{5}a_\nu^2b_\nu^2(tab)^2-
\frac{52}{5}H_\varrho^2a_\eta^4+
\frac53J\cdot i-\frac49i^3,\\
(ass')^4&=\frac{24}{5}a_\varrho^2a_{\varrho'}^2-
\frac{12}{5}t_\varrho^2+\frac53J\cdot i-\frac49i^3.
\end{aligned}
}
}

B) إن اختزال
\[
Z=(ss'u)^2=\theta(s)
\]
وما ينتج عنه من اختزال منظومة $s$ يحصل بإبدال الملازم $Z$ بالملازم الأدنى $g_\nu^2$ بحسب الصيغة (8.)b في \S\ 7، مع مراعاة العلاقة
\[
 s_\eta^2=s_\nu^2(\nu\nu_1x)^2-\frac13Z.
\]
للملازم $g_\nu^2$، بحسب الصيغة (13.) في \S\ 11، المختزل $s_\nu^2$ عاملا. وبحسب الصيغتين (8.) و(16.) نحصل على:
\[
 g_\nu^2=-Z+\frac{14}{5}ia_\eta^2+\frac{14}{15}i(asu)^2
 \modu{(\varrho,t)},
\]
وبحسب الصيغ (13.) و(14.) و(15.) و(16.):
\[
\begin{aligned}
g_\nu^2&=2[s_\vartheta^2]_{\nu^2}-\frac43 i\A_\nu^2
      =2s_\vartheta^2s_\nu^2-\frac65[s_\vartheta^2(\theta su)^2]\widehat{\nu x}^{2}
        -\frac43i\A_\nu^2\\
&=\frac45i\cdot a_\eta^2-\frac25i\cdot(asu)^2+\frac13J\cdot\theta
\modu{(\varrho,t)}.
\end{aligned}
\]
ومن ذلك ينتج
\srcnumdisplay{(23.)}{%
Z=2i\cdot a_\eta^2+\frac43i\cdot(asu)^2-\frac13J\cdot\theta
 \modu{(\varrho,t)}.
}

إذن تتحول المنظومة الكاملة نسبيا بالنسبة إلى المقياس $((ss'u)^4,\varrho,t)$ إلى منظومة كاملة نسبيا بالنسبة إلى المقياس $(\varrho,t)$، ومنها تنشأ المنظومة الكاملة مطلقا بعملية الاشتقاق والتقليص على المنظومة المعروفة لحدوديتين متجانستين من الدرجة الثانية. ولتكوين المنظومة الكاملة نسبيا بالنسبة إلى المقياس $(\varrho,t)$، علينا، بما أن منظومة $s$ تختزل بحسب الصيغة (23.) إلى الملازم $s$، إجراء عملية الاشتقاق والتقليص للمنظومة الكاملة نسبيا بالنسبة إلى المقياس $(s,\varrho,t)$ على $s$ وقوى $s$. وسيتبيّن أن قوى $s$ لا تدخل في الطيّ إلا مع المنظومة I.

نعرّف \glqq{}الملازمات الأعلى\grqq{}، بين الملازمات المتساوية في الرتبة والدرجة، بأنها تلك التي تنشأ من ملازمات أعلى في المنظومة الكاملة نسبيا بالنسبة إلى المقياس $(s,\varrho,t)$ بطيّها مع $s$، إذ نعدّ الطيّ مع $s$ مجرد أخذ الأشكال القطبية للملازمات الأصلية. وبذلك يتحدد ترتيب الملازمات المفردة تحديدا وحيدا (انظر \S\ 1. سلاسل الملازمات 2). فمثلا:
\[
(\theta_\eta(\theta Hu),s,u)^2\text{ أعلى من }s_\eta((\theta Hu)^2,s,u)\srcfnmark{*)},
\]
\[
(\theta_\eta(\theta Hu)^2,s,u)\text{ أعلى من }(\theta_\eta(\theta Hu),s,u)^2.
\]

نقسم المنظومة الناتجة إلى المنظومات الجزئية الأربع:
\[
\begin{array}{ll}
\text{المنظومة }s_I:&\text{ناتجة من طيّ }s\text{ وقوى }s\text{ مع المنظومة I،}\\
\text{المنظومة }s_{II}:&\text{ناتجة من طيّ }s\text{ مع المنظومة II،}\\
\text{المنظومة }s_{III}:&\text{ناتجة من طيّ }s\text{ مع الملازمات المفردة}\\
&\text{(دون الجداءات) من المنظومة III، ومع جداءات كل منها من}\\
&\text{ملازمين، أحدهما من المنظومة I والآخر من المنظومة II،}\\
\text{المنظومة }s_{IV}:&\text{ناتجة من طيّ }s\text{ مع جداءات ملازمين، أحدهما من}\\
&\text{المنظومة I والآخر من المنظومة III، أو أحدهما من المنظومة II والآخر من المنظومة III، أو كليهما من المنظومة III.}
\end{array}
\]
ونلاحظ أيضا أن جميع الصيغ من الآن فصاعدا تفهم بالنسبة إلى المقياس $(\varrho,t)$، وابتداء من الصيغة (27.) بالنسبة إلى المقياس $(\varrho,t,i,J,\text{ الملازمات الأعلى})$، دون التصريح بذلك في كل مرة.

للاختزال نجمع الصيغ التي حصلنا عليها سابقا:
\srcnumdisplay{(24.)}{%
\begin{aligned}
s_\nu^2&=\frac43i\,a_\nu^2+\frac16J\cdot u_x^2-\frac29i^2\cdot u_x^2,\\
s_\nu^3&=\frac13J\cdot u_x,
& s_\nu^4&=J,\\
g&=2s_\vartheta^2-\frac43i\Delta,\\
s_\vartheta^2(\theta su)&=0,\\
s_\vartheta^2(\theta su)^2&=\frac89i\cdot a_\nu^2-\frac4{27}i^2\cdot u_x^2,\\
Z&=2i\cdot a_\eta^2+\frac43i(asu)^2-\frac13J\cdot\theta,\\
g_\nu&=-s_\eta+u_x\left\{2ia_\eta^2-\frac23i(asu)^2+\frac23J\cdot\theta\right\},\\
g_\nu^2&=\frac45ia_\eta^2-\frac25i(asu)^2+\frac13J\cdot\theta,
&g_\nu^3&=0,\quad g_\nu^4=0.
\end{aligned}
}
\srcfntext{*)}{اختصارا للكتابة، ندرج حدّ عملية الاشتقاق والتقليص، بدلا من العملية الكاملة $[(\theta Hu)^2]_{\widehat{su}}s$، بوصفه ملازما من المنظومة. وعند وضع الصيغ تستبدل جميع عمليات الاشتقاق والتقليص بحدودها الابتدائية؛ فمثلا يستبدل $(\theta_\eta(\theta Hu),s,u)^2$ بـ$\theta_\eta(\theta Hu)(\theta su)^2$.}

\noindent\emph{النتائج:}
\[
 s_\nu^2,\quad s_\vartheta^2(\theta su),\quad s_\vartheta^2s_\nu,\quad s_\vartheta^2\theta_\nu
\]
هي مختزلات.

\begin{center}
\textbf{\S\ 18.}\quad\emph{المنظومة $s_I$.}
\end{center}

\[
\text{طيّ }\left\{\begin{array}{l}
s\text{ مع } f;\ \theta;\ K,j,\Delta;\ f\cdot j,N;\ \Delta\cdot j,\\
s^2\text{ مع }\theta,K.
\end{array}\right.
\]

\noindent\emph{الملازمات غير القابلة للاختزال.}

\noindent الرتبة الخامسة:
\[
(asu);\qquad (asu)^2;\qquad (asu)^3;\qquad (asu)^4.
\]

\noindent الرتبة السادسة:
\[
\begin{array}{cccc}
(\theta su)&(\theta su)^2&(\theta su)^3&(\theta su)^4\\
s_\vartheta&s_\vartheta(\theta su)&s_\vartheta(\theta su)^2&s_\vartheta(\theta su)^3\\
\cdot&s_\vartheta^2&\cdot&\cdot
\end{array}
\]

\noindent الرتبة السابعة:
\[
\begin{array}{c@{\quad}c@{\quad}c@{\quad}c@{\qquad}l}
\cdot&(Ksu)^2&(Ksu)^3&(Ksu)^4&\\
\text{A)}\ s_k&s_k(Ksu)&s_k(Ksu)^2&s_k(Ksu)^3&
\text{B)}\ s_j,\\
\cdot&s_k^2&\cdot&\cdot&
\text{C)}\ (\Delta su),\ (\Delta su)^2.
\end{array}
\]

\noindent الرتبة الثامنة:
\[
\begin{array}{ll}
\text{A)}\ s_j(asu),\ s_j(asu)^2,\ s_j(asu)^3,
&
\text{B)}\ \begin{array}[t]{cc}
\cdot&(Nsu)^2\\
s_n&s_n(Nsu).
\end{array}
\end{array}
\]

\noindent الرتبة العاشرة:
\[
\text{A)}\ s_j(\Delta su),\qquad \text{B)}\ s_\vartheta(\theta s'u)^4.
\]

\noindent الرتبة الحادية عشرة:
\[
\begin{array}{cc}
(Ks^2u)^6&\cdot\\
s_k(Ks^2u)^5&s_k(Ks^2u)^6.
\end{array}
\]

\noindent\emph{الاختزالات:}

في سلاسل الملازمات
\[
 s_\vartheta^2(\theta su),\quad s_k^2(Ksu),\quad s_j^2,\quad (\Delta su)^3,\quad (Nsu)^3
\]
يكون المختزل $s_\vartheta^2(\theta su)$ عاملا؛ أما $s_j^2$ و$(\Delta su)^3$ فينتجان بالرجوع من $j$ و$\Delta$ إلى ملازمات أدنى في سلسلة الملازمات:
\[
\begin{aligned}
s_\vartheta^2(\theta su)(a\theta u)^3&=-\frac12s_j^2,\\
s_\vartheta^2(\theta su)(a\theta u)^2&=\frac13(\Delta su)^3,\\
s_\vartheta^2(\theta su)^2(a\theta u)^2&=\frac13(\Delta su)^4+\frac13s_j^2.
\end{aligned}
\]
الملازمان $s_\vartheta^2s_{\vartheta'}$ و$s_\vartheta^2s_{\vartheta'}^2$: المختزلان $s_\vartheta^2(\theta su)$ و$\theta_{\vartheta'}$:
\[
\begin{aligned}
s_\vartheta^2s_{\vartheta'}&=s_\vartheta^2\theta_{\vartheta'}-s_\vartheta^2(\theta s\widehat{\vartheta'}x),\\
s_\vartheta^2s_{\vartheta'}^2&=s_\vartheta^2s_{\vartheta'}\theta_{\vartheta'}-s_\vartheta^2s_{\vartheta'}(\theta s\widehat{\vartheta'}x).
\end{aligned}
\]
سلسلة الملازمات $s_j(\Delta su)^2$: المختزلان $\Delta_j$ و$(\Delta su)^3$.

\noindent\emph{النتائج:}

لا يدخل $s$ بقواه في الطيّ مع $f,j,\Delta,N$. ولا تدخل $f,j,\Delta,N$ في الطيّ إلا في جداءات مع ملازمات متخالفة التدرج؛ غير أن هذه الملازمات يجوز، في الجداءات مع $f$، أن تبقى من الدرجة الثانية في $x$، وفي الجداءات مع $\Delta$ من الدرجة الأولى في $x$؛ أي ينبغي النظر في جداءات الملازمات الآتية:
\[
 f\cdot(0,\lambda),\quad f\cdot(1,\lambda),\quad f\cdot(2,\lambda),\quad
 j\cdot(\mu,0),\quad N,\quad
 \Delta\cdot(0,\lambda),\quad \Delta\cdot(1,\lambda)\quad(\lambda>0),
\]
حيث يدل $(\mu,\lambda)$ على ملازم درجته $\mu$ بالنسبة إلى $x$ ودرجته $\lambda$ بالنسبة إلى $u$.

لا تدخل $\theta$ ولا $K$، في قوى أو جداءات مع ملازمات اعتباطية، في الطيّ مع $s$ أو $s^2$. وبالنسبة إلى الجداءات مع المحددة اليعقوبية $K$، أي $K\cdot\varphi_x$ حيث $\varphi\ne f$ أو $\theta$، تسري العلاقة الآتية بين الملازمات القابلة للتحلل:
\[
K\cdot\varphi_x=(a\theta u)\varphi_x=f\cdot(\varphi\theta u)-\theta\cdot(\varphi a u).
\]
وبذلك يثبت المخطط أعلاه لطيّ المنظومة $s_I$.

\begin{center}
\textbf{\S\ 19.}\quad\emph{المنظومة $s_{II}$.}
\end{center}

طيّ $s$ مع $\nu;\ H;\ L;\ H^2$.

\noindent\emph{الملازمات غير القابلة للاختزال:}
\[
\begin{array}{c@{\qquad}c@{\qquad}c@{\qquad}c}
\text{الرتبة السادسة} & \text{الرتبة الثامنة} & \text{الرتبة العاشرة} & \text{الرتبة الثانية عشرة}\\[0.3em]
s_\nu & (Hsu),\ (Hsu)^2 & (Lsu)^2,\ (Lsu)^3 & (H^2su)^3\\
\cdot & s_\eta & s_l & \cdot\\
\cdot & \cdot & \cdot & \cdot\\
s_\nu^4=J & \cdot & \cdot & \cdot
\end{array}
\]

\noindent\emph{الاختزالات:}

سلسلتا الملازمات $s_\eta(Hsu)$ و$s_l(Lsu)$: المختزل $s_\nu^2$.

سلسلة الملازمات $s_\sigma$: المختزل $s_\nu^2$ المستخرج من $H$، حيث $s_\nu^2=\dfrac23s_\sigma$.

ينتج تحلل $(H^2su)^4$ من الصيغة (7.)a (انظر \S\ 11. A). ومن ذلك ينتج تحلل $(H\cdot L,s,u)^4$.

\noindent\emph{النتائج:}

لا يدخل $s^2$ في الطيّ مع المنظومة II.

لا يدخل $\nu$ في الطيّ إلا في جداءات مع ملازمات متخالفة التدرج؛ أما $H$ فلا يدخل في الطيّ، عند اعتباره ملازما موافقا للتغير، إلا في جداءات مع ملازمات $(1,\lambda)$ و$(2,\lambda)$ و$(3,\lambda)$، حيث $\lambda$ اعتباطي.

لا يدخل $\sigma$ ولا $(\nu\sigma x)$ في الطيّ مطلقا، ولا يدخل $L$، بوصفه محددة يعقوبية، في الطيّ ضمن جداءات (إذ تصبح عمليات الاشتقاق والتقليص الكاملة على الملازمات الموافقة للتغير في المنظومة III قابلة للاختزال).

وتبقى من جداءات المنظومتين I وII:
\[
f\cdot \nu,\qquad \Delta\cdot \nu.
\]

\begin{center}
\textbf{\S\ 20.}\quad\emph{ملازمات الرتبة السابعة (المنظومة $s_{III}$).}
\end{center}

طيّ $s$ مع $f\cdot \nu$ و$a_\nu$ و$a_\nu^2$.

\noindent\emph{الملازمات غير القابلة للاختزال:}
\[
\begin{gathered}
s_\nu(asu),\quad a_\nu(asu),\quad s_\nu(asu)^2,\quad a_\nu(asu)^2,\quad
s_\nu(asu)^3,\quad a_\nu(asu)^3,\\
a_\nu s_\nu,\qquad a_\nu s_\nu(asu),\qquad a_\nu^2(asu),\qquad a_\nu^2(asu)^2.
\end{gathered}
\]

\noindent\emph{الاختزالات:}

لن نذكر الاختزالات البديهية الناتجة من الطيّ المباشر مع المختزلين $s_\nu^2$ و$s_\eta(Hsu)$، ولا تحلل عمليات الاشتقاق والتقليص مع المحددات اليعقوبية.

الملازمان $a_\nu s_\nu(asu)^2$ و$a_\nu^2s_\nu(asu)^2$: انظر الصيغة (19.)، \S\ 12؛
\[
a_\nu s_\nu(asu)^3=a_\nu(a\widehat{\nu_1}\nu_2u)^3(\nu_1\nu_2\nu)=0.
\]
سلسلة الملازمات $a_\nu^2s_\nu$: المختزل $s_\vartheta^2(\theta su)$ بالرجوع بـ$a_\nu^2$ إلى ملازمات أدنى، أي بالاختزال المزدوج للعبارتين
\[
s_\vartheta^2(a\theta s)(a\theta u),\qquad
s_\vartheta^2(a\theta s)(a\theta u)^2
\]
وفقا للصيغة المفترضة الآتية:
\[
\begin{aligned}
s_\vartheta^2(a\theta s)(a\theta u)
&=[(a\theta u)^2]s^3+\frac15[a_\vartheta(a\theta u)^2]s^2+\frac1{60}[a_\vartheta^2(a\theta u)^2]s\\
&=\frac12a_\nu^2s_\nu-\frac23a_\nu s_\nu^2,\\
s_\vartheta^2(a\theta s)(a\theta u)^2
&=\frac45[a_\vartheta(a\theta u)^2]_{\widehat{us}}s^2+\frac{23}{180}[a_\vartheta^2(a\theta u)^2]_{\widehat{us}}s\\
&=-\frac12a_\nu^2s_\nu(asu).
\end{aligned}
\]
وفي ذلك تتحقق $a_\nu^2b_\nu(abu)=0$.

فنحصل على الصيغ:
\srcnumdisplay{(25.)}{%
\begin{array}{rcl@{\qquad}rcl}
\underline{a_\nu^2s_\nu}&=&\dfrac43i\cdot a_\nu^2b_\nu,
& \underline{a_\nu s_\nu(asu)^2}&=&\dfrac13i\cdot\theta_\nu^2-\dfrac1{18}iH,\\[0.7em]
\underline{a_\nu s_\nu(asu)^3}&=&0,
& \underline{a_\nu^2s_\nu(asu)}&=&0,\\[0.7em]
\underline{a_\nu^2s_\nu(asu)^2}&=&0.
\end{array}
}

\noindent\emph{النتائج:}

1) الملازمان $a_\nu s_\nu(asu)^2$ و$a_\nu^2s_\nu$ مختزلان.

2) تنتج قيم الملازمين $(\Delta su)^3$ و$(\Delta su)^4$ اللذين عُدّا قابلين للاختزال في \S\ 19؛ وكذلك قيم سلسلة الملازمات المناظرة $(agu)^2$، التي يفضي طيّها مع $\nu$ بواسطة المختزل $g_\nu$ إلى الاختزال المزدوج.
\srcnumdisplay{(26.)}{%
\begin{aligned}
\text{a)}\quad (\Delta su)^3&=-\frac65a_\nu^2(asu)+3a_\nu s_\nu(asu)+2i\theta_\nu(\vartheta \nu x),\\
\text{b)}\quad (\Delta su)^4&=-\frac25a_\nu^2(asu)^2+\frac43i\theta_\nu^2+\frac19iH.
\end{aligned}
}

من الصيغتين (24.) و(26.)، بترديد $(i,J)$ (قارن ص. 71):
\srcnumdisplay{(27.)}{%
\begin{aligned}
\text{a)}\quad (agu)^2&=\frac23(\Delta su)^2-\frac43a_\nu s_\nu+\frac35s\cdot a_\nu^2,\\
\text{b)}\quad (agu)^3&=2a_\nu s_\nu(asu)-a_\nu^2(asu)^2,\\
\text{c)}\quad (agu)^4&=-a_\nu^2(asu)^2.
\end{aligned}
}


\begin{center}
\textbf{\S\ 21.}\quad\emph{ملازمات الرتبة الثامنة (المنظومة $s_{III}$).}
\end{center}

طيّ $s$ مع ملازمات الرتبة الرابعة (المنظومة III). (\S\ 9.)

\noindent\emph{الملازمات غير القابلة للاختزال:}
\[
\begin{array}{c@{\quad}c|c|c}
\cdot
& (\vartheta \nu x)s_\nu,\ ((\vartheta \nu x)^2,s,u)
& ((\vartheta \nu x),s,u)^2,\ (\theta_\nu su)
& \bigstar\\[0.45em]
& (\vartheta \nu x)^2s_\nu
& ((\vartheta \nu x)^2,s,u)^2,\ \theta s_\nu
& \\[0.45em]
&& ((\vartheta \nu x)^2,s,u)s_\nu,\ (\theta_\nu(\vartheta \nu x)^2,s,u)
&
\end{array}
\]
\[
\begin{array}{c|c|c|c}
\bigstar
& ((\vartheta \nu x),s,u)^3,\ (\theta_\nu su)^2
& ((\vartheta \nu x),s,u)^4,\ (\theta_\nu su)^3
& \cdot\\[0.45em]
& ((\vartheta \nu x)^2,s,u)^3,\ (\theta_\nu(\vartheta \nu x),s,u)^2,\ (\theta_\nu^2su)
& (\theta_\nu(\vartheta \nu x),s,u)^3,\ (\theta_\nu^2su)^2
& (\theta_\nu(\vartheta \nu x),s,u)^4\\[0.45em]
& \cdot
& (\theta_\nu^3su)
& \cdot
\end{array}
\]

\noindent\emph{الاختزالات:}

يتحلل $((\vartheta \nu x),s,u)s_\nu$ بوصفه عملية اشتقاق وتقليص على المحددات اليعقوبية وفقا لـ\S\ 3.

سلسلة الملازمات $((\vartheta \nu x),s,u)^2s_\nu$: المختزلان $s_\nu^2$ و$s_\vartheta^2\theta_\nu$:
\[
(\widehat{\vartheta \nu su})s_\nu(\theta su)=s_\nu s_\vartheta(\theta su)=-\frac12(\widehat{\vartheta \nu su})^2(\theta su),
\]
\[
\theta_\nu(\widehat{\vartheta \nu su})s_\nu=\theta_\nu s_\nu s_\vartheta=-\frac12\theta_\nu(\widehat{\vartheta \nu su})^2.
\]

\noindent(قاعدة الضرب.)

سلسلة الملازمات $\theta_\nu s_\nu(\theta su)$: المختزل $a_\nu s_\nu(asu)^2$ بالاختزال المزدوج للعبارتين
\[
a_\nu s_\nu(asu)^2(abu)\quad\text{و}\quad a_\nu s_\nu(asu)^2(abu)(sbu);
\]
ويُختزل $\theta_\nu s_\nu(\theta su)^3$ بالاختزال المزدوج لـ$(agu)^3(abu)(bgu)^3$.

سلسلة الملازمات $\theta_\nu(\vartheta \nu x)s_\nu$: المختزل $a_\nu^2s_\nu$ من $\theta_\nu(\vartheta \nu x)s_\nu=a_\nu^2(abu)s_\nu$.

سلسلة الملازمات $(\theta_\nu(\vartheta \nu x)^2,s,u)$: الاختزال المزدوج لملازمات سلسلة الملازمات $\theta_\nu(\widehat{\vartheta \nu su})s_\nu$، التي تنتمي في الوقت نفسه إلى سلسلتي الملازمات $((\vartheta \nu x)su)^2s_\nu$ و$\theta_\nu(\vartheta \nu x)s_\nu$.

يتحلل الملازم $(\theta_\nu(\vartheta \nu x),s,u)$ بوصفه عملية اشتقاق وتقليص مفردة على المحددة اليعقوبية $\theta_\nu(\vartheta \nu x)=a_\nu^2(abu)$.

الملازم $((\vartheta \nu x)^2,s,u)^4$: الاختزال المزدوج للعبارة $(a\A u)^2(\A su)^4$، أو كذلك لـ$(\theta gu)^4$ من الصيغ (27.) وبالقياس على اختزال $(j\eta x)^4$ (\S\ 12 C).

ينتج عن الاختزال المزدوج الصيغ الآتية:
\srcnumdisplay{(28.)}{%
\begin{aligned}
0&=\theta_\nu s_\nu(\theta su)-\frac16(\widehat{\vartheta \nu su})^2(\theta su)-\frac1{12}(Hsu),\\
0&=\theta_\nu s(\theta su)^2-\frac14(\widehat{\vartheta \nu su})^2(\theta su)^2-\frac1{24}(Hsu)^2,\\
0&=(\widehat{\vartheta \nu su})^2(\theta su)^2+2\theta_\nu(\widehat{\vartheta \nu su})(\theta su)^2
  +3\theta_\nu^2(\theta su)^2-\frac13H(su)^2,\\
0&=\theta_\nu s_\nu(\theta su)^3+\frac12\theta_\nu(\widehat{\vartheta \nu su})(\theta su)^3,\\
0&=\theta_\nu s_\nu s_\vartheta-\frac14s_\eta,
\quad 0=\theta_\nu s_\nu s_\vartheta(\theta su),
\quad 0=\theta_\nu s_\nu s_\vartheta(\theta su)^2.
\end{aligned}
}
ويقابل الأخير $((\theta_\nu(\vartheta \nu x)^2,s,u)^2,\ldots)$.

\noindent\emph{النتائج:}

1) الملازمات $\theta_\nu(\vartheta \nu x)s_\nu$ و$(\widehat{\vartheta \nu su})(\theta su)s_\nu$ و$\theta_\nu(\widehat{\vartheta \nu su})^2$ و$(\widehat{\vartheta \nu su})^2(\theta su)^2$ مختزلات.

2) لا تدخل ملازمات الرتبة الرابعة $(5,\lambda)$ و$(6,\lambda)$ وما إلى ذلك في الطيّ مع $s^2$.

3) لسلسلة الملازمات $(\theta gu)^2$ تنتج، وفقا لـ(27.) ومع مراعاة اختزالات هذا القسم، الصيغ الآتية:
\srcnumdisplay{(29.)}{%
\begin{aligned}
\text{a)}\quad (\theta gu)^2&=\frac43\theta_\nu s_\nu+\frac23s_\nu s_\vartheta-\frac13s\cdot\theta_\nu^2-\frac19s\cdot H-\frac13f\cdot a_\nu^2(asu)^2+\frac13a_\nu^2\cdot(asu)^2,\\
\text{b)}\quad (\theta gu)^3&=-\frac13(\widehat{\vartheta \nu su})^2(\theta su)-\theta_\nu(\widehat{\vartheta \nu su})(\theta su)-\theta_\nu^2(\theta su),\\
\text{c)}\quad (\theta gu)^4&=2\theta_\nu^2(\theta su)^2-\frac13(Hsu)^2,\\
g\vartheta(\theta gu)&=(\widehat{\vartheta \nu su})(\vartheta \nu x)s_\nu-\theta_\nu(\vartheta \nu x)(\widehat{\vartheta \nu su}),\\
g\vartheta(\theta gu)^2&=\frac23s\cdot\theta_\nu^3,
\quad g\vartheta(\theta gu)^3=\frac13\theta_\nu^3(\theta su),
\quad g\vartheta(\theta gu)^4=0.
\end{aligned}
}

\begin{center}
\textbf{\S\ 22.}\quad\emph{ملازمات الرتبة التاسعة (المنظومة $s_{III}$).}
\end{center}

طيّ $s$ مع ملازمات الرتبة الخامسة (المنظومة III) ومع الجداء $\A\cdot \nu$ (\S\ 10).

\noindent\emph{الملازمات غير القابلة للاختزال:}
\[
\begin{array}{c@{\quad}c|c|c|c}
\text{A)}& \cdot & \cdot & (K_\nu su)^2 & \bigstar\\[0.45em]
& \cdot & ((k\nu x)^2,s,u)^2,\ K_\nu s_\nu & ((k\nu x)^2,s,u)^3 & \\[0.45em]
& (k\nu x)^2s_\nu,\ ((k\nu x)^3,s,u) & ((k\nu x)^3,s,u)^2 & \cdot & \cdot\\[0.45em]
& (k\nu x)^3s_\nu & ((k\nu x)^3,s,u)s_\nu,\ (K_\nu(k\nu x)^3,s,u) & \cdot & \cdot\\[0.45em]
& \bigstar\ (K_\nu su)^3 & (K_\nu su)^4 & \cdot & (K_\nu^2su)^4
\end{array}
\]
\[
\begin{array}{c@{\quad}c|c@{\qquad}c@{\quad}c}
\text{B)}& (j\nu x)s_\nu,\ ((j\nu x)^2,s,u) & ((j\nu x)^2,s,u)^2 && \\
& ((j\nu x)^3,s,u) & \cdot && \text{C)}\quad s_\nu(\A su),\ (\A_\nu su),
\end{array}
\]
\[
\begin{array}{c@{\quad}c|c|c}
\text{D)}& \cdot & ((aHu),s,u)^2,\ ((aHu)^2,s,u) & \bigstar\\[0.45em]
& (aHu)s_\eta,\ (a_\eta su) & (aHu)^2s_\eta,\ (a_\eta su)^2 & \\[0.45em]
& \cdot & (a_\eta^2su) & \\
\bigstar& ((aHu),s,u)^3,\ ((aHu)^2,s,u)^2 & ((aHu),s,u)^4 & \\
& (a_\eta su)^3,\ (a_\eta(aHu),s,u)^2,\ (a_\eta(aHu)^2,s,u) & (a_\eta(aHu),s,u)^3 &
\end{array}
\]

\noindent\emph{الاختزالات:}

A) تتبع الاختزالات مباشرة من اختزالات \S\ 21 بواسطة طيّ واحد. تتخذ سلسلة الملازمات $K_\nu s_\nu(Ksu)^2$ المختزلين $a_\nu s_\nu(asu)^2$ و$\theta_\nu s_\nu(\theta su)^2$ عاملين لها. ومن ذلك تحلل الملازمين الأعلى رتبة $K_\nu^2(Ksu)^2$ و$K_\nu^2(Ksu)^3$ وفقا للفرض الآتي:
\[
0=a_\nu s_\nu(asu)^2(a\theta u)=K_\nu s_\nu(Ksu)^2=\theta_\nu s_\nu(\theta su)^2(a\theta u).
\]

B) سلسلة الملازمات $(j\nu x)^2s_\nu$: المختزل $a_\nu^2s_\nu$ من
\[
(j\nu x)^2s_\nu=3a_\nu^2s_\nu(a\theta u)^2.
\]
ولـ$((j\nu x)^3,s,u)^2$: المختزل $a_\nu^2s_\nu$ من
\[
((j\nu x)^3,s,u)^2=-6a_\nu^2s_\nu(a\theta u)^2(\theta su).
\]

C) الملازمان $\A_\nu(\A su)^2$ و$s_\nu(\A su)^2$: المختزل $g_\nu$ من الصيغة (27.)a.

الملازم $\A_\nu s_\nu$: 1) المختزل $a_\nu^2s_\nu$ من $a_ga_\nu^2s_\nu=\frac23\A_\nu s_\nu$؛

2) يتحلل وفقا للصيغة (27.)a بالاختزال المزدوج للعبارة $(\A s\widehat{\nu}x)^2$.

ومن ذلك: تحلل الملازم الأعلى رتبة $a_\eta s_\eta$.

D) سلسلة الملازمات $(aHu)^2s_\eta(asu)$: المختزل $a_\nu s_\nu(asu)^2$ بالاختزال المزدوج لسلسلة الملازمات $[a_\nu s_\nu(asu)^2]_{\widehat{\nu_1}x}$؛ واختزال $(aHu)^2s_\eta(asu)^2$ من $a_\nu s_\nu(asu)^2$ و$a_\nu s_\nu(asu)^3$. ومن ذلك اختزال $(a_\eta,s,u)^4$.

سلسلة الملازمات $a_\eta(aHu)s_\eta$: المختزل $a_\nu^2s_\nu$؛ ويُختزل $a_\eta^2s_\eta$ بالاختزال المزدوج للعبارة $(\A s\widehat{\nu}x)^3$، لأن $a_\eta^2s_\eta$ لا ينشأ بالطيّ من الحد القابل للاختزال $a_\nu^2s_\nu(\nu\nu_1x)$ في الملازم $a_\eta(aHu)s_\eta$.

سلسلة الملازمات $(a_\eta^2su)^2$: المختزل $a_\nu^2s_\nu$ من
\[
a_\nu^2s_\nu s_{\nu_1}=-\frac12a_\eta^2(Hsu)^2.
\]
يتحلل الملازم $(a_\eta(aHu),s,u)$ بوصفه عملية اشتقاق وتقليص على محددة يعقوبية.

تنتج صيغ الاختزال الآتية:
\srcnumdisplay{(30.)}{%
\begin{aligned}
\text{a)}\quad 0&=(aHu)^2(asu)s_\eta-a_\eta(asu)(Hsu)^2-a_\eta(aHu)(asu)(Hsu),\\
\text{b)}\quad 0&=(aHu)^2(asu)^2s_\eta-a_\eta(aHu)(asu)^2(Hsu),\\
\text{c)}\quad 0&=a_\eta(asu)^2(Hsu)^2,\\
0&=a_\eta s_\eta+\frac14a_\nu^2\cdot g-\frac14s\cdot a_\eta^2,\\
0&=a_\eta(aHu)s_\eta-\frac12a_\eta^2(Hsu),
\quad 0=a_\eta^2(Hsu)^2,\ldots,
\quad 0=a_\eta^2s_\eta.
\end{aligned}
}

\noindent\emph{النتائج:}

1) الملازمات $(j\nu x)^2s_\nu$ و$\A_\nu s_\nu$ و$\A_\nu(\A su)^2$، ومن ثم $N_\nu(Nsu)^2$ و$(aHu)^2(asu)s_\eta$ و$a_\eta(aHu)s_\eta$ و$a_\eta^2(Hsu)^2$، مختزلات؛ أما $a_\eta s_\eta$ فهو ملازم متحلل يتحول عند جميع الطيات المفردة والمزدوجة إلى ملازمات متحللة.

2) لا تدخل ملازمات الرتبة الخامسة $(5,\lambda)$ و$(6,\lambda)$ وما إلى ذلك في الطيّ مع $s^2$.


\begin{center}\textbf{\S\ 23.}\quad\emph{ملازمات الرتبة العاشرة (المنظومة $s_{III}$).}\end{center}
طيّ $s$ مع ملازمات الرتبة السادسة (المنظومة III) (\S\ 11).

\noindent\emph{الملازمات غير القابلة للاختزال:}
\[
\begin{array}{r@{\quad}c}
\text{A)}&
\begin{array}{c|c}
(\vartheta^2\nu x)^3s_\nu&\cdot\\[0.45em]
\cdot&(\vartheta^2\nu x)^3s_\nu s_\vartheta,\ \theta_\nu(\vartheta^2\nu x)^3s_\vartheta
\end{array}
\end{array}
\]
\[
\begin{array}{r@{\quad}c}
\text{B)}&
\begin{array}{c|ccc}
\cdot&(a_\nu(j\nu x),s,u)^2&(a_\nu(j\nu x),s,u)^3&(a_\nu(j\nu x),s,u)^4\\[0.45em]
a_\nu(j\nu x)s_\nu&\cdot&(a_\nu^2(j\nu x),s,u)^2&(a_\nu^2(j\nu x),s,u)^3
\end{array}
\end{array}
\]
\[
\begin{array}{r@{\quad}c}
\text{C)}&
\begin{array}{c|c|c|c}
\cdot&\cdot&((\theta Hu),s,u)^2,\ ((\theta Hu)^2,s,u)&\bigstar\\[0.35em]
\cdot&((\vartheta\eta x),s,u)^2,\ (\theta Hu)s_\eta,\ (\theta_\eta su)&((\vartheta\eta x),s,u)^3,\ (\theta_\eta su)^2&\\[0.35em]
((\vartheta\eta x)^2,s,u)&\cdot&(\theta_\eta^2su)&\\[0.35em]
\cdot&(\theta_\eta(\vartheta\eta x)^2,s,u)&\cdot&
\end{array}
\end{array}
\]
\[
\begin{array}{c@{\quad}|c|c}
\bigstar&((\theta Hu),s,u)^3,\ ((\theta Hu)^2,s,u)^2&((\theta Hu),s,u)^4\\[0.45em]
&(H_\vartheta(\theta Hu),s,u)^2,\ (\theta_\eta(\theta Hu),s,u)^2,\ (\theta_\eta(\theta Hu)^2,s,u)&(\theta_\eta su)^4,\ (\theta_\eta(\theta Hu),s,u)^3\\[0.45em]
&\cdot&(\theta_\eta^2(\theta Hu),s,u)^2
\end{array}
\]

\noindent\emph{الاختزالات:}

A) يتحلل الملازمان $((\vartheta^2\nu x)^3,s,u)^2$ و$((\vartheta^2\nu x)^3,s,u)^3$:
\[
(\vartheta \nu x)(\widehat{\vartheta \nu su})(\widehat{\vartheta'\nu su})
=(\vartheta \nu x)s_\vartheta s_{\vartheta'}-(\vartheta \nu x)s_\nu s_\vartheta
=-2\theta(\vartheta \nu x)s_\nu s_\vartheta+(\vartheta \nu x)s_\vartheta^2.
\]
الملازمات الباقية: إما أن يكون مختزل عاملا فيها، وإما تنتج من طيّ واحد مع ملازمات متحللة.

B) المختزلان $a_\nu^2s_\nu$ و$(\widehat{j\nu su})s_\nu$.

C) 1. سلسلة الملازمات $(\theta Hu)^2s_\eta$: a) المختزل $(aHu)^2(asu)s_\eta$؛ b) الاختزال المزدوج لسلسلتي الملازمات $(\theta gu)^3g_\nu$ و$g_\nu(agu)^3(bgu)^2(abu)$.

ومن ذلك اختزال الملازمات الأعلى رتبة $(\theta_\eta su)^3$ و$(H_\vartheta(\theta Hu),s,u)^3$ و$(a_\nu\cdot b_{\nu_1}^2,s,u)^4$ [الملازم الأخير بدلا من $(H_\vartheta su)^4$].

2. $((\vartheta\eta x),s,u)^4$: المختزل $(a_\eta su)^4$.

3. $((\vartheta\eta x)^2,s,u)^3$: المختزل $s_\vartheta^2s_\nu$ من $(\widehat{\vartheta\eta}su)^2(Hsu)$.
تتحلل الملازمات $(\vartheta\eta x)s_\eta$ و$(\vartheta\eta x)^2s_\eta$ و$((\vartheta\eta x)^2,s,u)^2$ و$\theta_\eta s_\eta$ بالطيّ مع الملازم المتحلل $a_\eta s_\eta$.

4. سلاسل الملازمات $H_\vartheta(\theta Hu)s_\eta$ و$\theta_\eta(\theta Hu)s_\eta$ و$H_\vartheta(\vartheta\eta x)s_\eta$ و$\theta_\eta(\vartheta\eta x)s_\eta$:
المختزلان $a_\eta(aHu)s_\eta$ و$a_\eta^2s_\eta$.

سلسلة الملازمات $(\theta_\eta(\vartheta\eta x),s,u)^2$: المختزل $a_\eta^2(Hsu)^2$ [باستثناء $(\theta_\eta^2(\theta Hu),s,u)^2$، الذي ينشأ بالطيّ من الحد $a_\eta^2(bsu)(Hsu)$].

5. $(H_\vartheta(\vartheta\eta x),s,u)$ و$(H_\vartheta(\vartheta\eta x),s,u)^2$: الاختزال المزدوج لـ~$a_\eta(aHu)s_\eta(abu)$ و$g_\vartheta(\theta gu)g_\nu$.

سلسلة الملازمات $(H_\vartheta(\vartheta\eta x),s,u)^3$: المختزل $((\vartheta \nu x)^2,s,u)^2s_\nu$.

6. يتحلل $(H_\vartheta(\theta Hu),s,u)$ بطيّ واحد مع الملازم المتحلل $(\theta_\nu(\vartheta \nu x),s,u)$ وفقًا لـ~\S\ 3. 2)، c)؛ أي بالاختزال المزدوج للملازم المتحلل الأدنى رتبة $(H_\vartheta su)^2$:\srcfn{*)}{تعني العلامة $\equiv$ أن جداءات ملازمات من درجات أدنى محذوفة.}
\[
\begin{aligned}
0&\equiv \theta_\nu(\vartheta\nu\nu_1)(\theta su)+\frac23\theta_\nu^2(\vartheta \nu_1 x)(\theta su)+\theta_\nu(\vartheta \nu x)(\theta su)s_{\nu_1}\\
&\equiv -\frac12H_\vartheta(\theta Hu)(\theta su)+\theta_\eta(\theta su)(Hsu)+\frac12H_\vartheta(\theta su)(Hsu)\\
&\equiv -\frac12H_\vartheta(\theta Hu)(\theta su)+\theta_\eta(\theta su)(Hsu)+\frac12[H_\vartheta]_{\widehat{su}^2}-\frac12\cdot\frac35H_\vartheta(\theta Hu)(\theta su)\\
&\equiv -\frac35H_\vartheta(\theta Hu)(\theta su).
\end{aligned}
\]
تنتج بالاختزال المزدوج، وفقًا لـ1. و5.، الصيغ الآتية:
\[
\begin{array}{@{}r@{\quad}l@{}}
&\displaystyle
\begin{aligned}
0={}&\theta_\eta(\theta su)(Hsu)^2+\frac14H_\vartheta(\theta Hu)(\theta su)^2-\frac34\theta_\eta(\theta Hu)(\theta su)(Hsu)\\
&-\frac54\theta_\eta(\theta Hu)^2(\theta su)-(asu)^3\cdot b_\eta(bHu)^2-a_\nu(asu)^3\cdot b_{\nu_1}^2,
\end{aligned}\\[0.55em]
(31.)&\displaystyle 0=H_\vartheta(\theta Hu)(\theta su)^3-7\theta_\eta(\theta Hu)(\theta su)^2(Hsu),\\[0.55em]
&\displaystyle 0=a_\nu(asu)^2b_{\nu_1}^{2}(bsu)^2-\theta_\eta(\theta su)^2(Hsu)^2+6\theta_\eta(\theta Hu)(\theta su)^2(Hsu),\\[0.55em]
&\displaystyle 0\equiv (H_\vartheta(\vartheta\eta x),s,u),\qquad
0\equiv H_\vartheta(\widehat{\vartheta\eta su})(Hsu)-4\theta_\eta^2(Hsu).
\end{array}
\]

\noindent\emph{النتائج:}

1) الملازمات $(\theta Hu)^2s_\eta$ و$H_\vartheta(\theta Hu)s_\eta$ و$\theta_\eta(\theta Hu)s_\eta$ و$H_\vartheta(\vartheta\eta x)s_\eta$ و$\theta_\eta(\vartheta\eta x)s_\eta$ و$(H_\vartheta(\vartheta\eta x),s,u)$ و$(\theta_\eta(\vartheta\eta x),s,u)^2$ مختزلات.

2) لا يدخل $s^2$ في الطيّ مع ملازمات الرتبة السادسة $(5,\lambda)$ وما إلى ذلك.

\begin{center}\textbf{\S\ 24.}\quad\emph{ملازمات الرتبة الحادية عشرة (المنظومة $s_{III}$).}\end{center}
طيّ $s$ مع ملازمات الرتبة السابعة (المنظومة III) (\S\ 12).

\noindent\emph{الملازمات غير القابلة للاختزال:}
\[
\text{A)}\quad
\begin{array}{c|c|c|c|c}
\cdot&\cdot&\cdot&((KHu)^2su)^2&\bigstar\\[0.35em]
\cdot&\cdot&(K_\eta su)^2&\cdot&\\[0.35em]
\cdot&\cdot&\cdot&\cdot&\\[0.35em]
((k\eta x)^3,s,u)&\cdot&\cdot&\cdot&\\[0.35em]
\cdot&(K_\eta(k\eta x)^3,s,u)&\cdot&\cdot&
\end{array}
\]
\[
\begin{array}{c@{\quad}|c|c}
\bigstar&((KHu)^2,s,u)^3&((KHu)^2,s,u)^4\\[0.45em]
&\cdot&(K_\eta(KHu)^2,s,u)^3
\end{array}
\]
\[
\text{B)}\quad ((j\eta x),s,u)^2,\ (H_jsu)\ \big|\ ((j\eta x),s,u)^3,
\qquad \text{C)}\quad (\D_\eta su),
\]
\[
\text{D)}\quad
\begin{array}{c|c|c}
\cdot&((aLu)^2,s,u)^2,\ ((aLu)^3,s,u)&((aLu)^2,s,u)^3\\[0.5em]
(aLu)^2s_\eta,\ (a_lsu)^2&(a_lsu)^3&(a_l(aLu)^2,s,u)^2
\end{array}
\]

\noindent\emph{الاختزالات:}

A) الاختزال أو التحلل بطيّ واحد مع الملازمات المناظرة المشتملة على الرمزين $\theta\cdot H$.

يتحلل الملازمان $(K_\eta(KHu)^2,s,u)$ و$(K_\eta(KHu)^2,s,u)^2$ وفقًا لـ~\S\ 3. 2)، a)، بالطيّ مع $K_\nu^2(Ksu)$ و$K_\nu^2(Ksu)^2$.

يتحلل $(K_\eta(KHu),s,u)^4\equiv(a_\nu^2\cdot\theta_{\nu_1},s,u)^4$ وفقًا لـ~\S\ 3. 2)، b)، من الملازم القابل للتحلل $K_\nu^2(Ksu)^3$.

B) سلسلة الملازمات $(j\eta x)s_\eta$: المختزل $(j\nu x)^2s_\nu$.

الملازم $H_j(\widehat{j\eta su})(Hsu)$: المختزل $(j\nu x)(\widehat{j\nu su})^2$.

الملازم $H_j(j\eta x)(Hsu)$: يتحلل باختزال الملازم القابل للتحلل $((j\eta x)^2,s,u)^2$، وذلك بالمختزل $(j\nu x)^2s_\nu$ (الصيغة (18.)a):
\[
0=-2(j\nu x)^2s_\nu s_{\nu_1}=[(j\eta x)^2]_{\widehat{su}^2}+H_j(j\eta x)(Hsu).
\]
C) المختزلان $\D_\nu s_\nu$ و$\D_\nu(\D su)^2$.

D) الاختزال والتحلل بطيّ واحد مع الملازمات المناظرة المشتملة على الرمزين $f\cdot H$.

E) ينتج عن (30.)c اختزال سلسلة الملازمات القابلة للتحلل $(a_\sigma su)^2$:
\[
\begin{aligned}
0&=a_\eta(Hsu)^2(as\widehat{\nu}x)^2
=a_\eta(Hsu)^2(a\widehat{\nu}\eta u)^2+2a_\eta(a\widehat{\nu}\eta u)(\widehat{\nu}\eta su)(Hsu)^2
=\frac13a_\sigma(asu)^2,\\
0&=a_\eta a_\nu(Hsu)^2(as\widehat{\nu}x)(asu)
=a_\eta(a\widehat{\nu}\eta u)^2(Hsu)^2(asu)+a_\eta(a\widehat{\nu}\eta u)(\widehat{\nu}\eta su)(Hsu)^2(asu)
=\frac13a_\sigma(asu)^3.
\end{aligned}
\]
\noindent\emph{النتيجة:} لا يدخل $s^2$ في الطيّ مع ملازمات الرتبة السابعة.


\begin{center}\textbf{\S\ 25.}\quad\emph{ملازمات الرتب الثانية عشرة والثالثة عشرة والرابعة عشرة والخامسة عشرة (المنظومة $s_{III}$).}\end{center}
طيّ $s$ مع ملازمات الرتب الثامنة والتاسعة والعاشرة والحادية عشرة (المنظومة III) (\S\S\ 13--16).

\noindent\emph{الملازمات غير القابلة للاختزال:}

\noindent الرتبة الثانية عشرة.
\[
\begin{adjustbox}{max width=0.98\linewidth}
$\displaystyle
\begin{array}{c@{\qquad}c@{\qquad}c}
\text{A)}&
\begin{array}{c|c}
((\vartheta^2\eta x)^4,s,u)&\cdot\\[0.35em]
\cdot&\theta_\eta(\vartheta^2\eta x)^3s_\vartheta
\end{array}
&
\text{B)}\quad ((aHu)H_j,s,u)^2\;\big|\;((aHu)H_j,s,u)^3
\\[1.0em]
\text{C)}&
\begin{array}{c|c|c}
\cdot&((\theta Lu)^2,s,u)^2,\ ((\theta Lu)^3,s,u)&((\theta Lu)^2,s,u)^3\\[0.35em]
(\theta_lsu)^2&\cdot&(\theta_l(\theta Lu)^2,s,u)^2
\end{array}
&
\end{array}$
\end{adjustbox}
\]
\noindent\emph{الاختزالات:} لن تُذكر بعد الآن الاختزالات الناشئة عن الطيّ المباشر مع المختزلات أو الملازمات القابلة للتحلل.

B) يتحلل كلٌّ من $((aHu)H_j,s,u)$ و$(a_\eta(aHu)H_j,s,u)$ على صورة عملية الاشتقاق والتقليص على محددات يعقوبية.

يتحلل الملازم $(a_\eta(aHu)H_j,s,u)^2$ بالاختزال المزدوج للملازم القابل للتحلل $[\theta_\nu\cdot\theta_{\nu_1}^{\prime 3}]_{\widehat{su^3}}$ (\S\ 13. C):
\begin{align*}
0&\equiv[\theta_{\nu_1}^{\prime 3}\cdot\theta_\nu]_{\widehat{su^3}}
\equiv-\frac34\theta_\nu(\theta su)^2(\theta\theta'u)\theta_{\nu_1}^{\prime 3}\\
&\equiv-\frac98(\widehat{\theta\theta'}\widehat{\eta su})(\theta\theta'u)(\theta Hu)(\theta'Hu)\theta_\eta'(\theta su)
\equiv-\frac38a_\eta(aHu)H_j(asu)^2.
\end{align*}

C) ينتج عن الصيغ (31.) اختزال الملازمين القابلين للتحلل
\[
 [L_\vartheta(\theta Lu)]_{\widehat{su^3}}\quad\text{و}\quad [\theta_l(\theta Lu)]_{\widehat{su^3}},
\]
أي الجداءين $[\theta_\nu^2\cdot H]_{\widehat{su^3}}$ و$[a_\nu^2\cdot(bHu)^2]_{\widehat{su^3}}$:
\[
\begin{array}{@{}r@{\quad}l@{}}
&\displaystyle 0\equiv\theta_\nu^2(\theta su)^2(Hsu),\qquad
0\equiv [L_\vartheta(\theta Lu)]_{\widehat{su^4}}-7[\theta_l(\theta Lu)]_{\widehat{su^4}},\\[0.55em]
&\displaystyle 0\equiv a_\nu^2(asu)^2(bHu)^2(bsu)+6\theta_l(\theta Lu)^2(\theta su)(Lsu),\\[0.55em]
(32.)&\displaystyle 0\equiv\theta_\nu^3(\theta su)(Hsu)^2
\quad\text{من }0\equiv\theta_l^2(\theta Lu)(\theta su)(Lsu)^2\\[-0.05em]
&\displaystyle\hfill(\text{المختزل }a_\eta^2(Hsu)^2).
\end{array}
\]

\noindent\emph{الملازمات غير القابلة للاختزال.}

\noindent الرتبة الثالثة عشرة.
\[
\begin{array}{rl}
\text{A)}&((KLu)^3,s,u)^2;\ ((KLu)^3,s,u)^3,\\[0.25em]
\text{B)}&(L_jsu)^2,\\[0.25em]
\text{C)}&((aH^2u)^3,s,u)^2.
\end{array}
\]
\noindent الرتبة الرابعة عشرة.
\[
\begin{array}{rl}
\text{A)}&((aLu)^2L_j,s,u)^2,\\[0.25em]
\text{B)}&((\theta H^2u)^3,s,u)^2.
\end{array}
\]
\noindent الرتبة الخامسة عشرة.
\[
 ((KH^2u)^4,s,u)^2.
\]
\noindent\emph{الاختزالات:} يتحلل كلٌّ من $(K_l(KLu)^3,s,u)^2$، و$(H_\vartheta(\theta H^2u)^3,s,u)$، و$((K,H\cdot L,u)^5,s,u)$ وفقًا لـ~\S\ 3. 2)، c)؛ أي بالنظر في آن واحد إلى الملازمات القابلة للتحلل $(K_l(KLu)^2,s,u)^3$، و$(H_\vartheta(\theta H'u)^2,s,u)^2$، و$((K,H\cdot L,u)^4,s,u)^2$.

\noindent\emph{النتيجة:} لا يدخل $s^2$ في طيٍّ مع أيّ ملازم منفرد من المنظومة III.

\begin{center}\textbf{\S\ 26.}\quad\emph{المنظومة $s_{IV}$.}\end{center}
\noindent 1) \emph{طيّ $s$ مع المنظومة I$\cdot$III.}

\noindent\emph{الاختزالات:} بعد اختزالات الفقرات السابقة $(a_\nu^2s_\nu,\ (aHu)^2(asu)s_\eta$ وما إلى ذلك)، لا تتيح الملازمات $(0,\lambda)$ و$(1,\lambda)$ و$(2,\lambda)$ عمليتي الطيّ I وIII معًا؛ ولذلك لا تدخل $f,\D,N$، وفقًا لـ~\S\ 18، في الطيّ ضمن جداءات لتكوين المنظومة $s_{IV}$.

$j$ لا يدخل في الطيّ، إذ إن عمليات الطيّ المتخالفة التدرج بالنسبة إلى $j$، بعد الاختزالات نفسها،

\[
(K_\nu(k\nu x)^3,s,u),\qquad ((\vartheta^2\eta x)^4,s,u)\quad\text{إلخ.}
\]
تقابلها عمليات الطيّ المتوافقة التدرج بالنسبة إلى $j$:

\[
K_\nu(k\nu x)^2s_k,
\qquad (\vartheta^2\eta x)^3s_\vartheta\quad\text{إلخ.}
\]

\noindent 2) \emph{طيّ $s$ مع المنظومة II$\cdot$III}
أي طيّ $s$ مع $H\cdot(1,\lambda)$ و$H\cdot(2,\lambda)$ و$H\cdot(3,\lambda)$.

\noindent\emph{الملازمات غير القابلة للاختزال:}

\[
\begin{array}{rl}
\text{A)}&(a_\nu\cdot H,s,u)^4,
((aHu)^2\cdot H,s,u)^3,
((aHu)^2\cdot H,s,u)^4,\\[0.25em]
&((aLu)^2\cdot H,s,u)^4,
((aLu)^3\cdot H,s,u)^3,
((aH^2u)^3\cdot H,s,u)^4,\\[0.45em]
\text{B)}&(a_\nu^2\cdot H,s,u)^3,
(a_\eta(aHu)^2\cdot H,s,u)^3,
(a_l(aLu)^2\cdot H,s,u)^4,\\[0.45em]
\text{C)}&(\theta_\nu\cdot H,s,u)^4,
((\theta Hu)^2\cdot H,s,u)^3,
((\theta Hu)^2\cdot H,s,u)^4,\\[0.25em]
&((\theta Lu)^2\cdot H,s,u)^4,
((\theta Lu)^3\cdot H,s,u)^3,
((\theta H^2u)^3\cdot H,s,u)^4,\\[0.45em]
\text{D)}&(\theta_\eta(\theta Hu)^2\cdot H,s,u)^3,
(\theta_l(\theta Lu)^2\cdot H,s,u)^4,\\[0.45em]
\text{E)}&((KLu)^3\cdot H,s,u)^4,
((KH^2u)^4\cdot H,s,u)^4,\\[0.45em]
\text{F)}&((j\nu x)^2\cdot H,s,u)^3,
((j\nu x)^2\cdot H,s,u)^4,
((j\eta x)\cdot H,s,u)^4,\\[0.25em]
&(H_j\cdot H,s,u)^3,
(L_j\cdot H,s,u)^4.
\end{array}
\]
\noindent\emph{الاختزالات:} لا يدخل $\nu$ في الطيّ، على نحو مماثل لـ$j$.

B) وD). يتحلل كل من $(a_\nu^2\cdot H,s,u)^4$ و$(\theta_\nu^2\cdot H,s,u)^4$ استنادًا إلى الصيغتين (27.)c) و(29.)c)، مع مراعاة $(gHu)^2=-\frac13s\cdot\sigma$:

\[
(a_\nu^2\cdot H,s,u)^4\equiv\frac13(asu)^4\cdot\sigma,
\qquad
(\theta_\nu^2\cdot H,s,u)^4\equiv-\frac16(\theta su)^4\cdot\sigma;
\]
ومن ذلك ينتج تحلل $(a_\eta(aHu)\cdot H,s,u)^4$ و$(\theta_\eta(\theta Hu)\cdot H,s,u)^4$.

يتحلل كل من $(\theta_\nu^2\cdot H,s,u)^3$ و$(\theta_\nu^3\cdot H,s,u)^3$، ومن ثم يتحلل $(\theta_\eta^2(\theta Hu)\cdot H,s,u)^4$، وفقًا لـ~\S\ 25، ملازمات الرتبة الثانية عشرة C).
\noindent 3) \emph{طيّ $s$ مع المنظومة III$\cdot$III}، أي طيّ $s$ مع ملازمات المنظومة III:
\[
(3,\lambda)\cdot(3,\lambda),\quad (3,\lambda)\cdot(2,\lambda),\quad (3,\lambda)\cdot(1,\lambda),\quad
(2,\lambda)\cdot(2,\lambda),\quad (2,\lambda)\cdot(1,\lambda),\quad
(1,\lambda)\cdot(1,\lambda).
\]
\noindent\emph{الملازمات غير القابلة للاختزال:}
\[
\begin{array}{rl}
\text{A)}&(a_\nu^2\cdot a_\nu^2,s,u)^3,
(a_\nu^2\cdot a_\nu^2,s,u)^4,
(a_\nu^2\cdot a_\eta(aHu)^2,s,u)^3,\\[0.45em]
\text{B)}&(a_\nu\cdot H_j,s,u)^4,
((aHu)^2\cdot H_j,s,u)^3,
((aLu)^2\cdot H_j,s,u)^4,\\[0.25em]
&((aLu)^3\cdot H_j,s,u)^2,
((aH^2u)^3\cdot H_j,s,u)^3.
\end{array}
\]
\noindent\emph{الاختزالات:}

ينبغي أن تُعَدّ جداءات II$\cdot$III، عند تساوي الرتبة في الرمزين $\nu$ و$f$، أعلى من جداءات III$\cdot$III.

تنشأ الاختزالات جزئيًا من علاقات بين الجداءات (السيزيجيات)، وجزئيًا من إحلال الملازمات القابلة للتحلل الموافقة محل الجداءات واختزال طيّ هذه الملازمات مع $s$. (تُحذَف دائمًا الجداءات التي تحتوي على ملازمات متخالفة التدرج، لأنها تتحول إلى جداءات.)

A) الجداءات تربيعية في $f$، والرموز $\nu$ من رتبة اعتباطية. وفقًا لـ~\S\ 11، الصيغ (12.)، وبحساب مماثل، تصح الصيغ الآتية:
\begin{align*}
1)\quad 0&=a_\nu\cdot b_{\nu_1}+\frac12f\cdot(aHu)^2-\frac14\theta\cdot H+\frac12u_xH_\vartheta-\frac14u_x^2\cdot H_\vartheta^2,\\
2)\quad 0&=a_\nu\cdot b_{\nu_1}^2+f\cdot a_\eta(aHu)^2-\theta_\eta-\frac12H_\vartheta,\\
3)\quad 0&=a_\nu^2\cdot b_{\nu_1}^2-H_\vartheta^2+2\theta_\eta^2.
\end{align*}
ومن ذلك ينتج:
\begin{align*}
4)\quad 0&=a_\nu^2\cdot b_{\nu_1}^2(j\nu_1x)\quad(\text{من }3)),\\
5)\quad L_\vartheta(\theta Lu)&=-a_\nu\cdot b_\eta(bHu)^2+\frac34a_\nu^2\cdot(bHu)^2+\frac34\theta_\nu^2\cdot H\quad(\text{من }2)),\\
6)\quad L_\vartheta(\theta Lu)&=-6a_\nu\cdot b_\eta(bHu)^2+2a_\nu^2\cdot(bHu)^2-\frac12\theta_\nu^2\cdot H\quad(\text{من }1)),\\
7)\quad 0&=a_\nu\cdot(bHu)^2+f\cdot(aLu)^3+\theta_\nu\cdot H\quad(\text{من }1)),\\
\end{align*}
\[
\left.
\begin{aligned}
8)\quad 0&=a_\nu\cdot(bLu)^3+\frac34(\theta Hu)^2\cdot H\\
9)\quad 0&=(aHu)^2\cdot(bHu)^2+2(\theta Hu)^2\cdot H
\end{aligned}
\right\}\quad\text{(من }7)\text{)،}
\]
\begin{align*}
10)\quad 0&=a_\eta(aHu)^2\cdot b_\eta(bHu)^2\quad(\text{من }5)\text{ و}6)),\\
11)\quad 0&= [a_\nu^2\cdot b_\eta(bHu)]_{\widehat{su}^4}\quad(\text{من الصيغة }(30.)b)).
\end{align*}
إن اختزالات
\[
(H_\vartheta su)^4,
\qquad (L_\vartheta(\theta Lu),s,u)^3,
\qquad (L_\vartheta(\theta Lu),s,u)^4
\]
(الواردة في الصيغتين (31.) و(32.)) تعطي، مع الصيغ من 1) إلى 11)، اختزال جميع عمليات طيّ $s$ مع الجداءات التربيعية في $f$ وذات الرتبة الاعتباطية في $\nu$.

B) جداءات كلِّ اثنين من الرموز $f,\theta,j$؛ و$\nu$ من رتبة اعتباطية. وفقًا للصيغ (17.) و(18.) و(20.)، وبحساب مباشر، تصح العلاقات الآتية:
\[
\left.
\begin{aligned}
1)\quad 0&=a_\nu\cdot\theta_{\nu_1}+\frac14\theta\cdot(aHu)^2+\frac14f\cdot(\theta Hu)^2,\\
2)\quad 0&=\theta_\nu\cdot\theta_{\nu_1}+\frac12\theta\cdot(\theta Hu)^2
\end{aligned}
\right\}\quad\text{(بحساب مباشر على غرار }A\ 1)\text{).}
\]
ومن ذلك ينتج:
\[
\left.
\begin{aligned}
3)\quad 0&=3a_\nu\cdot(\theta Hu)^2+\theta\cdot(aLu)^3+2f\cdot(\theta Lu)^3\\
4)\quad 0&=3\theta_\nu\cdot(aHu)^2+f\cdot(\theta Lu)^3+2\theta\cdot(aLu)^3
\end{aligned}
\right\}\quad\text{(من }1)\text{ و}A\ 6)\text{)،}
\]
\begin{align*}
5)\quad 0&=a_\nu\cdot(\theta Lu)^3,\quad 0=\theta_\nu\cdot(aLu)^3,\quad 0=(aHu)^2\cdot(\theta Hu)^2
\quad\text{(من }3),\ 4)\text{ و}A\ 8)\text{)،}\\
6)\quad 0&=a_\nu\cdot\theta_\nu^2+\theta_\nu\cdot a_\nu^2+f\cdot\theta_\eta(\theta Hu)^2+\theta\cdot a_\eta(aHu)^2\quad(\text{الصيغة }(17.)),\\
7)\quad 0&=\theta_\nu\cdot\theta_\nu^2+\theta\cdot\theta_\eta(\theta Hu)^2+\frac16f\cdot H_j\\
&\qquad [\text{اشتقاق مماثل لاشتقاق }6)\text{ من }\theta_\nu^2(\theta\theta'\widehat{\nu}_1x)],\\
8)\quad 0&=a_\nu\cdot(j\nu x)^2+f\cdot H_j\quad(\text{من }6)),\\
9)\quad 0&=(aHu)^2\cdot(j\nu x)^2-2a_\nu\cdot H_j\quad(\text{من }8)),\\
10)\quad 0&=\theta_\nu\cdot(j\nu x)^2,\quad 0=\theta\cdot H_j\quad(\text{من }7)\text{ و}8)),
\end{align*}
\[
\left.
\begin{aligned}
11)\quad 0&=a_\nu\cdot\theta_\nu^3-\frac34(j\eta x)^2,\\
12)\quad 0&=a_\nu^2\cdot\theta_\nu^2-\frac13(j\eta x)^2-\frac13(\D Hu)^2
\end{aligned}
\right\}\quad\text{(الصيغة }(18.)\text{)،}
\]
\begin{align*}
13)\quad 0&=a_\nu^2\cdot\theta_\nu^3-\frac12a_\sigma\quad(\text{الصيغة }(20.)),\\
14)\quad 0&=\theta_\nu\cdot\theta_\nu^3,
&0&=a_\eta H_j\quad(\S\ 13. C)).
\end{align*}
تُتابَع الصيغ إلى الحدّ الذي تصبح فيه الجداءات قابلة للاستقطاب؛ أي إنه لا ينشأ بالطيّ إلا جداء جديد \emph{واحد}، وذلك بأن
\begin{center}
\begin{tabular}{rl}
a)&يصبح طيّ الجداء في نفسه قابلاً للاختزال،\\
b)&تصبح سائر الجداءات الناشئة بالطيّ قابلة للاختزال
\end{tabular}
\end{center}
أو تفضي إلى ملازمات متخالفة التدرج (قارن \S\ 3. 2)، a) وb)).

تعطي الصيغ من 1) إلى 5) اختزال جميع عمليات طيّ $s$ مع الجداءات التي تنشأ من $a_\nu\cdot\theta_\nu$ بالطيّ،

والصيغ من 6) إلى 10) اختزال الجداءات التي تنشأ من $a_\nu\cdot\theta_\nu^2$ بالطيّ.

وفقًا لـ~\S\ 24 A)، تتحلل، مع مراعاة الصيغ من 6) إلى 10)، عمليات طيّ $s$ مع الجداءات التي تنشأ من $a_\nu^2\cdot\theta_\nu$.

فنحصل على:
\begin{align*}
0&\equiv a_\nu^2(asu)^2\theta_{\nu_1}(\theta su)^2,
&0&\equiv a_\nu^2(asu)\theta_{\nu_1}(\theta su)^3-K_\eta(KHu)^2(Ksu)^3,\\
0&\equiv a_\nu^2(asu)^2\theta_{\nu_1}(\theta su),
&0&\equiv a_\nu^2(asu)\theta_{\nu_1}(\theta su)^2,\\
0&\equiv a_\nu^2(asu)\theta_{\nu_1}(\theta su).
\end{align*}
المختزلات:
\[
 ((j\eta x)^2,s,u)^3
 \ [\text{من }(\vartheta\eta\widehat{su})^2(Hsu).\ \S\ 23.\ C\ 3)],
\]
\[
(H_j(j\eta x),s,u)^2\ [\S\ 24,\ B],
\quad ((\D Hu)^2,s,u)^2\ [\S\ 24,\ C],
\]
\[
(a_\sigma su)^2\ [\S\ 24,\ E]
\]
تعطي هذه المختزلات، مع الصيغ من 11) إلى 14)، اختزال جميع الجداءات الناشئة من
$a_\nu\cdot\theta_{\nu_1}^3$، $a_\nu^2\cdot\theta_{\nu_1}^2$، و$a_\nu^2\cdot\theta_{\nu_1}^3$.

يمكن إجمالُ حساباتنا السابقة في
\emph{النتيجة الآتية}:

تتألّف المنظومة الكاملة نسبيًّا بترديد $\modu{(\R,t)}$ من الحدوديات المتجانسة والمنظومات الجزئية الآتية، وعددها \srcLTR{331}، وسنوردها بعد ذلك مرتّبةً أيضًا في جدول:
\[
\left.
\begin{array}{lr}
u_x&1\ \srcRTL{حدودية متجانسة}\\
\srcRTL{المنظومة}~\srcLTR{I}&6\ \srcRTL{حدوديات متجانسة}\\
\srcRTL{المنظومة}~\srcLTR{II}&5\ \srcRTL{حدوديات متجانسة}\\
\srcRTL{المنظومة}~\srcLTR{III}&117\ \srcRTL{حدوديات متجانسة}
\end{array}
\right\}
\quad
\begin{gathered}
\srcRTL{منظومة كاملة نسبيًّا بترديد}~(s,\varrho,t)\\
[\srcRTL{انظر الجدول}~\srcLTR{I}].
\end{gathered}
\]
\[
\left.
\begin{array}{lr}
s&1\ \srcRTL{حدودية متجانسة}\\
\srcRTL{المنظومة}~s_I&35\ \srcRTL{حدوديات متجانسة}\\
\srcRTL{المنظومة}~s_{II}&9\ \srcRTL{حدوديات متجانسة}\\
\srcRTL{المنظومة}~s_{III}&125\ \srcRTL{حدوديات متجانسة}\\
\srcRTL{المنظومة}~s_{IV}&32\ \srcRTL{حدوديات متجانسة}
\end{array}
\right\}
\quad [\srcRTL{انظر الجدول}~\srcLTR{II}].
\]

\clearpage
\begin{NoetherWideTable}{297mm}{420mm}
\thispagestyle{plain}
\begin{otherlanguage}{arabic}
\begin{center}\bfseries \srcRTL{الجدول}~\srcLTR{I}~\srcRTL{(انظر ص.}~\srcLTR{90}\srcRTL{).}\end{center}
\end{otherlanguage}
\vspace{-0.6em}
\begingroup
\setlength{\tabcolsep}{2pt}
\renewcommand{\arraystretch}{1.35}
\scriptsize
\begin{center}
\begin{adjustbox}{max width=0.98\linewidth,max totalheight=0.86\textheight,center}
\begin{tabular}{|c|c|c|c|c|c|c|c|c|c|c|c|c|c|c|c|c|c|c|c|c|c|c|}
\hline
{} & 0 & 1 & 2 & 3 & 4 & 5 & 6 & 7 & 8 & 9 & 10 & 11 & 12 & 13 & 14 & 15 & 16 & 17 & 18 & 19 & 20 & 21 \\ \hline
0 & \mcell{i} &  &  &  & \mcell{f} &  & \mcell{\A} &  & \mcell{K_{\nu}(k\nu x)^3} &  & \mcell{K_{\eta}(k\eta x)^3} &  &  &  & \mcell{(\vartheta^2\eta x)^4} & \mcell{H_k(k\cdot\vartheta,\eta x)^4} &  & \mcell{\theta_l(\vartheta\cdot k,lx)^5} &  &  &  & \mcell{(\vartheta^3lx)^6} \\ \hline
1 &  & \mcell{u_x} &  &  &  & \mcell{\theta_\nu(\vartheta\nu x)^2} &  & \mcell{\theta_\eta(\vartheta\eta x)^2} & \mcell{N} & \mcell{(k\nu x)^3} & \mcell{\theta_\nu(\vartheta^2\nu x)^3} & \mcell{(k\eta x)^3} & \mcell{H_{\vartheta}(\vartheta^2\eta x)^3\\ \theta_\eta(\vartheta^2\eta x)^3} &  & \mcell{\theta_l(\vartheta^2lx)^4} &  & \mcell{(\vartheta k,\eta,x)^4} &  & \mcell{(\vartheta k,l,x)^5} &  &  &  \\ \hline
2 &  &  & \mcell{a_\nu^2} &  & \mcell{\theta\\ a_\eta^2} &  & \mcell{(\vartheta\nu x)^2} & \mcell{K_\nu(k\nu x)^2} & \mcell{(\vartheta\eta x)^2} & \mcell{H_k(k\eta x)^2\\K_\eta(k\eta x)^2} &  & \mcell{(\vartheta^2\nu x)^3\\ K_l(klx)^3} &  & \mcell{(\vartheta^2\eta x)^3} &  & \mcell{(\vartheta^2lx)^4} &  &  &  &  &  &  \\ \hline
3 &  & \mcell{\theta_\nu^3} &  & \mcell{a_\nu} & \mcell{\theta_\nu(\vartheta\nu x)} & \mcell{\A_\nu\\a_\eta} & \mcell{K\\H_{\vartheta}(\vartheta\eta x)\\\theta_\eta(\vartheta\eta x)} & \mcell{\A_\eta} & \mcell{(k\nu x)^2\\\theta_l(\vartheta lx)^2} &  & \mcell{(k\eta x)^2} &  & \mcell{(klx)^3} &  &  &  &  &  &  &  &  &  \\ \hline
4 & \mcell{\nu} &  & \mcell{H\\\theta_\nu^2} & \mcell{(j\nu x)^3\\a_\eta(aHu)} & \mcell{\theta_\eta^2} & \mcell{(\vartheta\nu x)\\a_l^2} &  & \mcell{N_\nu\\(\vartheta\eta x)} &  & \mcell{H_n\\N_\eta\\(\vartheta lx)^2} &  &  &  &  &  &  &  &  &  &  &  &  \\ \hline
5 &  & \mcell{a_\eta(aHu)^2} & \mcell{\theta_\eta^2(\theta Hu)} & \mcell{\theta_\nu} & \mcell{K_\nu^2\\(aHu)} & \mcell{\theta_\eta} & \mcell{K_\eta^2\\(\A Hu)\\a_l} &  & \mcell{\A_l} &  &  &  &  &  &  &  &  &  &  &  &  &  \\ \hline
6 & \mcell{j\\\sigma} &  & \mcell{(j\nu x)^2\\(aHu)^2} & \mcell{L\\a_\nu^2(j\nu x)\\H_{\vartheta}(\theta Hu)\\\theta_\eta(\theta Hu)} &  & \mcell{K_\nu\\\theta_l^2} &  & \mcell{K_\eta} &  &  &  &  &  &  &  &  &  &  &  &  &  &  \\ \hline
7 &  & \mcell{\theta_\eta(\theta Hu)^2} & \mcell{H_j(j\eta x)\\a_l(aLu)^2} &  & \mcell{a_\nu(j\nu x)\\(\theta Hu)} &  & \mcell{\A_\nu(j\nu x)\\\theta_l} & \mcell{K_l^2} &  &  &  &  &  &  &  &  &  &  &  &  &  &  \\ \hline
8 & \mcell{H_j^2\\a_l(aLu)^3} & \mcell{(\nu\sigma x)\\(j\nu x)} & \mcell{(\theta Hu)^2} & \mcell{K_\eta(KHu)^2\\(j\eta x)\\(aLu)^2} &  &  &  &  &  &  &  &  &  &  &  &  &  &  &  &  &  &  \\ \hline
9 &  & \mcell{H_j\\(aLu)^3} & \mcell{a_\eta(aHu)H_j\\L_{\vartheta}(\theta Lu)^2\\\theta_l(\theta Lu)^2} &  & \mcell{(KHu)} &  &  &  &  &  &  &  &  &  &  &  &  &  &  &  &  &  \\ \hline
10 & \mcell{\theta_l(\theta Lu)^3} & \mcell{L_j^2\\(j\sigma x)\\a_\eta(aH^2u)^3} &  & \mcell{H_j(aHu)\\(\theta Lu)^2} &  &  &  &  &  &  &  &  &  &  &  &  &  &  &  &  &  &  \\ \hline
11 &  & \mcell{(\theta Lu)^3} & \mcell{K_l(KLu)^3\\L_j\\(aH^2u)^3} &  &  &  &  &  &  &  &  &  &  &  &  &  &  &  &  &  &  &  \\ \hline
12 & \mcell{(aH^2u)^4} & \mcell{a_l(aLu)^2L_j\\H_{\vartheta}(\theta H^2u)^3\\\theta_\eta(\theta H^2u)^3} &  & \mcell{(KLu)^3} &  &  &  &  &  &  &  &  &  &  &  &  &  &  &  &  &  &  \\ \hline
13 & \mcell{(\nu\sigma j)} &  & \mcell{(aLu)^2L_j\\(\theta H^2u)^3} &  &  &  &  &  &  &  &  &  &  &  &  &  &  &  &  &  &  &  \\ \hline
14 & \mcell{(\theta H^2u)^4} & \mcell{K_\eta(KH^2u)^4\\(a,H\!\cdot L,u)^4} &  &  &  &  &  &  &  &  &  &  &  &  &  &  &  &  &  &  &  &  \\ \hline
15 & \mcell{L_\vartheta(\theta,H\!\cdot L,u)^4} &  & \mcell{(KH^2u)^4} &  &  &  &  &  &  &  &  &  &  &  &  &  &  &  &  &  &  &  \\ \hline
16 &  & \mcell{(\theta,H\!\cdot L,u)^4} &  &  &  &  &  &  &  &  &  &  &  &  &  &  &  &  &  &  &  &  \\ \hline
17 & \mcell{K_\eta(K,H\!\cdot L,u)^5} &  &  &  &  &  &  &  &  &  &  &  &  &  &  &  &  &  &  &  &  &  \\ \hline
18 &  & \mcell{(K,H\!\cdot L,u)^5} &  &  &  &  &  &  &  &  &  &  &  &  &  &  &  &  &  &  &  &  \\ \hline
19 &  &  &  &  &  &  &  &  &  &  &  &  &  &  &  &  &  &  &  &  &  &  \\ \hline
20 &  &  &  &  &  &  &  &  &  &  &  &  &  &  &  &  &  &  &  &  &  &  \\ \hline
21 & \mcell{(KH^3u)^6} &  &  &  &  &  &  &  &  &  &  &  &  &  &  &  &  &  &  &  &  &  \\ \hline
\end{tabular}
\end{adjustbox}
\end{center}
\vspace{-0.2em}
\begin{otherlanguage}{arabic}
\begin{center}\footnotesize (\srcRTL{تدلّ الأعداد الواقعة في الصفّ الأفقي الأعلى على الدرجة في متغيّرات}~$x$\srcRTL{، وتدلّ الأعداد الواقعة في العمود الرأسي الأول على الدرجة في متغيّرات}~$u$\srcRTL{.})\end{center}
\end{otherlanguage}
\endgroup
\end{NoetherWideTable}
\clearpage
\begin{NoetherWideTable}{420mm}{594mm}
\thispagestyle{plain}
\begin{otherlanguage}{arabic}
\begin{center}\bfseries \srcRTL{الجدول}~\srcLTR{II}~\srcRTL{(انظر ص.}~\srcLTR{90}\srcRTL{).}\end{center}
\end{otherlanguage}
\vspace{-0.6em}
\begingroup
\setlength{\tabcolsep}{2pt}
\renewcommand{\arraystretch}{1.35}
\scriptsize
\begin{center}
\begin{adjustbox}{max width=0.98\linewidth,max totalheight=0.86\textheight,center}
\begin{tabular}{|c|c|c|c|c|c|c|c|c|c|c|c|c|c|c|c|c|c|}
\hline
{} & 0 & 1 & 2 & 3 & 4 & 5 & 6 & 7 & 8 & 9 & 10 & 11 & 12 & 13 & 14 & 15 & 16 \\ \hline
0 & \mcell{J} &  &  &  & \mcell{s} &  & \mcell{s_\vartheta^2} &  &  &  &  & \mcell{s_n} & \mcell{(k\nu x)^3s_\nu} & \mcell{(\vartheta^3\nu x)^3s_\nu s_\vartheta\\\theta_\nu(\vartheta^2\nu x)^3s_\vartheta} &  & \mcell{\theta_\eta(\vartheta^2\eta x)^3s_\vartheta} &  \\ \hline
1 &  &  &  &  &  &  & \mcell{(asu)} & \mcell{s_\vartheta} & \mcell{s_k^2\\(\Delta su)} & \mcell{s_n(Nsu)\\(\vartheta\nu x)^2s_\nu} & \mcell{((k\nu x)^3,s,u)s_\nu\\(K_\nu(k\nu x)^3,s,u)} &  & \mcell{(K_\eta(k\eta x)^3,s,u)} &  & \mcell{(\vartheta^2\nu x)^3s_\nu} &  & \mcell{((\vartheta^2\eta x)^4,s,u)} \\ \hline
2 &  &  &  &  & \mcell{(asu)^2} & \mcell{s_\vartheta(\theta su)} & \mcell{(\Delta su)^2\\a_\nu s_\nu} & \mcell{((\vartheta\nu x)^2,s,u)s_\nu\\(\theta_\nu(\vartheta\nu x)^2,s,u)} &  & \mcell{s_k\\(\theta_\eta(\vartheta\eta x)^2,s,u)} &  & \mcell{(k\nu x)^2s_\nu\\((k\nu x)^3,s,u)} &  & \mcell{((k\eta x)^3,s,u)} &  &  &  \\ \hline
3 &  &  & \mcell{(asu)^3} & \mcell{s_\vartheta(\theta su)^2\\s_\nu} & \mcell{a_\nu s_\nu(asu)\\a_\nu^2(asu)} & \mcell{s_\eta} & \mcell{(\theta su)\\(a_\eta^2su)} & \mcell{s_k(Ksu)} & \mcell{(Nsu)^2\\(\vartheta\nu x)s_\nu\\((\vartheta\nu x)^2,s,u)} & \mcell{((k\nu x)^3,s,u)^2} & \mcell{((\vartheta\eta x)^2,s,u)} &  &  &  &  &  &  \\ \hline
4 & \mcell{(asu)^4} & \mcell{s_\vartheta(\theta su)^3} & \mcell{a_\nu^2(asu)^2} & \mcell{(\theta_\nu^3su)} & \mcell{(\theta su)^2} & \mcell{s_k(Ksu)^2\\a_\nu(asu)\\s_\nu(asu)} & \mcell{((\vartheta\nu x)^2,s,u)^2\\\theta_\nu s_\nu} & \mcell{s_\nu(\Delta su)\\(\Delta_\nu su)\\(aHu)s_\eta\\(a_\eta su)} &  & \mcell{(\Delta_\eta su)} &  &  &  &  &  &  &  \\ \hline
5 &  &  & \mcell{(\theta su)^3} & \mcell{s_k(Ksu)^3,s_j,\\s_\vartheta(\theta s^\prime u)^4\\s_\nu(asu)^2\\a_\nu(asu)^2} & \mcell{(Hsu)\\((\vartheta\nu x)^2,s,u)^3\\(\theta_\nu(\vartheta\nu x),s,u)^2\\(\theta_\nu^2su)} & \mcell{((j\nu x)^3,s,u)\\(aHu)^3s_\eta\\(a_\eta su)^2} & \mcell{(Ksu)^2\\s_l\\(\theta_\eta^2su)} &  & \mcell{((k\nu x)^2,s,u)^2\\K_\nu s_\nu} &  &  &  &  &  &  &  &  \\ \hline
6 & \mcell{(\theta su)^4} & \mcell{s_\nu(asu)^3\\a_\nu(asu)^3} & \mcell{(Hsu)^2\\(\theta_\nu(\vartheta\nu x),s,u)^3\\(\theta_\nu^2su)^2} & \mcell{(a_\eta su)^3\\(a_\eta(aHu),s,u)^2\\(a_\eta(aHu)^2,s,u)} & \mcell{(Ksu)^3} & \mcell{s_j(asu)\\((\vartheta\nu x),s,u)^2\\(\theta_\nu su)} & \mcell{((k\nu x)^2,s,u)^3} & \mcell{s_j(\Delta su)\\a_\nu(j\nu x)s_\nu\\((\vartheta\eta x),s,u)^2\\(\theta Hu)s_\eta\\(\theta_\eta su)} &  &  &  &  &  &  &  &  &  \\ \hline
7 & \mcell{(\theta_\nu(\vartheta\nu x),s,u)^4} & \mcell{(a_\eta(aHu),s,u)^3} & \mcell{(Ksu)^4\\(\theta_\eta^2(\theta Hu),s,u)^2\\(a_\nu^2\cdot a_\nu^2,s,u)^3} & \mcell{s_j(asu)^2\\s_k(Ks^2u)^5\\((\vartheta\nu x),s,u)^3\\(\theta_\nu su)^2} & \mcell{(j\nu x)s_\nu\\((j\nu x)^2,s,u)\\((aHu),s,u)^2\\((aHu)^2,s,u)} & \mcell{((\vartheta\eta x),s,u)^3\\(\theta_\eta su)^2} & \mcell{(aLu)^2s_l\\(a_lsu)^2} &  &  &  &  &  &  &  &  &  &  \\ \hline

8 & \mcell{(a_\nu^2\cdot a_\nu^2,s,u)^4} & \mcell{s_j(asu)^3\\s_k(Ks^2u)^6\\((\vartheta\nu x),s,u)^4\\(\theta_\nu su)^3} & \mcell{((j\nu x)^2,s,u)^2\\((aHu),s,u)^3\\((aHu)^2,s,u)^2} & \mcell{(Lsu)^2\\(a_\nu^2(j\nu x),s,u)^2\\(H_\vartheta(\theta Hu),s,u)^2\\(\theta_\eta(\theta Hu),s,u)^2\\(\theta_\eta(\theta Hu)^2,s,u)} & \mcell{(a_lsu)^3} & \mcell{(K_\nu su)^2} &  & \mcell{(K_\eta su)^2} &  &  &  &  &  &  &  &  &  \\ \hline
9 & \mcell{(K_\nu^2su)^4\\((aHu),s,u)^4} & \mcell{(Lsu)^3\\(a_\nu^2(j\nu x),s,u)^3\\(\theta_\eta su)^4\\(\theta_\eta(\theta Hu),s,u)^3} & \mcell{(Ks^2u)^6\\(a_l(aLu)^2,s,u)^2\\(a_\nu^2\cdot H,s,u)^3} & \mcell{(K_\nu su)^3} & \mcell{(a_\nu(j\nu x),s,u)^2\\((\theta Hu),s,u)^2\\((\theta Hu)^2,s,u)} &  & \mcell{(\theta_lsu)^3} &  &  &  &  &  &  &  &  &  &  \\ \hline
10 &  & \mcell{(K_\nu su)^4\\(a_\nu^2\cdot a_\eta(aHu)^2,s,u)^3} & \mcell{(a_\nu(j\nu x),s,u)^3\\((\theta Hu),s,u)^3\\((\theta Hu),s,u)^2} & \mcell{((j\eta x),s,u)^2\\(H_jsu)\\((aLu)^2,s,u)^2\\((aLu)^3,s,u)} &  &  &  &  &  &  &  &  &  &  &  &  &  \\ \hline
11 & \mcell{(a_\nu(j\nu x),s,u)^4\\((\theta Hu),s,u)^4} & \mcell{(K_\eta(KHu)^2,s,u)^3\\((j\eta x),s,u)^3\\((aLu)^2,s,u)^3\\(a_\nu\cdot H,s,u)^4} & \mcell{(H^2su)^3\\(\theta_l(\theta Lu)^2,s,u)^2} &  & \mcell{((KHu)^2,s,u)^2} &  &  &  &  &  &  &  &  &  &  &  &  \\ \hline
12 &  & \mcell{(a_\eta(aHu)^2\cdot H,s,u)^3} & \mcell{((KHu)^2,s,u)^3} & \mcell{((aHu)H_j,s,u)^2\\((\theta Lu)^2,s,u)^2\\((\theta Lu)^3,s,u)} &  &  &  &  &  &  &  &  &  &  &  &  &  \\ \hline
13 & \mcell{((KHu)^2,s,u)^4} & \mcell{((aHu)H_j,s,u)^3\\((\theta Lu)^2,s,u)^3\\(\theta_\nu\cdot H,s,u)^4} & \mcell{(L_jsu)^2\\((aH^2u)^3,s,u)^2\\((j\nu x)^2\cdot H,s,u)^3\\((aHu)^2\cdot H,s,u)^3} &  &  &  &  &  &  &  &  &  &  &  &  &  &  \\ \hline
14 & \mcell{((j\nu x)^2\cdot H,s,u)^4\\((aHu)^2\cdot H,s,u)^4} & \mcell{(\theta_\eta(\theta Hu)^2\cdot H,s,u)^3} &  & \mcell{((KLu)^3,s,u)^2} &  &  &  &  &  &  &  &  &  &  &  &  &  \\ \hline
15 & \mcell{(a_l(aLu)^2\cdot H,s,u)^4} & \mcell{((KLu)^3,s,u)^3} & \mcell{((aLu)^2L_j,s,u)^2\\((\theta H^2u)^3,s,u)^2\\((\theta Hu)^2\cdot H,s,u)^3} &  &  &  &  &  &  &  &  &  &  &  &  &  &  \\ \hline
16 & \mcell{((\theta Hu)^2\cdot H,s,u)^4\\(a_\nu\cdot H_j,s,u)^4} & \mcell{((j\eta x)\cdot H,s,u)^4\\(H_j\cdot H,s,u)^3\\((aLu)^2\cdot H,s,u)^4\\((aLu)^3\cdot H,s,u)^3} &  &  &  &  &  &  &  &  &  &  &  &  &  &  &  \\ \hline
17 & \mcell{(\theta_l(\theta Lu)^2\cdot H,s,u)^4} &  & \mcell{((KH^2u)^4,s,u)^3} &  &  &  &  &  &  &  &  &  &  &  &  &  &  \\ \hline
18 &  & \mcell{((\theta Lu)^2\cdot H,s,u)^4\\((\theta Lu)^3\cdot H,s,u)^3\\((aHu)^2\cdot H_j,s,u)^3} &  &  &  &  &  &  &  &  &  &  &  &  &  &  &  \\ \hline
19 & \mcell{((aH^2u)^3\cdot H,s,u)^4\\(L_j\cdot H,s,u)^4} &  &  &  &  &  &  &  &  &  &  &  &  &  &  &  &  \\ \hline
20 &  & \mcell{((KLu)^3\cdot H,s,u)^4} & \mcell{((aLu)^3\cdot H_j,s,u)^2} &  &  &  &  &  &  &  &  &  &  &  &  &  &  \\ \hline
21 & \mcell{((\theta H^2u)^3\cdot H,s,u)^4\\((aLu)^2\cdot H_j,s,u)^4} &  &  &  &  &  &  &  &  &  &  &  &  &  &  &  &  \\ \hline
22 &  &  &  &  &  &  &  &  &  &  &  &  &  &  &  &  &  \\ \hline
23 & \mcell{((KH^2u)^4\cdot H,s,u)^4} & \mcell{((aH^2u)^3\cdot H,s,u)^3} &  &  &  &  &  &  &  &  &  &  &  &  &  &  &  \\ \hline
\end{tabular}
\end{adjustbox}
\end{center}
\vspace{-0.2em}
\begin{otherlanguage}{arabic}
\begin{center}\footnotesize (\srcRTL{تدلّ الأعداد الواقعة في الصفّ الأفقي الأعلى على الدرجة في متغيّرات}~$x$\srcRTL{، وتدلّ الأعداد الواقعة في العمود الرأسي الأول على الدرجة في متغيّرات}~$u$\srcRTL{.})\end{center}
\end{otherlanguage}
\endgroup
\end{NoetherWideTable}

\setcounter{footnote}{0}


\begin{center}
{\Large\bfseries \srcLTR{3.}~\srcRTL{في نظرية لامتغيرات الحدوديات المتجانسة ذات}~$n$~\srcRTL{متغيرات.}\srcfn{1)}{محاضرة أُلقيت في اجتماع علماء الطبيعة في سالزبورغ سنة \srcLTR{1909}.}}\par
\vspace{1.2em}
\srcRTL{بقلم \srcspaced{إيمي نوتر} في إرلنغن.}\par
\medskip
\srcLTR{Jahresbericht d. DMV 19 (1910), S. 101--104}
\end{center}

\vspace{1em}
في نظرية اللامتغيرات الإسقاطية للحدوديات المتجانسة ذات $n$ متغيرات حُلّت المسألة الرئيسة: فقد ثبت تناهي منظومة الملازمات. غير أن طائفة من المسائل الأخرى لم تُبحث، وهي التي تتعلق بطرائق دراسة العلاقات بين الحدوديات. وأود أن أعرض بإيجاز النتائج التي حصلت عليها في شأن هذه المسائل؛ غير أنني، توخيًا للوضوح، سأوجز أولًا صيغ المسائل في المجال الثلاثي، حيث هي معروفة.

وأعني أولًا المبرهنتين الأساسيتين للطريقة الرمزية: أولاهما أن جميع التكوينات اللامتغيرة يمكن تمثيلها تمثيلًا رمزيًا؛ وثانيتهما أن جميع العلاقات القائمة بين اللامتغيرات تُستخرج بالتطبيق المتعاقب لعدد قليل من المتطابقات الرمزية؛ ومن ثم تعطي الطريقة الرمزية جميع العلاقات من غير أن تتجاوز نطاق اللامتغيرات.\srcfn{1)}{\srcLTR{Study:} \emph{\srcLTR{Methoden zur Theorie der ternären Formen}}\srcLTR{. II. \S\ 6. (Leipzig, Teubner 1889.)}} وبناء على ذلك تأتي عمليات توليد الحدوديات المتجانسة: عملية الطي الرمزية، وعمليات الاشتقاق غير الرمزية، ونظرية التوسعات المتسلسلة، وما شابه ذلك. وفي هذا السياق أود أيضًا أن أذكر مبرهنة أثبتها في أطروحتي\srcfn{2)}{\emph{\srcLTR{Über die Bildung des Formensystems der ternären biquadratischen Form}}\srcLTR{. \S\ 1. (Crelle's Journal Bd. 134. 1908.)}}، مفادها أن جميع عمليات الطي يمكن توليدها بالتطبيق المتعاقب لعمليتَي الطي الممكنتَين اللتين تُسمَّيان \glqq{}متوافقتَي التدرج\grqq{}؛ وأعني بهما، في حاصل ضرب رمزي $s_x t_x u_\sigma u_\tau$، عمليتي الطي $(stu)$ و$(\sigma\tau x)$. وعلى هذه المبرهنة يمكن بناء نظرية المختزلات واختزال منظومات الملازمات، على غرار المجال الثنائي، كما بيّنت في أطروحتي.

وإذا انتقلنا الآن إلى المجال ذي $n$ متغيرات، لم تعد حتى المبرهنة الأولى للطريقة الرمزية تصح إلا على نحو مشروط. فمن المعلوم أنه تظهر، حتى في حالة الحدوديات المتجانسة رباعية المتغيرات، إلى جانب إحداثيات النقاط $x$ وإحداثيات المستويات $u$، إحداثيات المستقيمات $p_{ik}$، وهي محددات صغرى مكوّنة من صفين من إحداثيات النقاط. وبالمثل، لدينا في المجال ذي $n$ متغيرات صفوف المتغيرات $p_\rho$، التي تكون عناصرها محددات صغرى مأخوذة من $\rho$ من صفوف $x$، وصفوف الرموز $S^\rho$، التي تتصرف عناصرها كمحددات صغرى مأخوذة من $\rho$ من صفوف $s$ المتوافقة التدرج مع $u$. ومن المعلوم أن التمثيل الرمزي بواسطة المحددات لا يصح الآن إلا إذا فُككت صفوف الرموز $S^\rho$ إلى صفوف $s$، ثم وُزعت هذه الصفوف على محددات مختلفة؛ وعندئذ لا تكون الدلالة الفعلية إلا لتراكيب المحددات التي تعتمد على $S^\rho$. غير أنه لا يمكن تبيّن أي تراكيب المحددات تحقق ذلك؛ وبذلك تضيع سائر المبرهنات المذكورة في المجال الثلاثي.

أود الآن أن أبيّن مبدأ بسيطًا للحصول على \emph{تمثيل رمزي صريح بدلالة صفوف $S^\rho$}، ومن ثم استعادة سائر المبرهنات أيضًا. وهذا تطبيق للمبرهنة المعروفة في تقابل المصفوفات: \glqq{}يمكن أن يُقرن بكل صف متغيرات $p_\rho$ صف آخر $q_{n-\rho}$ بحيث يكون التقابل واحدًا لواحد، وتكون عناصر الصف الآخر محددات صغرى مأخوذة من $n-\rho$ من صفوف $u$، المرتبطة بـ$\rho$ من صفوف $x$ بالمعادلات $(u\mid x)=0$.\grqq{} وبالمثل، يُقرن بكل صف رموز $S^\rho$ صف آخر $S_{n-\rho}$ يتصرف كما لو كان مركبًا من $n-\rho$ من الصفوف المتوافقة التدرج مع $x$.

وتقبل المبرهنة صياغة ثانية ملائمة للتطبيقات: \glqq{}يقابل كل حاصل ضرب مصفوفي $(v^{(1)}\cdots v^{(\rho)}\mid y^{(1)}\cdots y^{(\rho)})$،\srcfn{1)}{أي مجموع حواصل ضرب محددات متناظرة، مأخوذة من مصفوفة ذات $\rho$ من صفوف $v$ ومن مصفوفة ذات $\rho$ من صفوف $y$.} حيث تكون صفوف المصفوفة الثانية متخالفة التدرج مع صفوف الأولى، محددٌ؛ وبالعكس يمكن تحويل كل محدد إلى حاصل ضرب مصفوفي ذي عدد صفوف معيّن سلفًا هو $\rho$.\grqq{} وبذلك يظهر المحدد وحاصل الضرب المصفوفي متكافئين تكافؤًا تامًا، وتأخذ المبرهنة الأولى للطريقة الرمزية الصورة الآتية: \glqq{}يمكن تمثيل جميع التكوينات اللامتغيرة تمثيلًا رمزيًا بحواصل ضرب مصفوفية.\grqq{} ويمكن الآن، بمساعدة صفوف المتغيرات، أن نكوّن بسهولة من صفّي الرموز $S^\sigma$ و$T^\tau$ حاصل ضرب مصفوفيًا يشتمل صراحة على صفي الرموز هذين؛ فهو إذن يمثّل تكوينًا لامتغيرًا من الصورة $(S^\sigma\mid p_\sigma)(T^\tau\mid p_\tau)$، أعني الآتي:
\srcnumdisplay[3.2em]{(1)}{%
  \pair{q_{n-\sigma-\lambda}S^\sigma}{T_{n-\tau}p_{\tau-\lambda}},
  \qquad (\lambda=1,2,\ldots,\tau\ \text{\srcRTL{أو}}\ n-\sigma).
}

والأهم هو العكس: \emph{يمكن تمثيل جميع التكوينات اللامتغيرة تمثيلًا صريحًا بحواصل ضرب مصفوفية}؛ ومن ثم فإن العملية السابقة تمثّل توسيعًا لعملية الطي في المجالين الثنائي والثلاثي. ولإثبات هذا العكس يلزم نقل المبرهنة الثانية للطريقة الرمزية إلى المجال ذي $n$ متغيرات؛ أي إقامة المجموعة الكاملة للمتطابقات المستقلة غير القابلة للتفكيك، التي تُعدّ فيها جميع صفوف $S^\rho$ و$p_\rho$ غير قابلة للتفكيك. وتُستمد هذه المتطابقات من المتطابقات المعروفة التي لا تشتمل إلا على صفوف $x$ و$u$،\srcfn{2)}{\srcLTR{E. Pascal in} \emph{\srcLTR{Memorie d. R. Acc. d. Lincei}}\srcLTR{. (4) 5 (1888), p. 375.}} وذلك باستبدال مبرهنة حاصل ضرب المحددات بمنظومة من مبرهنات حاصل ضرب المصفوفات، ثم ضم حواصل ضرب مصفوفية مختلفة في حاصل ضرب واحد بواسطة مبرهنات حاصل الضرب نفسها. وهكذا نصل إلى صيغتين متقابلتين وفق مبدأ الثنائية، لا نورد منهما إلا إحداهما:
\srcnumdisplay[3.2em]{(2)}{%
\begin{aligned}
 \pair{S_1^{\rho_1}S_2^{\rho_2}\cdots S_k^{\rho_k}}{x p_{\rho-1}}
 &=\sum_{i=1}^{k}\eps\,\bigl(S_1^{\rho_1}\cdots S_{i-1}^{\rho_{i-1}}S_{i+1}^{\rho_{i+1}}\cdots S_k^{\rho_k}q_{n-\rho+1},\,S_{n-\rho_i}x\bigr);\\
 \rho&=\sum_{i=1}^{k}\rho_i\le n;\qquad \eps=\pm1.
\end{aligned}
}

لا يمكن تجنب التخصيص الوحيد الذي تنطوي عليه هاتان الصيغتان، وهو دخول صف $x$؛ وعند أخذ صفوف الرموز والمتغيرات عامة تمامًا، نصل إلى \glqq{}متطابقة التفكيك\grqq{}، أي متطابقة يُفكك فيها صف $p_\rho$ إلى صفوف $x$ المفردة. وقد استخدم كليبش من قبل أبسط حالة خاصة لمتطابقة التفكيك في بحثه \glqq{}\srcLTR{Fundamentalaufgabe der Invariantentheorie}\grqq{} لاشتقاق الملازمات الأولية الموافقة للتغير في التوسعات المتسلسلة. وبمساعدة هذه المتطابقات، التي تتوسط الانتقال من عبارات المحددات غير الصريحة إلى حاصل الضرب المصفوفي، يمكن بالفعل البرهنة على أن التمثيل بالمصفوفات يحقق التمثيل الرمزي الصريح المنشود.

أما عملية الطي فتصح فيها مبرهنة مماثلة تمامًا لما في المجال الثلاثي، ويمكن أن تُبنى عليها مبرهنات الاختزال على النحو نفسه. فإذا سمّينا في \srcLTR{(1)} العدد $\lambda$ نقص الطي، وسمّينا عمليات الطي التي فيها $\sigma=\tau$ متوافقة التدرج، \emph{أمكن توليد جميع عمليات الطي بالتطبيق المتعاقب لعمليات الطي المتوافقة التدرج ذات النقص \srcLTR{1}}. ويعتمد برهان هذه المبرهنة أساسًا على إمكان استبدال أعم الحدوديات المتجانسة بـ\glqq{}صيغ نظامية\grqq{}، أي بحدوديات متجانسة تساوي الصفر تطابقيًا عند تطبيق جميع العمليات التفاضلية اللامتغيرة. وقد أثبت السيد \srcspaced{مرتنس} هذه الواقعة الأخيرة في المجال رباعي المتغيرات\srcfn{1)}{\srcLTR{Über invariante Gebilde quaternärer Formen.} \srcLTR{Wiener Berichte. Bd. 98. 1889.}}؛ ومعرفتنا بأعم العمليات التفاضلية، التي تنشأ كما هو معلوم من عمليات الطي باستبدال صفوف الرموز برموز تفاضلية، تتيح لنا، باستعمال متطابقة التفكيك، نقل برهان مرتنس إلى المجال ذي $n$ متغيرات نقلًا يكاد يكون حرفيًا.

\clearpage
\setcounter{footnote}{0}
\begin{center}
{\Large\bfseries \srcRTL{في نظرية لامتغيرات}\\
\srcRTL{الحدوديات المتجانسة ذات}~$n$~\srcRTL{متغيرات}\footnote{قارن بلاغًا موجزًا مقدّمًا إلى اجتماع علماء الطبيعة في سالزبورغ. \srcLTR{Jahresbericht d. deutsch. Math.-Vereinigung. Bd.\ 19, S.\ 101.}}}\par
\vspace{0.5em}
\srcRTL{بقلم الآنسة \emph{\srcspaced{إيمي نوتر}} في إرلنغن.}\par
\medskip
\rule{0.22\linewidth}{0.4pt}\par
\vspace{0.7em}
\srcLTR{Journal f.\ d.\ reine u.\ angew.\ Math.\ 139 (1911), S.\ 118--154}
\end{center}

\begin{center}
\textbf{مقدمة.}
\end{center}
في نظرية اللامتغيرات الإسقاطية للحدوديات المتجانسة ذات $n$ متغيرات، لم تكَد تُبحث، فيما وراء المجالين الثنائي والثلاثي، المسائل المتصلة بالمسألة الرئيسة -- إثبات تناهي منظومة الملازمات. وهذه هي مسائل الاعتماد المتبادل بين الحدوديات، ومسائل الطرائق التي تتيح الإحاطة ببنية العلاقات بينها. وتقوم هذه الطرائق أساسًا -- كما ثبت على وجه كامل في المجالين الثنائي والثلاثي -- على مبرهنتين سمّاهما السيد شتودي ``المبرهنتين الأساسيتين للطريقة الرمزية''، هما مبرهنة إمكان تمثيل الحدوديات تمثيلًا رمزيًا، ومبرهنة المتطابقات الرمزية.\footnote{\srcLTR{E. Study,} \emph{\srcLTR{Methoden zur Theorie der ternären Formen}} \srcLTR{(Leipzig, Teubner, 1889). II. \S\ 6.} -- ولورود هاتين المبرهنتين الأساسيتين أيضًا في نظرية لامتغيرات زمر تحويل خاصة، قارن \srcLTR{Study,} \emph{\srcLTR{Das Apollonische Problem}} \srcLTR{(Math. Ann. Bd. 49 [1897])} و\emph{\srcLTR{Zur Differentialgeometrie der analytischen Kurven}} \srcLTR{(Transactions of the American Math. Society, Bd. 10 [1909]).}} وتقرر المبرهنة الأخيرة أن جميع العلاقات القائمة بين التكوينات اللامتغيرة كثيرة الحدود تُستخرج بالتطبيق المتعاقب لعدد قليل من المتطابقات الرمزية؛ أي إن الطريقة الرمزية، وهي وحدها، تتيح معرفة بنية العلاقات بين الحدوديات من غير الخروج من نطاق اللامتغيرات.

وقد أشار السيد غوردان من قبل، في برنامجه الإرلنغني،\footnote{\srcLTR{P. Gordan,} \emph{\srcLTR{Über das Formensystem binärer Formen}}\srcLTR{. S. 46. (Leipzig, Teubner, 1875).}} إلى منشأ الصعوبة عند الانتقال إلى $n$ متغيرات؛ وهو أن المبرهنة الأولى للطريقة الرمزية لا تعود صحيحة في صيغتها العامة. فمن المعلوم أنه تظهر في المجال ذي $n$ متغيرات صفوف متغيرات من رتب أعلى $p_\rho$\footnote{في الاصطلاح، قارن \S\ 1.}، وبالمثل صفوف رموز من رتب أعلى $S^\rho$؛ ونصل إلى هذه الصفوف سواء انطلقنا، كما فعل كليبش أولًا،\footnote{\srcLTR{A. Clebsch, über eine Fundamentalaufgabe der Invariantentheorie. (Abh. der Ges. d. Wiss. zu Göttingen, Bd. XVII, 1872).}} من حدوديات متجانسة تشتمل على عدد اعتباطي من صفوف $p_\rho$، أم انطلقنا، على نهج \srcLTR{F. Deruyts} و\srcLTR{A. Capelli}، من حدوديات لا تشتمل إلا على عدد اعتباطي من صفوف $x$ المتوافقة التدرج.\footnote{قارن \srcLTR{A. Capelli,} \emph{\srcLTR{Lezioni sulla teoria delle forme algebriche}} \srcLTR{(Napoli 1902), \S\ 22--24,} والمراجع المذكورة هناك.} غير أنه لا يوجد تمثيل رمزي يشتمل صراحة\footnote{تعريف أدق لـ``التمثيل الصريح'' في \S\ 2.} على الصفوف $p_\rho$ و$S^\rho$؛ فلا بد من تفكيك $p_\rho$ و$S^\rho$ إلى فرادى صفوف $x$ وصفوف $s$ المتخالفة التدرج معها، ثم توزيع هذه الصفوف على محددات مختلفة. ومع أن الشروط التفاضلية (قارن الصيغة \srcLTR{(19.)}) تتيح التحقق من أن مجموعًا معيّنًا من المحددات لا يعتمد إلا على الصفوف $p_\rho$ و$S^\rho$، فإن هذه الطريقة لا تتيح تكوين نظرة شاملة إلى جميع التكوينات اللامتغيرة والعمليات التي تولدها.\footnote{كذلك فإن تطويرًا للحساب الرمزي وضعه السيد شتودي، وله قيمة حسابية كبيرة جدًا (نقله \srcLTR{F. Engel, Ber. d. K. Sächs. Ges. d. Wiss., math.-phys. Kl., 1900, S. 126,} وطبّقه السيد شتودي على المجال الرباعي في \emph{\srcLTR{Geometrie der Dynamen}}\srcLTR{, S. 135})، ليس تمثيلًا صريحًا بالمعنى المقصود هنا. ويُستبعد، مثلًا، أن يقود بسهولة إلى عمليات الاشتقاق غير الرمزية التي تولّد الحدوديات.}

ونبيّن الآن في \S\S\ 2 و6 كيف تُستعاد المبرهنة الأساسية الأولى في عمومها باستعمال ``مبرهنة تقابل المصفوفات''،\footnote{عدا ما ورد في نظرية الامتداد لغراسمان سنة \srcLTR{1862, Nr. 112}، ظهرت أولًا عند \srcLTR{A. Brill, über zwei Berührungsprobleme, \S\ 2 (Math. Ann. Bd. 4. 1871).}} وذلك بإحلال تمثيل بـ``حواصل ضرب مصفوفية'' محل التمثيل بالمحددات. ففي \S\ 2 يُبيَّن أن كل حاصل ضرب مصفوفي يشتمل صراحة على الصفوف $S^\rho,p_\rho$ هو تكوين لامتغير للحدودية الأساسية؛ أما العكس الوارد في \S\ 6، أي إمكان تمثيل كل تكوين لامتغير صراحة بحواصل ضرب مصفوفية، فلا يتحقق إلا بنقل المبرهنة الأساسية الثانية (\S\S\ 3--5).

وهذه المبرهنة معروفة للحدوديات المتجانسة التي لا تشتمل إلا على صفوف المتغيرات $x$.\footnote{\srcLTR{E. Pascal, Giorn. di mat. 26 (1888), S. 33, 102; Rendic. d. Accad. d. Linc. (4) 4 (1888), S. 119; ib. Mem. (4) 5 (1888), S. 375.}} وقد صاغ السيد شتودي المسألة العامة المتمثلة في تعيين جميع المتطابقات غير القابلة للتفكيك، مع عدّ جميع الصفوف $S^\rho,p_\rho$ غير قابلة للتفكيك؛\footnote{\srcLTR{``Methoden \dots'', S. 204, Anmerkung 17).}} وهو يرى كذلك في عدد المتطابقات التي تظهر مقياسًا لصعوبة الطرائق في المجال ذي $n$ متغيرات. وبما أنه أمكن جمع جميع المتطابقات في صيغة واحدة \srcLTR{(36.)}، فقد خُفِّضت هذه الصعوبة، على الأقل، إلى حدها الأدنى. غير أنه كثيرًا ما تظهر عند التطبيق حالات خاصة متميزة من \srcLTR{(36.)}، مثل \srcLTR{(20.)} و\srcLTR{(21.)}، تمتاز بأنها تتيح التمثيل الصريح من غير إحلال صفوف جديدة؛ أو الصيغة \srcLTR{(31.)} التي تسمى ``متطابقة التفكيك''، إذ يظهر فيها صف $p_\rho$ كما لو كان مفككًا إلى الصفوف $y$.

ويمكن الآن بناء سائر النظرية بسهولة على المبرهنتين الأساسيتين. فالمبرهنة الأولى تعطي مباشرة عمليات توليد الحدوديات: عملية الطي الرمزية وعمليات الاشتقاق غير الرمزية،\footnote{في بحث \srcLTR{``Invariante Gebilde von Nullsystemen''} \srcLTR{(Sitzungsber. d. Ak. d. Wiss. Wien, Bd. 97, Abt. IIa, 1888)} توصّل السيد مرتنس، بطريق أخرى لا تقبل التعميم المباشر، إلى عملية الطي في المجال الرباعي. واللامتغير المتناوب المؤلف من \srcLTR{6} من صفوف $S^2$ والوارد هناك يختلف عن اللامتغير الذي ينشأ لدينا من تطبيق الإحلال \srcLTR{(11.)} مرة واحدة بمجموع من اللامتغيرات الزوجية. وترد عملية الطي الرباعية، في حالة خاصة، أيضًا عند السيد \srcLTR{P. Gordan,} \emph{\srcLTR{Das simultane System von zwei quadratischen quaternären Formen}} \srcLTR{(Math. Ann. Bd. 56. 1903).}} في حين تستبعد متطابقة التفكيك من بين هذه العمليات كل ما يكافئ غيره (\S\ 6). وتتيح معرفة عمليات الاشتقاق، فضلًا عن ذلك، نقل مفهوم ``الصيغة النظامية'' إلى المجال ذي $n$ متغيرات؛ كما تفضي، باستعمال طريقة برهان قدّمها السيد مرتنس في المجال الرباعي،\footnote{\srcLTR{F. Mertens, Über invariante Gebilde quaternärer Formen. (Sitzber. d. Ak. d. Wiss. Wien. Bd. 98. Abt. IIa. 1889.)}} إلى المبرهنة القائلة إن منظومة من التكوينات اللامتغيرة يمكن أن تُستبدل بها منظومة مكافئة من الصيغ النظامية (\S\ 7). وتحت هذا الفرض الأخير بيّنت في أطروحتي\footnote{\srcLTR{Über die Bildung des Formensystems der ternären biquadratischen Form. (Dieses Journal, Bd. 134. 1908.)}} أن عملية الطي في المجال الثلاثي يمكن توليدها بالتطبيق المتعاقب لطيّتين أساسيتين. ثم طوّرت، بالاستناد إلى هذه المبرهنة ونظرية التوسعات المتسلسلة، أهم مبرهنات الاختزال لمنظومات الملازمات، بعد إدخال مفهوم ``سلسلة الملازمات''. وبالمثل تمامًا، يمكن في المجال ذي $n$ متغيرات توليد عملية الطي من $n-1$ طيّات أساسية (\S\ 8)، وصياغة مبرهنات الاختزال (\S\ 9).

\paragraph{ملاحظة لاحقة.} بمناسبة بلاغ سالزبورغ، نبّهني الأستاذ إ. مولر من فيينا إلى أن الحساب بحواصل الضرب المصفوفية الذي تقوم عليه هذه المقالة مطابق من حيث المبدأ لحساب نظرية الامتداد عند غراسمان.\footnote{\srcLTR{Hermann Graßmanns gesammelte mathematische und physikalische Werke. Ersten Bandes zweiter Teil: Die Ausdehnungslehre von 1862. In Gemeinschaft mit H. Graßmann d. Jüngeren herausgegeben v. Fr. Engel (Leipzig, Teubner 1896).}} والواقع أن ``الحاصل الضربي التوافقي'' لغراسمان $[S^{\rho_1}S^{\rho_2}\cdots S^{\rho_i}]$ يمكن أن يُفهم على أنه مصفوفة تعيّنها هذه الصفوف (قارن \S\ 2): فـ``حاصل الضرب التقدّمي'' هو مصفوفة الصفوف نفسها، و``حاصل الضرب التراجعي'' هو مصفوفة الصفوف المقترنة $S_{n-\rho_1},S_{n-\rho_2},\ldots,S_{n-\rho_i}$، أما المصفوفة المقابلة لـ``حاصل الضرب المختلط'' فتنشأ بضم عدة صفوف $S$، عن طريق الإحلال \srcLTR{(9.)}، في صفوف جديدة $P,Q$. وفي هذا التقابل، يقابل ``صفّنا المقترن'' ``متمّمة'' غراسمان، ويقابل ``حاصل الضرب المصفوفي'' عندنا ``حاصل ضرب مختلط من الدرجة \srcLTR{0}''.

وبموجب هذا التقابل، يصبح التوسّع الوارد في نظرية الامتداد لسنة \srcLTR{1862, Nr. 113} مطابقًا لصيغة ثنوية لصيغتنا \srcLTR{(32.)}. وفي الحالة الخاصة $\rho=1$، تتحول \srcLTR{(32.)} إلى \srcLTR{(17.)}، وتتحول الصيغة الثنوية إلى \srcLTR{(16.)}. وقد اشتق السيد مولر حديثًا،\footnote{\srcLTR{E. Müller, Beiträge zur Graßmannschen Ausdehnungslehre. 1. Mitteilung: Einige allgemeine Sätze. (Sitzungsber. d. Ak. d. Wiss. Wien; Bd. 118. Abt. IIa. Juli 1909), Satz 16 und 18.}} من هاتين الحالتين الخاصتين لتوسّع غراسمان، المتطابقتين \srcLTR{(20.)} و\srcLTR{(21.)} المذكورتين أعلاه.

ويتميّز اشتقاقنا عن اشتقاق غراسمان--مولر بأنه يقدّم في الوقت نفسه برهان اكتمال المتطابقات، وهو الأمر الأهم للمسائل المعنية هنا، كما أنه يتجاوز \srcLTR{(20.)} و\srcLTR{(21.)} إلى أعم المتطابقات غير القابلة للتفكيك \srcLTR{(36.)}، التي يبدو أنها لم تُلحظ حتى الآن.

\section*{\S\ 1. الترميز والتعريفات. خلاصة النتائج المعروفة.}
نسمّي، على المعتاد، صفَّ المتغيرات $x$ صفَّ العناصر $x_1,x_2,\ldots,x_n$، ونسمّي، بالمثل، صفَّ المتغيرات $p_\rho$ ($\rho=1,2,\ldots,n-1$) صفَّ العناصر البالغ عددها $\binom n\rho$، وهي $p_{i_1i_2\ldots i_\rho}$؛ حيث إن $p_{i_1i_2\ldots i_\rho}$ هي المحددات المكوَّنة من مصفوفة ذات $\rho$ من الأسطر $x^{(1)},x^{(2)},\ldots,x^{(\rho)}$، والبالغ عددها $\binom n\rho$، أي
\[
\left|\begin{array}{cccc}
 x_{i_1}^{(1)}&x_{i_2}^{(1)}&\cdots&x_{i_\rho}^{(1)}\\
 x_{i_1}^{(2)}&x_{i_2}^{(2)}&\cdots&x_{i_\rho}^{(2)}\\
 \vdots&\vdots&&\vdots\\
 x_{i_1}^{(\rho)}&x_{i_2}^{(\rho)}&\cdots&x_{i_\rho}^{(\rho)}
\end{array}\right|,
\qquad (i_1<i_2<\cdots<i_\rho=1,2,\ldots,n)
\]
وبموجب مبرهنة تقابل المصفوفات\footnote{قارن، مثلًا: \srcLTR{Gordan-Kerschensteiner,} \emph{\srcLTR{Vorlesungen über Invariantentheorie}}, \srcLTR{Bd. I, S. 99.}} يُقرَن بكل صف $p_\rho$ صفٌّ ثانٍ $q_{n-\rho}$ بحيث يكون التقابل واحدًا لواحد، وذلك بالعلاقة الآتية التي تصدق على كل عنصر من عناصر الصف:
\begin{equation}
 p_{i_1i_2\ldots i_\rho}=(-1)^{\frac{\rho(\rho+1)}{2}+i_1+i_2+\cdots+i_\rho}q_{j_1j_2\ldots j_{n-\rho}},
\tag{1}
\end{equation}
حيث إن فهارس $i$ وفهارس $j$ مرتبة تصاعديًا، كلٌّ على حدة، وتؤلف معًا تبديلًا كاملًا للأعداد من 1 إلى $n$. ويتألف الصف $q_{n-\rho}$ من المحددات البالغ عددها $\binom n\rho$، من الدرجة $(n-\rho)$، وهي $q_{j_1j_2\ldots j_{n-\rho}}$ المكوَّنة من مصفوفة ذات $n-\rho$ من الأسطر المتخالفة التدرج مع $x$، وهي $u^{(1)},u^{(2)},\ldots,u^{(n-\rho)}$. وتحقق هذه الأسطر معادلات الشرط البالغ عددها $\rho(n-\rho)$:
\begin{equation}
 (u^{(k)}\mid x^{(l)})=\sum_{\alpha=1}^{n}u_\alpha^{(k)}x_\alpha^{(l)}=0;
 \qquad (k=1,2,\ldots,n-\rho;\ l=1,2,\ldots,\rho).
\tag{2}
\end{equation}
وسنفترض هذه الأسطر مطبَّعةً (مع كليبش\footnote{مرجع كليبش المذكور آنفًا، \S\ 5.}) بحيث يُجعل عامل التناسب، الذي كان سيدخل لولا ذلك في \srcLTR{(1.)}، مساويًا للواحد.

يمكن، بحسب كليبش\footnote{المرجع نفسه، \S\ 12.} (وكذلك بحسب أبحاث \srcLTR{F. Deruyts} و\srcLTR{A. Capelli} اللاحقة؛ قارن المقدمة)، ردُّ أعمّ الحدوديات إلى حدوديات لا تشتمل إلا على صف واحد على الأكثر من كلٍّ من الصفوف البالغ عددها $n-1$ من النوع $p_\rho$؛ ومن ثم يمكن تمثيلها رمزيًا على النحو الآتي:
\begin{equation}
 (S^1\mid p_1)^{m_1}(S^2\mid p_2)^{m_2}\cdots(S^{n-1}\mid p_{n-1})^{m_{n-1}},
\tag{3}
\end{equation}
حيث وضعنا، قياسًا على \srcLTR{(2.)}:
\[
 (S^\rho\mid p_\rho)=\sum S^{i_1i_2\ldots i_\rho}p_{i_1i_2\ldots i_\rho}.
\]
ويمكن تطبيع صفوف الرموز $S^\rho$، بحكم المتطابقات غير الخطية القائمة بين عناصر صف $p_\rho$، بحيث تحقق تلك المتطابقات نفسها تمامًا؛ أي بحيث تتصرف مثل محددات مصفوفة ذات $\rho$ من الأسطر، أسطرها $s^{(1)},s^{(2)},\ldots,s^{(\rho)}$ متخالفة التدرج مع $x$.\footnote{كليبش، المرجع نفسه، \S\ 4.} ومن ثم يمكن أن يُقرَن بكل صف $S^\rho$ صفٌّ آخر $S_{n-\rho}$ بحيث يكون التقابل واحدًا لواحد، وذلك بالعلاقة:
\begin{equation}
 S_{i_1i_2\ldots i_{n-\rho}}=(-1)^{\frac{(n-\rho)(n-\rho+1)}2+i_1+i_2+\cdots+i_{n-\rho}}S^{j_1j_2\ldots j_\rho},
\tag{4}
\end{equation}
حيث تكون فهارس $i$ وفهارس $j$، مرة أخرى، مرتبة تصاعديًا، كلٌّ على حدة، وتستنفد معًا الأعداد من 1 إلى $n$ استنفادًا تامًا. وتتصرف عناصر $S_{i_1i_2\ldots i_{n-\rho}}$ من الصف $S_{n-\rho}$ مثل المحددات المكوَّنة من $(n-\rho)$ من الأسطر المتوافقة التدرج مع $x$، وهي $\sigma^{(1)},\sigma^{(2)},\ldots,\sigma^{(n-\rho)}$؛ وترتبط أسطر $s$ و$\sigma$ بمعادلات الشرط البالغ عددها $\rho(n-\rho)$:
\begin{equation}
 (s^{(k)}\mid \sigma^{(l)})=0,\qquad(k=1,2,\ldots,\rho;\ l=1,2,\ldots,n-\rho).
\tag{5}
\end{equation}
وبحكم العلاقتين \srcLTR{(1.)} و\srcLTR{(4.)}، يقبل كل عامل من عوامل \srcLTR{(3.)} تمثيلًا رمزيًا في أربع صيغ متكافئة:
\begin{equation}
 (S^\rho\mid p_\rho)=(S^\rho q_{n-\rho})=(S_{n-\rho}p_\rho)=(-1)^{\rho(n-\rho)}(q_{n-\rho}\mid S_{n-\rho})\footnotemark,
\tag{6}
\end{equation}
\footnotetext{قارن كليبش، المرجع نفسه، \S\ 5.}
وفي ذلك، كما هي الحال دائمًا فيما يأتي، يُستعمل $(S^\rho q_{n-\rho})$ و$(S_{n-\rho}p_\rho)$ اختصارًا للمحددتين $(s^{(1)}\cdots s^{(\rho)}u^{(1)}\cdots u^{(n-\rho)})$ و$(\sigma^{(1)}\cdots\sigma^{(n-\rho)}x^{(1)}\cdots x^{(\rho)})$؛ ويتبين، بنشر لابلاس، أن هاتين المحددتين لا تعتمدان إلا على الصفوف $S^\rho,q_{n-\rho}$ وعلى الصفوف $S_{n-\rho},p_\rho$، على الترتيب، وهو ما يُراد للترميز أعلاه أن يوحي به. أما التعبيران $(S^\rho\mid p_\rho)$ و$(q_{n-\rho}\mid S_{n-\rho})$ فيعنيان، بحسب ما تقدم، جداءَ مصفوفة أسطر $s$ في مصفوفة أسطر $x$، وجداءَ مصفوفة أسطر $u$ في مصفوفة أسطر $\sigma$، على الترتيب؛ والمقصود هنا بجداء المصفوفات مجموعُ حواصل ضرب المحددات المتناظرة، أي قيمةُ محدد مصفوفة الجداء.

نسمّي، مع كليبش، العدد المميّز $\rho(n-\rho)$ العائد إلى صف رموز (أو صف متغيرات) وزنَ الصف؛\footnote{المرجع نفسه، \S\ 11.} ونسمّي العدد $\rho$، الذي يبيّن عدد أسطر $s$، دليلَ الوزن العلوي، والعدد $n-\rho$، الذي يبيّن عدد أسطر $\sigma$، دليلَ الوزن السفلي. ويترتب من العلاقات \srcLTR{(4.)} أن صف الرموز الواحد يمكن أن يظهر في العلاقة نفسها مرةً بدليل وزن علوي ومرةً بدليل وزن سفلي؛ ويصدق الأمر نفسه على ظهور الصفين المقترنين $p_\rho$ و$q_{n-\rho}$. ونشير أيضًا إلى أننا نرمز، على وجه العموم، إلى صفوف المتغيرات من رتب أعلى التي تتألف مصفوفتها من أسطر $x$ بالحرف $p$، وإلى التي تتألف مصفوفتها من أسطر $u$ بالحرف $q$؛ وإلى صفوف الرموز المتوافقة التدرج مع $x$ بالحروف اليونانية الصغيرة: $\sigma,\tau,\alpha,\beta$؛ وإلى صفوف الرموز المتوافقة التدرج مع $u$ بالحروف اللاتينية الصغيرة: $s,t,a,b$؛ وإلى سائر صفوف الرموز بالحروف اللاتينية الكبيرة المزودة بدليل وزن: $S^\sigma$، $T^\tau$، $A^\rho$، أو، على الترتيب، $S_{n-\sigma}$، $T_{n-\tau}$، $A_{n-\rho}$. وتبعًا لكون دليل الوزن علويًا أو سفليًا، نكتب كذلك العناصر المفردة للصف بفهارس علوية أو سفلية، كما في \srcLTR{(4.)} مثلًا.

\section*{\S\ 2. التمثيل بجداءات المصفوفات.}
نقصد دائمًا فيما يأتي بـ``جداء المصفوفات'' مجموع حواصل ضرب المحددات المتناظرة لمصفوفتين لهما العدد نفسه من الأسطر، بحيث تكون أسطر المصفوفة الثانية متخالفة التدرج مع الأسطر المناظرة لها في الأولى. ويمكن في ذلك أن تكون أسطر كل مصفوفة: 1. معطاةً منفردة؛ 2. مجمَّعةً، كما في (6.)، في صف واحد $S^\rho$؛ 3. أو مجموعةً في صفوف مختلفة $S^{\rho_1},S^{\rho_2}\ldots$. ويرمز إلى جداء المصفوفات بالرمز $(\ \mid\ )$ المستعمل في (6.). والتمثيل (6.) تطبيق خاص لصيغة ثانية من مبرهنة تقابل المصفوفات: ``يقترن بكل جداء مصفوفات ذي $\rho$ من الأسطر $(v^{(1)}v^{(2)}\cdots v^{(\rho)}\mid y^{(1)}\cdots y^{(\rho)})$ محددٌ؛ وبالعكس، يمكن أن يقترن بكل محدد جداءُ مصفوفات ذي عدد من الأسطر محدد سلفًا، هو $\rho$ ($\rho=1,2,\ldots,n-1$)''.
\footnote{من المعلوم أن جداء المصفوفات ينعدم إذا كان $\rho>n$؛ أما إذا كان $\rho=n$ فيؤول إلى حاصل ضرب محددين.}
ولا يلزم، في الواقع، سوى جمع أسطر $y$ في صف $p_\rho$ (أو جمع أسطر $v$ في صف $q_\rho$) وإجراء الإحلالات المعطاة في (1.)، لكي ننتقل من جداء مصفوفات معطى إلى المحدد، وبالعكس. غير أن قيمة المحدد لا تُعطى هنا بعناصره المفردة، بل لا تُستخرج إلا بنشر لابلاس. ولما ثبت بذلك تكافؤ المحدد وجداء المصفوفات، يمكن أيضًا صياغة مبرهنة قابلية التمثيل الرمزي للحدوديات (المبرهنة الأساسية الأولى) كما يأتي:

\emph{المبرهنة \srcLTR{I}:} ``يمكن تمثيل جميع التكوينات اللامتغيرة رمزيًا بجداءات المصفوفات؛ وبالعكس، فإن جميع جداءات المصفوفات تكوينات لامتغيرة''.
\footnote{ومن البديهي أن التكوينات اللامتغيرة للحدوديات الأساسية المعطاة لا تكون إلا تلك التجميعات من جداءات المصفوفات التي تعتمد على صفوف الرموز $S^\rho\ldots$.}

يتيح التمثيل بجداءات المصفوفات تذليل الصعوبة الجوهرية الملازمة للتمثيل غير الصريح، وهي صعوبة يفرضها التمثيل بالمحددات.
\footnote{قارن المقدمة.}
ونقصد هنا بـ``التمثيل الصريح'' تمثيلًا لا تدخل فيه إلا الصفوف $S^\rho,T^\tau,\ldots$ نفسها، أو صفوف $R^\rho$ التي يمكن استنتاجها منها على نحو وحيد، ولا تعود الصفوف الجزئية المنفردة $s,t,\ldots$ تظهر فيه. وسنحصل في \S\ 6، لجميع التكوينات اللامتغيرة التي لا تعتمد إلا على صفوف الرموز $S^\rho,T^\tau,\ldots$، على تمثيل صريح بهذه الصفوف الرمزية، وبذلك ننقل عمليات التوليد، وهي عملية الطي وعمليات الاشتقاق، إلى المجال ذي $n$ متغيرات. ويتضح إمكان التمثيل الصريح من أن أبسط جداءات المصفوفات وحدها يمكن أن تقابل تمامًا المحددات التي تظهر في التمثيل بالمحددات؛ ومن ثم فإن جميع جداءات المصفوفات الأخرى تجمع عدة محددات مختلفة في تعبير واحد. ذلك أنه إذا حُلِّلت صفوف الرموز والمتغيرات إلى الصفوف المفردة $s,u,x$، فلا بد، بمقتضى المبرهنة الأساسية الأولى في صيغتها المعتادة، من أن ينشأ تجميع من المحددات. أما إثبات العكس، أي أن جميع تجميعات المحددات التي تمثل تكوينات لامتغيرة لحدوديات أساسية معينة سلفًا يمكن جمعها في جداءات مصفوفات صريحة، فيُستخلص باستعمال المبرهنة الأساسية الثانية (\S\ 6).

والآن، لكي نحصل من صفّي الرموز $S^\sigma,T^\tau$، بالاستعانة بصفوف المتغيرات، على تكوين لامتغير للحدودية الأساسية $(S^\sigma\mid p_\sigma)(T^\tau\mid p_\tau)$ يشتمل صراحةً على جميع الصفوف (وهي عملية الطي)، لا يلزمنا إلا أن نكوّن، وفق التعريف المعطى في أول هذا المبحث، جداء مصفوفات من النوع الثالث:
\begin{equation}
 \pair{q_{n-\sigma-\lambda}S^\sigma}{T_{n-\tau}p_{\tau-\lambda}},
 \qquad (\lambda=1,2,\ldots,\tau,n-\sigma\footnotemark).
\tag{7}
\end{equation}
\footnotetext{يمتد العدد $\lambda$ دائمًا إلى أصغر العددين المذكورين أخيرًا.}
وبنشر (7.) نحصل على:
\begingroup
\setlength{\arraycolsep}{2pt}
\begin{equation}
\begin{aligned}
&\pair{v^{(1)}\cdots v^{(n-\sigma-\lambda)}s^{(1)}\cdots s^{(\sigma)}}
       {\alpha^{(1)}\cdots\alpha^{(n-\tau)}y^{(1)}\cdots y^{(\tau-\lambda)}} \\
&=\sum
\left|\begin{array}{ccc}
 v_{i_1}^{(1)}&\cdots&v_{i_{n-\lambda}}^{(1)}\\
 \vdots&&\vdots\\
 v_{i_1}^{(n-\sigma-\lambda)}&\cdots&v_{i_{n-\lambda}}^{(n-\sigma-\lambda)}\\
 s_{i_1}^{(1)}&\cdots&s_{i_{n-\lambda}}^{(1)}\\
 \vdots&&\vdots\\
 s_{i_1}^{(\sigma)}&\cdots&s_{i_{n-\lambda}}^{(\sigma)}
\end{array}\right|
\left|\begin{array}{ccc}
 \alpha_{i_1}^{(1)}&\cdots&\alpha_{i_{n-\lambda}}^{(1)}\\
 \vdots&&\vdots\\
 \alpha_{i_1}^{(n-\tau)}&\cdots&\alpha_{i_{n-\lambda}}^{(n-\tau)}\\
 y_{i_1}^{(1)}&\cdots&y_{i_{n-\lambda}}^{(1)}\\
 \vdots&&\vdots\\
 y_{i_1}^{(\tau-\lambda)}&\cdots&y_{i_{n-\lambda}}^{(\tau-\lambda)}
\end{array}\right|\\
&=\sum_{i,k,l}\eps\,q_{k_1\ldots k_{n-\sigma-\lambda}}
 S^{k_{n-\sigma-\lambda+1}\ldots k_{n-\lambda}}
 T_{l_1\ldots l_{n-\tau}}
 p^{l_{n-\tau+1}\ldots l_{n-\lambda}},
\end{aligned}
\tag{8}
\end{equation}
\endgroup
أي إن الناتج تعبير لا يعتمد إلا على الصفوف $S,T,q,p$. وهنا $\eps=\pm1$\footnote{سيُحافظ على هذا الترميز في جميع ما يأتي.}، ويُفهم الجمع على أن فهارس كل صف من الصفوف $S\ldots$ مرتبة تصاعديًا فيما بينها، وأنه، لكل توفيقة $i_1,i_2,\ldots,i_{n-\lambda}$، تمر فهارس $k$ وفهارس $l$، كلٌّ على حدة، بجميع تباديل أعداد $i$ التي تظل ممكنة. ونسمي العدد $\lambda$ الداخل في (7.) نقص الطي في الحالة $\sigma\geqq\tau$؛ وسيتضح في \S\ 6 سبب هذا التقييد الأخير.

يمكن عدّ محددات المصفوفتين الداخلتين في (7.) عناصرَ في صفين جديدين $Q^{n-\lambda},P_{n-\lambda}$، ولْيكن الصفان المقترنان بهما $Q_\lambda,P^\lambda$. ويمكن، بحسب (8.)، التعبير عن الصفين $Q,P$ بدلالة الصفوف الأصلية $S$ و$q$، و$T$ و$p$، على الترتيب. ونشير إلى ذلك بالمعادلتين:
\begin{equation}
\begin{aligned}
 q_{n-\sigma-\lambda}S^\sigma&=Q^{n-\lambda}\sim Q_\lambda, &\quad\text{أو بعبارة أخرى}\quad Q_\lambda&=[q_{n-\sigma-\lambda}S^\sigma]_\lambda,\\
 T_{n-\tau}p_{\tau-\lambda}&=P_{n-\lambda}\sim P^\lambda, &\quad\text{أو بعبارة أخرى}\quad P^\lambda&=[T_{n-\tau}p_{\tau-\lambda}]^\lambda.
\end{aligned}
\tag{9}
\end{equation}
ومع مراعاة (4.)، يكون لجداء المصفوفات (7.)، قياسًا على (6.)، التمثيل الرمزي الرباعي المتكافئ الآتي:
\begin{equation}
\pair{q_{n-\sigma-\lambda}S^\sigma}{T_{n-\tau}p_{\tau-\lambda}}
=(q_{n-\sigma-\lambda}S^\sigma P^\lambda)=(Q_\lambda T_{n-\tau}p_{\tau-\lambda})=(-1)^{\lambda(n-\lambda)}\pair{P^\lambda}{Q_\lambda}.
\tag{10}
\end{equation}
ويمكن إجراء إحلال آخر لصفوف جديدة في (7.)، وذلك بوضع:
\begin{equation}
 \pair{q_{n-\sigma-\lambda}S^\sigma}{T_{n-\tau}p_{\tau-\lambda}}
 =\pair{R^{\sigma+\lambda}}{p_{\sigma+\lambda}}(R^{\tau-\lambda}p_{\tau-\lambda}).
\tag{11}
\end{equation}
وهنا أيضًا تُبيّن (8.) أن حواصل ضرب عناصر صفوف $R$ يمكن التعبير عنها بصورة وحيدة بدلالة الصفوف $S$ و$T$. وتُستعمل الإحلالات (9.) و(11.) خصوصًا عندما يُجرى، باستعمال (7.)، طيٌّ آخر مع صفوف أخرى.

يتيح لنا التمثيل بالمصفوفات أن نصوغ على نحو لامتغير العلاقات التربيعية بين عناصر صف $p_\rho$ (التي تترتب عليها، كما هو معلوم، جميع العلاقات الأخرى)، أي، بحسب \S\ 1، العلاقات الخطية في معاملات الحدودية الأساسية التي تحققها الصفوف $S^\rho$. وهذه هي المتطابقات:
\begin{equation}
 \pair{q_{n-\rho-\lambda}S^\rho}{S_{n-\rho}p_{\rho-\lambda}}=0\footnotemark,
 \qquad (\lambda=1,2,\ldots,\rho,n-\rho),
\tag{12}
\end{equation}
\footnotetext{قارن ملحقًا لـ\srcLTR{Study} في \srcLTR{Graßmanns Werke, Bd. I, Teil 2, S. 510} (شرط أن تكون الكمية $S^\rho$ بسيطة).}
حيث تُعَدّ $p$ و$q$ كميتين اعتباطيتين. وسنرى في \S\ 5 أن المتطابقات ذات عدد النقص الفردي $\lambda$ تترتب على المتطابقات ذات النقص الأعلى. والعلاقات (12.) نتيجة مباشرة للعلاقات (5.)؛ إذ لدينا:
\[
 \pair{q_{n-\rho-\lambda}S^\rho}{S_{n-\rho}p_{\rho-\lambda}}
 =(v^{(1)}\cdots v^{(n-\rho-\lambda)}s^{(1)}\cdots s^{(\rho)}\mid \sigma^{(1)}\cdots\sigma^{(n-\rho)}y^{(1)}\cdots y^{(\rho-\lambda)}),
\]
ويعطي تطبيق مبرهنة حاصل الضرب، بحكم (5.)، محددًا منعدمًا ذا $(n-\lambda)$ من الأسطر. وتقرر المتطابقات (12.) أنه يكفي، للعثور على جميع أنماط عمليات الطي، أن نطوي حاصل ضرب حدوديات خطية في كل صف من صفوف المتغيرات؛ أو، بعبارة مكافئة وفق \S\ 6، أن الحدودية المطبَّعة $(S^\rho\mid p_\rho)^m$ لا تملك أي تكوين لامتغير خطي في المعاملات.

والشروط (12.) كافية أيضًا لتوصيف الصفوف $p_\rho$. فبما أن خاصية كون العناصر المفردة من $p_\rho$ محدداتٍ لمصفوفة هي خاصية لامتغيرة، فلا بد أن يمكن التعبير عنها على نحو لامتغير؛ غير أن الشروط (12.) تمثل، بموجب المبرهنة \srcLTR{VI}، أكبر قيد لامتغير ممكن على $p_\rho$، وهو أن الحدودية الخطية المبنية من الصف $p_\rho$ بوصف عناصره معاملاتٍ لا تملك أي تكوينات لامتغيرة إطلاقًا.

\section*{\S\ 3. المتطابقات الرمزية\footnote{لصياغة المسألة ومراجعها، قارن المقدمة.}.}
علينا الآن أن ننقل المبرهنة الأساسية الثانية للطريقة الرمزية، أي أن نعيّن جميع المتطابقات غير القابلة للاختزال وغير القابلة للتفكيك التي تُعَدّ فيها جميع الصفوف $S^\rho,p_\rho$ غير قابلة للتفكيك. ونثبت في هذا الصدد الاصطلاح الآتي أيضًا: انسجامًا مع معنى صفوف الرموز وصفوف المتغيرات، تُعَدّ مركبةً تلك المتطابقات التي تنشأ حين نخصص صف المتغيرات $p_\rho$ إلى $p_{\rho-1}y$ ثم نطبق على التعابير الناتجة إحدى متطابقات المنظومة؛ أما المتطابقات التي تشتمل على $k$ من صفوف الرموز فتُعَدّ غير قابلة للتفكيك، حتى إذا كانت قد نتجت بالعملية المذكورة أعلاه من متطابقات تشتمل على عدد أقل من صفوف الرموز. ولا يلزم أن تدخل الصفوف $S^\rho,p_\rho$ صراحةً في المتطابقات؛ ولعدّها كميات غير قابلة للتفكيك يكفي إثبات أن حواصل ضرب المصفوفات الواردة لا تعتمد إلا على هذه الصفوف. غير أننا سنبين أنه يمكن في هذه الحالة دائمًا الحصول على تمثيل صريح، باستعمال الإحلال (11.) في بعض المواضع. ونحصل على المتطابقات العامة، بواسطة مبدأ التمثيل بالمصفوفات، من المتطابقات الخاصة التي وضعها السيد باسكال، والتي لا تشتمل إلا على الصفوف $x$ و$u$ وتمثل نقلًا مباشرًا للمتطابقات التي وضعها السيد شتودي للمجال الثلاثي:
\begin{align}
\sum_{i=1}^n(-1)^{i+1}\pair{a^{(i)}}{x}(a^{(1)}\cdots a^{(i-1)}a^{(i+1)}\cdots a^{(n)}u)&=(-1)^{n+1}\pair{u}{x}(a^{(1)}\cdots a^{(n)}),\tag{13}\\
\sum\pm\pair{a^{(1)}}{x^{(1)}}\pair{a^{(2)}}{x^{(2)}}\cdots\pair{a^{(n)}}{x^{(n)}}&=(a^{(1)}\cdots a^{(n)})(x^{(1)}\cdots x^{(n)}),\tag{14}\\
\sum_{i=1}^n(-1)^{i+1}\pair{u}{a^{(i)}}(a^{(1)}\cdots a^{(i-1)}a^{(i+1)}\cdots a^{(n)}x)&=(-1)^{n+1}\pair{u}{x}(a^{(1)}\cdots a^{(n)}).\tag{15}
\end{align}
وتنشأ متطابقتان أخريان من (13.) و(15.) بالإحلالين $(a^{(i)}\mid x)=(a^{(i)}q_{n-1})$ و$(u\mid a^{(i)})=(p_{n-1}a^{(i)})$، على الترتيب؛ ولذلك لا تُعَدّان، وفق اصطلاحنا، مختلفتين اختلافًا جوهريًا. وتمثل (14.) مبرهنة حاصل الضرب للمحددات؛ وسنبين كيف يمكن إحلال منظومة مبرهنات حاصل الضرب للمصفوفات، المصوغة على نحو مناسب، محل (13.) و(14.) (وبالثنوية محل (15.) و(14.)). وتكون مبرهنة حاصل الضرب، عند تطبيقها مرة على مصفوفات ذات $\rho$ من الأسطر ومرة على مصفوفات ذات $(\rho-1)$ من الأسطر، كما يأتي:
\[
\pair{a^{(1)}a^{(2)}\cdots a^{(\rho)}}{xp_{\rho-1}}
=(a^{(1)}a^{(2)}\cdots a^{(\rho)}\mid xy^{(2)}\cdots y^{(\rho)})
=\sum\pm\pair{a^{(1)}}{x}\pair{a^{(2)}}{y^{(2)}}\cdots\pair{a^{(\rho)}}{y^{(\rho)}},
\]
\[
\sum\pm\pair{a^{(2)}}{y^{(2)}}\cdots\pair{a^{(\rho)}}{y^{(\rho)}}
=(a^{(2)}\cdots a^{(\rho)}\mid y^{(2)}\cdots y^{(\rho)})
=(a^{(2)}\cdots a^{(\rho)}p_{\rho-1})=(a^{(2)}\cdots a^{(\rho)}q_{n-\rho+1});
\]
ومن هذين تنتج منظومة مبرهنات حاصل الضرب الآتية، حيث $\rho=2,3,\ldots,n$:
\begin{equation}
\begin{aligned}
\pair{a^{(1)}a^{(2)}\cdots a^{(\rho)}}{xp_{\rho-1}}
 &=\sum_{i=1}^{\rho}(-1)^{i+1}\pair{a^{(i)}}{x}(a^{(1)}\cdots a^{(i-1)}a^{(i+1)}\cdots a^{(\rho)}\mid p_{\rho-1})\\
 &=\sum_{i=1}^{\rho}(-1)^{i+1}\pair{a^{(i)}}{x}(a^{(1)}\cdots a^{(i-1)}a^{(i+1)}\cdots a^{(\rho)}q_{n-\rho+1}).\footnotemark
\end{aligned}
\tag{16}
\end{equation}
\footnotetext{يمكن أيضًا عدّ المتطابقات (16.) متطابقات (13.) المنقولة من مجال يضم $\rho$ من المتغيرات.}
عندما $\rho=n$ تتحول (16.) إلى (13.)؛ وإذا تابعنا التخصيص بوضع $p_{n-1}=y^{(2)}p_{n-2}$، و$p_{n-2}=y^{(3)}p_{n-3}$، وهكذا، وطبقنا المتطابقات (16.) على التعابير $(a^{(1)}\cdots a^{(i-1)}a^{(i+1)}\cdots a^{(n)}\mid y^{(2)}p_{n-2})$، و$(a^{(1)}\cdots a^{(i-1)}a^{(i+1)}\cdots a^{(k-1)}a^{(k+1)}\cdots a^{(n)}\mid y^{(3)}p_{n-3})$، وهكذا، انتقلنا من (13.) إلى (14.). ومن ثم تكون المتطابقة (14.)، بحسب تعريفنا، مركبةً من المتطابقات (16.)؛ أو تكون منظومة (16.) مكافئةً للمتطابقتين (13.) و(14.). وتحقق منظومة (16.) أحد الشروط المطلوبة من المتطابقات العامة؛ فهي تشتمل على الصف $p_{\rho-1}$ بوصفه كمية غير قابلة للتفكيك، وتشتمل عليه في الوقت نفسه بصورة صريحة. وهي المنظومة الوحيدة التي تحقق هذا المطلب وحده. ذلك أنه لا يمكننا تكوين محددات أو حواصل ضرب مصفوفية أخرى من الصفوف $a,x,p$؛ ولا يمكن أن توجد بين هذه المحددات والحواصل علاقات مستقلة عن (16.)، وإلا لنشأت، بحذف $(a^{(1)}\cdots a^{(\rho)}\mid xp_{\rho-1})$، علاقة بين $\rho$ من التعابير، مع الشرط $\rho\leqq n$، من الشكل $(a^{(i)}\mid x)(a^{(1)}\cdots a^{(i-1)}a^{(i+1)}\cdots a^{(\rho)}q_{n-\rho+1})$؛ في حين يلزم، بحسب (13.)، أن يدخل في العلاقة ما لا يقل عن $n+1$ من هذه التعابير. وتفضي اعتبارات مماثلة إلى المنظومة الآتية، المكافئة للمتطابقتين (15.) و(14.):
\begin{equation}
\begin{aligned}
\pair{u q_{\rho-1}}{a^{(1)}\cdots a^{(\rho)}}
 &=\sum_{i=1}^{\rho}(-1)^{i+1}\pair{u}{a^{(i)}}(q_{\rho-1}\mid a^{(1)}\cdots a^{(i-1)}a^{(i+1)}\cdots a^{(\rho)})\\
 &=\sum_{i=1}^{\rho}(-1)^{i+1}\pair{u}{a^{(i)}}(p_{n-\rho+1}a^{(1)}\cdots a^{(i-1)}a^{(i+1)}\cdots a^{(\rho)}).
\end{aligned}
\tag{17}
\end{equation}
وللوصول من (16.) إلى متطابقات بين صفوف اعتباطية $A^{\rho_1},B^{\rho_2},\ldots,x,p$، نلاحظ أن هذه المتطابقات يجب أن تصير مطابقةً لـ(16.) حالما نفكك $A^{\rho_1}=a^{(1)}a^{(2)}\cdots a^{(\rho_1)}$، و$B^{\rho_2}=b^{(1)}b^{(2)}\cdots b^{(\rho_2)}$، وهكذا؛ وأنه متى كانت التعابير المفردة من نمط التعابير الداخلة في (16.)، فلا يمكن جمعها في تعابير نهائية إلا بالعمليات التي تحددها (16.). وهكذا نحصل، إذا وضعنا أيضًا
\begin{equation}
 \rho=\rho_1+\rho_2+\cdots+\rho_k\leqq n
\tag{18}
\end{equation}
على:
\begin{align*}
&\pair{a^{(1)}\cdots a^{(\rho_1)}b^{(1)}\cdots b^{(\rho_2)}\cdots k^{(1)}\cdots k^{(\rho_k)}}{xp_{\rho-1}}
  =\pair{A^{\rho_1}B^{\rho_2}\cdots K^{\rho_k}}{xp_{\rho-1}} \\
&\quad=\sum_{i=1}^{\rho_1}(-1)^{i+1}\pair{a^{(i)}}{x}(a^{(1)}\cdots a^{(i-1)}a^{(i+1)}\cdots a^{(\rho_1)}B^{\rho_2}\cdots K^{\rho_k}q_{n-\rho+1})\\
&\qquad+(-1)^{\rho_1\rho_2}\sum_{i=1}^{\rho_2}(-1)^{i+1}\pair{b^{(i)}}{x}(b^{(1)}\cdots b^{(i-1)}b^{(i+1)}\cdots b^{(\rho_2)}A^{\rho_1}\cdots K^{\rho_k}q_{n-\rho+1})\\
&\qquad+\cdots\\
&\qquad+(-1)^{(\rho_1+\cdots+\rho_{k-1})\rho_k}\sum_{i=1}^{\rho_k}(-1)^{i+1}\pair{k^{(i)}}{x}(k^{(1)}\cdots k^{(i-1)}k^{(i+1)}\cdots k^{(\rho_k)}A^{\rho_1}\cdots (K-1)^{\rho_{k-1}}q_{n-\rho+1}).
\end{align*}
وكل مجموع من هذه المجاميع (لا مجرد جزء منه) يشتمل الآن فعلًا على الصفوف $A^{\rho_1},B^{\rho_2}\ldots$ بوصفها كميات غير قابلة للتفكيك؛ إذ يحقق الشروط الضرورية والكافية لذلك:\footnote{يورد \srcLTR{Capelli}، في المرجع السابق، \S\ 23، شروطًا كافية: $\bigl(a^{(k)}\mid \frac{\partial}{\partial a^{(i)}}\bigr)=0$، $(k>i,\ i=1,2,\ldots,\rho_1-1)$. ومن هذه تُستخلص الشروط الضرورية والكافية بأيسر طريق بتطبيق عملية قوس بواسون وفقًا لـ\srcLTR{Wellstein, Von den Differentialgleichungen der projektiven Invarianten, Math. Ann. 67 (1909), \S\ 3}.}
\begin{equation}
 \left(a^{(i+1)}\middle|\frac{\partial}{\partial a^{(i)}}\right)=0;\ \ldots\ ;
 \left(k^{(j+1)}\middle|\frac{\partial}{\partial k^{(j)}}\right)=0.
 \quad (i=1,2,\ldots,\rho_1-1;\ j=1,2,\ldots,\rho_k-1)
\tag{19}
\end{equation}
ونبين أنها تشتمل أيضًا على جميع الصفوف صراحةً. فإذا وضعنا $B^{\rho_2}\cdots K^{\rho_k}q_{n-\rho+1}=q_{n-\rho_1+1}$، حصلنا، بحسب (16.) و(10.)، على:
\begin{align*}
&\sum(-1)^{i+1}\pair{a^{(i)}}{x}(a^{(1)}\cdots a^{(i-1)}a^{(i+1)}\cdots a^{(\rho_1)}q_{n-\rho+1})\\
&=\pair{A^{\rho_1}}{xp_{\rho_1-1}}\\
&=(A_{n-\rho_1}xp_{\rho_1-1})=(-1)^{(\rho_1+1)(n+1)}\pair{q_{n-\rho_1+1}}{A_{n-\rho_1}x}\\
&=(-1)^{(\rho_1+1)(n+1)}\pair{B^{\rho_2}\cdots K^{\rho_k}q_{n-\rho+1}}{A_{n-\rho_1}x}.
\end{align*}
وإذا حسبنا سائر المجاميع على نحو مماثل، ثم وضعنا بعد ذلك $B^{\rho_2}=A^{\rho_2}$، و$C^{\rho_3}=A^{\rho_3},\ldots$، بما أن أدلة الوزن قد ميزت الصفوف بعضها من بعض، حصلنا على المتطابقات المطلوبة:
\begin{equation}
\begin{aligned}
\pair{A^{\rho_1}A^{\rho_2}\cdots A^{\rho_k}}{xp_{\rho-1}}
&=\sum_{i=1}^{k}(-1)^{\delta_i}\pair{A^{\rho_1}\cdots A^{\rho_{i-1}}A^{\rho_{i+1}}\cdots A^{\rho_k}q_{n-\rho+1}}{A_{n-\rho_i}x},\\
\delta_i&=(\rho_1+\rho_2+\cdots+\rho_{i-1})\rho_i+(\rho_i+1)(n+1).
\end{aligned}
\tag{20}
\end{equation}
ومن (17.) نحصل بالثنوية على:
\begin{equation}
\begin{aligned}
\pair{u q_{\sigma-1}}{A_{\sigma_1}A_{\sigma_2}\cdots A_{\sigma_k}}
&=\sum_{i=1}^{k}(-1)^{\eps_i}\pair{u A^{n-\sigma_i}}{p_{n-\sigma+1}A_{\sigma_1}\cdots A_{\sigma_{i-1}}A_{\sigma_{i+1}}\cdots A_{\sigma_k}},\\
\eps_i&=(\sigma_1+\sigma_2+\cdots+\sigma_{i-1})\sigma_i+(\sigma_i+1)(n+\sigma).
\end{aligned}
\tag{21}
\end{equation}
\emph{المبرهنة \srcLTR{II}.} تستنفد (20.) و(21.)، على الترتيب، جميع المتطابقات غير القابلة للتفكيك بين الصفوف $A^{\rho_i},x,p$ وبين الصفوف $A_{\sigma_i},u,q$؛ وتستنفدان في الوقت نفسه جميع المتطابقات غير القابلة للتفكيك التي تقبل تمثيلًا صريحًا من غير إجراء الإحلال (11.).

ينتج الجزء الأول من المبرهنة من الاعتبارات المتعلقة باشتقاق (20.) و(21.)؛ وينتج الجزء الثاني من استحالة وجود متطابقة بين صفين من صفوف الرموز:
\[
\pair{q_\rho A^\sigma}{B_{\rho+\sigma-\tau}p_\tau}
=\eps_1\pair{q_\rho B^{n-\rho-\sigma+\tau}}{A_{n-\sigma}p_\tau}
+\eps_2\pair{q_\rho q_{n-\tau}}{B_{\rho+\sigma-\tau}A_{n-\sigma}},
\]
حيث $\rho\leqq\tau\leqq\sigma$، ولا يكون أي من الصفوف ذا وزن $1\cdot(n-1)$. (وترد الفروض الأخرى المتعلقة بـ$\rho,\tau,\sigma$ إلى هذه الفروض بتبديل تسميات الصفوف.) وتنتج هذه الاستحالة، مثلا، من نشر (8.) عند التخصيص $q_\rho=lM^{\rho-1}$، و$q_{n-\tau}=lN^{n-\tau-1}$، إذا وضعنا $l_1=1$، و$M^{2,3,\ldots,\rho}=1$، و$A^{\rho+1,\ldots,\rho+\sigma}=1$، وجعلنا جميع العناصر الأخرى في الصفوف $l,M,A$ مساوية للصفر، مع إبقاء $N$ و$B$ اعتباطيين. ولما كانت المتطابقات بين صفوف اعتباطية تنتقل، بفك صف $p_\tau$ إلى $y^{(1)}y^{(2)}\cdots y^{(\tau)}$، إلى متطابقات مركبة من (20.)، فإن الجزء الثاني من الادعاء يتم برهانه متى أمكن توليد (20.) من متطابقات لا تشتمل إلا على صفين من صفوف الرموز. وبالفعل، نحصل من
\begin{equation}
\pair{A^{\rho_1}A^{\rho_2}}{xp_{\rho-1}}
=\eps\pair{A^{\rho_1}q_{n-\rho+1}}{A_{n-\rho_2}x}
+\pair{A^{\rho_1}}{xp_{\rho_1-1}},
\quad \text{حيث }p_{\rho_1-1}\sim A^{\rho_2}q_{n-\rho+1},
\tag{20a}
\end{equation}
بالتخصيص $A^{\rho_1}=B^{\sigma_1}B^{\sigma_2}$، أي بواسطة
\[
\pair{B^{\sigma_1}B^{\sigma_2}}{xp_{\rho_1-1}}
=\eps\pair{B^{\sigma_1}q_{n-\rho_1+1}}{B_{n-\sigma_2}x}
+\pair{B^{\sigma_1}}{xp_{\sigma_1-1}},
\quad \text{حيث } p_{\sigma_1-1}\sim B^{\sigma_2}q_{n-\rho_1+1},
\]
متطابقة من نمط (20.) في ثلاثة صفوف من صفوف الرموز؛ وبمزيد من التخصيصات لـ$B^{\sigma_1}=C^{\tau_1}C^{\tau_2}$، وهكذا، نحصل على المتطابقات العامة (20.). ومن ثم يتيح عدم وجود تمثيل صريح لمتطابقة بين صفين من صفوف الرموز وصفين اعتباطيين من صفوف المتغيرات الاستنتاج بأنه لا يوجد تمثيل صريح كهذا عند وجود عدد اعتباطي من صفوف الرموز، إذا لم يكن بين صفوف المتغيرات أي صف $x$ أو $u$.

\clearpage
\section*{\S\ 4. متطابقات التفكيك.}
ننظر الآن في الحالة الأعم للمتطابقات بين صفوف اعتباطية من صفوف الرموز والمتغيرات، وهي المتطابقات التي يظهر في تعبيراتها النهائية، من غير إجراء الإحلال (11.)، صف $p_\tau$ مفككا إلى صفوفه المفردة $y^{(1)}\cdots y^{(\tau)}$؛ ونسمّي هذه المتطابقات «متطابقات التفكيك». وبناء على الملاحظة الواردة في نهاية الفقرة السابقة، نحددها أولا في حالة صفين فقط من صفوف الرموز؛ أي إننا ننشر التعبير $\pair{q_\rho A^\sigma}{B_{\rho+\sigma-\tau}p_\tau}$. ونضع عندئذ:
\begin{equation}
\begin{aligned}
q_{\rho-\alpha}^{(i_1\ldots i_\alpha)}&\sim p_{n-\rho}\,y^{(i_1)}\cdots y^{(i_\alpha)},\qquad
p_{\tau-\alpha}^{(i_1\ldots i_\alpha)}=y^{(i_{\alpha+1})}\cdots y^{(i_\tau)},\qquad
p_\tau=y^{(1)}\cdots y^{(\tau)}.
\end{aligned}
\tag{22}
\end{equation}
وبالنسبة إلى الأعداد $\rho,\sigma,\tau$ ينبغي تمييز أربع حالات:
\[
\begin{array}{llll}
1)&\rho+\sigma\leqq n,& \tau\leqq\rho,& \tau\leqq\sigma;\\[2pt]
2)&\rho+\sigma\leqq n,& \tau\geqq\rho,& \tau\leqq\sigma;\\[2pt]
3)&\rho+\sigma\leqq n,& \tau\leqq\rho,& \tau\geqq\sigma;\\[2pt]
4)&\rho+\sigma\leqq n,& \tau\geqq\rho,& \tau\geqq\sigma,
\end{array}
\]
نختار من بينها الحالة الأولى لنعالجها، ثم نصف الحالات الباقية بإيجاز. فبحسب (20.) و(10.) و(9.)، إذا فككنا الصف $y^{(1)}y^{(2)}\cdots y^{(\tau)}B_{\rho+\sigma-\tau}$ إلى $y^{(1)}$ و$y^{(2)}\cdots y^{(\tau)}B_{\rho+\sigma-\tau}$، نحصل على:
\begin{equation}
\begin{aligned}
\pair{q_\rho A^\sigma}{y^{(1)}y^{(2)}\cdots y^{(\tau)}B_{\rho+\sigma-\tau}}
&=\pair{q_{\rho-1}^{(1)}A^\sigma}{y^{(2)}\cdots y^{(\tau)}B_{\rho+\sigma-\tau}}\\
&\quad+(-1)^\rho\pair{q_\rho[A_{n-\sigma}y^{(1)}]^{\sigma-1}}{y^{(2)}\cdots y^{(\tau)}B_{\rho+\sigma-\tau}}.
\end{aligned}
\tag{23}
\end{equation}
وللحد الأخير في (23.) صورة التعبير الأصلي نفسها تماما؛ غير أن $\sigma$ استبدلت بـ$\sigma-1$، و$\tau$ استبدلت بـ$\tau-1$، ومن ثم يمكن تطبيق المعادلة (23.) عليه من جديد. ولذلك، بما أن $\tau\leqq\sigma$، يمكننا تطبيق العملية (23.) $\tau$ مرة، فنصل، مع مراعاة (10.)، إلى العلاقة:
\begin{equation}
\begin{aligned}
\pair{q_\rho A^\sigma}{p_\tau B_{\rho+\sigma-\tau}}
&=(-1)^{\rho\sigma+(\sigma+\tau)(n+1)}
  \pair{q_\rho B^{n-\rho-\sigma+\tau}}{A_{n-\sigma}p_\tau}\\
&\quad+\sum_{i_1=1}^{\tau}(-1)^{\rho(i_1-1)}
 \pair{q_{\rho-1}^{(i_1)}[A_{n-\sigma}y^{(1)}\cdots y^{(i_1-1)}]^{\sigma-i_1+1}}
 {y^{(i_1+1)}\cdots y^{(\tau)}B_{\rho+\sigma-\tau}}.
\end{aligned}
\tag{24}
\end{equation}
وكل حد تحت علامة المجموع من نمط التعبير الأصلي من جديد؛ فقد استبدلت $\rho$ بـ$\rho-1$، و$\sigma$ بـ$\sigma-i_1+1$، و$\tau$ بـ$\tau-i_1$، وبذلك يتحقق الشرط 1) على نحو أقوى. ولذلك، بما أن $\tau\leqq\rho$، يمكننا تطبيق العملية (24.) $\tau$ مرة، فنحصل على متطابقة التفكيك المطلوبة:
\begin{equation}
\begin{aligned}
\pair{q_\rho A^\sigma}{p_\tau B_{\rho+\sigma-\tau}}
&=\eps_{i_0}\pair{q_\rho B^{n-\rho-\sigma+\tau}}{A_{n-\sigma}p_\tau}\\
&\quad+\sum_{i_1=1}^{\tau}\eps_{i_1}
 \pair{q_{\rho-1}^{(i_1)}B^{n-\rho-\sigma+\tau}}{A_{n-\sigma}p_{\tau-1}^{(i_1)}}\\
&\quad+\sum_{i_2=i_1+1}^{\tau}\eps_{i_2}
 \pair{q_{\rho-2}^{(i_1i_2)}B^{n-\rho-\sigma+\tau}}{A_{n-\sigma}p_{\tau-2}^{(i_1i_2)}}+\cdots\\
&\quad+\eps_{i_\tau}\pair{q_{\rho-\tau}^{(i_1\ldots i_\tau)}B^{n-\rho-\sigma+\tau}}{A_{n-\sigma}}\\
&=\sum_{\alpha=0}^{\tau}\sum_{i_\alpha=i_{\alpha-1}+1}^{\tau}\eps_{i_\alpha}
 \pair{q_{\rho-\alpha}^{(i_1\ldots i_\alpha)}B^{n-\rho-\sigma+\tau}}{A_{n-\sigma}p_{\tau-\alpha}^{(i_1\ldots i_\alpha)}}.
\end{aligned}
\tag{25}
\end{equation}
أما $\eps$ فنحصل من (24.) على:
\begin{equation}
\begin{aligned}
\eps_{i_\alpha}&=(-1)^{\delta_{i_\alpha}},\qquad
\delta_{i_\alpha}=\sum_{\beta=1}^{\alpha}(\rho-\beta+1)(i_{\beta-1}+i_\beta+1)\\
&\quad +(\rho-\alpha)(\sigma+i_\alpha+\alpha)+(\sigma+\tau+\alpha)(n+1),
\end{aligned}
\tag{26}
\end{equation}
حيث ينبغي وضع $i_0=0$. وينبغي فهم علامة المجموع في (25.) على النحو الآتي: تضاف إلى كل تركيبة $i_1i_2\ldots i_{\alpha-1}$، حيث $i_1<i_2<\cdots<i_{\alpha-1}$، جميع قيم $i_\alpha$ التي تحقق $i_{\alpha-1}<i_\alpha\leqq\tau$. وكما يسهل التحقق مع مراعاة (26.)، فإن المجاميع المفردة في (25.) تحقق المعادلات التفاضلية المناظرة لـ(19.):
\begin{equation}
 \left(\frac{\partial}{\partial y^{(i)}}\middle|y^{(i+1)}\right)=0,
 \qquad (i=1,2,\ldots,\tau-1),
\tag{27}
\end{equation}
ومن ثم لا تعتمد إلا على الصف $p_\tau$. وعلى ذلك تكون متطابقة التفكيك متطابقة غير قابلة للتفكيك بين الصفوف $A,B,q,p$؛ وسنعود إلى التمثيل الصريح في الفقرة التالية.

في الحالة 2) يمكن كذلك تطبيق العملية (23.) $\tau$ مرة، غير أن العملية (24.) تتوقف بعد المرة $\rho$. ولذلك ينبغي في التعبير النهائي من (25.) استبدال علامة المجموع الأولى بـ$\sum_{\alpha=0}^{\rho}$. أما في الحالتين 3) و4)، فلا يمكن إجراء العملية (23.) إلا $\sigma$ مرة؛ وتحل عندئذ محل (24.) العلاقة:
\begin{equation}
\begin{aligned}
\pair{q_\rho A^\sigma}{p_\tau B_{\rho+\sigma-\tau}}
&=(-1)^{\rho\sigma}
 \pair{q_{\rho-(\tau-\sigma)}^{(\sigma+1\ldots\tau)}B^{n-\rho+\tau-\sigma}}{A_{n-\sigma}p_{\tau-(\tau-\sigma)}^{(\sigma+1\ldots\tau)}}\\
&\quad+\sum_{i_1=1}^{\sigma}(-1)^{\rho(i_1-1)}
 \pair{q_{\rho-1}^{(i_1)}[A_{n-\sigma}y^{(1)}\cdots y^{(i_1-1)}]^{\sigma-i_1+1}}
 {y^{(i_1+1)}\cdots y^{(\tau)}B_{\rho+\sigma-\tau}};
\end{aligned}
\tag{28}
\end{equation}
وتكرار (28.) $(\tau-\sigma)$ مرة، بحيث تصير علامة المجموع $\sum_{i_\beta=i_{\beta-1}+1}^{\sigma+\beta-1}$ ــ كما يتبين من تطبيق (23.) $(\sigma-i_1+1)$ مرة على الحدود الواقعة تحت علامة المجموع في (28.) ــ يعطي جميع حدود المجموع المفرد المقابل للدليل $\alpha=\tau-\sigma$ في (25.)، مع الإشارة المناسبة. وعندئذ تنتقل الحالتان 3) و4) إلى الحالتين 1) و2)، لأن العددين المقابلين لـ$\sigma,\tau$ يصيران: $\sigma'=\sigma-i_{\tau-\sigma}+\tau-\sigma$، و$\tau'=\tau-i_{\tau-\sigma}=\sigma'$؛ ومع مواصلة الإجراء يصبح $\tau'<\sigma'$. ولذلك ينبغي في التعبير النهائي من (25.) استبدال علامة المجموع الأولى بـ$\sum_{\alpha=\tau-\sigma}^{\tau,\rho}$.

وبالمثل تنتج من (21.) متطابقة التفكيك الثنوية، التي يُفكك فيها الصف $q_\rho$ إلى الصفوف المفردة $v^{(1)},\ldots,v^{(\rho)}$. فإذا وضعنا
\begin{equation}
 \dot p_{\tau-\beta}^{(i_1\ldots i_\beta)}\sim v^{(i_1)}\cdots v^{(i_\beta)}q_{n-\tau};\qquad
 \dot q_{\rho-\beta}^{(i_1\ldots i_\beta)}=v^{(i_{\beta+1})}\cdots v^{(i_\rho)};\qquad
 q_\rho=v^{(1)}\cdots v^{(\rho)},
\tag{29}
\end{equation}
نحصل على:
\begin{equation}
 \pair{q_\rho A^\sigma}{p_\tau B_{\rho+\sigma-\tau}}
 =\sum_{\beta=\substack{0,\tau-\sigma}}^{\substack{\tau,\rho}}
   \sum_{i_\beta=i_{\beta-1}+1}^{\rho}\eps
\pair{\dot q_{\rho-\beta}^{(i_1\ldots i_\beta)}B^{n-\rho-\sigma+\tau}}{A_{n-\sigma}\dot p_{\tau-\beta}^{(i_1\ldots i_\beta)}} ,
\tag{30}
\end{equation}
حيث يمتد دليل الجمع $\beta$ من الأكبر بين الحدين السفليين إلى الأصغر بين الحدين العلويين. وجميع حدود كل مجموع مفرد في (25.) و(30.) أشكال قطبية لحدودية واحدة، نشأت بعمليتي أخذ الأشكال القطبية الآتيتين:
\[
\left(\frac{\partial}{\partial p_{\tau-\alpha}}\middle|p_{\tau-\alpha}^{(i_1\ldots i_\alpha)}\right),\qquad
\left(q_{\rho-\alpha}^{(i_1\ldots i_\alpha)}\middle|\frac{\partial}{\partial q_{\rho-\alpha}}\right).
\]
وللتعبير عن ذلك، نرمز بـ$\Pi$ إلى تركيب من عمليات أخذ الأشكال القطبية هذه، مضروبة في ثوابت، فنحصل بدلا من (25.) و(30.) على العلاقة:
\begin{equation}
 \pair{q_\rho A^\sigma}{p_\tau B_{\rho+\sigma-\tau}}
 =\sum_{\alpha=\substack{0,\tau-\sigma}}^{\substack{\tau,\rho}}
 \Pi\pair{q_{\rho-\alpha}B^{n-\rho-\sigma+\tau}}{A_{n-\sigma}p_{\tau-\alpha}}.
\tag{31}
\end{equation}

\section*{\S\ 5. متطابقات التفكيك في صورة صريحة.}
إذا خصصنا في (25.)، في الحالة 4)، $\tau$ إلى $\sigma+\rho$، نحصل على:
\begin{equation}
 \pair{q_\rho A^\sigma}{y^{(1)}\cdots y^{(\rho+\sigma)}}
 =\sum_{i_\rho=i_{\rho-1}+1}^{\tau}
 \eps_{i_\rho}\pair{q_\rho}{y^{(i_1)}\cdots y^{(i_\rho)}}
 \pair{A^\sigma}{y^{(i_{\rho+1})}\cdots y^{(i_{\rho+\sigma})}},
\tag{32}
\end{equation}
حيث تصبح $\eps_{i_\rho}=(-1)^{\delta_{i_\rho}}$، و$\delta_{i_\rho}=\rho(\rho+1)/2+i_1+\cdots+i_\rho$. وهذه صيغة أعم لمبرهنة حاصل الضرب للمصفوفات من (16.) و(17.)، وتنتج بنشر لابلاس لمحدد مصفوفة الجداء
\[
 \sum\pm(v^{(1)}y^{(1)})\cdots(v^{(\rho)}y^{(\rho)})(a^{(1)}y^{(\rho+1)})\cdots(a^{(\sigma)}y^{(\rho+\sigma)}).
\]
فمبرهنة حاصل الضرب الأعم مؤلفة إذن من المتطابقتين (20.) و(21.)، كل بحسب الحال، وبذلك يتضح استخلاص الصورتين الخاصتين (16.) و(17.).

توفر العلاقة (32.) الآن وسيلة التمثيل الصريح لمتطابقات التفكيك. ولهذا نعدّ الحدوديات الواردة في (31.) حدوديات أساسية؛ أي نجري الإحلال (11.)، وهو لا يغير شيئا في التمثيل الصريح بدلالة الصفين $A,B$، لأن الصفوف المستحدثة $R$، بحسب (8.)، يمكن التعبير عنها بالصفين $A,B$ نفسيهما، لا بصفوفهما الجزئية $a^{(1)}a^{(2)}\ldots$ و$b^{(1)}b^{(2)}\ldots$. فإذا كان
\begin{equation}
 \pair{q_{\rho-\alpha}B^{n-\rho-\sigma+\tau}}{A_{n-\sigma}p_{\tau-\alpha}}
 =\bigl(R_{\rho-\alpha}p_{n-\rho+\alpha}\bigr)
  \pair{R^{\tau-\alpha}}{p_{\tau-\alpha}},
\tag{33}
\end{equation}
فإن المجموع المفرد من (25.) المقابل للدليل $\alpha$ يصير:
\begin{equation}
\begin{aligned}
&\sum \eps_{i_\alpha}\bigl(R_{\rho-\alpha}p_{n-\rho}y^{(i_1)}\cdots y^{(i_\alpha)}\bigr)
 \pair{R^{\tau-\alpha}}{y^{(i_{\alpha+1})}\cdots y^{(i_\tau)}}\\
&\quad=\sum \eps_{i_\alpha}\bigl([R_{\rho-\alpha}p_{n-\rho}]^\alpha y^{(i_1)}\cdots y^{(i_\alpha)}\bigr)
 \pair{R^{\tau-\alpha}}{y^{(i_{\alpha+1})}\cdots y^{(i_\tau)}}\\
&\quad=\eps\pair{[R_{\rho-\alpha}p_{n-\rho}]^\alpha R^{\tau-\alpha}}{p_\tau}
 =\eps\pair{R^{\tau-\alpha}q_{n-\tau}}{R_{\rho-\alpha}p_{n-\rho}}.
\end{aligned}
\tag{34}
\end{equation}
وبذلك صار لدينا بالفعل تعبير ختامي صريح؛ وتدل الصورة المتناظرة التي يدخل فيها $q_\rho$ و$p_\tau$ على أن (30.) يقود إلى التعبير الختامي نفسه، فهو إذن مستقل عن التفكيك الأصلي. ولا يبقى فرق بين (25.) و(30.) إلا ما دامت الصفوف $y$ و$v$ ذات دلالة مستقلة، أي ما دامت متطابقة التفكيك تنتقل، بحسب تعريفنا، إلى متطابقة قابلة للتفكيك. وللوصول إلى متطابقات ذات عدد أكبر من صفوف الرموز، لا يلزم، بحسب خاتمة \S\ 3، إلا تفكيك صف $q_\rho$ إلى $q_{\rho_1}C^{\sigma_1}$، أو صف $p_\tau$ إلى $p_{\tau_1}C_{\tau-\tau_1}$. والتعبيران الناشئان:
\begin{equation}
 \pair{q_{\rho_1}C^{\sigma_1}}{[R^{\tau-\alpha}q_{n-\tau}]_\alpha R_{\rho_1+\sigma_1-\alpha}},
 \qquad
 \pair{[R_{\rho-\alpha}p_{n-\rho}]^\alpha R^{\tau-\alpha}}{p_{\tau_1}C_{\tau-\tau_1}},
\tag{35}
\end{equation}
هما الآن من النمط نفسه الذي يظهر عند وجود صفين من صفوف الرموز؛ ولذلك يقبلان، بعد إحلال مناظر لـ(33.)، تمثيلا صريحا من جديد، ومن ثم يفضي التكرار إلى تمثيل صريح عند عدد اعتباطي من الصفوف. ويجدر أيضا أن نلاحظ أن المجموع الأول والمجموع الأخير في (25.) و(30.)، كل بحسب الحال، يعطيان التمثيل الصريح، من غير تطبيق (33.)، بالصيغة (32.) وحدها، أو يقتصران أصلا على حد واحد يشتمل على جميع الصفوف صراحة. ويمكن عدّ العلاقة الناشئة من تفكيك $q_\rho$ إلى $q_{\rho_1}C^{\sigma_1}$، وهكذا، في (32.) مبرهنة حاصل الضرب في أعم صورها.

وخلاصة ما سنبينه هي:

\emph{المبرهنة \srcLTR{III}:} بالمتطابقات
\begin{equation}
 \pair{q_\rho A^\sigma}{p_\tau B_{\rho+\sigma-\tau}}
 =\sum_{\alpha=\substack{0,\tau-\sigma}}^{\substack{\tau,\rho}}
 \eps\pair{R^{\tau-\alpha}q_{n-\tau}}{R_{\rho-\alpha}p_{n-\rho}},
\tag{36}
\end{equation}
حيث تتحدد الصفوف $R$ بـ(33.)، وبالمتطابقات الناشئة منها بتفكيك $q_\rho,p_\tau$ إلى $q_{\rho_1}C^{\sigma_1}$، وهكذا، تُستنفد جملة المتطابقات غير القابلة للتفكيك.

والمتطابقات الناشئة هي بالفعل أعم المتطابقات من نوعها، لأن الصفوف المفردة لم تعد تخضع لأي قيد من حيث الوزن أو العدد، بخلاف ما كان عليه الحال في (20.) و(21.). كذلك لا توجد متطابقات أخرى بين هذه الصفوف؛ لأنه عند تفكيك صف $p_\tau$ إلى $y^{(1)}\cdots y^{(\tau)}$ تكون أعم المتطابقات مؤلفة من (20.)، أي، بحسب خاتمة \S\ 3، من (20a.)؛ والتركيب الذي نوقش عقب (16.) هو بعينه ما أُجري في \S\ 4. أما الانتقال إلى عدد اعتباطي من الصفوف فينتج، كما تبين من مناقشة (20a.)، بتطبيق العملية نفسها دائما على التعبيرات الناشئة من تفكيك $q_\rho$ إلى $q_{\rho_1}C^{\sigma_1}$، وهكذا. وينتج من ذلك أيضا أن الانتقال المباشر من (20.) بدلا من (20a.) لا يعطي شيئا جديدا؛ ويمكن التحقق من هذا حسابيا بسهولة بأن نخصص مرة صف $q_\rho$ في (25.)، ونجري مرة أخرى الانتقال من (20.) بطريقة \S\ 4؛ ففي المرتين تنشأ المجاميع نفسها، وتنتقل هذه المجاميع، بتطبيق مبرهنة حاصل الضرب العامة (انظر أعلاه)، إلى المتطابقات الناتجة من (36.). وتقابل المتطابقتان (20.) و(21.) الحالتين الخاصتين $\tau=1$ و$\rho=1$؛ فلا يظهر عندئذ في كل حالة إلا مجموعان مفردان، ومن ثم نحصل، وفقا للملاحظة السابقة، على تمثيل صريح من غير الإحلال (33.)، وهو ما يوافق ما سبق.

وفي الختام، نذكر بإيجاز حالتين خاصتين أخريين لمتطابقة التفكيك. فإذا وضعنا في (31.): $\tau=\sigma$، و$\rho=n-\sigma$، وسمّينا الصف $p_\tau$ باسم $p'_\sigma\sim q'_{n-\sigma}$، حصلنا على:
\begin{equation}
\begin{aligned}
\pair{A^\sigma}{p_\sigma}\pair{B^\sigma}{p'_\sigma}
-\pair{A^\sigma}{p'_\sigma}\pair{B^\sigma}{p_\sigma}
=\sum_{\alpha=1}^{\sigma,n-\sigma}
\Pi\pair{q_{n-\sigma-\alpha}B^\sigma}{A_{n-\sigma}p_{\sigma-\alpha}}.
\end{aligned}
\tag{37}
\end{equation}
وهذه هي صيغة كليبش\footnote{كليبش، المرجع نفسه، ص. 25--32. يثبت كليبش أن الحدود المفردة في الأطراف اليمنى لا تعتمد إلا على الصفين $A,B$؛ وينتج التمثيل الصريح من تطبيق (32.) على عباراته $\Phi_h$.} لتوسعة حدودية ذات صفين متوافقي التدرج $p_\sigma,p'_\sigma$ بدلالة الملازمات الأولية الموافقة للتغير. وعندئذ تكون الملازمات الأولية الموافقة للتغير العائدة إلى الحدودية $\pair{A^\sigma}{p_\sigma}^m\pair{B^\sigma}{p'_\sigma}^n$ هي الآتية:
\begin{equation}
 \left\{\sum_{\alpha=1}^{\sigma,n-\sigma}
 \pair{q_{n-\sigma-\alpha}B^\sigma}{A_{n-\sigma}p_{\sigma-\alpha}}
 \right\}^{\rho}
 \pair{A^\sigma}{p_\sigma}^{m-\rho}\pair{B^\sigma}{p'_\sigma}^{n-\rho},
 \qquad (\rho=0,1,\ldots,m,n).
\tag{38}
\end{equation}
وهي تنشأ من الحدودية الأساسية بما سنسمّيه اختصارًا \glqq{}الطي المتوافق التدرج ذي النقص $\alpha$\grqq{} (قارن \S\ 2).

وثانيًا، إذا وضعنا: $\tau=\sigma-\lambda$، و$\rho=n-\sigma-\lambda$؛ و$B^\sigma=A^\sigma$، حصلنا على:
\begin{equation}
\begin{aligned}
\pair{q_{n-\sigma-\lambda}A^\sigma}{A_{n-\sigma}p_{\sigma-\lambda}}
&=(-1)^\lambda\pair{q_{n-\sigma-\lambda}A^\sigma}{A_{n-\sigma}p_{\sigma-\lambda}}\\
&\quad+(-1)^{(n-\sigma)(\sigma-\lambda)}
 \sum_{\alpha=1}^{\sigma-\lambda,n-\sigma-\lambda}
 \Pi\pair{q_{n-\sigma-\lambda-\alpha}A^\sigma}{A_{n-\sigma}p_{\sigma-\lambda}}.
\end{aligned}
\tag{39}
\end{equation}
فتترتب، كما لوحظ في \S\ 2، الشروط (12.) المميِّزة لصفوف الرموز، عند كون النقص $\lambda$ فرديًا، على الشروط ذات النقص الأعلى. وفي حالة الحدودية الخطية $\pair{A^\sigma}{p_\sigma}\equiv\pair{B^\sigma}{p_\sigma}$، يترتب من (39.) أن تكويناتها اللامتغيرة غير القابلة للاختزال، التربيعية في المعاملات، تُرَدّ إلى تكوينات ذات نقص زوجي؛ وهذا يوافق الواقعة المعروفة أن الحدودية الخطية في المجالين الثنائي والثلاثي لا تملك تكوينات لامتغيرة.

ولنلاحظ أن المتطابقات العامة اشتُقت في الأصل على افتراض أن الصفوف المفردة تحقق الشروط (12.). غير أنه لما كانت هذه الصفوف المفردة لا تدخل في المتطابقات إلا خطيًا، فإن هذه المتطابقات تصح أيضًا للصفوف الأعمّ التي لا تحقق الشروط (12.).

\section*{\S\ 6. التمثيل الصريح للتكوينات اللامتغيرة.\\عمليات الطي وعمليات الاشتقاق.}
تتيح لنا نتائج الفقرات السابقة أن نُتمّ المبرهنة التي صيغت في \S\ 2 بشأن قابلية التمثيل الرمزي (المبرهنة الأساسية الأولى)، وذلك بإضافة قابلية التمثيل الصريح، إذ نبرهن ما يأتي:

\emph{المبرهنة \srcLTR{IV}:} ``يمكن تمثيل جميع التكوينات اللامتغيرة لحدوديات أساسية اعتباطية تمثيلًا صريحًا بجداءات المصفوفات؛ أي إن التعابير النهائية لا تدخل فيها، عدا صفوف المتغيرات $p_\rho$، إلا الصفوف $S^\sigma,\ldots,T^\tau$ للحدوديات الأصلية، أو الصفوف $P,Q,R$ المشتقة منها بالتعويض (9.) أو (11.).''

لإثبات ذلك، نفترض أن تكوينًا لامتغيرًا يعتمد، إلى جانب صفوف المتغيرات، على الصفوف $S^\sigma,\ldots,T^\tau$؛ ولتكن هذه الصفوف قد فُككت بما يكفي إلى صفوف مفردة: $S^\sigma=S^{\sigma_1}\cdots S^{\sigma_i},\ldots,T^\tau=T^{\tau_1}\cdots T^{\tau_k}$، بحيث يتيسر تمثيلها بالمحددات. ولتكن الصفوف مرتبة وفق ترتيب $S^\sigma,\ldots,T^\tau$ اعتباطي في ذاته، لكنه مثبت على نحو معين. ونختار الآن تجميعًا من المحددات يعتمد على الصف الكلي الأول $S^\sigma=S^{\sigma_1}\cdots S^{\sigma_i}$، لا على صفوف جزئية منفردة منه فحسب. غير أن هذا الاعتماد، بحسب \S\ 3--\S\ 5، نتيجة للمتطابقات العامة. فإذا فككنا صف المتغيرات $p_\rho$ إلى الصفين $y,v$، أو فككنا أحد الصفوف اللاحقة $T$ إلى الصفوف المفردة $t$ أو $\tau$، فإن أعم المتطابقات تكون مركبة من (20.) و(21.). لكن هذه المتطابقات الأخيرة تؤدي دائمًا إلى تعبير يشتمل صراحة على الصف $S^\sigma=S^{\sigma_1}\cdots S^{\sigma_i}$، هو الطرف الأيسر من (20.) و(21.)؛ إذ يتبين من (19.) أن المجموع الكامل القائم في الطرف الأيمن هو وحده الذي يملك خاصية الاعتماد على $S^\sigma$ وحده، لا جزء منه يقابل حدًا من (20a.). ومع متابعة الإجراء، تظل جميع الصفوف التي دخلت صراحةً مرة واحدة ظاهرةً صراحةً؛ ولا يمكن أن يقتضي تطبيق المتطابقات إلا، ربما، جمع عدة صفوف في صف واحد، أي التعويض (9.). فإذا لم يبق في النهاية إلا صف واحد $T^\tau$ ينبغي إظهاره في الصيغة الصريحة، أمكن وقوع حالتين. فقد يبقى صف متغيرات حر، أي غير مرتبط بصفوف أخرى بالتعويض (9.)؛ وعندئذ يمكن إتمام الإجراء بتفكيكه إلى الصفوف الجزئية $y$ أو $v$، فتنشأ، كما في (31.)، أشكال قطبية لحدودية واحدة أو لعدة حدوديات تشتمل صراحة على جميع الصفوف. أما إذا لم يعد يظهر صف متغيرات حر، فإننا نتعرف، من كيفية نشوء العبارات، في مجموع العبارات التي تتضمن الصفوف المفردة $t$ أو $\tau$ ولكنها تعتمد على الصف $T^\tau$، حدود متطابقة تفكيك؛ وهذه تتيح دائمًا، بحسب \S\ 5 وباستخدام التعويض (11.)، التمثيل الصريح في جميع الصفوف. وبذلك ثبتت المبرهنة السابقة على وجه العموم. ونلاحظ أيضًا أنه حين لا يظهر سوى صفين من صفوف الرموز، لا يقع التعويض (11.) قط؛ ذلك أنه -- بسبب $\sigma,n-\sigma,\tau,n-\tau<n$ -- لا يمكن ألا تدخل صفوف متغيرات إلا إذا جُمع صفا الرموز في محدد واحد، فيظهران إذن صراحةً؛ أما حالما تدخل صفوف المتغيرات، فإن (34.) و(33.) تبينان أن التعويض (11.) يغدو زائدًا.

ينتج مما تقدم أنه يكفي، لتوليد جميع التكوينات اللامتغيرة، ضم صفوف الرموز، مع إضافة صفوف متغيرات، لتكوين جميع جداءات المصفوفات الممكنة. وأعم جداء للمصفوفات هو الآن من الصورة:
\begin{equation}
 \pair{P^{\mu_1}\cdots P^{\mu_i}}{Q_{\nu_1}\cdots Q_{\nu_k}};\qquad \sum\mu=\sum\nu,
\tag{40}
\end{equation}
حيث يمكن أن تكون $P$ و$Q$ صفوفًا للحدوديات الأصلية، أو صفوفًا نشأت من الصفوف الأصلية بتتابع من التعويضات (9.) أو (11.)، أو تكون أخيرًا صفوف متغيرات. غير أن العبارات (40.) يمكن توليدها بالضم المتعاقب لكل صفين من صفوف الرموز في جداء مصفوفات، مع الاستعانة بصفوف متغيرات؛ إذ نحصل من
\[
 \pair{q_{n-\mu-\lambda}P^\mu}{Q_{\nu_1}p_{n-\nu_1-\lambda}}
 =\bigl(R_\lambda Q_{\nu_1}p_{n-\nu_1-\lambda}\bigr)
\]
بضم الصف $Q_{\nu_2}$ إليه على عبارة
\[
 \pair{[R_\lambda Q_{\nu_1}]^{n-\lambda-\nu_1}}{Q_{\nu_2}p_{n-\nu_1-\nu_2-\lambda}}
 =\pair{q_{n-\mu-\lambda}P^\mu}{Q_{\nu_1}Q_{\nu_2}p_{n-\nu_1-\nu_2-\lambda}},
\]
ومن ثم نحصل، بمتابعة الإجراء، على عبارة من الصورة (40.). وإذا لم تكن $P$ و$Q$ صفوف الحدوديات الأصلية، فإن (9.) و(11.)، بالاقتران بما برهن للتو، تبينان أنها نشأت هي أيضًا من تتابع من عمليات ضم كل صفين من صفوف الرموز. ومن ذلك ينتج:

\emph{المبرهنة \srcLTR{V}:} تتألف عملية توليد جميع التكوينات اللامتغيرة، أي عملية الطي، من الضم المتعاقب لكل صفين من صفوف الرموز في جداء مصفوفات، مع الاستعانة بصفوف متغيرات.

من صفَّي الرموز $S^\sigma,T^\tau$، حيث $\sigma\geqq\tau$، يمكن الآن تكوين جداءات المصفوفات الآتية:
\begin{align}
&\pair{q_{n-\sigma-\lambda}S^\sigma}{T_{n-\tau}p_{\tau-\lambda}}, &&(\lambda=1,2,\ldots,\tau,n-\sigma),\tag{41}\\
&\pair{q_{n-\tau-\mu}T^\tau}{S_{n-\sigma}p_{\sigma-\mu}}, &&(\mu=\sigma-\tau+\lambda),\tag{42}\\
&\pair{q_{n-\tau-\nu}T^\tau}{S_{n-\sigma}p_{\sigma-\nu}}, &&(\nu=1,2,\ldots,\sigma-\tau).
\tag{43}
\end{align}
غير أن (31.) يفضي، بالنسبة إلى (42.) و(43.)، إلى القيم الآتية:
\begin{align}
&\pair{q_{n-\sigma-\lambda}S^\sigma}{T_{n-\tau}p_{\tau-\lambda}}
 +\sum_{\alpha}\Pi\pair{q_{n-\sigma-\lambda-\alpha}S^\sigma}{T_{n-\tau}p_{\tau-\lambda-\alpha}},\tag{42a}\\
&\Pi(q_{n-\sigma}S^\sigma)(T_{n-\tau}p_\tau)
 +\sum_{\beta}\Pi\pair{q_{n-\sigma-\beta}S^\sigma}{T_{n-\tau}p_{\tau-\beta}}.
\tag{43a}
\end{align}
ومن ثم يمكن التعبير خطيًا عن الحدوديات (42.) بواسطة (41.) والأشكال القطبية لـ(41.)، وعن الحدوديات (43.) بواسطة الأشكال القطبية للحدودية الأساسية وللحدوديات (41.). كذلك ينتج من (31.) اعتماد الحدوديات (41.) خطيًا على (42.) وعلى الأشكال القطبية لـ(42.). وبناءً على ذلك تكون الحدوديات (42.)، وفق تعريف كليبش\footnote{كليبش، المرجع نفسه، \S\ 7.}، مكافئة للحدوديات (41.)، في حين يُرَدّ (43.) إلى الحدودية الأساسية. إذن تُعطى عملية التوليد، على قدم المساواة، بـ(41.) أو (42.)؛ وسنعرّف العملية (41.)، الأبسط من حيث العدد $\lambda$، بأنها عملية الطي. والعدد $\lambda$، أي نقص الطي (قارن \S\ 2)، يبيّن عدد الصفوف الناقصة لإتمام جداء المحددات، وذلك في جداء المصفوفات الذي يشتمل، عند طيٍّ معيّن، على أكبر عدد من الصفوف. ولما كانت $\lambda$ لا تتجاوز أصغر القيمتين $\tau,n-\sigma$، حيث $\sigma\geqq\tau$، فإن:
\begin{equation}
 \lambda\leqq\frac n2,\ n\text{ زوجي};\qquad
 \lambda\leqq\frac{n-1}{2},\ n\text{ فردي}.
\tag{44}
\end{equation}
ويبقى ردّ (42.) و(43.) إلى (41.) صحيحًا أيضًا حين تكون الصفوف $p,q$ قد تحولت، بطيٍّ إضافي، إلى صفوف رموز؛ إذ ينتج وفقًا لـ(36.):
\[
 \pair{A^{n-\tau-\varkappa}T^\tau}{S_{n-\sigma}B_{\sigma-\varkappa}}
 =\sum_{\alpha=0,\varkappa-\sigma+\tau}^{\tau,n-\sigma}
 \eps\pair{R^{\tau-\alpha}B^{n-\sigma+\varkappa}}{A_{\tau+\varkappa}R_{n-\sigma-\alpha}},
\]
حيث تكون الصفوف $R$ قد نشأت، وفقًا لـ(33.)، من $S$ و$T$. ومن ثم تُحصل هذه الحدوديات، التي لا تحتوي إلا على صفوف رموز، بطيّ $\pair{q_\tau A^{n-\tau-\varkappa}}{B_{\sigma-\varkappa}p_{n-\sigma}}$ أو $\pair{q_{\tau-\alpha}B^{n-\sigma+\varkappa}}{A_{\tau+\varkappa}p_{n-\sigma-\alpha}}$، على الترتيب ــ بحسب كون $\sigma-\tau\gtreqless 2\varkappa$ ــ مع الحدودية الأساسية والحدوديات (41.).

تُستخرج من عملية الطي الرمزية، على النحو المعروف، عمليات الاشتقاق اللامتغيرة، إذا استُبدل بصفَّي الرموز $S^\sigma,T_{n-\tau}$ صفّا رموز الاشتقاق $\frac{\partial}{\partial p_\sigma}$، $\frac{\partial}{\partial q_{n-\tau}}$؛ ومن المعلوم أن لكل جداء $\frac{\partial}{\partial p_{i_1\ldots i_\sigma}}\cdot\frac{\partial}{\partial q_{j_1\ldots j_{n-\tau}}}$ الدلالة الفعلية $\frac{\partial^2}{\partial p_{i_1\ldots i_\sigma}\partial q_{j_1\ldots j_{n-\tau}}}$. ولما كان صفّا $\frac{\partial}{\partial p_\sigma}$، $\frac{\partial}{\partial q_{n-\tau}}$ متوافقَي التدرج مع الصفَّين $S^\sigma,T_{n-\tau}$، فإنهما يخضعان للمتطابقات نفسها التي يخضع لها هذان، وبخاصة المتطابقات القائمة بين (42.) و(42a.)، وبين (43.) و(43a.). وبالتلخيص نحصل على ما يأتي:

\emph{المبرهنة \srcLTR{VI}:} تنشأ جميع التكوينات اللامتغيرة لحدوديات أساسية اعتباطية منها، بالتطبيق المتكرر لعملية الطي الرمزية (41.)، أو كذلك لعمليات الاشتقاق اللامتغيرة:
\begin{equation}
 \left(q_{n-\sigma-\lambda}\frac{\partial}{\partial p_\sigma}\middle|\frac{\partial}{\partial q_{n-\tau}}p_{\tau-\lambda}\right),
 \qquad(\sigma\geqq\tau,\ \lambda=1,2,\ldots,\tau,n-\sigma).
\tag{45}
\end{equation}
ويبقى أن نلاحظ أنه عند كل طيٍّ (41.) ينقص الوزن الكلي للحدودية، أي مجموع أوزان صفوف المتغيرات المفردة (قارن \S\ 1)، بمقدار العدد الموجب $2(\lambda^2+\lambda(\sigma-\tau))$. وينتج من ذلك، بمعزل عن برهان التناهي العام، تناهي سلسلة الحدوديات، أي مجموعة التكوينات اللامتغيرة الخطية في معاملات حدودية أساسية معيّنة.\footnote{قارن كليبش، المرجع نفسه، \S\ 11 و12: تناهي منظومة الملازمات الأولية الموافقة للتغير.}

ولنلاحظ كذلك أن المبرهنة \srcLTR{VI} تظل صحيحة أيضًا للحدوديات التي لا تستوفي الشروط (12.). إذ يمكن استبدال كل عامل من (3.)، وهو $\pair{S^\rho}{p_\rho}^{m_\rho}$، بجداء من $m_\rho$ عوامل خطية $\pair{A^\rho}{p_\rho}\pair{B^\rho}{p_\rho}\cdots\pair{M^\rho}{p_\rho}$؛ فتتحول التكوينات اللامتغيرة بذلك إلى تكوينات تخص منظومة من الحدوديات الخطية، لا يلزم مساواة بعضها ببعض إلا لاحقًا. أما كل حدودية خطية بمفردها فتستوفي دائمًا الشروط (12.)، التي لا تتضمن عند إدخال رموز الاشتقاق إلا مشتقات من الرتبة الثانية.

\section*{\S\ 7. توسيع الحدوديات العامة بحسب الصيغ النظامية.}
سنستخلص فيما يأتي (\S\ 8) خاصية أساسية لعملية الطي (41.)، هي قابليتها للتركيب من طيات أساسية؛ غير أننا نحتاج لذلك إلى نقل مبرهنة مهمة صاغها السيد مرتنس\footnote{\srcLTR{F. Mertens, \emph{Über invariante Gebilde quaternärer Formen} (انظر المقدمة)، ويُحال إليها باسم ``Mertens``.}} في المجال الرباعي إلى المجال ذي $n$ متغيرات. تقول هذه المبرهنة إنه يمكن استبدال كل منظومة منتهية اعتباطية من الحدوديات بمنظومة مكافئة من الصيغ النظامية. ونعني هنا بـ``الصيغة النظامية'' ــ تعميمًا للتعريف الثنائي والثلاثي\footnote{\srcLTR{Study, a. a. O. S. 54.}} ــ حدودية تنعدم تطابقيًا جميع تكويناتها اللامتغيرة الخطية في المعاملات، باستثناء الحدودية الأساسية نفسها. ومن ثم، وفقًا للمبرهنة السادسة (\S\ 6)، تحقق الصيغة النظامية $\varphi$ منظومة المعادلات التفاضلية:
\begin{equation}
 \left(q_{n-\sigma-\lambda}\frac{\partial}{\partial p_\sigma}\middle|\frac{\partial}{\partial q_{n-\tau}}p_{\tau-\lambda}\right)\varphi=0,
 \qquad(\sigma\geqq\tau,\ \lambda=1,2,\ldots,\tau,n-\sigma),
\tag{46}
\end{equation}
حيث تُعامل الصفوف $q_{n-\sigma-\lambda},p_{\tau-\lambda}$ بوصفها كميات اعتباطية. وتتطابق هذه المعادلات التفاضلية في المجال الرباعي ــ مع مراعاة (39.) عندما $\sigma=\tau=2$ ــ تطابقًا تامًا مع المعادلات التفاضلية العشر التي أعطاها السيد مرتنس من غير إدخال الصفوف الاعتباطية $p,q$ (الموضع المذكور، الصيغة 28))؛ وتتيح لنا معرفتها، مقترنة بالمبرهنة السادسة، أن ننقل إلى $n$ متغيرات طريقة البرهان الواردة هناك نقلًا يكاد يكون حرفيًا. ولا يلزم أن نضيف إلا البرهان، المستند إلى متطابقة التفكيك (30.)، على أن الحدوديات المنشأة صيغ نظامية؛ ففي المجال الرباعي ينتج هذا البرهان من غير أي تحويل إضافي. ولذلك تكفي أيضًا خلاصة موجزة للبرهان، ونقتصر طلبًا للبساطة على الحالة الوحيدة التي ستعنينا فيما يأتي، وهي حالة حدودية أساسية واحدة $f$ وسلسلة ملازماتها (قارن خاتمة \S\ 6)، لأن البرهان يبقى هو نفسه للتكوينات اللامتغيرة من درجة اعتباطية في معاملات أي عدد من الحدوديات الأساسية.

الفكرة الأساسية للبرهان هي أن نستبدل، بإدخال المعاملات المحوَّلة، بحدودية ذات $n-1$ من الصفوف $p_\rho$ حدودياتٍ ذات $n$ من الصفوف المتوافقة التدرج مع $x$، وهي صفوف $\xi$ لكميات التحويل. وتُشتق من هذه الحدوديات، بواسطة عمليات اشتقاق لامتغيرة، الصيغ النظامية المطلوبة مباشرة. ويجري توسيع الحدوديات العامة بحسب الصيغ النظامية بوساطة توسعة متسلسلة لحدوديات ذات $\rho$ ($\rho\leqq n$) من الصفوف $\xi$ المتوافقة التدرج؛ ثم تعيد عمليات الاشتقاق اللامتغيرة إلى الحدوديات في الصفوف الأصلية $p_\rho$.

لتكن $\xi^{(1)},\ldots,\xi^{(n)}$ صفوف كميات التحويل، بحيث يكون:
\begin{equation}
\begin{aligned}
x_i&=\xi_i^{(1)}x'_1+\xi_i^{(2)}x'_2+\cdots+\xi_i^{(n)}x'_n,\\
p_{i_1\ldots i_\rho}&=\sum_j(\xi_{i_1}^{(j_1)}\cdots \xi_{i_\rho}^{(j_\rho)})p'_{j_1\ldots j_\rho}.
\end{aligned}
\tag{47}
\end{equation}
ولنرمز إلى المعاملات المحوَّلة للحدوديات $f,\varphi,\ldots$ بالرموز $(f),(\varphi),\ldots$، ولتكن
\[
 f\equiv\pair{S^1}{p_1}^{m_1}\pair{S^2}{p_2}^{m_2}\cdots\pair{S^{n-1}}{p_{n-1}}^{m_{n-1}}.
\]
عندئذ يكون:
\begin{equation}
\begin{aligned}
(f)&=\pair{S^1}{\xi^{(1)}}^{m_{11}}\cdots\pair{S^1}{\xi^{(n)}}^{m_{1n}}
\pair{S^2}{\xi^{(1)}\xi^{(2)}}^{m_{21}}\cdots
\pair{S^{n-1}}{\xi^{(2)}\cdots\xi^{(n)}}^{m_{n-1,n}}\\
&=\pair{s^{(1)}}{\xi^{(1)}}^{k_1}\pair{s^{(2)}}{\xi^{(2)}}^{k_2}\cdots\pair{s^{(n)}}{\xi^{(n)}}^{k_n},
\end{aligned}
\tag{48}
\end{equation}
حيث $m_{11}+\cdots+m_{1n}=m_1$، وهكذا. ومن كل معامل $(f)$ يتحقق فيه
\begin{equation}
 k_1\geqq k_2\geqq k_3\geqq\cdots\geqq k_n,
\tag{49}
\end{equation}
ينتج، بواسطة عملية الاشتقاق التي تستنفد صفوف $\xi$ استنفادًا تامًا:
\begin{equation}
\left(\frac{\partial}{\partial\xi^{(1)}}\middle|p_1\right)^{k_1-k_2}
\left(\frac{\partial}{\partial\xi^{(1)}}\frac{\partial}{\partial\xi^{(2)}}\middle|p_2\right)^{k_2-k_3}
\cdots
\left(\frac{\partial}{\partial\xi^{(1)}}\cdots\frac{\partial}{\partial\xi^{(n-1)}}\middle|p_{n-1}\right)^{k_{n-1}-k_n}
\left(\frac{\partial}{\partial\xi^{(1)}}\cdots\frac{\partial}{\partial\xi^{(n)}}\right)^{k_n}
\tag{50}
\end{equation}
الحدودية
\begin{equation}
 \varphi=\bigl(s^{(1)}\mid p_1\bigr)^{k_1-k_2}
 \bigl(s^{(1)}s^{(2)}\mid p_2\bigr)^{k_2-k_3}\cdots
 \bigl(s^{(1)}\cdots s^{(n-1)}\mid p_{n-1}\bigr)^{k_{n-1}-k_n}
 \bigl(s^{(1)}s^{(2)}\cdots s^{(n)}\bigr)^{k_n}.
\tag{51}
\end{equation}
والحدوديات $\varphi$ هي الصيغ النظامية المطلوبة؛ إذ تتصف بالخاصيتين الآتيتين اللتين تعرّفان الصيغ النظامية:

1. إنها تكوينات لامتغيرة للحدودية الأساسية $f$، خطية في المعاملات. وينتج ذلك من كونها تنشأ من $f$ بعمليات الاشتقاق اللامتغيرة $\left(\frac{\partial}{\partial p_\rho}\middle|\xi^{(i_1)}\xi^{(i_2)}\cdots\xi^{(i_\rho)}\right)$ وبالعملية (50.). ولذلك يمكن (وفقًا لخاتمة \S\ 6) التعبير عنها خطيًا بدلالة الأشكال القطبية لسلسلة الملازمات المنتهية للحدودية $f$. ونحصل على هذه الأشكال القطبية بأن نعدّ المتغيرات الداخلة في $\varphi$، وهي $p_\rho$، معاملات لحدودية متحللة إلى حدوديات خطية، ذات صفوف المتغيرات $\mathfrak{p}_\sigma$:
\begin{equation}
 H=(q_{n-1}\mid \mathfrak{p}_{n-1})^{k_1-k_2}(q_{n-2}\mid \mathfrak{p}_{n-2})^{k_2-k_3}\cdots(q_1\mid \mathfrak{p}_1)^{k_{n-1}-k_n}.
\tag{52}
\end{equation}
وتعطي اللامتغيرات المتزامنة لسلسلتي ملازمات $f$ و$H$ جميع الأشكال القطبية الداخلة في الاعتبار، لأن درجة هذه الأخيرة بالنسبة إلى متغيرات كل صف مفرد $p_\rho$ يجب أن تساوي درجة الحدودية $\varphi$ ذاتها بالنسبة إلى المتغيرات نفسها.

2. إنها تحقق المعادلات التفاضلية (46.). فإذا طبقنا على $\varphi$ عملية الاشتقاق (45.) حصلنا على:
\[
 \varphi'=\bigl(q_{n-\sigma-\lambda}s^{(1)}\cdots s^{(\sigma)}\mid [s^{(1)}\cdots s^{(\tau)}]_{n-\tau}p_{\tau-\lambda}\bigr),\qquad (\sigma\geqq\tau).
\]
ومن متطابقة التفكيك (30.) ينتج، إذا وحّدنا الصفين المنشأين من صفوف $s$، أي $q_\sigma=s^{(1)}\cdots s^{(\sigma)}$ و$p_{n-\tau}=[s^{(1)}\cdots s^{(\tau)}]_{n-\tau}$، مع الصفين $q_\rho,p_\tau$ الواردين هناك، أن:
\[
 \varphi'=\sum_{\beta=\sigma-\tau+\lambda}^{n-\tau,\sigma}\sum_{i_\beta}\eps\,
 \pair{\dot{q}_{\sigma-\beta}^{(i_1\ldots i_\beta)}q_{n-\tau+\lambda}}{p_{\sigma+\lambda}\dot{p}_{n-\tau-\beta}^{(i_1\ldots i_\beta)}}.
\]
ويكون عندئذ:
\[
 \dot{p}_{n-\tau-\beta}^{(i_1\ldots i_\beta)}\sim s^{(1)}\cdots s^{(\tau)}s^{(i_1)}\cdots s^{(i_\beta)},
\]
حيث تدل $i_1,\ldots,i_\beta$ على أيّ $\beta$ من الأعداد من 1 إلى $\sigma$، مختارة بحيث تكون متميزة بعضها من بعض. ولما كان $\beta>\sigma-\tau$، فلا بد أن يوجد بين الأعداد $i_1,\ldots,i_\beta$ عدد واقع بين 1 و$\tau$؛ وبذلك ينعدم $\dot{p}_{n-\tau-\beta}^{(i_1\ldots i_\beta)}$ تطابقيًا، ومن ثم ينعدم كل جداء مصفوفات مفرد، وتنعدم معه أيضًا $\varphi'$.

ولإتمام تكافؤ منظومة الصيغ النظامية مع سلسلة الملازمات $f$، لا يبقى الآن إلا أن نبيّن أن جميع حدوديات سلسلة الملازمات يمكن توسيعها بدلالة الأشكال القطبية للصيغ النظامية بالمعنى المحدد في 1. ولتكن $\Pi$ تركيبًا من عمليات الاستقطاب
\srcnumdisplay[4.2em]{(53.)}{%
 \left(\frac{\partial}{\partial\xi^{(i)}}\middle|\xi^{(k)}\right),
 \qquad (i\ne k;\ i=1,2,\ldots,\rho;\ k=1,2,\ldots,\rho)}
ولتكن $\Pi^\sigma$ تركيبًا من العمليات الخاصة (53.) التي يكون فيها $i=\sigma$ و$k<\sigma$. فإذا كانت $F$ حدودية في الصفوف الـ$\rho$، أي $\xi^{(1)},\ldots,\xi^{(\rho)}$، وكانت في الصف $\xi^{(\rho)}$ من الدرجة $k_\rho$، فتصح التوسعة الآتية\srcfn{*)}{\srcLTR{F. Mertens, \emph{Über eine Formel der Determinantentheorie}. Sitzungsber. d. Ak. d. Wiss., Wien, math.-naturw. Kl., Bd. 91. 2. Abt. (1885). Formel 20).}}:
\srcnumdisplay[4.2em]{(54.)}{%
 F=\sum_{\alpha=0}^{k_\rho}\Pi F_\alpha,
 \qquad(\rho\leqq n)}
حيث تكون $F_\alpha$ في الصف الأخير $\xi^{(\rho)}$ من الدرجة $\alpha$، وتنشأ من $F$ بعمليات الاشتقاق اللامتغيرة
\srcnumdisplay[4.2em]{(55.)}{%
 \left(\frac{\partial}{\partial\xi^{(1)}}\cdots\frac{\partial}{\partial\xi^{(\rho)}}\middle|\xi^{(1)}\cdots\xi^{(\rho)}\right)^\alpha}
وبـ$\Pi^\rho$. فإذا استُبدلت، عند إنشاء $F_\alpha$، العملية (55.) بـ
\srcnumdisplay[4.2em]{(56.)}{%
 \left(\frac{\partial}{\partial\xi^{(1)}}\cdots\frac{\partial}{\partial\xi^{(\rho)}}\middle|p_\rho\right)^\alpha,}
نتجت حدودية $F_\alpha(p_\rho)$ لا تضم إلا $\rho-1$ من الصفوف $\xi^{(1)},\ldots,\xi^{(\rho-1)}$، ويمكن تطبيق (54.) عليها من جديد. ويؤدي استمرار هذه الطريقة إلى حدوديات $G(p_\rho,p_{\rho-1},\ldots,p_1)$. وللحصول منها على توسعة $F$ وفق (54.)، ينبغي استبدال كل صف $p_\sigma$ بالصفوف $\xi^{(1)},\ldots,\xi^{(\sigma)}$ وتطبيق عملية الاستقطاب $\Pi$ (53.) على الحدوديات الناتجة؛ وهذا يعطي (قارن (48.)) مجموعًا من المعاملات المحوَّلة $(G)$. وعندما $\rho=n$ تبقى العملية (55.) في الخطوة الأولى. وإذا رمز $c$ إلى ثوابت عددية، ووضعنا للاختصار أيضًا
\srcnumdisplay[4.2em]{(57.)}{%
 (\xi^{(1)}\xi^{(2)}\cdots\xi^{(n)})=\Delta;
 \qquad
 \left(\frac{\partial}{\partial\xi^{(1)}}\frac{\partial}{\partial\xi^{(2)}}\cdots\frac{\partial}{\partial\xi^{(n)}}\right)=\nabla,}
حصلنا، في حالة $\rho=n$، على:
\srcnumdisplay[4.2em]{(58.)}{%
 F=\sum c\Delta^\alpha(G),}
أما عندما $\rho<n$ فلا يدخل في الطرف الأيمن إلا $\alpha=0$، أي $\sum c(G)$.

إذا طبقنا (58.) على معامل محوَّل $(f)$ من (48.)، فإن قابلية عمليات الاستقطاب $\Pi^\sigma$ للتبادل مع جميع العمليات (56.) التي يكون فيها $\rho\geqq\sigma$، وكون الأشكال القطبية للمعاملات المحوَّلة تعطي بدورها معاملات محوَّلة، يؤديان إلى:
\[
\begin{aligned}
 G={}&\left(\frac{\partial}{\partial\xi^{(1)}}\middle|p_1\right)^{\beta_1}
 \left(\frac{\partial}{\partial\xi^{(1)}}\frac{\partial}{\partial\xi^{(2)}}\middle|p_2\right)^{\beta_2}
 \cdots
 \left(\frac{\partial}{\partial\xi^{(1)}}\cdots\frac{\partial}{\partial\xi^{(n-1)}}\middle|p_{n-1}\right)^{\beta_{n-1}}
 \nabla^{\beta_n}\Pi(f)\\
 ={}&\sum c\varphi
\end{aligned}
\]
ومن (58.) ينتج:
\srcnumdisplay[4.2em]{(59.)}{%
 (f)=\sum c\Delta^\alpha(\varphi)\qquad\text{(مرتنس، الصيغة 23)).}}
وللانتقال من (59.) إلى توسعة $f$ نفسها، نعود فنعدّ المتغيرات $p_\rho$ معاملات لحدودية $H$ في (52.)، ذات أعداد الدرجات المناظرة لـ$f$، وهي $m_1,\ldots,m_{n-1}$، ولتكن $m=m_1+\cdots+m_{n-1}$. وعندئذ، وفقًا لتعريف اللامتغيرات، يكون:
\[
 f\cdot\Delta^m=\sum(f)(H)=\sum c\Delta^\alpha(\varphi)(H),
\]
ويعطي تطبيق العملية التي تستنفد صفوف $\xi$ استنفادًا تامًا، وهي $\nabla^m$:
\srcnumdisplay[4.2em]{(60.)}{%
 f=\sum c\Phi,}
حيث تُشتق الحدوديات $\Phi$ من $\varphi$ و$H$ بعمليات اشتقاق لامتغيرة بالنسبة إلى المتغيرات، فتكون لامتغيرات متزامنة لسلسلتي الملازمات $\varphi$ و$H$\srcfn{*)}{بدلًا من الاستدلال المباشر، يُدخل مرتنس، في \S\ 7، سلسلة من عمليات الاشتقاق الإضافية التي ترمي أساسًا إلى تكوين سلسلة الملازمات $H$، وهو ما تؤديه مبرهنتنا السادسة.}. ولكن لما كانت $\varphi$ صيغة نظامية، اقتصرت سلسلة الملازمات $\varphi$ على الحدودية $\varphi$ نفسها؛ أما سلسلة الملازمات $H$ فتحتوي جميع التراكيب اللامتغيرة الممكنة للمتغيرات؛ ومن ثم فإن اللامتغيرات المتزامنة لـ$\varphi$ و$H$ هي الأشكال القطبية لـ$\varphi$ بالمعنى المحدد في 1؛ أي إن (60.) تمثل توسعة $f$ بدلالة الأشكال القطبية للصيغ النظامية. ولا يزال علينا الآن إثبات التوسعة نفسها لحدودية اعتباطية $\theta$ من سلسلة ملازمات $f$. ولتكن $\Gamma$ الصيغ النظامية المشتقة من $(\theta)$ بواسطة (50.)، بحيث يكون:
\srcnumdisplay[4.2em]{(61.)}{%
 (\theta)=\sum c\Delta^\beta(\Gamma).}
وتتميز الحدوديات $\theta$ بالعلاقة الآتية التي تصح لكل معامل محوَّل:
\[
 (\theta)\Delta^\sigma=\sum c(f)\qquad\text{(مرتنس، الصيغة 9)).}
\]
ولكن كل حدودية $F$ في الصفوف الـ$n$، أي $\xi$، إذا احتوت العامل $\Delta^\sigma$ فإنها تحتوي أيضًا العامل الثاني، أي $\nabla^\sigma\Pi F$ (مرتنس، الصيغة 20))، ومن ثم ينتج:
\[
 (\theta)=\sum c\nabla^\sigma(f),\qquad\text{إذن}\qquad \Gamma=\sum c\varphi,
\]
ومن ذلك، وفقًا لـ(61.)، نحصل على العلاقة المناظرة لـ(59.):
\[
 (\theta)=\sum c\Delta^\beta(\varphi),
\]
وهي تستلزم للحدودية $\theta$ التوسعة (60.) بدلالة الأشكال القطبية للصيغ النظامية $\varphi$. وخلاصة ما تقدم:

\emph{المبرهنة السابعة. يمكن استبدال كل سلسلة ملازمات $f$ بمنظومة مكافئة من الصيغ النظامية.}

\begin{center}
\emph{\S\ 8. ردُّ عملية الطي إلى عمليات الطي الأساسية.}
\end{center}
استكمالًا للمبرهنة السابعة، ننفّذ الآن ردَّ عملية الطي إلى عمليات الطي الأساسية، الذي ذُكر في مطلع \S\ 7. ونسمّي في ذلك (قارن \S\ 5) كلَّ طيٍّ لصفَّين متوافقي التدرج $S^\rho,T^\rho$ «\emph{طيًّا متوافقَ التدرج}»، ونسمّي، على وجه الخصوص، الطيَّ المتوافقَ التدرج ذا النقص 1 لصفَّين $S^\rho,T^\rho$ «\emph{طيًّا أساسيًّا من النوع $\rho$}». ثم نبرهن، إتمامًا للمبرهنة السادسة، المبرهنةَ الآتية:

\emph{المبرهنة الثامنة: يمكن تركيب عملية الطي العامة (41.) من عمليات الطي الأساسية البالغ عددها $(n-1)$ من النوع $\rho$، حيث $\rho=1,2,\ldots,n-1$. وبعبارة أخرى: يكفي لتوليد جميع التكوينات اللامتغيرة التطبيق المتعاقب لعمليات الطي الأساسية البالغ عددها $n-1$.}

في المجال الثنائي تتطابق عملية الطي (41.) مباشرةً مع عملية الطي الأساسية الوحيدة الممكنة؛ وفي المجال الثلاثي برهنتُ المبرهنة الثامنة في \S\ 1 من رسالتي للدكتوراه (قارن المقدمة). وهناك، بحكم (44.)، تظل عمليات الطي المتوافقة التدرج العامة مطابقة لعمليات الطي الأساسية، ومن الجدير بالذكر أن المبرهنة الثامنة، في حالة $n$ من المتغيرات، تردّ المسألة إلى عمليات الطي الأساسية، لا إلى عمليات الطي المتوافقة التدرج فحسب، وهو ردّ أضعف بكثير. ويحاكي البرهان، من حيث المبدأ، برهان الحالة الثلاثية؛ إذ تُنشأ أعمّ عمليات الطي من عمليات الطي الأساسية باستعمال المتطابقات الرمزية. وسننشئ:
\begin{enumerate}
\renewcommand{\labelenumi}{\arabic{enumi})}
\item عمليات طي اعتباطية ذات نقص 1 بواسطة عمليات طي متوافقة التدرج ذات نقص 1، أي بواسطة عمليات الطي الأساسية (نظير الحالة الثلاثية)،
\item عمليات طي اعتباطية ذات نقص $\lambda$ بواسطة عمليات طي متوافقة التدرج ذات نقص $\lambda$ وبواسطة عمليات طي اعتباطية ذات نقص 1،
\item عمليات طي متوافقة التدرج ذات نقص $\lambda$ بواسطة عمليات طي متوافقة التدرج ذات نقص $\lambda-1$ وبواسطة عمليات طي اعتباطية ذات نقص 1.
\end{enumerate}
وفي ذلك نفترض، وهذا أمر جوهري، وفق المبرهنة السابعة، أن الحدوديتين $S$ و$T$ اللتين سيُجرى عليهما الطي صيغتان نظاميتان. ومن ثم تصح، بحسب (46.)، العلاقات:
\srcnumdisplay[4.2em]{\raisebox{-1.7ex}{(62.)}}{%
\begin{aligned}
\pair{q_{n-\sigma_1-\lambda}S^{\sigma_1}}{S_{n-\sigma_2}p_{\sigma_2-\lambda}}&=0,
&& (\sigma_1\geqq\sigma_2,\ \lambda=1,2,\ldots,\sigma_2,n-\sigma_1),\\
\pair{q_{n-\tau_1-\lambda}T^{\tau_1}}{T_{n-\tau_2}p_{\tau_2-\lambda}}&=0,
&& (\tau_1\geqq\tau_2,\ \lambda=1,2,\ldots,\tau_2,n-\tau_1).
\end{aligned}}
ويكفي أيضًا، وفقًا لخاتمة \S\ 2، أن نفترض الحدّيتين $S$ و$T$ خطيتين في جميع صفوف المتغيرات؛ أي إن الحدوديات المنشأة تنتمي إلى سلسلة الملازمات الآتية:
\[
 f=\pair{S^1}{p_1}\pair{S^2}{p_2}\cdots\pair{S^{n-1}}{p_{n-1}}
   \pair{T^1}{p_1}\pair{T^2}{p_2}\cdots\pair{T^{n-1}}{p_{n-1}}.
\]
\emph{1) توليد $\pair{q_{n-\sigma-1}S^\sigma}{T_{n-\tau}p_{\tau-1}}$، حيث $\sigma>\tau$، من طيّات أساسية من النوع $\tau+\alpha$؛ $\alpha=0,1,\ldots,\sigma-\tau$.} ومن الطيّ الأساسي من النوع $\tau$
\[
 \pair{q_{n-\tau-1}S^\tau}{T_{n-\tau}p_{\tau-1}}
\]
نحصل، بإجراء الإحلال
\[
 w\sim T_{n-\tau}p_{\tau-1}
\]
على حدودية من سلسلة ملازمات $f$، ذات ثلاثة صفوف دليل الوزن العلوي لكلٍّ منها (قارن \S\ 1) هو $\tau+1$:
\srcnumdisplay[4.2em]{(63.)}{%
 \pair{wS^\tau}{p_{\tau+1}}\pair{S^{\tau+1}}{p_{\tau+1}}\pair{T^{\tau+1}}{p_{\tau+1}}.}
وينتج من تطبيق طيّ أساسي من النوع $\tau+1$ على (63.) ما يأتي:
\[
 \pair{q_{n-\tau-2}wS^\tau}{S_{n-\tau-1}p_\tau};
\]
ومن ذلك، بحسب (31.) وبعد إجراء الإحلال $q_{n-\tau-2}w=q'_{n-\tau-1}$، ينتج المقدار الآتي:
\[
 \eps\bigl(q_{n-\tau-2}wS^{\tau+1}\bigr)\pair{S^\tau}{p_\tau}
 +\sum_{\alpha=1}^{\tau,n-\tau-1}\Pi\pair{q'_{n-\tau-1-\alpha}S^{\tau+1}}{S_{n-\tau}p_{\tau-\alpha}}.
\]
وينعدم التعبير الواقع تحت علامة الجمع تطابقيًا بمقتضى (62.)؛ ومن ثم نكون قد وصلنا إلى الحدودية:
\[
 \pair{wS^{\tau+1}}{p_{\tau+2}}\pair{S^{\tau+2}}{p_{\tau+2}}\pair{T^{\tau+2}}{p_{\tau+2}},
\]
وهي تنتج من (63.) بإبدال $\tau$ بـ$\tau+1$. ومن ثم فإن تكرار العملية $(\sigma-\tau)$ مرةً يفضي إلى الطيّ المطلوب:
\[
 \pair{wS^\sigma}{p_{\sigma+1}}=\eps\pair{q_{n-\sigma-1}S^\sigma}{T_{n-\tau}p_{\tau-1}}.
\]
ونرمز بإيجاز إلى العملية التي أجريناها بالصيغة:
\srcnumdisplay[4.2em]{(64.)}{%
 T^\tau S^\tau,\quad S^{\tau+1},\quad S^{\tau+2},\ldots,S^\sigma;}
ونرى أن المستعمل ليس إلا خاصية الصيغة النظامية لـ$S$، لا خاصية $T$. ويمكن إجراء عملية ثنوية لـ(64.)، باستعمال خاصية الصيغة النظامية لـ$T$ وحدها، وبالترتيب المعكوس؛ وهي العملية الآتية:
\srcnumdisplay[4.2em]{(65.)}{%
 S^\sigma T^\sigma,\quad T^{\sigma-1},\quad T^{\sigma-2},\ldots,T^\tau}
بحسب العلاقات:
\begin{align*}
\pair{q_{n-\sigma-1}S^\sigma}{T_{n-\sigma}p_{\sigma-1}}&=\eps\bigl(T_{n-\sigma}y p_{\sigma-1}\bigr),\quad \text{حيث }y\sim q_{n-\sigma-1}S^\sigma,\\
\pair{q_{n-\sigma}T^{\sigma-1}}{T_{n-\sigma}yp_{\sigma-2}}&=\eps\bigl(T_{n-\sigma+1}y p_{\sigma-2}\bigr),\\
&\vdots\\
\pair{q_{n-\tau-1}T^\tau}{T_{n-\tau-1}yp_{\tau-1}}&=\eps\pair{q_{n-\sigma-1}S^\sigma}{T_{n-\tau}p_{\tau-1}}.
\end{align*}
ويجدر أيضًا التنبيه إلى الصلة بنقصان الوزن الكلي المذكور في خاتمة \S\ 6؛ فقيمته في الطيّات ذات النقص 1 هي $2(1+(\sigma-\tau))$، وفي الطيّات الأساسية هي $2\cdot1$. فإذا كان التركيب من الطيّات الأساسية ممكنًا أصلًا، لزم له بالضرورة $\sigma-\tau+1$ طيًّا من هذا القبيل، كما أُجري أعلاه.

\emph{2) توليد الطيات العامة من الطيات المتوافقة التدرج والأساسية.} يبلغ نقصان الوزن في طي عام (41.): $2(\lambda^2+\lambda(\sigma-\tau))=2[\lambda^2+(\sigma-\tau)(1+(\lambda-1))]$؛ ومن ثم يُظهر إمكان تفكيكه إلى طي متوافق التدرج ذي نقص $\lambda$ (نقصان الوزن $2\lambda^2$)، وإلى $\sigma-\tau$ طيات ذات نقص 1، والفرق بين دليلي الوزن فيها $\lambda-1$ (نقصان الوزن $2\{1+(\lambda-1)\}$). ويطابق هذا التفكيك، في حالة $\lambda=1$، التفكيك المبين في 1)؛ وسنبين أنه قابل للتحقيق بالفعل.

ومن الطي المتوافق التدرج ذي النقص $\lambda$
\[
 \pair{q_{n-\tau-\lambda}S^\tau}{T_{n-\tau}p_{\tau-\lambda}}
\]
نحصل، بإجراء الإحلال $R^\lambda\sim T_{n-\tau}p_{\tau-\lambda}$، على حدودية ذات صفين، دليلا وزنهما العلويان هما، على الترتيب، $\tau+\lambda$ و$\tau+1$:
\srcnumdisplay[4.2em]{(66.)}{%
 \pair{R^\lambda S^\tau}{p_{\tau+\lambda}}\pair{S^{\tau+1}}{p_{\tau+1}}.}
والصف ذو دليل الوزن الأقل $\tau+1$ ينتمي إلى صيغة نظامية؛ ومن ثم تقبل (66.) طيًا ذا نقص 1 مركبًا من $\lambda$ طيات أساسية، على ترتيب (65.):
\[
 (R^\lambda S^\tau)S^{\tau+\lambda},\quad S^{\tau+\lambda-1},\quad S^{\tau+\lambda-2},\ldots,S^{\tau+1}.
\]
وبمراعاة (31.) و(62.) ينتج من ذلك:
\[
 \pair{q_{n-\tau-\lambda-1}R^\lambda S^\tau}{S_{n-\tau-1}p_\tau}=\eps\pair{R^\lambda S^{\tau+1}}{p_{\tau+\lambda+1}},
\]
أي حدودية تنتج من (66.) بتحويل $\tau$ إلى $\tau+1$، على نحو مماثل لما كان عليه الحال في 1) مع (63.). ويؤدي تكرار العملية $(\sigma-\tau)$ مرة إلى الطي المطلوب:
\[
 \pair{R^\lambda S^\sigma}{p_{\sigma+\lambda}}=\eps\pair{q_{n-\sigma-\lambda}S^\sigma}{T_{n-\tau}p_{\tau-\lambda}}.
\]
ولم يُستعمل في العملية كلها إلا خاصية الصيغة النظامية لـ$S$؛ أما اعتبار $T$ صيغة نظامية فيفضي إلى العملية الثنوية، مع قلب الترتيب قلبًا تامًا، وفق الصيغ:
\begin{align*}
\pair{q_{n-\sigma-\lambda}S^\sigma}{T_{n-\sigma}p_{\sigma-\lambda}}&=\eps\bigl(T_{n-\sigma}R_\lambda p_{\sigma-\lambda}\bigr),\quad \text{حيث }R_\lambda\sim q_{n-\sigma-\lambda}S^\sigma,\\
\pair{q_{n-\sigma}T^{\sigma-1}}{T_{n-\sigma}R_\lambda p_{\sigma-\lambda-1}}&=\eps\bigl(T_{n-\sigma+1}R_\lambda p_{\sigma-\lambda-1}\bigr),\\
&\vdots\\
\pair{q_{n-\tau-1}T^\tau}{T_{n-\tau-1}R_\lambda p_{\tau-\lambda}}&=\eps\pair{q_{n-\sigma-\lambda}S^\sigma}{T_{n-\tau}p_{\tau-\lambda}}.
\end{align*}
\emph{3) توليد الطيات المتوافقة التدرج ذات النقص $\lambda$ من الطيات المتوافقة التدرج ذات النقص $\lambda-1$ ومن الطيات الأساسية.} يبلغ نقصان الوزن في الطيات المتوافقة التدرج ذات النقص $\lambda$: $2\lambda^2=2((\lambda-1)^2+2\lambda-1)$؛ ومن ثم يتيح إمكان تفكيك طي كهذا إلى طي متوافق التدرج ذي نقص $\lambda-1$ وإلى $2\lambda-1$ طيًا أساسيًا. ونجد هذا التفكيك على أبسط وجه إذا أجرينا أولًا التخصيص: $n=2\lambda$. وعندئذ يكون الطي المتوافق التدرج الوحيد الممكن، ذو النقص $\lambda$، هو: $\pair{S^{n/2}}{T_{n/2}}=\pair{S^\lambda}{T_\lambda}$؛ وهو يشتمل على صف واحد أقل من كل من $u$ و$x$، قياسًا إلى الطي ذي النقص $\lambda-1$ للصفين نفسيهما $\pair{uS^\lambda}{T_\lambda x}$. وعلى العكس، نركب الطي ذا النقص $\lambda$ من طيات أساسية تُفضي إلى اختفاء صف واحد من كل من $u$ و$x$ (نقصان الوزن $2(n-1)=2(2\lambda-1)$)، وتولد في الوقت نفسه صف الرموز الجديد $R^\lambda$، ومن طي متوافق التدرج ذي النقص $\lambda-1$. وأبسط طي يحقق ذلك هو الطي الآتي ذو النقص 1:
\[
 \pair{sT^\lambda}{\sigma p_\lambda}=\eps\pair{S^{2\lambda-1}}{R_{\lambda-1}p_\lambda}=\eps\pair{R^\lambda}{p_\lambda},
\]
حيث نضع:
\[
 R^{\lambda+1}=sT^\lambda,
 \quad s=S^1,
 \quad \sigma=S_1,
 \quad R^\lambda\sim\sigma R_{\lambda-1}.
\]
ويتألف هذا الطي، وفقًا لـ(65.) و(64.)، من $2\lambda-1$ طية أساسية:
\srcnumdisplay[4.2em]{(67.)}{%
 T^\lambda S^\lambda,S^{\lambda-1},\ldots,S^1;
 \qquad R^{\lambda+1}S^{\lambda+1},S^{\lambda+2},\ldots,S^{2\lambda-1}.}
ويعطي طيٌّ ذو نقص $\lambda-1$، مع مراعاة (21.) و(62.)، ما يأتي:
\[
 \pair{uS^\lambda}{\sigma R_{\lambda-1}x}=\eps\pair{R^{\lambda+1}}{S_\lambda x}
 =\eps\pair{sT^\lambda}{S^\lambda x}=\eps\pair{S^\lambda}{T_\lambda}.
\]
ويمكن الآن تعميم البناء مباشرةً من الحالة الخاصة إلى الحالة العامة، لصفين $S^\rho,T^\rho$ و$n$ اعتباطي. وتحل محل (67.) الطيات الأساسية الآتية، وعددها $2\lambda-1$:
\srcnumdisplay[4.2em]{(68.)}{%
 T^\rho S^\rho,S^{\rho-1},\ldots,S^{\rho-\lambda+1};
 \qquad R^{\rho+1}S^{\rho+1},S^{\rho+2},\ldots,S^{\rho+\lambda-1}.}
يعطي النصف الأول من الطيات الصيغة الآتية:
\[
 \pair{q_{n-\rho-1}T^\rho}{S_{n-\rho+\lambda-1}p_{\rho-\lambda}},
\]
ومنها، بالإحلالين
\[
 S_{n-\rho+\lambda-1}p_{\rho-\lambda}\sim w,
 \qquad T^\rho w=R^{\rho+1}
\]
وبتطبيق النصف الثاني من الطيات، تنتج الصيغة الآتية:
\[
 \pair{q_{n-\rho-\lambda}S^{\rho+\lambda-1}}{R_{n-\rho-1}p_\rho}.
\]
فإذا أجرينا الإحلال:
\[
 q_{n-\rho-\lambda}S^{\rho+\lambda-1}\sim y,
\]
نتج عن طيّ متوافق التدرج ذي نقص $\lambda-1$ ما يأتي:
\srcnumdisplay[4.2em]{(69.)}{%
 \pair{q_{n-\rho-\lambda+1}S^\rho}{yR_{n-\rho-1}p_{\rho-\lambda+1}}.}
ونبين أن هذا هو الطي المطلوب ذو النقص $\lambda$. فإذا أجرينا الإحلال
\[
 R_{n-\rho-1}p_{\rho-\lambda+1}=R_{n-\lambda}
\]
ولاحظنا أن فكّ $y$
\[
 \pair{q_{n-\rho-\lambda+1}R^\lambda}{S_{n-\rho}y}
 =\eps\pair{q_{n-\rho-\lambda}S^{\rho+\lambda-1}}{S_{n-\rho}p'_{\rho-1}}=0
\]
ينتج، فإن (20.) يعطي لـ(69.) القيمة:
\[
 \eps\pair{S^\rho R^\lambda}{p_{\rho+\lambda-1}y}
 =\eps\pair{q_{n-\rho-\lambda}S^{\rho+\lambda-1}}{[S^\rho R^\lambda]_{n-\rho-\lambda}p_{\rho+\lambda-1}}.
\]
غير أن هذا الطي من النمط (43.)، فهو إذن مكافئ لـ
\[
 \pair{q_{n-\rho-\lambda}S^\rho}{R_{n-\rho-1}p_{\rho-\lambda+1}}
 =\eps\pair{wT^\rho}{R_\lambda p_{\rho-\lambda+1}},
 \qquad \text{حيث }R_\lambda\sim q_{n-\rho-\lambda}S^\rho.
\]
وهنا يكون $R_\lambda$ قد اتخذ بالفعل الصيغة الختامية المطلوبة؛ ويطابق هذا التعبير الصيغة $\pair{sT^\lambda}{S_\lambda x}$ في الحالة الخاصة. وإذا لاحظنا الآن أن فكّ $w$ و$R_\lambda$ يعطي:
\[
 \pair{wR^{n-\lambda}}{T_{n-\rho}p_{\rho-\lambda+1}}
 =\eps\pair{q'_{\lambda-1}R^{n-\lambda}}{S_{n-\rho+\lambda-1}p_{\rho-\lambda}}
 =\eps\pair{q''_{n-\rho-1}S^\rho}{S_{n-\rho+\lambda-1}p_{\rho-\lambda}}=0,
\]
فإن (21.) يعطي القيمة:
\[
 \pair{wq_{n-\rho+\lambda-1}}{T_{n-\rho}R_\lambda}
 =\eps\pair{q_{n-\rho+\lambda-1}[T_{n-\rho}R_\lambda]^{n-\lambda}}{S_{n-\rho+\lambda-1}p_{\rho-\lambda}}
 \quad\text{مكافئ}\quad
 \pair{q_{n-\rho-\lambda}S^\rho}{T_{n-\rho}p_{\rho-\lambda}}.
\]
ولم تُستعمل في العملية كلها إلا خاصية الصيغة النظامية لـ$S$، لا خاصية $T$؛ ومن ثم يمكن استبدال الطي المتوافق التدرج ذي النقص $\lambda-1$ المستخدم لتوليد (69.) بـ$2\lambda-3$ طية أساسية وبطي متوافق التدرج ذي نقص $\lambda-2$، وهو بدوره يقبل التفكيك إلى طيات أساسية، وهكذا دواليك. وبالمثل يتم التوليد عند استعمال خاصية الصيغة النظامية لـ$T$ وتطبيق العملية الثنوية، أي بتطبيق الطيات الأساسية الآتية:
\[
 S^\rho T^\rho,\quad T^{\rho+1},\ldots,T^{\rho+\lambda-1};
 \qquad R^{\rho-1}T^{\rho-1},\quad T^{\rho-2},\ldots,T^{\rho-\lambda+1},
\]
وبطيّ ثنويّ لـ(69.). وينشأ الآن الطي المكافئ للطي السابق:
\[
 \pair{q_{n-\rho-\lambda}T^\rho}{S_{n-\rho}p_{\rho-\lambda}}.
\]

وإذا أضفنا التفكيك المتحصل عليه في 2)، نكون قد حصلنا بالفعل على \emph{توليد أعم الطيات من الطيات الأساسية}. ويبقى أن نثبت أن هذا التوليد، بصرف النظر عن تبديل الترتيب، \emph{وحيد}؛ أي إن لكل طي قيمةً واحدةً لا غير لـ$\alpha_\rho$، تمثل عدد الطيات الأساسية من النمط $\rho$، حيث $\rho=1,2,\ldots,n-1$.

ولتكن $m_\rho$ درجات صفوف المتغيرات $p_\rho$ في الصيغة الابتدائية $f$، ولتكن $l_\rho$ درجات الصيغة الناتجة عن الطي. ويحول كل طي أساسي من النمط $\rho$ الدرجات $m_{\rho-1},m_\rho,m_{\rho+1}$، على الترتيب، إلى $m_{\rho-1}+1,m_\rho-2,m_{\rho+1}+1$. ومن ثم تحقق قيم $\alpha$ منظومة المعادلات الآتية:
\srcnumdisplay[4.2em]{(70.)}{%
 m_\rho-l_\rho=-\alpha_{\rho-1}+2\alpha_\rho-\alpha_{\rho+1},
 \qquad (\rho=1,2,\ldots,n-1)}
حيث ينبغي جعل $\alpha_0$ صفرًا، وكذلك $\alpha_n=0$، ويكون محدد هذه المنظومة مساويًا لـ$n$\srcfn{*)}{إذا رمز $\Delta_n$ إلى المحدد ذي الرتبة $n$، فإن صيغة العودية هي: $\Delta_n=2\Delta_{n-1}-\Delta_{n-2}$؛ أي $\Delta_n-\Delta_{n-1}=\Delta_{n-1}-\Delta_{n-2}=\cdots=\Delta_2-\Delta_1=1$، لأن $\Delta_2=3$ و$\Delta_1=2$.}. وبذلك تتعين قيم $\alpha$ تعينًا وحيدًا؛ وهي، وفقًا للمبرهنة الثامنة، أعداد صحيحة $\ge0$. وهذا يفرض شروطًا على الفروق $m_\rho-l_\rho$.

تصح نتائج هذا المبحث \emph{بموجب العلاقات (62.)}؛ ولكن إذا لم تصح العلاقات (62.)، فلا تكون التعابير الختامية متكافئة مع التعابير الابتدائية. ذلك أن تعريف التكافؤ المعطى في \S\ 6 يقبل أيضًا الصياغة الآتية\srcfn{**)}{قارن في ذلك بالمناقشة في ختام المبحث التالي.}: \glqq{}تكون صيغتان متكافئتين إذا وفقط إذا كان اختلافهما مؤلفًا من صيغ ذات طي أعلى، أي من صيغ لها نقصان أكبر في الوزن.\grqq{} ويتحقق هذا النقصان الأكبر في الوزن دائمًا عندما يكون الصفان الداخلان في الطي $q_{n-\sigma-\lambda},p_{\tau-\lambda}$ صفَّي متغيرات بالفعل، كما كان الحال في (41.) و(42.)، وفي الطيات التي وردت في هذا المبحث على أنها متكافئة، وكذلك في الصيغ المتكافئة في المبرهنة السابعة. أما إذا اشتقت هذه الصفوف بالتعويض من صفوف الرموز، فلا يلزم أن يحدث هذا النقصان، كما كان الحال في صيغ هذا المبحث التي أمكن التعبير عن بعضها ببعض بموجب (62.). بل إن بعض الصيغ المنعدمة، كما يسهل أن نرى، ذات وزن كلي أعلى.

\begin{center}
\emph{\S\ 9. سلاسل الملازمات. مبرهنات الاختزال}\srcfn{*)}{قارن الاعتبارات المناظرة في المجال الثلاثي. الأطروحة، \S\S\ 1--3.}.
\end{center}
نعرّف الآن \emph{سلسلة الملازمات}، التي أُدخلت في ختام \S\ 6، تعريفًا أدق قليلًا بأنها \glqq{}مجموع التكوينات اللامتغيرة الخطية في معاملات صيغة ابتدائية معطاة؛ أي مجموع الصيغ الناشئة من الصيغة الابتدائية بطيها في نفسها، \emph{مرتبةً بحسب الصيغ الأعلى}\grqq{}.

ولتوضيح الصيغ الأعلى، لنفترض، وفقًا للمبرهنة السابعة، أن الصيغة الابتدائية معطاة بوصفها حاصل ضرب صيغتين نظاميتين $S\cdot T$؛ وعندئذ نسمي الصيغ ذات الطي الأعلى في نفسها \emph{صيغًا أعلى}؛ ونسمي كذلك، من بين الصيغ ذات الطي نفسه في نفسها، الصيغ المعيّنة بتطبيع خاص صيغًا أعلى. وبعبارة أدق: لنفترض أن الصيغة $A$ نشأت بعدد $\alpha_\rho$، والصيغة $B$ بعدد $\beta_\rho$، من الطيات الأساسية ذات النمط $\rho$. تكون $B$ أعلى من $A$:
\begin{enumerate}
\renewcommand{\labelenumi}{\arabic{enumi})}
\item إذا تحقق في منظومة المتباينات $\beta_\rho\geqq\alpha_\rho$، $\rho=1,2,\ldots,n-1$، رمز المتباينة عند قيمة واحدة على الأقل من $\rho$\srcfn{**)}{إذا ظهر في منظومة المتباينات الرمز $>$ في بعضها والرمز $<$ في بعضها الآخر، فلا تكون الصيغ قابلة للمقارنة فيما بينها؛ إذ إن الصيغة الناشئة من صيغة أخرى بالطي تقترن دائمًا بمنظومة موجبة من قيم الطيات الأساسية، أي، في حالتنا، بقيم موجبة أو منعدمة لـ$\beta_\rho-\alpha_\rho$.}؛
\item إذا تحقق رمز المساواة عند جميع قيم $\rho$، ولكن جُمعت صفوف رموز أقل في جداء مصفوفات واحد. ويتبين أن هذا التطبيع ضروري بسبب الاختزالات في \S\ 8. وبناءً عليه تصبح، مثلًا، $\pair{S^{n-1}}{T_{n-1}}$ أعلى من $\pair{S^{n-1}}{[S^1T^\rho]_{n-\rho-1}p_\rho}$.
\item إذا تحقق رمز المساواة عند جميع قيم $\rho$ وكان عدد صفوف الرموز المجموعة في جداء مصفوفات واحد متساويًا، عُدَّت $B$ أعلى من $A$ بواسطة تطبيع اعتباطي في ذاته، لكنه مثبت مرة واحدة ونهائيًا. وهذا التطبيع ممكن دائمًا بسبب العدد المنتهي للصيغ العائدة إلى منظومة قيم $\alpha$؛ ويمكن تحقيقه، مثلًا، بتمييز الصيغة $S$ (عدد أكبر من صفوف الرموز $S$، أو مجموع أعلى لمؤشرات الوزن العلوية لصفوف $S$، وهكذا).
\end{enumerate}
يمكننا، على نحو ملائم، أن نتصور سلسلة الملازمات مرتبةً في هيئة تجمّع شبكي من نقاط الشبكة ذي أبعاد عددها $(n-1)$، بحيث تقابل الاتجاهات الـ$n-1$ الطيات الأساسية الـ$n-1$. وعندئذ تقترن بكل صيغة منظومة قيم واحدة لا غير هي $\alpha$، أي نقطة شبكية ذات إحداثيات صحيحة $\alpha_\rho\ge0$؛ أما الصيغ المختلفة العائدة إلى منظومة قيم $\alpha$ فيعيّن التطبيع السابق ترتيبها تعيينًا وحيدًا. وكل صيغة اعتباطية في سلسلة الملازمات تكون الصيغة الابتدائية لسلسلة ملازمات جديدة محتواة بوصفها جزءًا من السلسلة الكلية؛ وهي محتواة تحديدًا في الصيغ التي نحصل عليها بالانطلاق من الصيغة الابتدائية الجديدة والتقدم في الاتجاهات الموجبة الـ$n-1$.

وبمقتضى المبرهنة الثامنة والنظرية العامة للتوسعات المتسلسلة، تصح لسلاسل الملازمات الآن بالضبط مبرهنات الاختزال المعروفة في المجالين الثنائي والثلاثي (قارن الأطروحة، \S\ 3)، وسنشتق بإيجاز مبرهنتها الرئيسة. ونقدم أولًا التعريف:

\glqq{}لتكن، في منظومة معطاة من الصيغ، مجموعة منتهية معيّنة من الصيغ قد ضُمّت في مقياس؛ ونسمي صيغةً \emph{قابلةً للاختزال} إذا أمكن التعبير عنها بصيغ تكون اللامتغيرات عوامل فيها، أو بـ\glqq{}صيغ أعلى\grqq{}؛ أي بصيغ أعلى من سلسلة الملازمات الكلية التي تنتمي إليها الصيغة، أو بصيغ تحتوي رموز المقياس برتبة أعلى. ونسمي سلسلة الملازمات القابلة للاختزال \emph{مختزلًا}.\grqq{} وعندئذ تصح

\emph{المبرهنة التاسعة: إذا كانت الصيغة الابتدائية لسلسلة ملازمات قابلةً للاختزال من حيث إنها نشأت بالطي مع مختزل، فإن سلسلة الملازمات كلها، الناشئة من هذه الصيغة الابتدائية، قابلة للاختزال.}

وللبرهان، نلاحظ أن طي الصيغة $f$ بصيغة ثانية $g$ يولّد صيغة $h$ يمكن، بموجب (45.)، أن تُعَدَّ صيغة مشتملةً على عدة صفوف متغيرات متوافقة التدرج $p_\rho,p'_\rho$، هي الصيغة $f$. غير أن مثل هذه الصيغ يمكن، وفقًا للنظرية العامة للتوسعات المتسلسلة، تمثيلها بأشكال قطبية للملازمات الأولية الموافقة للتغير للصيغة $f$، وذلك بتطبيق التوسعة المتسلسلة تباعًا على كل زوج من صفوف المتغيرات المتوافقة التدرج $p_\rho,p'_\rho$. وتكون الملازمات الأولية الموافقة للتغير التي تظهر عندئذ، بحسب (38.)، هي الصيغ المنشأة بالطي المتوافق التدرج. ولما كان ترتيب تطبيق التوسعة المتسلسلة لا يزال اعتباطيًا، أمكننا، على وجه الخصوص، اختيار ترتيب يوافق توليد أعم الطيات من الطيات الأساسية. وعندئذ تصبح منظومة الملازمات الأولية الموافقة للتغير الناتجة مطابقة لسلسلة الملازمات $f$؛ وتندرج الصيغ الناشئة بالطي المتوافق التدرج ذي النقص الأعلى ضمن الصيغ الناشئة بالطيات الأساسية. غير أننا نصل أيضًا إلى الأشكال القطبية لسلسلة الملازمات $f$ مع أي ترتيب اعتباطي للتوسعة المتسلسلة، وليقابل هذا الترتيب منظومة ثانية من الملازمات الأولية الموافقة للتغير. فلما كانت سلسلة الملازمات تستوعب جميع التكوينات اللامتغيرة، أمكن التعبير خطيًا عن كل صيغة في المنظومة الثانية بدلالة الأشكال القطبية لسلسلة الملازمات، وتحديدًا الأشكال القطبية للصيغ التي يظهر فيها نقصان مساوٍ أو أكبر في الوزن. وينتج هذا الأخير من أن الأشكال القطبية يمكن، كما في \S\ 7، أن تُعَدَّ لامتغيرات متزامنة لسلسلة ملازمات $H$ في (52.)، ذات درجات مناظرة للصيغة $h$، ولسلسلة الملازمات $f$. ويقابل كل طي لـ$H$ في نفسه نقصانٌ في الوزن ينتقل، بعد إجراء التعويض (11.)، إلى صفوف الرموز ومن ثم إلى صيغ $f$\srcfn{*)}{يقوم هذا الاعتبار أيضًا مقام الأساس للصياغة الثانية لمفهوم التكافؤ (ختام \S\ 8). وتوجد تفصيلات أخرى عن المنظومات المختلفة من الملازمات الأولية الموافقة للتغير في المجال الثلاثي عند شتودي، الموضع المذكور، II. \S\ 7، المبرهنتان 7) و8). وبمقتضى مبرهنتنا السابعة، تصح التفصيلات نفسها أيضًا لـ$n$ متغيرات.}.

كل صيغة اعتباطية من سلسلة الملازمات $h$ تنشأ بطي $h$ في نفسه؛ أي، من جهة، بطي أعلى لـ$g$ مع $f$، ومن جهة أخرى بطي $f$ أو $g$ في نفسه. وفي الحالتين تفضي التوسعة المتسلسلة، مثلما بُيّن أعلاه للصيغة الابتدائية $h$، إلى الأشكال القطبية لسلسلة الملازمات $f$. وإذا كان $f$، على وجه الخصوص، مختزلًا، نتجت توسعة بدلالة الأشكال القطبية للمختزل؛ ومن ثم تكون سلسلة الملازمات $h$ قابلة للاختزال.

وتنتج تطبيقات المبرهنة على اختزال منظومات الصيغ الكاملة على نحو مماثل لما في المجالين الثنائي والثلاثي.


\clearpage
\setcounter{footnote}{0}
\providecommand{\Kfield}{\mathfrak{K}}
\providecommand{\Ssys}{\mathfrak{S}}
\providecommand{\Mfield}{\mathfrak{M}}
\providecommand{\tq}{\tau}
\providecommand{\ds}{\displaystyle}

\begin{center}
{\Large\bfseries \srcLTR{5.}~\srcRTL{حقول الدوال النسبية.}\srcfn{1)}{محاضرة أُلقيت في اجتماع علماء الطبيعة في فيينا سنة \srcLTR{1913}.}}\par
\vspace{1.2em}
\srcRTL{بقلم \srcspaced{إيمي نوتر} في إرلنغن.}\par
\medskip
\srcLTR{J. Ber. d. DMV 22 (1913), S. 316--319}
\end{center}

\vspace{2em}

ترجع المسائل الآتية في أصلها إلى محادثات مع إ. فيشر. وقد سبق لإ. شتاينتز أن طرح بعض هذه المسائل وحلّها للحالة الخاصة لمتغير واحد، بل تحت فروض أعم على نطاق المعاملات.\srcfn{2)}{\srcLTR{E. Steinitz: Algebraische Theorie der Körper} (ولا سيما \S\ 24). \srcLTR{Crelles Journal 137. 1910.}}

أفهم بـ\glqq{}حقل دوال نسبية\grqq{} حقلًا عناصره دوال نسبية في $n$ من المتغيرات غير المعيّنة، ومعاملاتها من حقل أعداد مفروض قد يشمل، على وجه الخصوص، جميع الأعداد العقدية. ومن أمثلة هذه الحقول حقل الدوال المتناظرة في $n$ من الكميات، أو، على وجه أعم، مجموع الدوال النسبية في $n$ من المتغيرات غير المعيّنة التي تقبل تبديلات زمرة معيّنة (حقول النوع اللاغرانجية)؛ وكذلك حقل اللامتغيرات.

وثمة فرق مميّز بين الحقلين المذكورين أولًا والحقل الأخير: فالأولان يحتويان \emph{$n$ دوال مستقلة جبريًا}، هي الدوال المتناظرة الأولية؛ في حين أن عدد اللامتغيرات المستقلة جبريًا أصغر دائمًا من عدد المتغيرات غير المعيّنة. لذلك من المهم أن يكون من الممكن \emph{عمومًا إقامة تقابل واحد لواحد بين مثل هذه الحقول من النوع الثاني والحقول من النوع الأول}، وذلك باستبدال بعض المتغيرات غير المعيّنة بأعداد. ولهذا سأقتصر فيما يأتي \emph{على الحقول من النوع الأول، أي الحقول ذات $n$ دوال مستقلة جبريًا}؛ وبفضل ذلك التقابل تصح النتائج بعدئذ بوجه عام.

تنتظم المسائل حول ثلاثة مفاهيم أساسية: القاعدة النسبية، والقاعدة الدنيا، والقاعدة التكاملية.

1. أفهم بـ\glqq{}القاعدة النسبية\grqq{} عددًا منتهيًا من دوال الحقل بحيث يمكن تمثيل كل دالة في الحقل بتركيب نسبي من هذا العدد المنتهي، بمعاملات من نطاق المعاملات المفروض. وتُثبت، باعتبارات بسيطة ــ توجد أيضًا عند إ. شتاينتز في الموضع المذكور لحالة متغير واحد ــ قاعدة نسبية لكل حقل دوال نسبية. فمثلًا، وفقًا لمبرهنة لاغرانج، تؤلف الدوال المتناظرة الأولية، مع دالة تنتمي إلى الزمرة، قاعدة نسبية لحقول النوع اللاغرانجية المذكورة آنفًا.

2. يجب أن تحتوي كل قاعدة نسبية (للحقول من النوع الأول) على $n$ دوال على الأقل؛ فإن احتوت $n$ دوال بالضبط، وكانت لذلك مستقلة جبريًا بالضرورة، سميتها \glqq{}قاعدة دنيا\grqq{}. ويمكن الإجابة جزئيًا عن مسألة وجود القاعدة الدنيا بواسطة مبرهنات لوروت وكاستلنووفو وإنريكيس.\srcfn{1)}{\srcLTR{J. Lüroth: Beweis eines Satzes über rationale Kurven. Math. Ann. 9. 1875. --- G. Castelnuovo: Sulla razionalità delle involuzioni piane. Math. Ann. 44. 1893. --- F. Enriques: Sopra una involuzione non razionale dello spazio. Rend. Acc. Linc. vol. 21. 21. Jan. 1912.}} وبموجبها توجد دائمًا قاعدة دنيا لحقول متغير واحد: إنها دالة لوروت للحقل، المحددة تحديدًا وحيدًا حتى تحويل خطي كسري (قارن شتاينتز، \S\ 24). أما لحقول متغيرين، فإذا سُمح بتمديد جبري لنطاق المعاملات وُجدت قاعدة دنيا دائمًا، وبوجه عام لا توجد إلا عند السماح بهذا التمديد؛ في حين لا توجد عمومًا أي قاعدة دنيا لثلاثة متغيرات فأكثر. ولم تُبذل بعد محاولة لتوصيف الحقول الخاصة التي لها قاعدة دنيا.

وقد واصلت الآن بحث مسألة القاعدة الدنيا في حقول النوع اللاغرانجية. وهي تكتسب هنا المعنى الآتي: \glqq{}إذا كان حقل النوع اللاغرانجي المنتمي إلى الزمرة $G$ يملك قاعدة دنيا، أمكن إنشاء مجموع المعادلات ذات الزمرة $G$ المقررة إنشاءً نسبيًا بتمثيل بارامتري؛ وذلك لكل حقل أعداد يحتوي نطاق معاملات القاعدة الدنيا ــ أي، على وجه الخصوص، لكل حقل أعداد كيفي إذا كان نطاق المعاملات هذا هو حقل الأعداد النسبية.\grqq{}

تقدم الزمرة التناظرية أبسط مثال؛ فالدوال المتناظرة الأولية تؤلف هنا قاعدة دنيا ذات معاملات من الأعداد النسبية. وينتج من ذلك، بوصفه تمثيلًا بارامتريًا، المعادلة ذات المعاملات غير المعيّنة فحسب. أما كيفية الانتقال منها، في كل حقل أعداد، إلى عدد كيفي من المعادلات ذات زمرة غالوا التناظرية، فتبيّنها مبرهنة هيلبرت في عدم القابلية للاختزال.

وينتج الآن مباشرة من مبرهنتي لوروت وكاستلنووفو وجود القاعدة الدنيا لجميع الزمر التي تظهر في معادلات الدرجتين الثالثة والرابعة؛ كما يتبين أيضًا أن معاملات هذه القاعدة الدنيا من الأعداد النسبية. ومن ثم يمكن، لكل حقل أعداد كيفي، إنشاء مجموع معادلات الدرجتين الثالثة والرابعة ذات الزمرة المقررة إنشاءً نسبيًا. وتوجد للزمر الدورية والديهدريالية قاعدة دنيا يكون نطاق معاملاتها حقل جذور الوحدة من الدرجة $n$؛ ويصح الأمر نفسه لبعض الزمر الفوق دورية الأخرى. أما مسألة ما إذا كان وجود القاعدة الدنيا خاصية مميزة لبعض الزمر أصلًا، فلا بد أن تبقى غير محسومة.

3. للوصول إلى مفهوم القاعدة الأخير، أنظر في مجموع كثيرات الحدود الموجودة في الحقل. وأفهم بـ\glqq{}القاعدة التكاملية\grqq{} عددًا منتهيًا من كثيرات الحدود هذه بحيث يمكن تمثيل كل كثيرة حدود في الحقل بتركيب كثير حدود من هذا العدد المنتهي، بمعاملات من نطاق المعاملات المفروض. ومسألة وجود القاعدة التكاملية ــ التي تشمل، مثلًا، تناهي منظومة اللامتغيرات بوصفه حالة خاصة ــ طرحها هيلبرت، بصياغة مختلفة قليلًا، في \glqq{}مسائله الرياضية\grqq{}\srcfn{1)}{\srcLTR{D. Hilbert: Mathematische Probleme. Vortrag Paris 1900. Problem 14. Göttinger Nachr. 1900.}} بوصفها مسألة الدوال الصحيحة نسبيًا.

لا أستطيع في الوقت الحاضر أن أذكر إلا طائفة واحدة من الحقول التي توجد لها قاعدة تكاملية. فإذا احتوى الحقل منظومة من $n$ كثيرات حدود بحيث كانت المحصلة المتجانسة للأجزاء المتجانسة ذات أعلى درجة في كثيرات الحدود هذه غير منعدمة، وُجدت قاعدة تكاملية؛ وتُعطى هذه القاعدة بمعاملات المعادلة غير القابلة للاختزال في الحقل للشكل الخطي الآتي، بالنسبة إلى جميع متغيرات $u$:
\[
  u_0=u_1x_1+u_2x_2+\cdots+u_nx_n.\srcfnmark{2)}
\]
\srcfntext{2)}{ترد هذه المعادلة غير القابلة للاختزال، في حالة متغير واحد، عند إ. شتاينتز في برهان مبرهنة لوروت.}
وإذا كانت قيمة المحصلة، على وجه الخصوص، مساوية لوحدة من نطاق المعاملات، كانت تكاملية التمثيل مضمونة أيضًا. ومن أمثلة هذه الطائفة من الحقول حقول النوع اللاغرانجية، لأن محصلة الدوال المتناظرة الأولية تساوي $\pm 1$. ومثال آخر (من دون تكاملية) هو جميع الحقول التي لا تحتوي إلا كثيرة حدود واحدة مستقلة جبريًا؛ فالقاعدة التكاملية هنا مطابقة لدالة لوروت لأصغر حقل جزئي يحتوي جميع كثيرات الحدود. وهكذا، كما هو معلوم، تكون جميع اللامتغيرات الإسقاطية النسبية كثيرة الحدود لصيغة تربيعية ذات $n$ متغيرات دوال كثيرة حدود في المميّز.

الشرط المذكور كافٍ فحسب، وليس ضروريًا بأي حال لوجود القاعدة التكاملية. ويمكن بسهولة إنشاء حقول تنعدم فيها محصلة كل $n$ كثيرات حدود، ومع ذلك تملك قاعدة تكاملية تعطيها معاملات تلك المعادلة غير القابلة للاختزال. وهنا ينشأ السؤال عما إذا كانت معاملات تلك المعادلة تعطي القاعدة التكاملية حتى في أعم حالة.

\clearpage

\begin{center}
{\Large\bfseries \srcLTR{6.}~\srcRTL{الحقول ومنظومات الدوال النسبية}}\par
\vspace{0.6em}
\srcLTR{Emmy Noether (Erlangen).}\par
\vspace{1.2em}
\srcLTR{Math. Ann. 76 (1915), S. 161--196}
\end{center}

\vspace{2em}

يعالج هذا البحث مسائل القواعد في منظوماتٍ أيًّا كانت من الدوال النسبية والدوال النسبية الصحيحة (كثيرات الحدود)؛ وتُظهر مناهج نظرية الحقول المستعملة أن جوهر المسألة هو حل هذه الأسئلة في الحقول المؤلفة من دوال نسبية ــ أي حقول الدوال النسبية\srcfn{*)}{قارن بلاغًا تمهيديًا بمناسبة اجتماع علماء الطبيعة في فيينا: «حقول الدوال النسبية» ــ \srcLTR{Jahresber. d. D. Math.-Ver. 22 (1913)} ــ، وفيه عرض عام للمسائل المطروحة وللنتائج الخاصة بحقول الدوال.} ــ، في حين ينتج تعميم النتائج على منظوماتٍ أيًّا كانت بوصفه نتيجة مترتبة.

لم يُعرف حتى الآن من مسائل القواعد للمنظومات العامة إلا وجود قاعدة المودول (\srcLTR{Modulbasis}) الذي تكفله مبرهنة هيلبرت (\srcLTR{Math. Ann. 36}). وفيما يأتي تُحل مسألة القابلية للتمثيل النسبي حلًّا تامًّا بإثبات وجود القاعدة النسبية لكل منظومة أيًّا كانت (\S\ 7)؛ أما القاعدة النسبية للحقول فتحصل بالفعل في \S\ 4. ويتيح وجود القاعدة النسبية أن ننطلق في جميع المواضع من الحقل أو المنظومة المعرَّفين تعريفًا مجردًا، فنتجنب بذلك الصعوبات التي لا تنشأ عن المنظومة ذاتها بل عن الاختيار الخاص للقاعدة النسبية، مثل ظهور مقامات خاصة أو نقاط أساس لدوال القاعدة.

وتُبحث للحقول أيضًا مسألة القاعدة الدنيا، أي قاعدة نسبية مؤلفة من دوال مستقلة جبريًا (\S\ 6). وتفضي مناهج نظرية الحقول كذلك إلى قاعدة نسبية مميّزة، هي قاعدة الإنفولوشن (\srcLTR{Involutionsbasis})\srcfn{**)}{اختير هذا الاسم لأن حقل الدوال النسبية يمثل، في التأويل الهندسي، أعم إنفولوشن في فضاء خطي.} (\S\ 5)، وتغدو هذه القاعدة جوهرية بخاصة في المناطق التكاملية المؤلفة من كثيرات حدود (\S\ 8). فهي تعطي هنا تمثيلًا ذا مقام ثابت (قوة لدالة تعيّنها المنطقة التكاملية)، على نحو معروف مثلًا في الحالة الخاصة للتمثيل النمطي للامتغيرات.

ومن القابلية للتمثيل النسبي تُستخلص الآن أوسع النتائج الممكنة للقابلية للتمثيل الصحيح النسبي، أي لمسألة التناهي بالمعنى الأضيق.

تعتمد جميع مبرهنات التناهي المعروفة حتى الآن على أن مبرهنة هيلبرت في قاعدة المودول تضمن التناهي لجميع المنظومات التي يستتبع فيها تمثيل
\[
  F=A_1f_1+A_2f_2+\cdots+A_kf_k
\]
تمثيلًا ثانيًا
\[
  F=B_1f_1+B_2f_2+\cdots+B_kf_k
\]
حيث يكون كل $B_i$ من عناصر المنظومة، ودرجته أقل من درجة $F$.

أما مبرهنة القاعدة النسبية ومناهج نظرية الحقول فتتيح، على النقيض من ذلك، استنتاج التناهي تحت فروض من نوع مختلف تمامًا.

وهكذا تنتج، في إجابة جزئية عن إحدى مسائل هيلبرت (\srcLTR{Math. Probleme 14})، طائفة من الدوال الصحيحة نسبيًا (مناطق صحيحة نسبيًا من النوع الأول) تشكل مناطق تكاملية منتهية التوليد، وتوفر قاعدة الإنفولوشن المذكورة آنفًا قاعدتها التكاملية (\S\ 10). وبالقياس إلى برهان هيلبرت لمبرهنة كرونيكر في المنظومة الأساسية للكميات الصحيحة جبريًا (\srcLTR{Volle Invariantensysteme, \S\ 2; Math. Ann. 42})، يمكن كذلك إثبات أن جميع المنظومات المنتظمة من كثيرات الحدود ــ أي جميع المنظومات التي يمكن تحويلها إلى منظومات خالية من نقاط الأساس ــ تملك قاعدة تكاملية (\S\ 12)، في حين يسهل إعطاء منظومات غير منتظمة لا قاعدة تكاملية لها. وكمثال بسيط لهذه الحقيقة نذكر المبرهنة الآتية: \glqq{}في كل منظومة كيفما كانت من عدد غير منتهٍ من كثيرات الحدود في متغير غير معيّن واحد، يوجد عدد منتهٍ من كثيرات الحدود هذه بحيث يمكن تمثيل كل كثيرة حدود في المنظومة بتركيب كثير حدود من ذلك العدد المنتهي.\grqq{} وتوجد أمثلة أخرى في \S\ 13.

وتأتي الفقرتان الأخيرتان (\S\S\ 14 و15)، أخيرًا، بتقوية لمبرهنات القواعد من جهة التكاملية.

ويوفر مبدأ النقل الوارد في \S\ 3 أداة أساسية للبحث؛ فبموجبه يكفي أن تُبحث جميع مسائل القواعد المطروحة هنا في المنظومات التي يتساوى فيها عدد المتغيرات غير المعيّنة مع عدد الدوال المستقلة جبريًا في المنظومة.

كان الدافع إلى هذا البحث محادثات مع السيد إ. فيشر، ولا سيما سؤال منه عن القاعدة الدنيا لحقول النوع اللاغرانجية. أما البحوث المتصلة بذلك، التي أفضت إلى إنشاء معادلات ذات زمرة مقررة (قارن البلاغ المذكور عن \glqq{}حقول الدوال النسبية\grqq{}، الجزء 2)، فقد أُرجئت إلى منشور لاحق.

يكمن التوسيع الجوهري في البحث الحاضر، بالقياس إلى البلاغ الأصلي، في أن مبرهنات القواعد التي لم تُذكر هناك إلا للحقول أو لمناطق تكاملية خاصة قد نُقلت هنا إلى منظومات أيًّا كانت. ويتم هذا النقل على نحو بسيط بواسطة مفهوم أصغر حقل حاوٍ لمنظومة كيفما كانت، بوصفه تعميمًا لحقل القسمة لمنطقة تكاملية (\S\ 7). وقد حملني على صياغة هذا المفهوم أن السيد ك. هنتزلت استطاع، بعد اطلاعه على القاعدة النسبية للحقول، أن يثبت وجود القاعدة النسبية للمنظومات كيفما كانت بطريق غير مباشر يمر بقاعدة المودول الهيلبرتية. كما قادتني إلى مبرهنة القاعدة التكاملية للمنظومات المنتظمة ملاحظة ك. هنتزلت أن الاستدلالات التي طبقتها على المناطق التكاملية تصح للمنظومات كيفما كانت متى خضعت للفروض نفسها؛ وأدين للسيد هنتزلت أيضًا ببعض الملاحظات الصغرى المتفرقة.

وأخيرًا تجدر الإشارة إلى أن القاعدة النسبية والقاعدة الدنيا للحقول، في حالة متغير غير معيّن واحد، موجودتان عند إ. شتاينتز في «النظرية الجبرية للحقول» (\srcLTR{Algebraische Theorie der Körper})، \S\ 24 (\srcLTR{Crelles Journ. 137})؛ وقد افترض هناك أن نطاق المعاملات حقل بأعم معنى. ولم يُبحث فيما يأتي إلى أي مدى تبقى المبرهنات المقدمة هنا صحيحة تحت هذه الفروض الأعم؛ وعلى كل حال تتطلب البراهين تعديلًا، لأن أدوات نظرية المحددات الدالية، مثلًا، تتعطل بالفعل في الحقول ذات المميّز $p$.

\section*{\S\ 1. حقول $\Kfield_{n\rho}$ ومنظومات $\Ssys_{n\rho}$.}

يتناول ما يأتي بحث منظومات من الدوال النسبية في $n$ من المتغيرات غير المعيّنة، وبخاصة حقول الدوال النسبية.

ومن ثم يُراد دائمًا بـ$f(x_1\cdots x_n)$، و$g(x_1\cdots x_n)$، \ldots\ أو اختصارًا بـ$f(x)$، و$g(x)$، \ldots\ دوال نسبية في $x_1\cdots x_n$؛ وتُفترض هذه الدوال في صيغة مختزلة ــ أي إن البسط والمقام أوليان فيما بينهما. ويُرمز إلى الدوال النسبية الصحيحة (كثيرات الحدود) بأحرف كبيرة: $F(x)$، و$G(x)$، \ldots\ ويُفترض نطاق المعاملات $\Omega$ حقل أعداد كيفما كان، ومن ثم يمكنه، بخاصة، أن يشمل جميع الأعداد العقدية؛ كما يجوز أن يحتوي $\Omega$ أيضًا على عدد منتهٍ من الوسائط.

وتُعدّ الكميات من $\Omega$، أي جميع الدوال من الدرجة صفر، منتميةً إلى المنظومة دائمًا.

وبعد هذه الاصطلاحات يمكن تعريف الحقل المؤلف من دوال نسبية ــ حقل الدوال النسبية ــ تعريفًا مجردًا.

\emph{التعريف \srcLTR{I}: تُسمى منظومة من الدوال النسبية حقلًا إذا استوفت الشروط الآتية:}
\begin{enumerate}
\item \emph{إذا احتوت $f(x)$، فإنها تحتوي أيضًا، لكل كمية $c$ من $\Omega$، على $c\cdot f(x)$.}
\item \emph{إذا احتوت $f(x)$ و$g(x)$، فإنها تحتوي دائمًا أيضًا على $f(x)+g(x)$، و$f(x)\cdot g(x)$، وعلى حاصل القسمة ــ إذا كان $g(x)\ne0$ ــ $f(x):g(x)$.}\srcfn{*)}{يجوز هنا أن تكون $f(x)$ و$g(x)$ أيضًا من الدرجة صفر، أي من كميات $\Omega$.}
\end{enumerate}
ومن الحقول الخاصة ما ينشأ بإلحاق عدد منتهٍ من الدوال النسبية
\[
  f_1(x),\ f_2(x),\ldots, f_k(x)
\]
بـ$\Omega$؛ ويُرمز إليها بـ
\[
  \Omega\bigl(f_1(x)\cdots f_k(x)\bigr)
  \quad\text{أو}\quad
  \Omega(f_1\cdots f_k)
\]
.\srcfn{**)}{يتبين في \S\ 4 أن هذه \glqq{}الحقول الخاصة\grqq{} تستنفد بالفعل مجموع الحقول كلها.} ومن بين هذه الحقول الخاصة نذكر أيضًا الحقل الناشئ بإلحاق $x_1,\ldots,x_n$ بـ$\Omega$، وهو
\[
  \Omega(x_1\cdots x_n)
\]
المطابق لمجموع الدوال النسبية في $x_1\cdots x_n$ ذات المعاملات من $\Omega$.

وندخل أيضًا (اتباعًا لإ. شتاينتز) مفهوم الحقل الوسيط:
\begin{quote}
\glqq{}يقع حقل $\Mfield$ بين $\Omega_1$ و$\Omega_2$ إذا كان يحتوي $\Omega_1$ وكان محتوىً في $\Omega_2$\grqq{}؛
\end{quote}
وعندئذ يمكن أيضًا صياغة التعريف \srcLTR{I} على النحو الآتي:

حقل الدوال النسبية في $n$ من المتغيرات غير المعيّنة هو حقل وسيط بين $\Omega$ و$\Omega(x_1\cdots x_n)$، تظهر فيه المتغيرات غير المعيّنة الـ$n$ ظهورًا فعليًا في دالة واحدة على الأقل.

وإلى جانب عدد المتغيرات غير المعيّنة، توجد لمنظومات الدوال النسبية كمية مميّزة أخرى، هي عدد الدوال المستقلة جبريًا، أو الرتبة الجبرية، ويحددها التعريف الآتي:

\emph{التعريف \srcLTR{II}: تكون لمنظومة من الدوال النسبية الرتبة الجبرية $\rho$ إذا أمكن تعيين $\rho$ دالة من المنظومة مستقلة جبريًا، بينما تكون كل مجموعة من $(\rho+1)$ دالة من المنظومة معتمدة جبريًا.}\srcfn{***)}{تُسمى $\rho$ دالة $f_1(x)\cdots f_\rho(x)$ مستقلة جبريًا إذا استلزمت كل علاقة
\[
F\bigl(f_1(x)\cdots f_\rho(x)\bigr)=0\quad\text{مطابقًا في }x_1\cdots x_n
\]
أن يكون
\[
F(\lambda_1\cdots \lambda_\rho)=0\quad\text{مطابقًا في }\lambda_1\cdots \lambda_\rho.
\]
وفي الحالة المقابلة تُسمى هذه الدوال معتمدة جبريًا.}

ويُرمز بالفهرس المزدوج إلى منظومات مؤلفة من عدد منتهٍ أو غير منتهٍ من الدوال النسبية في $n$ من المتغيرات غير المعيّنة، ورتبتها الجبرية $\rho$، على الصورة
\[
  \Ssys_{n\rho}
\]
. ويُقصد، على وجه الخصوص، بـ
\[
  \Kfield_{n\rho}
\]
حقل في $n$ من المتغيرات غير المعيّنة، ورتبته الجبرية $\rho$. وتصح المتباينة:
\[
  1\leq \rho\leq n.
\]

ونورد أيضًا بعض الأمثلة على الحقول $\Kfield_{nn}$ و$\Kfield_{n\rho}$:

1. من الحقول $\Kfield_{nn}$ حقول النوع اللاغرانجية، ويمكن تعريفها بأنها
\begin{quote}
\glqq{}مجموع الدوال النسبية (والدوال النسبية الصحيحة) في $x_1\cdots x_n$ التي تقبل تبديلات زمرة تبديلات $G$ على $x_1\cdots x_n$.\grqq{}
\end{quote}
وهي تحتوي دائمًا حقل الدوال المتناظرة، ومن ثم تحتوي أيضًا الدوال المتناظرة الأولية الـ$n$ المستقلة جبريًا.

2. من الحقول $\Kfield_{n\rho}$ حقول اللامتغيرات الإسقاطية، وهي مؤلفة من مجموع اللامتغيرات النسبية (والنسبية الصحيحة) لشكل أساسي واحد أو لعدة أشكال أساسية.\srcfn{*)}{من الملائم ألا تُبحث إلا التحويلات ذات المحدِّد $1$؛ فعندئذ يقبل كل تركيب نسبي لعدد كيفما كان من اللامتغيرات التحويلَ أيضًا، ونحصل بالفعل على حقل. أما اللامتغيرات المتجانسة المتساوية الوزن، إذا أُخذت وحدها، فلا تؤلف حقلًا، لكنها تؤلف منظومة $\Ssys_{n\rho}$.} ويكون $\rho<n$، لأن عدد اللامتغيرات المستقلة جبريًا أصغر دائمًا من عدد المعاملات.

ويصح نظير ذلك لحقول اللامتغيرات بالنسبة إلى الزمر الجزئية من الزمرة الإسقاطية.

\section*{\S\ 2. قضية مساعدة في الدوال المعتمدة جبريًا.}

سيتبين في \S\ 3، بفضل القضية المساعدة في هذه الفقرة، أن الرتبة الجبرية لمنظومة هي العدد المميّز لها على وجه الدقة. فالقضية المساعدة تبيّن أن تعيين قيم عددية على التعاقب لبعض المتغيرات غير المعيّنة، على نحو تبقى معه $\rho$ دالة مستقلة جبريًا على استقلالها الجبري، لا يمكن أن يجعل أي دالة معتمدة جبريًا على هذه الدوال الـ$\rho$ تنعدم انعدامًا مطابقًا أو تصير غير معيّنة.

لتكن إذن المنظومة
\[
\begin{array}{@{}r@{\quad}l@{}}
(1)& g_1(x),\ g_2(x),\ldots,g_\rho(x)
\end{array}
\]
ذات الرتبة الجبرية $\rho$ --- حيث $\rho<n$ --- ولتكن $h(x)$ دالة نسبية معتمدة جبريًا على دوال (1). ومن مبرهنات معلومة أن رتبة المصفوفة اليعقوبية للمنظومة (1) تساوي $\rho$. فلنسمِّ المتغيرات غير المعيّنة، مثلًا، بحيث
\[
\begin{array}{@{}r@{\quad}l@{}}
(2)& \displaystyle
\frac{\partial(g_1\ g_2\cdots g_\rho)}{\partial(x_1\ x_2\cdots x_\rho)}
=\frac{\Gamma(x)}{G(x)^{\rho+1}}\ne0,
\end{array}
\]
حيث يرمز $G(x)$، كما هو معلوم، إلى المقام المشترك لدوال (1).

\emph{لنختر الآن العدد $a$ بحيث لا يكون حاصل الضرب $G(x)\cdot\Gamma(x)$ قابلًا للقسمة على $(x_i-a)$}، على أن يكون $i$ واحدًا من الأعداد $\rho+1,\rho+2,\ldots,n$. وعليه، لا يمكن للمقام المشترك لدوال (1) أن ينعدم عند $x_i=a$، كما أن الدوال الـ$\rho$
\[
\begin{array}{@{}r@{\quad}l@{}}
(3)& [g_1(x)]_{x_i=a},\ldots,[g_\rho(x)]_{x_i=a}
\end{array}
\]
مستقلة جبريًا هي أيضًا. وسنقول إن الرتبة الجبرية للمنظومة (1) تبقى محفوظة بالتعويض $x_i=a$.

ولتكن المعادلة غير القابلة للاختزال التي تحققها $h(x)$ بالنسبة إلى دوال (1) معطاة كما يأتي:
\[
\begin{array}{@{}r@{\quad}l@{}}
(4)& A_0(g_1\cdots g_\rho)\cdot h(x)^\tau
 +A_1(g_1\cdots g_\rho)\cdot h(x)^{\tau-1}
 +\cdots+A_\tau(g_1\cdots g_\rho)=0,
\end{array}
\]
حيث لا ينعدم $A_0(g_1\cdots g_\rho)$ ولا $A_\tau(g_1\cdots g_\rho)$ انعدامًا مطابقًا. وللوصول من ذلك إلى متطابقة بين كثيرات الحدود، نضع:
\[
\begin{array}{@{}r@{\quad}l@{}}
(5)& g_i(x)=G_i(x):G(x);\qquad h(x)=H_1(x):H_2(x),
\end{array}
\]
ونلاحظ أنه --- بسبب افتراض كون تمثيل $h(x)$ في صيغة مختزلة --- يكون $H_1(x)$ و$H_2(x)$ أوليين فيما بينهما. وبمقتضى (5) تتحول (4) إلى:
\[
\begin{array}{@{}r@{\quad}l@{}}
(6)& \begin{aligned}[t]
&B_0(G\,G_1\cdots G_\rho)\cdot G(x)^{\sigma_0}\cdot H_1(x)^\tau
+B_1(G\,G_1\cdots G_\rho)\cdot G(x)^{\sigma_1}\cdot H_1(x)^{\tau-1}H_2(x)\\
&\quad+\cdots+B_\tau(G\,G_1\cdots G_\rho)\cdot G(x)^{\sigma_\tau}\cdot H_2(x)^\tau=0,
\end{aligned}
\end{array}
\]
حيث تكون $B_k$ أشكالًا متجانسة في $G,G_i$، ويلزم أن يكون واحد على الأقل من الأسس غير السالبة $\sigma_k$ مساويًا للصفر.

ولنفترض الآن أن $H_1(x)$ قابل للقسمة على $(x_i-a)$؛ فعندئذ، بحسب (6)، يكون أيضًا
\[
  B_\tau(G\,G_1\cdots G_\rho)\cdot G(x)^{\sigma_\tau}\cdot H_2(x)^\tau
\]
قابلًا للقسمة على $(x_i-a)$. وهذا مستبعد بالفرض بالنسبة إلى $G(x)$؛ ومن ثم فإن قابلية $B_\tau$ للقسمة تستلزم أيضًا قابلية
\[
  A_\tau=B_\tau:G^\lambda
\]
للقسمة، فتقوم بين دوال (3) العلاقة $A_\tau=0$، خلافًا للاستقلال الجبري المفترض.
وأخيرًا، لا يمكن كذلك أن يكون $H_2(x)$، لكونه أوليًا فيما بينه وبين $H_1(x)$، قابلًا للقسمة على $(x_i-a)$\srcfn{*)}{ما عاد هذا الاستنتاج الأخير، المستند إلى الصيغة المختزلة، ممكنًا عند إجراء تعويض واحد متزامن: $x_i=a$ و$x_k=b$، مع بقاء سائر الفروض. فقد ينعدم $H_1(x)$ و$H_2(x)$ معًا، فيصير خارج القسمة غير معيّن بالتعويض $=\frac00$؛ كما في المثال
\[
 h(x)=\frac{x_3(\alpha x_1+\beta x_2)+x_4(\gamma x_1+\delta x_2)}
 {x_3(\alpha' x_1+\beta' x_2)+x_4(\gamma' x_1+\delta' x_2)}
\]
عند التعويضين $x_3=0$ و$x_4=0$.}؛ ومن ثم يؤدي افتراض قابلية $H_1(x)$ للقسمة على $(x_i-a)$ إلى تناقض. ويُبيَّن بالمثل أن $H_2(x)$ أيضًا لا يمكن أن يكون قابلًا للقسمة على $(x_i-a)$. وهكذا لا يمكن للدالة $[h(x)]_{x_i=a}=k(x)$ أن تنعدم انعدامًا مطابقًا ولا أن تصير غير معيّنة.

يمكن تطبيق الاستنتاج السابق مرة أخرى على المنظومة (3) والدالة $k(x)$، التي يُفترض من جديد أنها مكتوبة في صيغة مختزلة؛ ويؤدي استمرار هذا الإجراء في النهاية إلى

\emph{القضية المساعدة: لتكن كل من المنظومتين}
\[
\begin{array}{c}
g_1(x),\ g_2(x),\ldots,\ g_\rho(x),\\
g_1(x),\ g_2(x),\ldots,\ g_\rho(x),\ h(x)
\end{array}
\]
\emph{ذات رتبة جبرية $\rho$ --- حيث $\rho<n$ --- ولنفرض، مثلًا، أن}
\[
 \frac{\partial(g_1\cdots g_\rho)}
 {\partial(x_1\cdots x_\rho)}
 =\frac{\Gamma(x)}{G(x)^{\rho+1}}\ne0.
\]
\emph{ولتكن الأعداد}
\[
 a_n,\ a_{n-1},\ldots,\ a_{\rho+1}
\]
\emph{مختارة بحيث لا ينعدم، بالتعويض}
\[
 x_n=a_n,\quad x_{n-1}=a_{n-1},\ldots,\quad x_{\rho+1}=a_{\rho+1}
\]
\emph{حاصل الضرب $G(x)\cdot\Gamma(x)$ --- أي بحيث تبقى الرتبة الجبرية للمنظومة (1) محفوظة. وفي التعويضات المتعاقبة}
\[
\begin{array}{c}
[h(x)]_{x_n=a_n}=h^{(n-1)}(x);\quad
[h^{(n-1)}(x)]_{x_{n-1}=a_{n-1}}=h^{(n-2)}(x);\quad\cdots;\\[0.25em]
[h^{(\rho+1)}(x)]_{x_{\rho+1}=a_{\rho+1}}=h^*(x_1,x_2,\ldots,x_\rho)
\end{array}
\]
\emph{لنفترض أن كلًا من الدوال $h(x),h^{(n-1)}(x),\ldots,h^{(\rho+1)}(x)$ قد كُتبت في صيغة مختزلة. في ظل هذه الفروض لا يمكن لبسط الدالة $h^*(x_1\cdots x_\rho)$ ولا مقامها أن ينعدم انعدامًا مطابقًا، كما لا يمكن للدالة أن تصير $\frac00$.}\srcfn{**)}{يعطي التعويض المتعاقب، مثلًا، في الدالة $h(x)$ من الحاشية السابقة:
\[
 [h(x)]_{x_4=0}=h^{(3)}(x)=\frac{\alpha x_1+\beta x_2}{\alpha' x_1+\beta' x_2},\qquad
 h^*(x_1,x_2)=\frac{\alpha x_1+\beta x_2}{\alpha' x_1+\beta' x_2}.
\]}

\section*{\S\ 3. تقابل واحد لواحد بين المنظومات $\Ssys_{n\rho}$ والمنظومات $\Ssys^*_{\rho\rho}$.}

سنستخرج مباشرة من القضية المساعدة تقابلًا بين المنظومات $\Ssys_{n\rho}$ والمنظومات $\Ssys^*_{\rho\rho}$، وذلك بإحلال أعداد محل بعض المتغيرات غير المعيّنة؛ وبذلك تظهر الرتبة الجبرية $\rho$ بوصفها العدد المميّز بحقّ.

لما كانت المنظومة $\Ssys_{n\rho}$ ذات رتبة جبرية $\rho$، توجد $\rho$ دالة من دوال $\Ssys_{n\rho}$، هي
\[
\begin{array}{@{}l@{\quad}c@{}}
(1)& g_1(x)\cdots g_\rho(x),
\end{array}
\]
وهي مستقلة جبريًا. ولتكن، إضافة إلى ذلك، $f(x)$ دالة كيفما كانت من $\Ssys_{n\rho}$؛ فعندئذ تكون المنظومة
\[
\begin{array}{@{}l@{\quad}c@{}}
(2)& g_1(x)\cdots g_\rho(x),\ f(x)
\end{array}
\]
ذات رتبة جبرية $\rho$ أيضًا؛ وتحقق المنظومتان (1) و(2) شروط القضية المساعدة.\srcfn{*)}{تُعاد ترقيم المتغيرات غير المعيّنة عند اللزوم.}

لذلك نعيّن الأعداد
\[
 a_n,\ a_{n-1},\ldots,\ a_{\rho+1}
\]
وفق القضية المساعدة، بحيث تبقى الرتبة الجبرية $\rho$ للمنظومة (1) محفوظة تحت التعويض ــ وهذا ممكن على وجوه كثيرة كيفما شئنا ــ؛ فعندئذ لا يمكن، في الدالة المعرّفة بواسطة
\[
\begin{array}{c}
[f(x)]_{x_n=a_n}=f^{(n-1)}(x);\quad
[f^{(n-1)}(x)]_{x_{n-1}=a_{n-1}}=f^{(n-2)}(x);\quad\cdots;\\[0.25em]
[f^{(\rho+1)}(x)]_{x_{\rho+1}=a_{\rho+1}}=f^*(x_1\cdots x_\rho)
\end{array}
\]
أي في الدالة
\[
\begin{array}{@{}l@{\quad}c@{}}
(3)& f^*(x_1\cdots x_\rho)
\end{array}
\]
أن ينعدم البسط أو المقام انعدامًا مطابقًا.

والآن لتجرِ $f(x)$ على جميع دوال $\Ssys_{n\rho}$، سواء كان عددها منتهيًا أم غير منتهٍ. وننظر في المنظومة المؤلفة من جملة جميع الدوال (3)، ولها الخصائص الآتية:

1. رتبتها الجبرية $\rho$؛ إذ تحتوي الدوال الناشئة من (1)،
\[
g_1^*(x_1\cdots x_\rho),\ldots,g_\rho^*(x_1\cdots x_\rho),
\]
وهي مستقلة جبريًا بسبب اختيار $a_n,a_{n-1},\ldots,a_{\rho+1}$. ولذلك نسمّي هذه المنظومة
\[
\Ssys^*_{\rho\rho}
\]
.\srcfn{**)}{يُراد على الدوام، تبعًا لذلك، من $f^*(x)$ و$g^*(x)$ و\ldots{} دوال الصورة، أي دوال في $x_1\cdots x_\rho$.}

2. التقابل بين الدالتين $f(x)$ و$f^*(x)$ واحد لواحد بلا استثناء.

ذلك أنه إذا كانت $f_1(x)$ و$f_2(x)$ تنتميان إلى $\Ssys_{n\rho}$، فإن الفرق
\[
d(x)=f_1(x)-f_2(x)
\]
يكون أيضًا معتمدًا جبريًا على المنظومة (1)؛ وكذلك لدينا:
\[
d^*(x)=f_1^*(x)-f_2^*(x).
\]
ولما كانت $d(x)$ مع المنظومة (1) تحقق شروط القضية المساعدة، فإنه يلزم من
\[
d(x)\ne0
\]
بالضرورة:
\[
d^*(x)\ne0;
\]
أي إن الدوال المختلفة من $\Ssys_{n\rho}$ تقابلها دوال مختلفة من $\Ssys^*_{\rho\rho}$.

3. من كل علاقة بين دوال من $\Ssys^*_{\rho\rho}$:
\[
F\bigl(f_1^*(x),f_2^*(x),\ldots,f_k^*(x)\bigr)=0
\quad\text{مطابقًا في }x_1\cdots x_\rho
\]
تترتب، للدوال المناظرة لها على نحو وحيد من $\Ssys_{n\rho}$، العلاقة:
\[
F\bigl(f_1(x),f_2(x),\ldots,f_k(x)\bigr)=0
\quad\text{مطابقًا في }x_1\cdots x_n.
\]
ذلك أن كل تركيب نسبي لهذه الدوال، شأنه شأن $f_1(x),f_2(x),\ldots,f_k(x)$، معتمد جبريًا على المنظومة (1).

وينتج كذلك أن العلاقة الأولى الآتية تستلزم الثانية:
\[
\begin{array}{@{}l@{\quad}l@{}}
(4)& h(x)=F\bigl(f_1(x)\cdots f_k(x)\bigr):\\
& h^*(x)=F\bigl(f_1^*(x)\cdots f_k^*(x)\bigr).
\end{array}
\]
وبحسب القضية المساعدة، ينتج من تحقق شرط عدم الصفر الأول في العرض الآتي أن يتحقق الثاني:
\[
h(x)\ne0\quad\text{أيضًا}\quad h^*(x)\ne0;
\]
وهذا ما كان مطلوبًا إثباته.

وخلاصة ذلك أننا نحصل على

\emph{المبرهنة I: يمكن أن نقابل كل منظومة $\Ssys_{n\rho}$ مؤلفة من عدد منتهٍ أو غير منتهٍ من الدوال النسبية ــ على وجوه كثيرة كيفما شئنا ــ بمنظومة $\Ssys^*_{\rho\rho}$ تقابلًا واحدًا لواحد، بحيث تظل جميع العلاقات القائمة بين دوال $\Ssys^*_{\rho\rho}$ قائمة أيضًا في $\Ssys_{n\rho}$، والعكس بالعكس. وعلى وجه الخصوص، يقابل الحقل $\Kfield_{n\rho}$ وفق (4) حقلًا $\Kfield^*_{\rho\rho}$ أيضًا.}

ونستخلص من المبرهنة I أيضًا النتيجة الآتية التي تبسّط جميع البراهين اللاحقة: \emph{إن خواص المنظومات $\Ssys_{n\rho}$ التي تتميّز بعلاقات، منتهية العدد أو غير منتهية العدد، بين دوال المنظومة، تصح للمنظومات $\Ssys_{n\rho}$ متى صحت لمنظومة صورة $\Ssys^*_{\rho\rho}$.}

ونسمي هذه النتيجة اختصارًا \glqq{}مبدأ النقل\grqq{}.

وتقدم نظرية اللامتغيرات الجبرية أبسط مثال لهذا التقابل. ففي الواقع يمكن دومًا، بتحويل خطي، تثبيت بعض معاملات الشكل الأساسي عدديًا على نحو تظل معه جميع العلاقات الصحيحة للشكل الخاص صحيحة أيضًا في الحالة العامة.

\section*{\S\ 4. القاعدة النسبية للحقول $\Kfield_{n\rho}$.}

في البحوث الآتية، المتمحورة حول مفاهيم القاعدة، نستعمل على الدوام \glqq{}مبدأ النقل\grqq{} الوارد في \S\ 3. ونثبت أولًا وجود القاعدة للحقول، غير أننا نصوغ التعريفات على وجه عام:

\emph{التعريف الثالث: يُسمّى عدد منتهٍ من دوال $\Ssys_{n\rho}$ قاعدة نسبية إذا أمكن تمثيل كل دالة من $\Ssys_{n\rho}$ على صورة تركيب نسبي من هذا العدد المنتهي من الدوال، بمعاملات من $\Omega$.}

لا بد أن تشتمل القاعدة النسبية على $\rho$ من الدوال المستقلة جبريًا، لأن الرتبة الجبرية لمنظومة مشتقة من قاعدة نسبية تساوي الرتبة الجبرية للقاعدة النسبية.

نبين أولًا أن برهان وجود القاعدة النسبية يقبل تطبيق مبدأ النقل. ذلك أن التعريف الثالث يقبل الصياغة الآتية:

دوال $\Ssys_{n\rho}$
\[
g_1(x),\ g_2(x),\ldots,g_k(x)
\]
تكوّن قاعدة نسبية إذا وفقط إذا تحققت العلاقات اللامتناهية الآتية:
\[
\begin{array}{@{}l@{\quad}c@{}}
(1)& f(x)=\gamma\bigl(g_1(x)\cdots g_k(x)\bigr),
\end{array}
\]
حيث تجري $f(x)$ على جميع دوال $\Ssys_{n\rho}$، وتدل $\gamma$ على دالة نسبية ــ تتغير بتغير $f$ ــ في $k$ من الوسائط.

فإذا كانت في منظومة الصورة $\Ssys^*_{\rho\rho}$ قاعدة نسبية معطاة بالدوال
\[
g_1^*(x),\ g_2^*(x),\ldots,g_k^*(x),
\]
وهو ما تترتب عليه العلاقات
\[
f^*(x)=\gamma\bigl(g_1^*(x)\cdots g_k^*(x)\bigr)
\]
فإن المبرهنة الأولى تقضي بتحقق العلاقات (1) في $\Ssys_{n\rho}$؛ أي إن \emph{مبدأ النقل للقاعدة النسبية} ينص على ما يأتي:

\emph{من وجود قاعدة نسبية لمنظومة صورة $\Ssys^*_{\rho\rho}$ ينتج وجود قاعدة نسبية للمنظومة الأصلية $\Ssys_{n\rho}$.}\srcfn{*)}{ويصاغ مبدأ النقل على النحو المناظر لكل مسألة من مسائل القاعدة المعالجة هنا.}

ولإثبات وجود القاعدة النسبية لجميع الحقول $\Kfield_{n\rho}$، يمكننا إذن أن نقتصر على الحقول $\Kfield^*_{\rho\rho}$؛ وهنا يوصل إلى الغاية تعميم مباشر لطريقة استدلال استعملها إ. شتاينتس.\srcfn{*)}{إ. شتاينتس، \srcLTR{Algebraische Theorie der Körper}، \S\ 24، \srcLTR{Crelles Journ. 137 (1909).}}

لتكن
\[
\begin{array}{@{}l@{\quad}c@{}}
(2)& g_1^*(x_1\cdots x_\rho),\ldots,g_\rho^*(x_1\cdots x_\rho)
\end{array}
\]
منظومة مؤلفة من $\rho$ من الدوال المستقلة جبريًا من $\Kfield^*_{\rho\rho}$. وبحذف $x_1\cdots x_\rho$ نحصل على $\rho$ من العلاقات:
\[
\begin{array}{@{}l@{\quad}c@{}}
(3)& F_1\bigl(g_1^*(x)\cdots g_\rho^*(x),x_1\bigr)=0,\ldots,
F_\rho\bigl(g_1^*(x)\cdots g_\rho^*(x),x_\rho\bigr)=0,
\end{array}
\]
بحيث يظهر في كل علاقة $F_i=0$ المتغير غير المعيّن $x_i$ فعليًا؛ إذ لولا ذلك لوجد، خلافًا للفرض، اعتماد جبري بين الدوال (2). وتفيد العلاقات (3) أن
\[
x_1,x_2,\ldots,x_\rho
\]
عناصر جبرية بالنسبة إلى
\[
\Omega\bigl(g_1^*(x)\cdots g_\rho^*(x)\bigr)
\]
وأن الحقل الناشئ بإلحاق $x_1\cdots x_\rho$، وهو
\[
\Omega\bigl(g_1^*\cdots g_\rho^*,x_1\cdots x_\rho\bigr)=\Omega(x_1\cdots x_\rho)
\]
تمديد جبري منتهٍ لـ$\Omega(g_1^*\cdots g_\rho^*)$. وبحسب مبرهنة معروفة\srcfn{**)}{انظر مثلًا: \srcLTR{Weber, Lehrbuch d. Algebra 1}، \S\ 144 (الطبعة الأولى)، أو \S\ 55 (الطبعة المختصرة).} يكون $\Kfield^*_{\rho\rho}$ نفسه، بوصفه حقلًا وسيطًا بين $\Omega(g_1^*\cdots g_\rho^*)$ و$\Omega(x_1\cdots x_\rho)$، تمديدًا جبريًا لـ$\Omega(g_1^*\cdots g_\rho^*)$، في حين يكون $\Omega(x_1\cdots x_\rho)$ تمديدًا جبريًا لـ$\Kfield^*_{\rho\rho}$. ومن ثم ينشأ $\Kfield^*_{\rho\rho}$ من $\Omega(g_1^*\cdots g_\rho^*)$ بإلحاق عدد منتهٍ من عناصر $\Kfield^*_{\rho\rho}$، أو كذلك بإلحاق \emph{دالة واحدة} مختارة اختيارًا مناسبًا؛ أي إن
\[
\Kfield^*_{\rho\rho}
=\Omega(g_1^*\cdots g_\rho^*,g_{\rho+1}^*\cdots g_k^*)
=\Omega(g_1^*\cdots g_\rho^*g_0^*).
\]

وبمبدأ النقل نحصل من ثم على

\emph{المبرهنة الثانية: لكل حقل $\Kfield_{n\rho}$ قاعدة نسبية رتبتها الجبرية $\rho$؛ ويمكن بوجه خاص اختيار القاعدة النسبية بحيث لا تشتمل على أكثر من $\rho+1$ دالة.}

ولتكن منظومة كيفما كانت من $\rho$ من الدوال المستقلة جبريًا من $\Kfield_{n\rho}$ معطاة بالدوال $h_1(x)\cdots h_\rho(x)$؛ ويمكن بالطريقة السابقة استكمالها إلى قاعدة نسبية
\[
h_1(x)\cdots h_\rho(x);\ h_{\rho+1}(x)\cdots h_\sigma(x)
\]
وبحسب الخاصية التعريفية توجد العلاقات:
\[
\begin{array}{ll}
h_i(x)=\chi_i\bigl(g_1(x)\cdots g_k(x)\bigr)&(i=1,2,\ldots,\sigma),\\
g_j(x)=\psi_j\bigl(h_1(x)\cdots h_\sigma(x)\bigr)&(j=1,2,\ldots,k),
\end{array}
\]
حيث تدل $\chi_i$ و$\psi_j$ على دوال نسبية في وسائطها؛ أي:

ومن أي قاعدة نسبية كيفما كانت تنتج كل قاعدة نسبية أخرى بتحويل ثنائي نسبي، أي بتحويل نسبي ذي معكوس نسبي؛ وكذلك:

إذا كانت
\[
h_1(x)\cdots h_\rho(x)
\]
منظومة كيفما كانت من $\rho$ من الدوال المستقلة جبريًا من $\Kfield_{n\rho}$، فإن $\Kfield_{n\rho}$ تمديد جبري منتهٍ لـ
\[
\Omega\bigl(h_1(x)\cdots h_\rho(x)\bigr).
\]

\section*{\S\ 5. صيغة الإنفولوشن وقاعدة الإنفولوشن لحقول $\Kfield_{n\rho}$.}

كان يعيب كل قاعدة نسبية شُيّدت في \S\ 4 أنها تعتمد على الدوال الابتدائية المختارة اعتباطيًا؛ أما الآن فنثبت وجود \emph{قاعدة نسبية يعيّنها الحقل تعيينًا فريدًا}.

بحسب \S\ 4، يكون $\Omega(x_1\cdots x_\rho)$ حقلًا جبريًا ــ درجته، مثلًا، $i$ ــ فوق $\Kfield^*_{\rho\rho}$، وينشأ، بحكم
\[
\Kfield^*_{\rho\rho}(x_1\cdots x_\rho)=\Omega(x_1\cdots x_\rho)
\]
بإلحاق العناصر الجبرية بالنسبة إلى $\Kfield^*_{\rho\rho}$، وهي $x_1\cdots x_\rho$. ولذلك ننظر إلى كثيرة الحدود غير القابلة للاختزال بالنسبة إلى $\Kfield^*_{\rho\rho}$، وهي $\Phi^*(z,u)$، والتي لها الجذر
\[
z=u_1x_1+u_2x_2+\cdots+u_\rho x_\rho=u(x),
\]
ودرجة هذه الكثيرة هي $i$ في $z$. وإضافة إلى ذلك، فإن $\Phi^*(z,u)$، بوصفها حاصل ضرب الكميات المترافقة مع $[z-u(x)]$، شكل متجانس من الدرجة $i$ في $z,u_1\cdots u_\rho$؛ ومن ثم يكون معامل $z^i$ فيها دالة مستقلة عن $u$ ومنتمية إلى $\Kfield^*_{\rho\rho}$. فإذا جعلنا هذا المعامل الأعلى مساويًا لـ$1$ بالقسمة، تعيّنت كثيرة الحدود غير القابلة للاختزال $\Phi^*(z,u)$ تعيينًا فريدًا. وهذا الشكل المتجانس المعرّف على هذا النحو،
\[
\begin{array}{@{}l@{\quad}l@{}}
(1)& \displaystyle \Phi^*(z,u)=z^i+\sum{}'g^*_{\alpha\alpha_1\ldots\alpha_\rho}(x_1\cdots x_\rho)\cdot
z^\alpha u_1^{\alpha_1}\cdots u_\rho^{\alpha_\rho}\\
& \displaystyle\qquad(\alpha+\alpha_1+\cdots+\alpha_\rho=i)
\end{array}
\]
نسميه \emph{صيغة الإنفولوشن}\srcfn{*)}{يتضح سبب هذه التسمية في ختام الفقرة؛ وتدل $\sum'$ على أن الجمع يمتد إلى جميع الحدود ما عدا $z^i$.} للحقل $\Kfield^*_{\rho\rho}$؛ وينتج من (1)
\[
\begin{array}{@{}l@{\quad}l@{}}
(2)& \displaystyle \Phi(z,u)=z^i+\sum{}'g_{\alpha\alpha_1\ldots\alpha_\rho}(x)\cdot
z^\alpha u_1^{\alpha_1}\cdots u_\rho^{\alpha_\rho}\\
& \displaystyle\qquad(\alpha+\alpha_1+\cdots+\alpha_\rho=i)
\end{array}
\]
بوصفها \emph{صيغة إنفولوشن} للحقل $\Kfield_{n\rho}$.



ستتبيّن معاملات صيغة الإنفولوشن $\Phi^*(z,u)$ قاعدةً نسبيةً للحقل $\Kfield^*_{\rho\rho}$؛ أي إن العلاقة الآتية قائمة:
\[
\begin{array}{@{}l@{\quad}c@{}}
(3)& \Omega\bigl(g^*_{\alpha\alpha_1\ldots\alpha_\rho}(x_1\cdots x_\rho)\bigr)=\Kfield^*_{\rho\rho}\qquad(\alpha+\alpha_1+\cdots+\alpha_\rho=i).
\end{array}
\]
ولإثبات ذلك نلاحظ أن معاملات $\Phi^*(z,u)$ تنتمي إلى $\Omega(g^*_{\alpha\alpha_1\ldots\alpha_\rho}(x))$، وأن عدم قابليتها للاختزال بالنسبة إلى $\Kfield^*_{\rho\rho}$ يجعلها، من باب أولى، غير قابلة للاختزال بالنسبة إلى هذا الحقل الجزئي. ومن ثم تمثل $\Phi^*(z,u)$ المعادلة غير القابلة للاختزال لـ$u(x)$ بالنسبة إلى $\Omega(g^*_{\alpha\alpha_1\ldots\alpha_\rho}(x))$. وينتج من ذلك أن $\Omega(x_1\cdots x_\rho)$ حقل جبري ــ من الدرجة $i$ ــ فوق $\Omega(g^*_{\alpha\alpha_1\ldots\alpha_\rho}(x))$. ولما كان $\Kfield^*_{\rho\rho}$ واقعًا بين $\Omega(g^*_{\alpha\alpha_1\ldots\alpha_\rho}(x))$ و$\Omega(x_1\cdots x_\rho)$، فإن تساوي الدرجات يستلزم، كما هو معلوم، تساوي الحقول؛ فتنتج العلاقة (3). وبمبدأ النقل، تؤلف معاملات $\Phi(z,u)$، على النحو المناظر، قاعدة نسبية للحقل $\Kfield_{n\rho}$؛ أي إن لدينا

\emph{المبرهنة الثالثة: تؤلف معاملات صيغة الإنفولوشن قاعدة نسبية يعيّنها الحقل، هي قاعدة الإنفولوشن؛ وتكون هذه القاعدة معيّنة تعيينًا فريدًا للحقول $\Kfield^*_{\rho\rho}$.}\srcfn{*)}{أما للحقول $\Kfield_{n\rho}$ فقد تعتمد صيغة الإنفولوشن على التصوير؛ ويمكن تحصيل الوحدانية بتصوير ذي معاملات غير معيّنة.}

ونقدم أيضًا تفسيرًا غير نسبي لهذه الواقعة، يبيّن صلتها بالنظرية الهندسية للإنفولوشن.

لنفترض أن تحليل $\Phi^*(z,u)$ إلى عوامل خطية في حقل تمديد معطى على الصورة الآتية:
\[
\begin{array}{@{}l@{\quad}l@{}}
(4)& \Phi^*(z,u)=(z-u_1x_1-\cdots-u_\rho x_\rho)(z-u_1x_1^{(1)}-\cdots-u_\rho x_\rho^{(1)})\cdots\\[0.25em]
& \qquad\qquad\quad (z-u_1x_1^{(i-1)}-\cdots-u_\rho x_\rho^{(i-1)}).
\end{array}
\]
وتبين المقارنة مع (1) أن الدوال
\[
g^*_{\alpha\alpha_1\ldots\alpha_\rho}(x_1\cdots x_\rho)\qquad(\alpha+\alpha_1+\cdots+\alpha_\rho=i)
\]
تطابق الدوال المتناظرة الأولية لصفوف الكميات الآتية:
\[
\begin{array}{@{}l@{\quad}cccc@{}}
(5)&x_1&x_2&\cdots&x_\rho\\
&x_1^{(1)}&x_2^{(1)}&\cdots&x_\rho^{(1)}\\
&\vdots&\vdots&&\vdots\\
&x_1^{(i-1)}&x_2^{(i-1)}&\cdots&x_\rho^{(i-1)}.
\end{array}
\]
وعليه، فإن كل دالة من $\Kfield^*_{\rho\rho}$ تركيب نسبي لهذه الدوال الأولية، متناظر في صفوف الكميات (5)؛ وبالعكس، تنتمي إلى $\Kfield^*_{\rho\rho}$ كل دالة متناظرة في صفوف الكميات (5).

وبناءً على ذلك، إذا جرت
\[
f^*(x)=\frac{F^*(x)}{H^*(x)}
\]
على جميع دوال $\Kfield^*_{\rho\rho}$، قامت المنظومة الآتية من العلاقات اللامتناهية:
\[
\begin{array}{@{}l@{\quad}c@{}}
(6)&\displaystyle \frac{F^*(x)}{H^*(x)}=\frac{F^*(x^{(1)})}{H^*(x^{(1)})}=\cdots=\frac{F^*(x^{(i-1)})}{H^*(x^{(i-1)})};
\end{array}
\]
أو، باختصار:
\[
F^*(x^{(k)})H^*(x^{(l)})-F^*(x^{(l)})H^*(x^{(k)})=0\qquad(k,l=0,1,\ldots,(i-1));
\]
ولا ينعدم $H^*(x^{(k)})$ انعدامًا مطابقًا، وينطبق ذلك أيضًا على $F^*(x^{(k)})$ إذا كان $f^*\ne0$؛ فهذان التعبيران مترافقان مع $H^*(x)$ و$F^*(x)$، على الترتيب، ومتناظران صوريًا.

وبالعكس، فإن كل صف كميات $\xi$ يحقق العلاقات اللامتناهية
\[
\begin{array}{@{}l@{\quad}c@{}}
(7)&F^*(x)H^*(\xi)-H^*(x)F^*(\xi)=0,\qquad H^*(\xi)\ne0
\end{array}
\]
ينتمي إلى صفوف الكميات (5).

فإذا كان $G^*(x)$ يدل على المقام المشترك لمعاملات $\Phi^*(z,u)$، فإن دالة كثيرة الحدود الآتية، وهي كثيرة حدود في $\xi$ أيضًا، لا تنعدم انعدامًا مطابقًا بموجب الفرض (7):
\[
G^*(\xi)\cdot\Phi^*(z,u;\xi)
\]
ــ أي إن معاملات $\Phi^*(z,u;\xi)$ لا تصير غير معيّنة ــ ولها الجذر
\[
z=u_1\xi_1+u_2\xi_2+\cdots+u_\rho\xi_\rho.
\]
ومن ثم، بحكم (7)، يكون لـ$\Phi^*(z,u;x)$ أيضًا الجذر $z=u(\xi)$؛ أي إن $\xi$ ينتمي إلى صفوف الكميات (5). ويسري الأمر المناظر، بحسب (6)، إذا حقق $\xi$ منظومة معادلات بالنسبة إلى صف كميات مترافق:
\[
F^*(x^{(k)})H^*(\xi)-H^*(x^{(k)})F^*(\xi)=0;\qquad H^*(\xi)\ne0,
\]
أي إن صفوف الكميات (5) تُعرَّف بأنها الحلول $\xi$ لمنظومة المعادلات
\[
F^*(x)H^*(\xi)-H^*(x)F^*(\xi)=0;\qquad H^*(\xi)\ne0,
\]
حيث تجري $\frac{F^*(x)}{H^*(x)}$ على جميع دوال $\Kfield^*_{\rho\rho}$؛ ويجوز في ذلك استبدال صف $x$ بأي صف مترافق.

وهذا يعني هندسيًا، إذا فُسّر كل صف $x$ على أنه نقطة، أن حلول (7) تمثل في فضاء خطي ذي $\rho$ من الأبعاد $\infty^\rho$ مجموعات، في كل منها $i$ من النقاط، بحيث تتعيّن كل مجموعة بعينها تعيينًا فريدًا بأي نقطة من نقاطها. وتسمى مثل هذه المنظومة من مجموعات النقاط \glqq{}\emph{إنفولوشن في فضاء خطي}\grqq{}.\srcfn{*)}{تسمى الحلول المستبعدة من (7)، التي يتحقق عندها $F^*(\xi)=0$ و$H^*(\xi)=0$ ــ وكذلك الحلول المستبعدة من منظومة المعادلات المناظرة لقاعدة الإنفولوشن ــ النقاط الأساسية للإنفولوشن؛ ويجب أن تُعدّ معها أيضًا النقاط الناتجة من التجانس والواقعة \glqq{}في اللانهاية\grqq{}.}

وعلى النحو المناظر، يميّز الحقل $\Kfield_{n\rho}$ إنفولوشنًا في فضاء خطي مؤلفًا من منحنيات وسطوح وما إلى ذلك.

وبما أن الدرجة $i$ للحقل $\Omega(x_1\cdots x_\rho)$ بالنسبة إلى $\Kfield^*_{\rho\rho}$ تساوي بحسب ما تقدم عدد منظومات حلول (7) الخاضعة للشرط الإضافي $H^*(\xi)\ne0$ ــ وهي التي يمكن أن نسميها منظومات حلول $\Kfield^*_{\rho\rho}$ ــ فإن الاستدلال المستعمل لإثبات وجود قاعدة الإنفولوشن يقبل أيضًا الصياغة غير النسبية الآتية:

\glqq{}إذا كان $\Lsys^*$ حقلًا جزئيًا من $\Kfield^*_{\rho\rho}$، وكانت كل منظومة حلول للحقل $\Lsys^*$ منظومة حلول أيضًا لـ$\Kfield^*_{\rho\rho}$، فإن $\Lsys^*$ و$\Kfield^*_{\rho\rho}$ متطابقان.\grqq{}

ويسري الأمر المناظر، بمبدأ النقل، على الحقول الجزئية $\Lsys$ من $\Kfield_{n\rho}$.

\section*{\S\ 6. القاعدة الدنيا للحقول $\Kfield_{n\rho}$.}

تقدم هذه الفقرة عرضًا إجماليًا لمبرهنات معروفة، ولكن في صياغة موسعة أتاحها مبدأ النقل ووجود القاعدة النسبية.

\emph{التعريف الرابع: تسمى القاعدة النسبية للحقل $\Kfield_{n\rho}$ قاعدة دنيا إذا اقتصرت على العدد الأدنى من الدوال، أي $\rho$ دالة، وهي عندئذ يجب أن تكون مستقلة جبريًا.}

تنتج من مبرهنات لوروت وكاستلنووفو وإنريكيس\srcfn{*)}{\srcLTR{J. Lüroth: Beweis eines Satzes über rationale Kurven, Math. Ann. 9 (1875). --- G. Castelnuovo: Sulla razionalità delle involuzioni piane, Math. Ann. 44 (1893). --- F. Enriques: Sopra una involuzione non razionale dello spazio, Rend. Acc. Linc, Vol. 21, 21. Jan. 1912.}} الحقائق الآتية:

I. يملك كل حقل $\Kfield_{n1}$ قاعدة دنيا، هي دالة لوروت للحقل، المحددة تحديدًا وحيدًا حتى تحويل خطي كسري.

II. توجد للحقول $\Kfield_{n2}$ قاعدة دنيا دائمًا إذا سُمح بتمديد جبري لنطاق المعاملات $\Omega$، وبوجه عام لا توجد إلا إذا سُمح بذلك.

III. لا توجد عمومًا قاعدة دنيا للحقول $\Kfield_{n\rho}$ ذات $\rho\ge3$. ولم تُبذل بعد محاولة لتوصيف الحقول الخاصة $\Kfield_{n\rho}$ التي لها قاعدة دنيا.

ونشير بإيجاز إلى برهان مبرهنة لوروت للحقول $\Kfield^*_{11}$ الذي قدمه إ. شتاينتز\srcfn{**)}{في الموضع المذكور، \S\ 24.}:

لتكن $\Phi^*(z,u)$ صيغة الإنفولوشن للحقل $\Kfield^*_{11}$، ولتكن $\psi^*(x)$ معاملًا كيفيًا يظهر فيه المتغير $x$ فعلًا؛ أي معاملًا من معاملات $\Phi^*(z,u)$. ويتبين أن درجة $\Omega(x)$ بالنسبة إلى $\Omega(\psi^*(x))$ تساوي درجة $\Omega(x)$ بالنسبة إلى $\Kfield^*_{11}$، ومن ذلك يلزم تطابق الحقلين. ثم إذا كانت $\Omega(\chi^*(x))$ مساوية لـ$\Omega(\psi^*(x))$، فإن الاعتماد النسبي المتبادل بين $\chi^*(x)$ و$\psi^*(x)$ يستلزم أن ترتبط الدالتان بتحويل خطي كسري.

وبمقتضى مبدأ النقل، تقبل مبرهنة لوروت أيضًا الصياغة الآتية:

\emph{تتألف القاعدة الدنيا لكل حقل $\Kfield_{n1}$ من أي دالة من قاعدة الإنفولوشن؛ وترتبط كل دالتين من قاعدة الإنفولوشن للحقل $\Kfield_{n1}$ إحداهما بالأخرى بتحويل خطي كسري.}

يثبت ج. كاستلنووفو الحكم II بواسطة النظرية الهندسية للإنفولوشن للحقول $\Kfield^*_{22}$ المشتقة من قاعدة نسبية مؤلفة من ثلاث دوال. وبما أن مبدأ النقل يبقى صحيحًا أيضًا عند تمديد جبري منته لنطاق المعاملات $\Omega$، فإن وجود القاعدة النسبية يجعل البرهان بذلك تامًا لجميع الحقول $\Kfield_{n2}$.

وفي شأن III ينبغي ملاحظة أن إنريكيس أنشأ إنفولوشنًا غير نسبي في فضاء ثلاثي الأبعاد؛ أي حقلًا $\Kfield_{33}$ لا يملك قاعدة دنيا.

\section*{\S\ 7. منظومة كيفما كانت $\Ssys_{n\rho}$، وأسرة خطية $\Lsys_{n\rho}$، ومنطقة تكاملية $\Jdom_{n\rho}$ من الدوال النسبية. وجود القاعدة النسبية.}

ينتج الآن من وجود القاعدة النسبية للحقول $\Kfield_{n\rho}$ وجود القاعدة النسبية للمنظومات كيفما كانت من الدوال النسبية والصحيحة نسبيًا (أي كثيرات الحدود). ونرتب هذه المنظومات على النحو الذي تنشأ به إحداها من الأخرى تباعًا بالعمليات الأولية: الجمع والضرب والقسمة.

I. لتدل $\Ssys_{n\rho}$، كما في كل موضع، على منظومة كيفما كانت من الدوال النسبية (أو الصحيحة نسبيًا) في $n$ من المتغيرات غير المعيّنة، رتبتها الجبرية $\rho$. وتُعدّ المقادير من $\Omega$ من جملة عناصر المنظومة.

II. تسمى المنظومة $\Lsys_{n\rho}$ \emph{أسرة خطية} إذا استوفت الشرطين الآتيين:

1. تشتمل $\Lsys_{n\rho}$، مع $f(x)$، دائمًا على $c\cdot f(x)$ أيضًا، حيث $c$ مقدار كيفما كان من $\Omega$؛

2. وتشتمل $\Lsys_{n\rho}$، مع $f(x)$ و$g(x)$، دائمًا على $f(x)+g(x)$ أيضًا.

III. تسمى الأسرة الخطية $\Jdom_{n\rho}$ \emph{منطقة تكاملية} إذا استوفت، فوق ذلك، الشرط الآتي:

3. تشتمل $\Jdom_{n\rho}$، مع $f(x)$ و$g(x)$، دائمًا على $f(x)\cdot g(x)$ أيضًا.

IV. تسمى المنطقة التكاملية $\Kfield_{n\rho}$ \emph{حقلًا} إذا استوفت، فوق ذلك، الشرط الآتي:

4. تشتمل $\Kfield_{n\rho}$، مع $f(x)$ و$g(x)$ ــ حيث $g(x)\ne0$ ــ دائمًا على خارج القسمة $f(x):g(x)$ أيضًا.

الشروط من 1 إلى 4 هي الشروط المذكورة للحقل في \S\ 1.\srcfn{*)}{يجوز أيضًا أن تكون $f(x)$ و$g(x)$ من الدرجة صفر، أي مقدارين من $\Omega$.}

وسنشتق الآن هذه البنى بعضها من بعض تباعًا.

a) تُوسَّع منظومة كيفما كانت $\Ssys_{n\rho}$ إلى أسرة خطية $\Lsys_{n\rho}$ باشتراط ما يأتي:

تشتمل $\Lsys_{n\rho}$، إلى جانب $\Ssys_{n\rho}$، على جميع الدوال $h(x)$، ولا تشتمل إلا عليها، التي يمكن تمثيلها تركيبًا خطيًا، بمعاملات من $\Omega$، لعدد منتهٍ من دوال $\Ssys_{n\rho}$:
\[
h(x)=c_1f_1(x)+c_2f_2(x)+\cdots+c_kf_k(x),
\]
حيث تجري $f_1(x)\cdots f_k(x)$ على جميع دوال $\Ssys_{n\rho}$، وتجري $c_1\cdots c_k$ على جميع المقادير من $\Omega$.

ذلك أن إضافة تراكيب نسبية من دوال المنظومة لا تغير رتبتها الجبرية. وتستوفي $\Lsys_{n\rho}$ الشرطين 1 و2، فهي إذن أسرة خطية. ولما كانت كل أسرة خطية تشتمل على $\Ssys_{n\rho}$ لا بد، بمقتضى 1 و2، أن تشتمل أيضًا على $\Lsys_{n\rho}$، كانت $\Lsys_{n\rho}$ أصغر أسرة خطية تشتمل على $\Ssys_{n\rho}$.

b) وعلى المنوال نفسه، تنشأ من أسرة خطية $\Lsys_{n\rho}$ أصغر منطقة تكاملية $\Jdom_{n\rho}$ تشتمل عليها، باشتراط ما يأتي:

تشتمل $\Jdom_{n\rho}$، إلى جانب $\Lsys_{n\rho}$، على جميع الدوال $h(x)$، ولا تشتمل إلا عليها، التي يمكن تمثيلها \emph{تركيبًا نسبيًا صحيحًا} ــ أي تركيب كثير حدود ــ بمعاملات من $\Omega$، لعدد منتهٍ من دوال $\Lsys_{n\rho}$، $f_1(x)\cdots f_k(x)$، حيث تجري $f_1(x)\cdots f_k(x)$ على جميع دوال $\Lsys_{n\rho}$.

c) وأخيرًا، تنشأ من منطقة تكاملية $\Jdom_{n\rho}$ أصغر حقل $\Kfield_{n\rho}$ يشتمل عليها، باشتراط ما يأتي: يشتمل $\Kfield_{n\rho}$، إلى جانب $\Jdom_{n\rho}$، على جميع الدوال $h(x)$، ولا يشتمل إلا عليها، التي يمكن تمثيلها على صورة خارج قسمة دالتين من $\Jdom_{n\rho}$.

وإذا جُمعت التوسيعات الثلاثة في خطوة واحدة، نتج ما يأتي:

\emph{يمكن توسيع كل منظومة $\Ssys_{n\rho}$ إلى أصغر حقل $\Kfield_{n\rho}$ يشتمل عليها، باشتراط ما يأتي:}

\emph{يشتمل $\Kfield_{n\rho}$، إلى جانب $\Ssys_{n\rho}$، على جميع الدوال $h(x)$، ولا يشتمل إلا عليها، التي يمكن تمثيلها تركيبًا نسبيًا ــ بمعاملات من $\Omega$ ــ لعدد منتهٍ من دوال $\Ssys_{n\rho}$، $f_1(x)\cdots f_k(x)$، حيث تجري $f_1(x)\cdots f_k(x)$ على جميع دوال $\Ssys_{n\rho}$.}

ومن الحقيقة الأخيرة ينتج مباشرة وجود القاعدة النسبية للبنى كيفما كانت:

ليكن $\Kfield_{n\rho}$ أصغر حقل يشتمل على $\Ssys_{n\rho}$، ولتكن قاعدة نسبية لـ$\Kfield_{n\rho}$ معطاة بالدوال
\[
\begin{array}{@{}l@{\quad}c@{}}
(1)&h_1(x),\ h_2(x),\ldots,h_\sigma(x).
\end{array}
\]
وبحسب تعريف $\Kfield_{n\rho}$ توجد العلاقات الآتية:
\[
\begin{array}{@{}l@{\quad}c@{}}
(2)&h_1(x)=r_1(g_1(x)\cdots g_\lambda(x));\quad\ldots;\quad h_\sigma(x)=r_\sigma(g_\mu(x)\cdots g_\nu(x)),
\end{array}
\]
حيث $r_1\cdots r_\sigma$ دوال نسبية في وسائطها، و
\[
\begin{array}{@{}l@{\quad}c@{}}
(3)&g_1(x)\cdots g_\lambda(x),\ldots,g_\mu(x)\cdots g_\nu(x)
\end{array}
\]
تنتمي جميعها إلى $\Ssys_{n\rho}$. وبمقتضى العلاقات (2)، تمثل (3)، مثل (1)، قاعدة نسبية لـ$\Kfield_{n\rho}$؛ وبسبب انتماء الدوال (3) إلى $\Ssys_{n\rho}$، تكون كذلك قاعدة نسبية لـ$\Ssys_{n\rho}$. ومن ثم ثبتت

\emph{المبرهنة الرابعة: لكل منظومة كيفما كانت $\Ssys_{n\rho}$ من الدوال النسبية (أو الصحيحة نسبيًا) قاعدة نسبية مؤلفة من عدد منتهٍ من دوال المنظومة، ورتبتها الجبرية $\rho$.}

ولتكن الآن $\Lsys_{n\rho}$، على وجه الخصوص، أسرة خطية، ولتكن
\[
\begin{array}{@{}l@{\quad}c@{}}
(4)&g_1(x)\cdots g_\rho(x);\ g_{\rho+1}(x)\cdots g_\nu(x)
\end{array}
\]
قاعدة نسبية لـ$\Lsys_{n\rho}$. ولنفترض، مثلًا، أن
\[
g_1(x)\cdots g_\rho(x)
\]
مستقلة جبريًا. فإذا وضعنا
\[
g(x)=c_1g_{\rho+1}(x)+c_2g_{\rho+2}(x)+\cdots+c_{\nu-\rho}g_\nu(x),
\]
أمكن اختيار $c_i$ بوصفها أعدادًا من $\Omega$ على نحو يجعل
\[
\begin{array}{@{}l@{\quad}c@{}}
(5)&g_1(x)\cdots g_\rho(x),\ g(x)
\end{array}
\]
قاعدة نسبية لأصغر حقل مشتمل، وهو $\Kfield_{n\rho}$. ولكن لما كانت $\Lsys_{n\rho}$ أسرة خطية، انتمت $g(x)$ إلى $\Lsys_{n\rho}$، ومثلت (5) قاعدة نسبية لـ$\Lsys_{n\rho}$؛ أي ثبتت

\emph{المبرهنة الخامسة: يمكن، على وجه الخصوص، اختيار القاعدة النسبية لكل أسرة خطية $\Lsys_{n\rho}$ من الدوال النسبية (أو الصحيحة نسبيًا)، ومن ثم القاعدة النسبية لكل منطقة تكاملية $\Jdom_{n\rho}$ ولكل حقل $\Kfield_{n\rho}$، بحيث تشتمل على الأكثر على $\rho+1$ دالة (وعلى الأقل على $\rho$ دالة).}

وهكذا، فإن الحد العددي لدوال القاعدة النسبية الذي وجدناه للحقل في \S\ 4 يرتكز على خاصية الحقل بوصفه أسرة خطية؛ أما الاقتصار على $\rho$ دالة عند وجود القاعدة الدنيا فيستعمل جميع فروض الحقل.

ونلاحظ أيضًا أنه، وفقًا لـ\S\ 3 (4)، تُحفظ عند التصوير جميع الخواص المميزة للمنظومات ولأصغر البنى المشتملة عليها.


\section*{\S\ 8. قاعدة الإنفولوشن للمناطق التكاملية $\Jdom_{n\rho}$ المؤلفة من كثيرات الحدود.}

تتناول الفقرات التالية منظومات $\Ssys_{n\rho}$ عناصرها دوال صحيحة نسبيًا (كثيرات حدود)؛ وتتناول أولًا المناطق التكاملية $\Jdom_{n\rho}$ المؤلفة من كثيرات الحدود. وينبغي بادئ ذي بدء تثبيت مفاهيم القابلية للقسمة لهذه المناطق $\Jdom_{n\rho}$.

لتكن $F(x),G(x)$ كثيرتي حدود من $\Jdom_{n\rho}$، ولتكن $L(x)$ قاسمًا مشتركًا لهما، بحيث يكون:
\[
F(x)=L(x)\cdot F_1(x);\qquad G(x)=L(x)\cdot G_1(x).
\]
وتسمى $L(x)$ حينئذٍ، وحينئذٍ فقط، \emph{قاسمًا مشتركًا في $\Jdom_{n\rho}$} إذا كانت كثيرات الحدود الثلاث $L(x),F_1(x),G_1(x)$ تنتمي إلى $\Jdom_{n\rho}$.

تسمى كثيرتا حدود من $\Jdom_{n\rho}$ \emph{أوليين فيما بينهما في $\Jdom_{n\rho}$} إذا لم يكن لهما في $\Jdom_{n\rho}$ قاسم مشترك غير وحدي.

وليكن $\Kfield_{n\rho}$ أصغر حقل يشتمل على $\Jdom_{n\rho}$. فعندئذ ينتج من التعريف، لكل دالة من $\Kfield_{n\rho}$، تمثيل من الصورة:
\[
f(x)=F_1(x):F_2(x),
\]
حيث تنتمي $F_1(x),F_2(x)$ إلى $\Jdom_{n\rho}$ وهما أوليتان فيما بينهما في $\Jdom_{n\rho}$. وعلى نحو مماثل، يمكن توحيد مقامات عدة دوال من $\Kfield_{n\rho}$ على مقام مشترك في $\Jdom_{n\rho}$:
\[
f_1(x)=\frac{F_1(x)}{F(x)}\cdots f_k(x)=\frac{F_k(x)}{F(x)}.
\]
وتسمى $F(x)$ حينئذٍ، وحينئذٍ فقط، \emph{مقامًا مشتركًا أصغر في $\Jdom_{n\rho}$} إذا كانت كثيرات الحدود من $\Jdom_{n\rho}$، وهي $F(x),F_1(x),\ldots,F_k(x)$، أوليةً فيما بينها في $\Jdom_{n\rho}$.

وقد يكون مثل هذا التمثيل ممكنًا بطرائق مختلفة، لأن قوانين التحليل الوحيد إلى دوال غير قابلة للاختزال لا تكون نافذة بوجه عام في $\Jdom_{n\rho}$.

ونبحث الآن دلالة قاعدة الإنفولوشن للحقل $\Kfield_{n\rho}$ بالنسبة إلى $\Jdom_{n\rho}$، بأن ننقل مفهومي صيغة الإنفولوشن وقاعدة الإنفولوشن إلى المناطق التكاملية. ولتكن
\[
\begin{gathered}
\Phi(z,u)=z^i+\sum{}' g_{\alpha\alpha_1\ldots\alpha_\rho}(x)z^\alpha u_1^{\alpha_1}\cdots u_\rho^{\alpha_\rho}\\
(\alpha+\alpha_1+\cdots+\alpha_\rho=i)
\end{gathered}
\]
صيغة إنفولوشن للحقل $\Kfield_{n\rho}$، ولتكن $G(x)$ مقامًا مشتركًا أصغر في $\Jdom_{n\rho}$ لمعاملات $\Phi(z,u)$. وبالضرب في $G(x)$ تتحول $\Phi(z,u)$ إلى صيغة $\Psi(z,u)$ تكون معاملاتها كثيرات حدود من $\Jdom_{n\rho}$ وأوليةً فيما بينها في $\Jdom_{n\rho}$:
\[
\begin{gathered}
G(x)\cdot\Phi(z,u)=\Psi(z,u)=G(x)z^i+\sum{}'G_{\alpha\alpha_1\ldots\alpha_\rho}(x)z^\alpha u_1^{\alpha_1}\cdots u_\rho^{\alpha_\rho}\\
(\alpha+\alpha_1+\cdots+\alpha_\rho=i).
\end{gathered}
\]

\emph{كل صيغة $\Psi(z,u)$ تنشأ من صيغة إنفولوشن $\Phi(z,u)$ للحقل $\Kfield_{n\rho}$ بضربها في كثيرة حدود من $\Jdom_{n\rho}$، على نحو تكون معه معاملات $\Psi(z,u)$ كثيرات حدود من $\Jdom_{n\rho}$ وأوليةً فيما بينها في $\Jdom_{n\rho}$، تسمى صيغة إنفولوشن لـ$\Jdom_{n\rho}$؛ وتسمى معاملات $\Psi(z,u)$ قاعدة إنفولوشن مناظرة لها.}

ويتضح أولًا من معنى قاعدة الإنفولوشن للحقل $\Kfield_{n\rho}$ أن كل قاعدة إنفولوشن لـ$\Jdom_{n\rho}$ هي في الوقت نفسه قاعدة نسبية. ولأجل مواصلة البحث، ننظر في منطقة تصوير $\Jdom^*_{\rho\rho}$ يكون لأصغر حقل يشتمل عليها، $\Kfield^*_{\rho\rho}$، معادلة غير قابلة للاختزال ذات الجذر
\[
z=u_1x_1+\cdots+u_\rho x_\rho
\]
معطاة بالمعادلة $\Phi^*(z,u)=0$.

بحسب \S\ 5، تمثل معاملات
\[
\begin{gathered}
\Phi^*(z,u)=z^i+\sum{}'g^*_{\alpha\alpha_1\ldots\alpha_\rho}(x)z^\alpha u_1^{\alpha_1}\cdots u_\rho^{\alpha_\rho}\\
(\alpha+\alpha_1+\cdots+\alpha_\rho=i)
\end{gathered}
\]
الدوال المتناظرة الأولية لصفوف المقادير المترافقة، بالنسبة إلى $\Kfield^*_{\rho\rho}$،
\[
x,\ x^{(1)}\cdots x^{(i-1)}
\]
ولْتكن $F^*(x)$ كثيرة حدود كيفما كانت من $\Kfield^*_{\rho\rho}$؛ فعندئذ ينتج من
\[
F^*(x)=\frac1i\left(F^*(x)+F^*(x^{(1)})+\cdots+F^*(x^{(i-1)})\right),
\]
أن $F^*(x)$ دالة متناظرة صحيحة، أي كثيرة حدود، في صفوف المقادير هذه؛ ومن ثم يمكن التعبير عنها تركيبًا نسبيًا صحيحًا ــ أي تركيب كثير حدود ــ للدوال المتناظرة الأولية، أي للدوال
\[
g^*_{\alpha\alpha_1\ldots\alpha_\rho}(x)
\]
فإذا كانت الآن
\[
\begin{gathered}
G^*(x)\cdot\Phi^*(z,u)=\Psi^*(z,u)=G^*(x)z^i+
\sum{}'G^*_{\alpha\alpha_1\ldots\alpha_\rho}(x)z^\alpha u_1^{\alpha_1}\cdots u_\rho^{\alpha_\rho}\\
(\alpha+\alpha_1+\cdots+\alpha_\rho=i)
\end{gathered}
\]
صيغة إنفولوشن لـ$\Jdom^*_{\rho\rho}$، فإن كل كثيرة حدود من $\Kfield^*_{\rho\rho}$، وبخاصة كل كثيرة حدود من $\Jdom^*_{\rho\rho}$، تكون بناءً على ذلك تركيبًا نسبيًا صحيحًا لخوارج القسمة
\[
G^*_{\alpha\alpha_1\ldots\alpha_\rho}(x):G^*(x).
\]

ومع مراعاة مبدأ النقل، ينتج بذلك ما يأتي:

\emph{المبرهنة السادسة: لكل منطقة تكاملية من كثيرات الحدود $\Jdom_{n\rho}$ توجد قاعدة نسبية متميزة واحدة على الأقل، هي قاعدة الإنفولوشن، بحيث لا يظهر في مقام التمثيل النسبي لجميع كثيرات الحدود من $\Jdom_{n\rho}$ بدوال قاعدة الإنفولوشن إلا قوة دالة واحدة بعينها (هي المعامل الأعلى لصيغة الإنفولوشن المناظرة).}

\section*{\S\ 9. المناطق الصحيحة نسبيًا $\Gdom_{n\rho}$ المؤلفة من كثيرات الحدود.}

نبحث فيما يلي مناطق تكاملية خاصة مؤلفة من كثيرات الحدود، تشكل مرحلة وسطى بين المنطقة التكاملية والحقل.

تسمى منطقة تكاملية مؤلفة من كثيرات الحدود $\Gdom_{n\rho}$ «منطقة صحيحة نسبيًا» إذا تحقق الشرط الآتي:

\emph{4أ) تحتوي $\Gdom_{n\rho}$، إلى جانب $F(x)$ و$G(x)$ --- حيث $G(x)\ne0$ ---، دائمًا على خارج القسمة $F(x):G(x)$ إذا وفقط إذا كان هذا الخارج صحيحًا في المتغيرات غير المعينة، أي كثيرة حدود.}

ينبغي أولًا إثبات تطابق المناطق الصحيحة نسبيًا مع مجموع كثيرات الحدود في حقل $\Kfield_{n\sigma}$ ($\sigma\ge \rho$) من الدوال النسبية.

ليكن $\Kfield_{n\rho}$ أصغر حقل يحتوي $\Gdom_{n\rho}$. لكل دالة $f(x)$ من $\Kfield_{n\rho}$ تمثيل على صورة خارج قسمة كثيرتي حدود من $\Gdom_{n\rho}$:
\[
f(x)=F_1(x):F_2(x),
\]
حيث يُسمح أيضًا بكثيرات الحدود من الدرجة صفر. ومتى كان $f(x)$ كثيرة حدود انتمى، بمقتضى الشرط 4أ، إلى $\Gdom_{n\rho}$؛ ولما كانت $\Gdom_{n\rho}$ لا يمكن أن تحتوي أي كثيرات حدود أخرى، فإن \emph{كل منطقة صحيحة نسبيًا $\Gdom_{n\rho}$ تطابق مجموع كثيرات الحدود في أصغر حقل يحتويها.}

ولْيكن، بالعكس، $\Kfield_{n\sigma}$ حقلًا اعتباطيًا من الدوال النسبية (أي ليس أصغر حقل حاوٍ لمنظومة معطاة)، ولتكن $\mathfrak J$ مجموع كثيرات الحدود المحتواة في $\Kfield_{n\sigma}$. تحقق $\mathfrak J$ الشروط 1 و2 و3 للمنطقة التكاملية والشرط 4أ، فهي إذن منطقة صحيحة نسبيًا؛ أي إن \emph{مجموع كثيرات الحدود المحتواة في حقل $\Kfield_{n\sigma}$ يؤلف منطقة صحيحة نسبيًا $\Gdom_{n\rho}$ ($\rho\le\sigma$).}\srcfn{*)}{يجوز أيضًا أن يكون $\rho=0$؛ أي إن مجموع كثيرات الحدود من $\Kfield_{n\sigma}$ قد يتألف من نطاق المعاملات $\Omega$ وحده.}

ونبين كذلك تطابق المناطق الصحيحة نسبيًا مع «المناطق التكاملية المؤلفة من دوال صحيحة نسبيًا» التي أدخلها هيلبرت.\srcfn{**)}{د. هيلبرت: مسائل رياضية. محاضرة باريس 1900. المسألة 14 (أخبار غوتينغن 1900).}

لتكن
\[
\begin{array}{@{}l@{\quad}c@{}}
(1)& G_1(x)\cdots G_k(x)
\end{array}
\]
قاعدة نسبية لـ$\Gdom_{n\rho}$؛ فبمقتضى الشروط 1 و2 و3 و4أ ينتمي إلى $\Gdom_{n\rho}$ كل تركيب نسبي لكثيرات الحدود (1) يكون صحيحًا في المتغيرات غير المعينة، أي يصير كثيرة حدود. وبالعكس، لما كان يمكن، بحسب مفهوم القاعدة النسبية، التعبير نسبيًا عن كل كثيرة حدود من $\Gdom_{n\rho}$ بواسطة كثيرات الحدود (1)، فإن $\Gdom_{n\rho}$ \emph{تطابق مجموع تلك التركيبات النسبية لكثيرات الحدود (1) التي تكون صحيحة في المتغيرات غير المعينة، أي كثيرات حدود.} وهكذا نحصل على تعريف هيلبرت المنطلق من القاعدة النسبية الخاصة، في حين لا يستعمل تعريفنا إلا الخصائص المجردة للمنطقة.

في المناطق الصحيحة نسبيًا يتبسط تعريف القاسم المشترك المعطى في \S\ 8 للمناطق التكاملية العامة:

لتكن $F(x),G(x)$ كثيرتي حدود من $\Gdom_{n\rho}$، وليكن $L(x)$ قاسمًا مشتركًا لهما. عندئذ يسمى $L(x)$ \emph{قاسمًا مشتركًا في $\Gdom_{n\rho}$} إذا وفقط إذا كان $L(x)$ ينتمي إلى $\Gdom_{n\rho}$.

ذلك أن انتماء خارجَي القسمة $F(x):L(x)$ و$G(x):L(x)$ ينتج عندئذ من الشرط 4أ.

ومن الخواص الجوهرية للمناطق الصحيحة نسبيًا كذلك أن مفهوم «الصحيح جبريًا بالنسبة إلى $\Gdom_{n\rho}$» يمكن تعريفه بواسطة معادلة اعتباطية على النحو نفسه تمامًا المعروف بالنسبة إلى حقول الأعداد الجبرية أو إلى الحقل $\Omega(x_1\cdots x_n)$.

\emph{التعريف V: تسمى كمية $\xi$ صحيحة جبريًا بالنسبة إلى $\Gdom_{n\rho}$ --- أو صحيحة جبريًا نسبيًا أيضًا --- إذا حققت معادلة ما معاملها الأعلى هو الوحدة، ومعاملاتها الباقية كثيرات حدود من $\Gdom_{n\rho}$.}

ولتكن الكمية الصحيحة جبريًا نسبيًا $\xi$ محققة، مثلًا، للمعادلة
\[
L(z)=z^\lambda+G_1(x)z^{\lambda-1}+\cdots+G_\lambda(x)=0
\]
ولتكن $L(z)$ قابلة للقسمة على الدالة الصحيحة غير القابلة للاختزال بالنسبة إلى أصغر حقل حاوٍ $\Kfield_{n\rho}$
\[
M(z)=z^\mu+H_1(x)z^{\mu-1}+\cdots+H_\mu(x)
\]
لكن قابلية القسمة هذه تستلزم، بحسب التعميم المعروف لمبرهنة غاوس\srcfn{*)}{انظر مثلًا ي. كونيغ: مدخل إلى النظرية العامة للكميات الجبرية (لايبزيغ، ب. غ. تويبنر، 1903)، الفصل التاسع، \S\ 2، أو فيبر (الطبعة الصغيرة)، \S\ 20.}، أن الدوال النسبية $H_1(x)\cdots H_\mu(x)$ كثيرات حدود من $\Gdom_{n\rho}$، ومن ثم تنتمي إلى $\Gdom_{n\rho}$. وبذلك يُستخلص تعريف الكمية الصحيحة جبريًا نسبيًا $\xi$ من المعادلة غير القابلة للاختزال. وبوجه خاص، تكون كل كثيرة حدود $F(x)$ من $\Gdom_{n\rho}$ صحيحة جبريًا أيضًا بالنسبة إلى $\Gdom_{n\rho}$، لأنها تحقق المعادلة $z-F(x)=0$.

ولنذكر أيضًا أنه، كما يُعرّف أصغر حقل حاوٍ، يُعرّف كذلك لأية منطقة تكاملية $\Jdom_{n\rho}$ مؤلفة من كثيرات الحدود أصغر منطقة صحيحة نسبيًا حاوية $\Gdom_{n\rho}$ بالشرط الآتي:

تحتوي $\Gdom_{n\rho}$، إلى جانب $\Jdom_{n\rho}$، جميع كثيرات الحدود $H(x)$، وهذه وحدها، التي يمكن تمثيلها على صورة خارج قسمة كثيرتي حدود من $\Jdom_{n\rho}$.

وينتج كذلك من هذا التعريف أن كل صيغة إنفولوشن وقاعدة إنفولوشن لـ$\Jdom_{n\rho}$ تولد صيغة وقاعدة مناظرتين لـ$\Gdom_{n\rho}$، وذلك بفصل قاسم مشترك واحد على الأكثر في $\Gdom_{n\rho}$ من جميع كثيرات حدود القاعدة.

ذلك أن $\Jdom_{n\rho}$ و$\Gdom_{n\rho}$ لهما أصغر حقل حاوٍ واحد هو $\Kfield_{n\rho}$؛ وكل صيغة إنفولوشن لـ$\Jdom_{n\rho}$ تنشأ من صيغة إنفولوشن لـ$\Kfield_{n\rho}$ بضربها في كثيرة حدود من $\Jdom_{n\rho}$، فتكون معاملاتها من $\Gdom_{n\rho}$، غير أنها قد تقبل قاسمًا مشتركًا في $\Gdom_{n\rho}$.\srcfn{*)}{فعلى سبيل المثال، للمنطقة التكاملية المؤلفة من جميع التركيبات النسبية الصحيحة لـ$x_1^2$ و$x_1x_2$، وهي $\Jdom_{22}$، تكون صيغة الإنفولوشن:
\[
\Psi(z,u)=x_1^2z^2-x_1^4u_1^2-2x_1^3x_2u_1u_2-x_1^2x_2^2u_2^2.
\]
وبما أن $x_1^2$ ينتمي إلى $\Jdom_{22}$، ومن باب أولى إلى أصغر منطقة صحيحة نسبيًا حاوية $\Gdom_{22}$، ومن ثم هو قاسم مشترك في $\Gdom_{22}$، فإن صيغة الإنفولوشن لـ$\Gdom_{22}$ تصبح:
\[
\Psi(z,u)=z^2-x_1^2u_1^2-2x_1x_2u_1u_2-x_2^2u_2^2.
\]}

وأخيرًا نلاحظ أن \emph{منطقة تصوير} منطقة صحيحة نسبيًا $\Gdom_{n\rho}$، التي تكون بدورها بحسب \S\ 7 منطقة تكاملية $\Jdom^*_{\rho\rho}$ مؤلفة من كثيرات الحدود، لا يلزم أن تكون منطقة صحيحة نسبيًا. فإذا كان $F(x)$ و$G(x)$ كثيرتي حدود من $\Gdom_{n\rho}$ وكان خارج قسمتهما غير كثيرة حدود، وبالتالي لا ينتمي إلى $\Gdom_{n\rho}$، فإن خارج القسمة $F^*(x):G^*(x)$ لا ينتمي هو أيضًا إلى $\Jdom^*_{\rho\rho}$. غير أن هذا الخارج الناشئ من تخصيص بعض المتغيرات غير المعينة قد يكون بالفعل كثيرة حدود؛ وعندئذ لا يتحقق الشرط 4أ في $\Jdom^*_{\rho\rho}$.\srcfn{**)}{لتكن $\Gdom_{32}$، مثلًا، مؤلفة من جميع كثيرات الحدود التي هي تراكيب نسبية لـ$F=x_1^2x_2$ و$G=x_1^2(1+x_3)$. يعطي التخصيص المسموح $x_3=0$ العلاقة $F^*:G^*=x_2$، في حين أن $F:G$ ليست كثيرة حدود. ومن ثم فإن $\Jdom^*_{22}$ ليست منطقة صحيحة نسبيًا.}

\section*{\S\ 10. المناطق الصحيحة نسبيًا من النوع الأول وقاعدتها التكاملية.}

إن مفهوم «الصحيح جبريًا نسبيًا» الذي أُدخل في \S\ 9 (التعريف الخامس) يتيح تعيين فئة من المناطق الصحيحة نسبيًا تكون مناطق تكاملية منتهية التوليد؛ وبذلك يتيح الإجابة بالإيجاب، لهذه الفئة على الأقل، عن سؤال طرحه د. هيلبرت (المسائل الرياضية، المسألة 14). ونورد

\emph{التعريف السادس: يسمى عدد منتهٍ من كثيرات الحدود في $\Gdom_{n\rho}$ قاعدةً تكاملية إذا أمكن تمثيل كل كثيرة حدود في $\Gdom_{n\rho}$ بتركيب نسبي صحيح من ذلك العدد المنتهي، بمعاملات من $\Omega$. وتسمى المنطقة التكاملية المؤلفة من كثيرات الحدود، إذا كانت لها قاعدة تكاملية، منطقةً تكاملية منتهية التوليد.}

ونسمي فئة المناطق الصحيحة نسبيًا التي نريد النظر فيها، والتي سيتبين أن قاعدة الإنفولوشن فيها قاعدة تكاملية، مناطقَ صحيحة نسبيًا من النوع الأول، وفق

\emph{التعريف السابع: تسمى المنطقة الصحيحة نسبيًا $\Gdom_{n\rho}$ من النوع الأول إذا كان لها، بوصفها منطقة تصوير، منطقة صحيحة نسبيًا $\Gdom^*_{\rho\rho}$ بحيث تعتمد كل كثيرة حدود كيفما كانت في $x_1\cdots x_\rho$ (بمعاملات من $\Omega$) اعتمادًا صحيحًا جبريًا على $\Gdom^*_{\rho\rho}$.}

فبمقتضى هذا التعريف تعتمد المتغيرات غير المعينة $x_1\cdots x_\rho$، ومن ثم الشكل الخطي $u(x)=u_1x_1+\cdots+u_\rho x_\rho$ أيضًا، اعتمادًا صحيحًا جبريًا على $\Gdom^*_{\rho\rho}$. ولذلك فإن الدالة الصحيحة غير القابلة للاختزال ذات الجذر $z=u(x)$ --- وهي صيغة الإنفولوشن لأصغر حقل حاوٍ $\Kfield^*_{\rho\rho}$ --- معاملها الأعلى الوحدة ومعاملاتها الأخرى كثيرات حدود من $\Gdom^*_{\rho\rho}$؛ ومن ثم فهي في الوقت نفسه صيغة إنفولوشن لـ$\Gdom^*_{\rho\rho}$، ولا يمكن أن توجد في $\Gdom^*_{\rho\rho}$ صيغة إنفولوشن أخرى. وبحسب مبدأ النقل، تملك $\Gdom_{n\rho}$ أيضًا صيغة إنفولوشن معاملها الأعلى الوحدة؛ ولما كان، بحسب المبرهنة السادسة، لا يدخل في مقام التمثيل النسبي لجميع كثيرات الحدود في $\Gdom_{n\rho}$ بواسطة قاعدة الإنفولوشن المقابلة سوى هذا المعامل الأعلى، تصير قاعدة الإنفولوشن قاعدة تكاملية؛ أي:

\emph{المبرهنة السابعة: كل منطقة صحيحة نسبيًا من النوع الأول منطقة تكاملية منتهية التوليد؛ وقاعدتها التكاملية هي قاعدة الإنفولوشن المستخرجة من منطقة التصوير التي تميزها. وعلى وجه الخصوص، تُعطى القاعدة التكاملية للمناطق الصحيحة نسبيًا من النوع الأول $\Gdom_{\rho\rho}$ بقاعدة الإنفولوشن المحددة تحديدًا وحيدًا لـ$\Gdom_{\rho\rho}$.}

ونورد أيضًا معيارًا للمناطق الصحيحة نسبيًا من النوع الأول. لتكن $\Gdom_{n\rho}$ من النوع الأول، ولتكن قاعدة تكاملية لـ$\Gdom_{n\rho}$ معطاة بـ
\[
\begin{array}{@{}l@{\quad}c@{}}
(1)& F_1(x),F_2(x),\ldots,F_k(x).
\end{array}
\]
عندئذ توجد، بموجب مبرهنة مساعدة لهيلبرت\srcfn{*)}{د. هيلبرت: \srcLTR{Über die vollen Invariantensysteme}، \srcLTR{Math. Ann. 42 (1893)}، \S\ 1. (المبرهنة المساعدة التي لم تُذكر هناك إلا لكثيرات الحدود المتجانسة تصح كذلك لكثيرات الحدود غير المتجانسة.)}، $\rho$ كثيرة حدود مستقلة جبريًا من $\Gdom_{n\rho}$:
\[
\begin{array}{@{}l@{\quad}c@{}}
(2)& G_1(x),G_2(x),\ldots,G_\rho(x)
\end{array}
\]
بحيث تعتمد كثيرات الحدود (1)، ومن ثم كل كثيرة حدود في $\Gdom_{n\rho}$، اعتمادًا صحيحًا جبريًا على كثيرات الحدود (2). وينطبق الأمر نفسه على كل كثيرة حدود في منطقة التصوير $\Gdom^*_{\rho\rho}$ بالنسبة إلى كثيرات الحدود المقابلة:
\[
\begin{array}{@{}l@{\quad}c@{}}
(3)& G_1^*(x),G_2^*(x),\ldots,G_\rho^*(x).
\end{array}
\]


وبحسب التعريف السابع، تعتمد إذن كل كثيرة حدود كيفما كانت في $x_1\rcdots x_\rho$ اعتمادًا صحيحًا جبريًا على كثيرات الحدود (3).

ولنفترض الآن، عكسًا، أن منطقة تصوير $\Jdom^*_{\rho\rho}$ لمنطقة صحيحة نسبيًا $\Gdom_{n\rho}$ كيفما كانت تحتوي $\rho$ كثيرة الحدود (3)، بحيث تعتمد عليها اعتمادًا صحيحًا جبريًا كل كثيرة حدود في $x_1\rcdots x_\rho$. عندئذ نبين أن $\Jdom^*_{\rho\rho}$ تنتج من ذلك بوصفها منطقة صحيحة نسبيًا $\Gdom^*_{\rho\rho}$، وأن $\Gdom^*_{\rho\rho}$، ومن ثم $\Gdom_{n\rho}$، من النوع الأول.

وبالفعل، لتكن $F^*(x)$ كثيرة حدود في أصغر حقل حاوٍ لـ$\Jdom^*_{\rho\rho}$، وهو $\Kfield^*_{\rho\rho}$؛ فهي، بحسب الفرض، تعتمد اعتمادًا صحيحًا جبريًا على كثيرات الحدود (3). ولما كانت الدالة المقابلة في $\Kfield_{n\rho}$ تعتمد عندئذ اعتمادًا صحيحًا جبريًا على $\rho$ كثيرة الحدود في $\Gdom_{n\rho}$، فهي كثيرة حدود $F(x)$، ومن ثم تنتمي إلى $\Gdom_{n\rho}$؛ ولذلك تنتمي $F^*(x)$ إلى $\Jdom^*_{\rho\rho}$. وهكذا تشتمل منطقة التصوير $\Jdom^*_{\rho\rho}$ على جميع كثيرات الحدود في $\Kfield^*_{\rho\rho}$، فهي منطقة صحيحة نسبيًا $\Gdom^*_{\rho\rho}$. أما $\Gdom^*_{\rho\rho}$، ومن ثم $\Gdom_{n\rho}$ أيضًا، فينتج مباشرة من الفرض، بحسب التعريفين الخامس والسابع، أنهما من النوع الأول؛ أي:

\emph{تتميز المناطق الصحيحة نسبيًا من النوع الأول $\Gdom_{n\rho}$ بأن في منطقة تصوير $\Jdom^*_{\rho\rho}$ توجد $\rho$ كثيرة حدود تعتمد عليها اعتمادًا صحيحًا جبريًا كل كثيرة حدود كيفما كانت في $x_1\rcdots x_\rho$. وتنتج منطقة التصوير من هذا الفرض بوصفها منطقة صحيحة نسبيًا.}\srcfn{*)}{مثلًا، في المنطقة $\Gdom_{32}$ المذكورة في الملاحظة الختامية لـ\S\ 9، يعطي التخصص الجائز $x_1=1$ كثيرتي الحدود $F^*=x_2$ و$G^*=1+x_3$، اللتين تعتمد عليهما اعتمادًا صحيحًا جبريًا كل كثيرة حدود كيفما كانت في $x_2,x_3$. ومن ثم فإن $\Gdom_{32}$ من النوع الأول، وتكون $F(x)$ و$G(x)$ قاعدة تكاملية لها.}

\section*{\S\ 11. مبرهنة مساعدة في الاعتماد الصحيح جبريًا.}

نورد الآن معيارًا يضمن لكل كثيرة حدود كيفما كانت في $x_1\rcdots x_\rho$، بالنسبة إلى $\rho$ كثيرة الحدود الآتية في $x_1\rcdots x_\rho$، أن
\[
\begin{array}{@{}l@{\quad}c@{}}
(1)& G_1^*(x)\rcdots G_\rho^*(x)
\end{array}
\]
تعتمد عليها اعتمادًا صحيحًا جبريًا.

إذا كانت كثيرات الحدود (1)، على وجه الخصوص، صيغًا متجانسة، فإن جميع الصيغ المعتمدة عليها اعتمادًا صحيحًا جبريًا تنعدم معًا عند كل منظومة قيم معيّنة تُعدم (1)؛ ومن ثم تنعدم، بوجه خاص، $x_1\rcdots x_\rho$ أيضًا.

فلا يمكن إذن للصيغ (1) أن تنعدم عند أي منظومة قيم غير
\[
x_1=0,\ldots,x_\rho=0
\]
؛ ولا بد أن تكون محصلتها غير مساوية للصفر.

وعكسًا، هذا الشرط كافٍ أيضًا، ويقبل تعميمًا إلى كثيرات الحدود غير المتجانسة وفق المبرهنة المساعدة الآتية:

\emph{مبرهنة مساعدة: شرط كافٍ لأن تعتمد كل كثيرة حدود كيفما كانت في $x_1\rcdots x_\rho$ اعتمادًا صحيحًا جبريًا على $\rho$ كثيرة حدود في $x_1\rcdots x_\rho$
\[
\begin{array}{@{}l@{\quad}c@{}}
(1)& G_1^*(x)\rcdots G_\rho^*(x)
\end{array}
\]
هو أن تكون محصلة الأجزاء المتجانسة ذات أعلى درجة، المأخوذة بالنسبة إلى $x_1\rcdots x_\rho$، غير منعدمة}
\[
H_1^*(x)\rcdots H_\rho^*(x)
\]
\emph{لكثيرات الحدود (1):}
\[
\begin{array}{@{}l@{\quad}c@{}}
(2)& \displaystyle R_0=\left[\begin{array}{c} H_1^*\rcdots H_\rho^*\\ x_1\rcdots x_\rho\end{array}\right]\ne0.
\end{array}
\]
\emph{أما للصيغ المتجانسة (1) فالشرط (2) لازم.}\srcfn{*)}{في نظرية المحصلة، قارن ف. مرتنس: \srcLTR{Zur Theorie der Elimination (I. u. II. Teil)}، \srcLTR{Sitzungsber. d. k. Ak. d. Wiss. Wien, Math.-naturw. Kl. Abt. IIa, Bd. 108 (1899)}. (أُعيد عرضها جزئيًا عند ي. كونيغ، \srcLTR{Theorie d. algebr. Größen}، الفصل السادس).}

وللبرهان نكوّن كثيرات الحدود
\[
\begin{array}{@{}l@{\quad}c@{}}
(3)& \begin{gathered}
\bar G_1^*(x)=G_1^*(x)-\lambda_1,\ldots,\bar G_\rho^*(x)=G_\rho^*(x)-\lambda_\rho,\\
\bar\Phi^*(x)=\Phi^*(x)-\mu,
\end{gathered}
\end{array}
\]
حيث تدل $\lambda_1\rcdots\lambda_\rho,\mu$ على متغيرات غير معيّنة، وتدل $\Phi^*(x)$ على كثيرة حدود كيفما كانت في $x_1\rcdots x_\rho$. ولتكن محصلة كثيرات الحدود (3) بالنسبة إلى $x_1\rcdots x_\rho,1$\srcfn{**)}{أي المحصلة المتجانسة بالنسبة إلى $x_1\rcdots x_\rho x_0$، إذا جُعلت كثيرات الحدود (3) متجانسة بإدخال متغير غير معيّن إضافي $x_0$، بحيث تبقى الدرجة في $x$ محفوظة.} معطاةً بـ
\[
\begin{array}{@{}l@{\quad}c@{}}
(4)& \displaystyle R(\lambda,\mu)=\left[\begin{array}{cc} \bar G_1^*\rcdots \bar G_\rho^*\bar\Phi^*\\ x_1\rcdots x_\rho\quad 1\end{array}\right];
\end{array}
\]
فهي دالة نسبية صحيحة في $\lambda_1\rcdots\lambda_\rho,\mu$. ويمكن الآن، بحسب ف. مرتنس، تعيين المعامل الأعلى تعيينًا صريحًا في توسعة (4) بحسب قوى $\mu$\srcfn{***)}{الموضع نفسه، \S\ 3 (برهان المبرهنة الأولى)، وبصورة صريحة \S\ 6 (تعيين مجموع جميع حدود المحصلة $R$ التي تحتوي حاصل القوى المعطى $a_{\lambda\lambda}^{p_\lambda}a_{\mu\mu}^{p_\mu}\cdots a_{\sigma\sigma}^{p_\sigma}$ عاملًا). وفي حالتنا يوضع $p_\lambda=0$، و$p_\mu=0,\ldots,p_\sigma=\tau$، و$a_{\sigma\sigma}=(-\mu)$. أُعيد عرضه عند كونيغ، الفصل السادس، \S\ 9.}؛ فنحصل على
\[
\begin{array}{@{}l@{\quad}c@{}}
(5)& (-1)^\tau R(\lambda,\mu)=R_0^\sigma\cdot\mu^\tau+A_1(\lambda)\mu^{\tau-1}+\cdots+A_\tau(\lambda),
\end{array}
\]
حيث تدل $A_1(\lambda)\rcdots A_\tau(\lambda)$ على دوال كثيرة حدود ذات معاملات صحيحة في $\lambda_1\rcdots\lambda_\rho$ وفي معاملات كثيرات الحدود (3). وتبيّن هذه التوسعة أولًا أنه، تحت الفرض (2) $R_0\ne0$، لا يمكن أن تنعدم $R(\lambda,\mu)$ انعدامًا تطابقيًا أيضًا. وأما علاقة التطابق في $x_1\rcdots x_\rho,\lambda_1\rcdots\lambda_\rho\mu$:
\[
R(\lambda,\mu)\equiv0\pmod{\bar G_1^*(x),\ldots,\bar G_\rho^*(x),\bar\Phi^*(x)}
\]
فتتحول كذلك، بالتعويض:
\[
\lambda_i=G_i^*(x),\qquad \mu=\Phi^*(x)
\]
إلى:
\[
\begin{array}{@{}l@{\quad}l@{}}
(6)& \displaystyle R(G_i^*(x),\Phi^*(x))=R_0^\sigma\cdot\Phi^{*\tau}(x)+A_1(G_i^*(x))\Phi^{*\tau-1}(x)+\cdots\\
& \displaystyle\qquad +A_\tau(G_i(x))=0,
\end{array}
\]
وبذلك تثبت المبرهنة المساعدة، لأن $R_0$، بحسب الفرض، عدد من $\Omega$.

\section*{\S\ 12. القاعدة التكاملية للمنظومات المنتظمة $\Ssys_{n\rho}$ من كثيرات الحدود.}

تعطي المبرهنة المساعدة في \S\ 11، مع خاتمة \S\ 10، الحقيقة الآتية مباشرةً:

\emph{إذا وُجدت في منطقة تصوير $\Jdom^*_{\rho\rho}$ لـ$\Gdom_{n\rho}$ كثيرات حدود عددها $\rho$، بحيث لا تنعدم المحصلة المتجانسة للأجزاء ذات أعلى درجة انعدامًا تطابقيًا --- $R_0\ne0$ ---، فإن $\Gdom_{n\rho}$ من النوع الأول، ومن ثم هي منطقة تكاملية منتهية التوليد، وتكون قاعدة الإنفولوشن فيها قاعدةً تكاملية.}

ويتيح الشرط $R_0\ne0$ الآن، بفضل المبرهنة المساعدة، برهانًا ثانيًا لوجود القاعدة التكاملية، يرجع في جوهره إلى د. هيلبرت\srcfn{*)}{\srcLTR{Über die vollen Invariantensysteme}، \S\ 2. عند هيلبرت، وبافتراض أن وجود القاعدة التكاملية لحقل الامتغيرات قد ثبت من طريق آخر، لا يتعلق الأمر إلا بتفسير هذه القاعدة بواسطة منظومة كرونيكر الأساسية للكميات الصحيحة جبريًا في حقل.}؛ وهو، وإن لم يُظهر قاعدة الإنفولوشن من حيث هي كذلك، يسري في المقابل على نحو موحد على جميع المنظومات $\Ssys_{n\rho}$ ذات الخاصية المذكورة.

ولتكن، على وجه التحديد، كثيرات الحدود الـ$\rho$ من $\Jdom^*_{\rho\rho}$ التي يتحقق لها $R_0\ne0$، والتي تكون لذلك مستقلة جبريًا، معطاةً بـ
\[
\begin{array}{@{}l@{\quad}c@{}}
(1)& G_1^*(x)\rcdots G_\rho^*(x).
\end{array}
\]
بحسب المبرهنة المساعدة، تعتمد كل كثيرة حدود كيفما كانت في $x_1\rcdots x_\rho$ اعتمادًا صحيحًا جبريًا على كثيرات الحدود (1)؛ وينطبق ذلك، بوجه خاص، على كل كثيرة حدود تنتمي إلى أصغر حقل حاوٍ لـ$\Jdom^*_{\rho\rho}$، أي $\Kfield^*_{\rho\rho}$. فإذا ناظرت هذه الكثيرات في $\Ssys_{n\rho}$ كثيرات الحدود:
\[
\begin{array}{@{}l@{\quad}c@{}}
(2)& G_1(x)\rcdots G_\rho(x),
\end{array}
\]
فإن كل كثيرة حدود في أصغر حقل حاوٍ لـ$\Ssys_{n\rho}$، أي $\Kfield_{n\rho}$، وبالتالي كل كثيرة حدود من $\Ssys_{n\rho}$ على وجه الخصوص، تعتمد بحسب مبدأ النقل اعتمادًا صحيحًا جبريًا على كثيرات الحدود (2). وبالعكس، فإن كل دالة من $\Kfield_{n\rho}$ تعتمد اعتمادًا صحيحًا جبريًا على كثيرات الحدود (2) هي أيضًا كثيرة حدود. فإذا اعتبرنا، وفق \S\ 4، $\Kfield_{n\rho}$ حقلًا جبريًا فوق $\Omega(G_1\rcdots G_\rho)$، كانت جملة كثيرات الحدود في $\Kfield_{n\rho}$ --- أي المنطقة الصحيحة نسبيًا $\Gdom_{n\rho}$ --- مطابقةً لجملة الكميات الصحيحة جبريًا في هذا الحقل. ولتكن $G(x)$ كثيرة حدود من $\Kfield_{n\rho}$ تُكمل (2) إلى قاعدة نسبية، ولتكن $D(x)$ مميّز المعادلة غير القابلة للاختزال ذات الدرجة $k$ التي تحققها $G(x)$ بالنسبة إلى $\Omega(G_1\rcdots G_\rho)$. عندئذ تقبل كل كثيرة حدود في $\Kfield_{n\rho}$، وبوجه خاص كل كثيرة حدود $F(x)$ من $\Ssys_{n\rho}$، بوصفها كمية صحيحة جبريًا، التمثيل الآتي:
\[
\begin{array}{@{}l@{\quad}c@{}}
(3)& \displaystyle F(x)=\frac{\Gamma_1(G_i)\cdot G(x)^{k-1}+\Gamma_2(G_i)\cdot G(x)^{k-2}+\cdots+\Gamma_k(G_i)}{D(x)}.
\end{array}
\]

ولْتجرِ $F(x)$ الآن على جميع كثيرات الحدود من $\Ssys_{n\rho}$. عندئذ يبيّن تطبيق المبرهنة الأولى لهيلبرت في قاعدة المودول (\srcLTR{Math. Ann. 36}) على المنظومة التي تُكوَّن من الدوال المناظرة $v_1\Gamma_1(G_i)+v_2\Gamma_2(G_i)+\cdots+v_k\Gamma_k(G_i)$ وجود عدد منتهٍ من كثيرات الحدود من $\Ssys_{n\rho}$:
\[
\begin{array}{@{}l@{\quad}c@{}}
(4)& F_1(x),F_2(x),\ldots,F_\lambda(x)
\end{array}
\]
بحيث تقبل كل كثيرة حدود من $\Ssys_{n\rho}$ تمثيلًا من الصورة:
\[
\begin{array}{@{}l@{\quad}c@{}}
(5)& F(x)=A_1(G_i)F_1(x)+A_2(G_i)\cdot F_2(x)+\cdots+A_\lambda(G_i)F_\lambda(x),
\end{array}
\]
حيث تدل $A_1\rcdots A_\lambda$ على كثيرات حدود في $G_i(x)$. ومن ثم تكوّن كثيرات الحدود (2) و(4) قاعدة تكاملية لـ$\Ssys_{n\rho}$؛ أي:

\emph{المبرهنة VIII: إذا وُجدت في منظومة تصوير $\Jdom^*_{\rho\rho}$ للمنظومة $\Ssys_{n\rho}$ المؤلفة من كثيرات حدود، كثيرات حدود عددها $\rho$ بحيث لا تنعدم محصلة الأجزاء ذات أعلى درجة انعدامًا تطابقيًا --- $R_0\ne0$ ---، فإن $\Ssys_{n\rho}$ تملك قاعدة تكاملية.}

وتقبل هذه المبرهنة تعميمًا يسيرًا. فيكفي أن تنتمي كثيرات الحدود الـ$\rho$ في (1)، التي يتحقق لها $R_0\ne0$، إلى أصغر منطقة تكاملية حاوية لـ$\Ssys^*_{\rho\rho}$، أي $\Jdom^*_{\rho\rho}$. والواقع أن الاستدلالات المؤدية إلى المبرهنة VIII تبيّن أن التمثيل (5) يظل صحيحًا في هذه الحالة أيضًا. غير أن كل كثيرة حدود $L(x)$ في أصغر منطقة تكاملية حاوية $\Jdom_{n\rho}$ يمكن، بحسب التعريف، تمثيلها بتركيب كثير حدود من عدد منتهٍ --- يتغير بتغير $L(x)$ --- من كثيرات الحدود من $\Ssys_{n\rho}$؛ ولذلك توجد كثيرات حدود عددها $\sigma$ من $\Ssys_{n\rho}$
\[
\begin{array}{@{}l@{\quad}c@{}}
(6)& K_1(x),K_2(x),\ldots,K_\sigma(x)
\end{array}
\]
بحيث تصير كثيرات الحدود $G_i(x)$ تراكيب كثيرة حدود من (6). وهكذا تعطي (6) و(4) قاعدة تكاملية. وبإدخال التعريف الآتي:

\emph{التعريف VIII: تسمى منظومة $\Ssys_{n\rho}$ من كثيرات الحدود منتظمةً إذا وُجدت في منطقة تصوير $\Jdom^*_{\rho\rho}$ لأصغر منطقة تكاملية حاوية كثيرات حدود عددها $\rho$ بحيث لا تنعدم محصلة الأجزاء ذات أعلى درجة انعدامًا تطابقيًا --- $R_0\ne0$ ---}

ينتج:

\emph{المبرهنة IX: كل منظومة منتظمة $\Ssys_{n\rho}$ من كثيرات الحدود تملك قاعدة تكاملية.}\srcfn{*)}{أثبت هيلبرت بمثال وجود منظومات غير منتظمة لا قاعدة تكاملية لها (\srcLTR{Gött. Nachr. 1891, S. 233}). والمثال الآتي أبسط قليلًا:
\[
\Ssys_{22}=x_1,
x_1x_2,
x_1x_2^2,
\ldots,
x_1x_2^\nu,
\ldots\quad\text{in inf.}
\]
ويكافئ وجود قاعدة تكاملية هنا وجود عدد صحيح موجب $\nu$ بحيث تتحقق المتطابقات:
\[
x_1x_2^{\nu+\sigma}=\sum c\cdot x_1^{i_0}(x_1x_2)^{i_1}(x_1x_2^2)^{i_2}\cdots(x_1x_2^\nu)^{i_\nu},
\qquad(\sigma=1,2,\ldots).
\]
غير أن مقارنة الدرجتين في $x_1,x_2$ تؤدي، حتى عند $\sigma=1$، إلى معادلتي الشرط الآتيتين للأسس:
\[
1=i_0+i_1+i_2+\cdots+i_\nu,
\qquad
\nu+1=i_1+2i_2+\cdots+\nu i_\nu,
\]
ومن الواضح أنهما لا تقبلان حلًا في الأعداد الصحيحة غير السالبة. ومن ثم لا تملك $\Ssys_{22}$ قاعدة تكاملية.}

ونعطي أيضًا، على غرار ما فعلناه للإنفولوشن في \S\ 5، تأويلًا هندسيًا للمنظومات المنتظمة. ولهذا ننظر، كما في (7) هناك، في منظومة المعادلات
\[
\begin{array}{@{}l@{\quad}c@{}}
(7)& L^*(x)-L^*(\xi)=0,
\end{array}
\]
حيث تجري $L^*(x)$ على جميع كثيرات الحدود من $\Jdom^*_{\rho\rho}$. وإذا جُعلت (7) متجانسةً بواسطة $x_0,\xi_0$، فلتنتقل إلى منظومة المعادلات الآتية بين صيغ متجانسة:
\[
\begin{array}{@{}l@{\quad}c@{}}
(8)& \xi_0^m\cdot M^*(x)-x_0^m\cdot M^*(\xi)=0,
\end{array}
\]
وعندما يدل $H^*(x)$ على الجزء المتجانس ذي أعلى درجة من $L^*(x)$، تتحقق العلاقة:
\[
\begin{array}{@{}l@{\quad}c@{}}
(9)& H^*(x)\equiv M^*(x)\pmod{x_0}.
\end{array}
\]

نسمي \emph{النقاط الأساسية} لـ$\Jdom^*_{\rho\rho}$ (قارن \S\ 5، الحاشية) حلول (8) المستقلة عن $x$ --- أي المختلفة عن الصفوف المترافقة مع $x$. فالنقاط الأساسية هي إذن منظومات القيم غير المنظومة $\xi_1=0,\ldots,\xi_\rho=0$، والتي يتحقق فيها معًا:
\[
\xi_0^m=0,
\qquad M^*(\xi)=0
\]
أو، بحكم (9):
\[
\xi_0=0,
\qquad H^*(\xi)=0.
\]

فإذا وُجدت صيغ عددها $\rho$ من $H^*(\xi)$ يتحقق لها $R_0\ne0$، أي لا تملك منظومة حلول مشتركة، تعذر أن تملك $\Jdom^*_{\rho\rho}$ نقطة أساسية. وإذا كانت $\Ssys^*_{\rho\rho}$، على وجه الخصوص، مؤلفةً من صيغ متجانسة، تعذر أيضًا، عند $R_0\ne0$، أن تملك $\Ssys^*_{\rho\rho}$ نقطة أساسية؛ لأن هذه النقطة ستكون في الوقت نفسه نقطة أساسية لـ$\Jdom^*_{\rho\rho}$. غير أن العكس صحيح أيضًا. ولهذا ننظر في المنظومة المكوَّنة من جملة الصيغ $H^*(x)$، أي $\Jdom(H^*)$، وهي تكوّن بدورها منطقة تكاملية.

ينتج لذلك من المبرهنة الأولى لهيلبرت في قاعدة المودول ومن مبرهنة هيلبرت المساعدة (قارن \S\ 10، الإحالة) وجود صيغ عددها $\rho$ من $\Jdom(H^*)$، يستتبع انعدامها انعدام جميع الصيغ من $\Jdom(H^*)$. فإذا لم تكن لـ$\Jdom^*_{\rho\rho}$ نقطة أساسية، لم يمكن لهذه الصيغ الـ$\rho$ أن تنعدم معًا، ومن ثم تكون محصلتها $R_0\ne0$.

إذن:

\emph{تتميز المنظومات المنتظمة $\Ssys_{n\rho}$ من كثيرات الحدود بأن منطقة تصوير $\Jdom^*_{\rho\rho}$ لأصغر منطقة تكاملية حاوية لا تملك نقطة أساسية. وإذا كانت $\Ssys^*_{\rho\rho}$ منظومة منتظمة من صيغ متجانسة، فإن المنظومة نفسها لا تملك نقطة أساسية.}\srcfn{*)}{تملك المنظومة غير المنتظمة $\Ssys_{22}$ المذكورة في الحاشية السابقة النقطة الأساسية:
\[
\xi_1:\xi_2:\xi_0=0:1:0.
\]}

\section*{\S\ 13. أمثلة على مناطق صحيحة نسبيًا من النوع الأول وعلى منظومات منتظمة.}

1. بوصفه المثال الأول لمنطقة صحيحة نسبيًا من النوع الأول $\Gdom_{nn}$، ننظر في \emph{مجموع كثيرات الحدود في $x_1\rcdots x_n$ التي تقبل زمرة تبديلات معيّنة $G$} --- أي مجموع كثيرات الحدود في حقل من النوع اللاغرانجي. وبما أن $x_1\rcdots x_n$، بموجب المطابقة
\[
\begin{array}{@{}l@{\quad}c@{}}
(1)& x_k^n-\sigma_1(x)x_k^{n-1}+\sigma_2(x)x_k^{n-2}+\cdots\pm\sigma_n(x)=0\quad(k=1,2,\ldots,n)
\end{array}
\]
تعتمد اعتمادًا صحيحًا جبريًا على الدوال المتناظرة الأولية التابعة إلى $\Gdom_{nn}$، ومن ثم، بحسب التعريف V، تعتمد اعتمادًا صحيحًا جبريًا على $\Gdom_{nn}$، فإن $\Gdom_{nn}$ من النوع الأول. وتؤلف $\Gdom_{nn}$، فضلًا عن ذلك، بوصفها منطقة من النوع الأول، منظومة منتظمة، لأن عناصرها صيغ متجانسة ولأن $n=\rho$ (بحسب خاتمة \S\ 10). وبالفعل، لا يمكن، بحكم (1)، أن تنعدم محصلة الدوال المتناظرة الأولية $R_0$ انعدامًا تطابقيًا؛ وتُحسب $R_0$، وفق قواعد الحساب في نظرية المحصلة، فتساوي $\pm1$. أما \emph{صيغة الإنفولوشن} لـ$\Gdom_{nn}$ فهي --- لما كانت $\Gdom_{nn}$ من النوع الأول وكان $n=\rho$ --- محددة تحديدًا وحيدًا بوصفها صيغة الإنفولوشن لحقل النوع اللاغرانجي المقابل. وبما أن المقادير المترافقة مع $x_1\rcdots x_n$ بالنسبة إلى هذا الحقل تنشأ من تبديلات الزمرة $G$، فإن صيغة الإنفولوشن هذه معطاة بـ
\[
\begin{array}{@{}l@{\quad}c@{}}
(2)&
\begin{aligned}
  \Psi(z,u)&=\prod_G
  (z-u_1x_{i_1}-u_2x_{i_2}-\cdots-u_nx_{i_n})\\
  &=z^i+\sum{}'G_{\alpha\alpha_1\ldots\alpha_n}(x)
  z^\alpha u_1^{\alpha_1}\cdots u_n^{\alpha_n},
\end{aligned}
\end{array}
\]
حيث يمتد الجداء على جميع تبديلات الزمرة $G$. وتمثل (2) \glqq{}صيغة غالوا للزمرة $G$\grqq{}، ومن ثم تصح الحقيقة الآتية:

\emph{كل كثيرات الحدود في $x_1\rcdots x_n$ التي تقبل زمرة تبديلات $G$ هي تراكيب كثيرة حدود في المعاملات $G_{\alpha\alpha_1\ldots\alpha_n}(x)$ لصيغة غالوا للزمرة.}

2. وبوصفها أمثلة على \emph{المنظومات المنتظمة}، ننظر في المنظومات $\Ssys_{n1}$ المؤلفة من كثيرات حدود، أي المنظومات التي لا تحتوي إلا كثيرة حدود واحدة مستقلة جبريًا.

وبالفعل، لتكن $\Ssys^*_{11}$ منظومة تصوير لـ$\Ssys_{n1}$، ولتكن
\[
  F^*(x)=a_0x^k+a_1x^{k-1}+\cdots+a_k\qquad(a_0\ne0),
\]
كثيرة حدود من $\Ssys^*_{11}$؛ فعندئذ يكون
\[
\begin{array}{@{}l@{\quad}c@{}}
(3)& R_0=a_0\ne0\text{\srcfnmark{*)}}
\end{array}
\]
\srcfntext{*)}{ميرتنس، في الموضع المذكور، \S\ 5.}
ومن ثم تكون $\Ssys_{n1}$ منتظمة بحسب التعريف VIII، وبالتالي يصح:

\emph{كل منظومة $\Ssys_{n1}$ من كثيرات الحدود، وبخاصة كل منظومة من عدد غير منته من كثيرات الحدود في متغير غير معيّن واحد، تملك قاعدة تكاملية.}\srcfn{**)}{إن أبسط المنظومات غير المنتظمة الممكنة التي لا قاعدة تكاملية لها هي إذن منظومات $\Ssys_{22}$. ويدل على وجودها فعلًا مثال هيلبرت والمثال الوارد في ملاحظة \S\ 12.}

3. نخصّص الآن المنظومات $\Ssys_{n1}$ إلى مناطق صحيحة نسبيًا $\Gdom_{n1}$، وهي، بما أنها منظومات منتظمة، مناطق صحيحة نسبيًا من النوع الأول بحسب \S\ 12. ومن ثم فإن قاعدتها التكاملية معطاة بقاعدة إنفولوشن، هي في الوقت نفسه قاعدة إنفولوشن لأصغر حقل حاوٍ $\Kfield_{n1}$. وبحسب \S\ 6، فإن دالة اعتباطية من قاعدة الإنفولوشن لـ$\Kfield_{n1}$ تكون في الوقت نفسه قاعدة دنيا --- وقد سميناها دالة لوروت للحقل --- وتعتمد كل دالة من قاعدة الإنفولوشن اعتمادًا خطيًا كسريًا على دالة لوروت هذه. وفي حالتنا تكون دالة لوروت كثيرة حدود، لأنها تنتمي إلى $\Gdom_{n1}$؛ ومن ثم تعتمد كل كثيرة حدود أخرى من قاعدة الإنفولوشن اعتمادًا خطيًا كسريًا، وبحكم (3) اعتمادًا صحيحًا جبريًا، فتكون بالتالي كثيرة حدود خطية في دالة لوروت هذه $L(x)$، التي تصير بذلك قاعدةً تكاملية. وتكون صيغة الإنفولوشن لـ$\Gdom^*_{11}$، بعد وضع $(u_1=1)$ أيضًا، كما يأتي:
\[
  \Psi^*(z)=L^*(z)-L^*(x),
\]
ويتضح منها مرة أخرى أن $L(x)$ قاعدة تكاملية. ومن ثم يصح:

\emph{المناطق الصحيحة نسبيًا $\Gdom_{n1}$ من النوع الأول؛ وتتألف قاعدتها التكاملية من كثيرة حدود واحدة، هي دالة لوروت لأصغر حقل حاوٍ.}

4. وبوصفه مثالًا مخصوصًا آخر لمنطقة صحيحة نسبيًا $\Gdom_{n1}$، نذكر اللامتغيرات الإسقاطية كثيرة الحدود لصيغة تربيعية في $s$ من المتغيرات. واتساقًا مع ما تقدم، يمكن، كما هو معلوم، تمثيلها بتركيب كثير الحدود في محددة الصيغة $|a_{ik}|$؛ وهذه المحددة هي في الوقت نفسه دالة لوروت لحقل اللامتغيرات المقابل.

\section*{\S\ 14. منظومات من كثيرات الحدود الصحيحة.}

نريد الآن تشديد مبرهنات القاعدة التي حصلنا عليها، وذلك من جهة صحة المعاملات.

ولذلك نفترض فيما يأتي أن نطاق المعاملات $\Omega$ حقل أعداد جبري (منتهٍ)، ولتدل $[\Omega]$ على مجموعة الأعداد الصحيحة الجبرية في $\Omega$، التي سنسميها اختصارًا أعدادًا صحيحة. ونقصد بكثيرة حدود صحيحة كثيرةَ حدود معاملاتُها من $[\Omega]$؛ ومن الآن فصاعدًا لا ترمز $F(x),G(x),\ldots$ إلا إلى كثيرات حدود صحيحة.

وننظر، قياسًا على \S\ 7، في مختلف المنظومات المؤلفة من كثيرات حدود صحيحة:

تسمى كل منظومة كيفما كانت من كثيرات الحدود الصحيحة \emph{منظومة صحيحة المعاملات} $\Ssys_{n\rho}$.

الأسرة الخطية الصحيحة المعاملات $\Lsys_{n\rho}$ هي منظومة صحيحة المعاملات تشتمل، مع $F_1(x)$ و$F_2(x)$، دائمًا على $c_1F_1(x)+c_2F_2(x)$ أيضًا، حيث $c_1,c_2$ عددان صحيحان كيفما كانا من $[\Omega]$.

المنطقة التكاملية الصحيحة المعاملات $\Jdom_{n\rho}$ هي أسرة خطية صحيحة المعاملات تشتمل، مع $F(x)$ و$G(x)$، دائمًا على $F(x)\cdot G(x)$ أيضًا.

المنطقة الصحيحة نسبيًا ذات المعاملات الصحيحة $\Gdom_{n\rho}$ هي منطقة تكاملية تشتمل، مع $F(x)$ و$G(x)$ ــ حيث $G(x)\ne0$ ــ على خارج القسمة $F(x):G(x)$ دائمًا حين يكون هذا الخارج كثيرة حدود صحيحة، وحينئذٍ فقط.

وكما في \S\ 7، تتألف أصغر منطقة تكاملية صحيحة المعاملات $\Jdom_{n\rho}$ تشتمل على منظومة صحيحة المعاملات $\Ssys_{n\rho}$ من جميع كثيرات الحدود $H(x)$ التي هي تراكيب كثيرة حدود بمعاملات صحيحة لعدد منتهٍ من كثيرات الحدود في $\Ssys_{n\rho}$؛

وتتألف أصغر منطقة صحيحة نسبيًا ذات معاملات صحيحة $\Gdom_{n\rho}$ تشتمل عليها من جميع كثيرات الحدود الصحيحة $H(x)$ التي هي تراكيب نسبية لعدد منتهٍ من كثيرات الحدود في $\Ssys_{n\rho}$.

وننتقل الآن إلى نقل مبرهنات القاعدة. وفيما يتعلق بالقاعدة النسبية، ينبغي ملاحظة أن مفاهيم الحقل وأصغر حقل مشتمل والقاعدة النسبية، التي لا تقوم إلا على العمليات «النسبية»، لا يطرأ عليها تغيير. ويمكن، فوق ذلك، اختيار منظومة التصوير $\Ssys^*_{\rho\rho}$ لمنظومة صحيحة المعاملات $\Ssys_{n\rho}$، بطرائق كثيرة كيفما شئنا، على هيئة منظومة صحيحة المعاملات من جديد. ولما كانت الاستدلالات المؤدية إلى الحد العددي تبقى صحيحة أيضًا للأسرة الخطية الصحيحة المعاملات، ثبت ما يأتي:

\emph{المبرهنة الخامسة$'$: لكل منظومة صحيحة المعاملات $\Ssys_{n\rho}$ قاعدة نسبية؛ ويمكن، على وجه الخصوص، اختيار القاعدة النسبية لأسرة خطية صحيحة المعاملات $\Lsys_{n\rho}$ بحيث تشتمل على الأكثر على $\rho+1$ كثيرة حدود.}

ونبحث الآن النظير لقاعدة الإنفولوشن للمناطق التكاملية المؤلفة من كثيرات الحدود (\S\ 8). ولهذا ينبغي توسيع المبرهنة المتعلقة \emph{بالدوال المتناظرة لصفوف المقادير} من جهة \emph{صحة المعاملات} بالمبرهنة المساعدة الآتية:

\emph{المبرهنة المساعدة. إن الدوال المتناظرة كثيرة الحدود الصحيحة المعاملات في $n$ من صفوف المقادير $(yz\cdots t)$:}
\[
\begin{array}{@{}l@{\quad}c@{}}
(1)&
\begin{array}{cccc}
  y_1&z_1&\cdots&t_1\\
  y_2&z_2&\cdots&t_2\\
  \vdots&\vdots&&\vdots\\
  y_n&z_n&\cdots&t_n
\end{array}
\end{array}
\]
\emph{هي دوال كثيرة الحدود صحيحة المعاملات في عدد منتهٍ منها، هي الدوال المتناظرة الأولية}
\[
\begin{array}{@{}l@{\quad}c@{}}
(2)&
\begin{gathered}
  \sigma_1(y),\sigma_2(y),\ldots,\sigma_n(y);\quad
  \sigma_1(z),\sigma_2(z),\ldots,\sigma_n(z);\quad\ldots;\\
  \sigma_1(t),\ldots,\sigma_n(t);
\end{gathered}
\end{array}
\]
\emph{والدوال}
\[
\begin{array}{@{}l@{\quad}c@{}}
(3)&
  \chi_1(y,z,\ldots,t),\ldots,\chi_k(y,z,\ldots,t),
\end{array}
\]
\emph{التي يكون اعتمادها على الدوال (2) صحيحًا جبريًا وصحيحَ المعاملات.}

ففي الواقع، تعتمد المقادير الـ$n$، وهي $y_1,
\ldots,y_n$، على $\sigma_1(y),
\ldots,\sigma_n(y)$ اعتمادًا صحيحًا جبريًا وصحيحَ المعاملات؛ وينطبق الأمر نفسه على $z_1,
\ldots,z_n$ بالنسبة إلى $\sigma_1(z),
\ldots,\sigma_n(z)$، وهكذا. ومن ثم يكون حقل الدوال المتناظرة لصفوف المقادير (1) حقلًا جبريًا فوق الحقل الذي تحدده الدوال (2)، على نحو يجعل المقادير الصحيحة جبريًا والصحيحة المعاملات في هذا الحقل مطابقةً للدوال المتناظرة كثيرة الحدود الصحيحة المعاملات في صفوف المقادير (1). فإذا عُيّنت الدوال (3) منظومةَ كرونيكر الأساسية، ثبتت المبرهنة المساعدة.

ولتكن الآن $\Jdom_{n\rho}$ منطقة تكاملية صحيحة المعاملات، و$\Jdom^*_{\rho\rho}$ منطقة تصوير صحيحة المعاملات، ولتكن
\[
  \Psi^*(z,u)=G^*(x)\cdot z^i+
  \sum{}'G^*_{\alpha\alpha_1\ldots\alpha_\rho}(x)\cdot
  z^\alpha u_1^{\alpha_1}\cdots u_\rho^{\alpha_\rho}
\]
صيغة إنفولوشن لـ$\Jdom^*_{\rho\rho}$؛ ويجوز هنا أن يكون $G^*(x)$ من البعد صفر أيضًا، أي عددًا من $[\Omega]$. وعندئذ تكون الدوال المتناظرة الأولية المناظرة للدوال (2) معطاة ببعض خوارج القسمة الآتية:
\[
\begin{array}{@{}l@{\quad}c@{}}
(4)&
  G^*_{\alpha\alpha_1\ldots\alpha_\rho}(x):G^*(x).
\end{array}
\]
ولما كان اعتماد الدوال (3)، بحسب المبرهنة المساعدة، على (2) صحيحًا جبريًا وصحيحَ المعاملات، فإن الدوال المناظرة لـ(3) من $\Kfield^*_{\rho\rho}$ تكون ــ بحسب توسيع مبرهنة غاوس إلى كثيرات الحدود ذات المعاملات الصحيحة الجبرية ــ معطاةً بـ
\[
\begin{array}{@{}l@{\quad}c@{}}
(5)&
  K_\lambda^*(x):(G^*(x))^\sigma,
\end{array}
\]
حيث $K_\lambda^*(x)$ كثيرات حدود صحيحة المعاملات من $\Kfield^*_{\rho\rho}$. ومن ثم فإن $K_\lambda^*(x)$، عندما يتعلق الأمر بمناطق صحيحة نسبيًا، تنتمي إلى المنطقة؛ وهي في الحالة العامة خوارج قسمة كثيرات حدود من $\Jdom^*_{\rho\rho}$.

وإذا كانت $F^*(x)$ كثيرة حدود صحيحة كيفما كانت من $\Jdom^*_{\rho\rho}$، كان لدينا
\[
  F^*(x)=\frac1i\Bigl(F^*(x)+F^*(x^{(1)})+\cdots+F^*(x^{(i-1)})\Bigr)
\]
ومن ثم تكون $i\cdot F^*(x)$ دالة متناظرة كثيرة الحدود صحيحة المعاملات في صفوف المقادير $x,x^{(1)},\ldots,x^{(i-1)}$.

وبحسب المبرهنة المساعدة ومبدأ النقل، يثبت من ثم ما يأتي:

\emph{المبرهنة السادسة$'$: لكل منطقة تكاملية صحيحة المعاملات $\Jdom_{n\rho}$ تحدد كل صيغة إنفولوشن قاعدة نسبية متميزة، بحيث تقبل كل كثيرة حدود $F(x)$ من $\Jdom_{n\rho}$ تمثيلًا من الصورة:}
\[
  F(x)=\frac{\Gamma(H(x),H_1(x),\ldots,H_\sigma(x))}{i\cdot H(x)^k},
\]
\emph{حيث ترمز $\Gamma$ إلى دالة كثيرة الحدود صحيحة المعاملات. وللمناطق الصحيحة نسبيًا ذات المعاملات الصحيحة $\Gdom_{n\rho}$، إذا كانت منطقة التصوير $\Gdom^*_{\rho\rho}$ أيضًا صحيحة نسبيًا وذات معاملات صحيحة، يكون المقام $H(x)$ هو المعامل الأعلى لصيغة الإنفولوشن المناظرة.}

\section*{\S\ 15. مناطق صحيحة نسبيًا ذات معاملات صحيحة من النوع الأول ومنظومات منتظمة صحيحة المعاملات.}

تحلّ القاعدة التكاملية الصحيحة المعاملات، في المنظومات الصحيحة المعاملات، محل القاعدة التكاملية؛ أي إنها عدد منتهٍ من كثيرات الحدود الصحيحة المعاملات في المنظومة، بحيث يمكن تمثيل كل كثيرة حدود صحيحة المعاملات في المنظومة على صورة تركيب كثير حدود صحيح المعاملات في تلك الكثيرات الحدود المنتهية العدد.

تسير المبرهنات المتعلقة بالقاعدة التكاملية الصحيحة المعاملات، في جوهرها، موازيةً للمبرهنات السابقة، ولا تتطلب إلا تعديلات طفيفة في البراهين.

ويمكن أولًا نقل التعريف الخامس (\S\ 9) نقلًا مباشرًا.

\emph{التعريف الخامس$'$: تسمى كمية $\xi$ صحيحةً جبريًا وصحيحةَ المعاملات بالنسبة إلى المنطقة الصحيحة نسبيًا ذات المعاملات الصحيحة $\Gdom_{n\rho}$ إذا حققت معادلةً ما معاملها الأعلى وحدةٌ من $[\Omega]$، ومعاملاتها الباقية كثيرات حدود من $\Gdom_{n\rho}$.}

ذلك أن توسيع مبرهنة غاوس إلى كثيرات الحدود ذات المعاملات الصحيحة الجبرية يبين، كما في \S\ 9، أنه يمكن من ذلك تحصيل التعريف بواسطة المعادلة غير القابلة للاختزال.

وينتج من التعريف الخامس$'$ كذلك نقل التعريف السابع.

\emph{التعريف السابع$'$: تسمى المنطقة الصحيحة نسبيًا ذات المعاملات الصحيحة $\Gdom_{n\rho}$ من النوع الأول إذا كان لها، بوصفها منطقة تصوير، منطقة صحيحة نسبيًا $\Gdom^*_{\rho\rho}$ بحيث تعتمد كل كثيرة حدود صحيحة المعاملات كيفما كانت في $x_1\rcdots x_\rho$ اعتمادًا صحيحًا جبريًا وصحيحَ المعاملات على $\Gdom^*_{\rho\rho}$.}

ولإثبات أن هذه المناطق من النوع الأول منتهية التوليد، نلاحظ أولًا أنه، كما في \S\ 10، يترتب من التعريف أن المعامل الأعلى لصيغة الإنفولوشن وحدةٌ من $[\Omega]$.

ومن ثم، فبحسب المبرهنة السادسة$'$ (\S\ 14)، تقبل كل كثيرة حدود $F(x)$ من $\Gdom_{n\rho}$ تمثيلًا على الصورة الآتية:
\[
  F(x)=\frac{\Gamma(H_1(x)\rcdots H_\sigma(x))}{i}.
\]
فإذا طبقنا هنا مبرهنة هيلبرت الثانية في القاعدة المودولية الصحيحة المعاملات (Math. Ann. 36) --- وهي المبرهنة المصوغة لمعاملات من الأعداد الصحيحة النسبية ---، بعد توسيعها إلى كثيرات الحدود ذات المعاملات الصحيحة الجبرية، على المنظومة المؤلفة من كثيرات الحدود الصحيحة المعاملات $\Gamma(H_1\rcdots H_\sigma)$\srcfn{*)}{إذا كانت $w_1\rcdots w_k$ منظومة أساسية للأعداد الصحيحة الجبرية من $\Omega$، ووضعنا
\[
  \Gamma(H_1\rcdots H_\sigma)=\Gamma_1(H)w_1+\cdots+\Gamma_k(H)w_k,
\]
كانت $\Gamma_1(H)\rcdots \Gamma_k(H)$ ذات معاملات من الأعداد الصحيحة النسبية. ويبين تطبيق المبرهنة الثانية على المنظومة المؤلفة من الصيغ $v_1\Gamma_1(H)+\cdots+v_k\Gamma_k(H)$ مباشرةً صحة هذا التوسيع.}، اتضح من ذلك وجود القاعدة التكاملية الصحيحة المعاملات؛ أي:

\emph{المبرهنة السابعة$'$: لكل منطقة صحيحة نسبيًا ذات معاملات صحيحة من النوع الأول قاعدة تكاملية صحيحة المعاملات.}

وننقل كذلك تعريف المنظومات المنتظمة:

\emph{التعريف الثامن$'$: تسمى المنظومة الصحيحة المعاملات $\Ssys_{n\rho}$ منتظمةً صحيحةَ المعاملات إذا وُجدت في منطقة تصوير $\Jdom^*_{\rho\rho}$ لأصغر منطقة تكاملية صحيحة المعاملات حاوية كثيرات حدود عددها $\rho$ بحيث تكون محصلة الأجزاء ذات أعلى درجة وحدةً من $[\Omega]$ --- $R_0=\eps$.}

ولنقل برهان وجود القاعدة التكاملية، نلاحظ أن المبرهنة المساعدة في \S\ 11 تقبل الصيغة الأقوى الآتية --- وهي صيغة تنتج من المعادلتين (5) و(6) في \S\ 11 ومن كون $A_1(\lambda)\rcdots A_\tau(\lambda)$ دوال كثيرة الحدود صحيحة المعاملات في $\lambda_1\rcdots\lambda_\rho$:

\emph{شرط كافٍ لكي تعتمد كل كثيرة حدود صحيحة المعاملات كيفما كانت في $x_1\rcdots x_\rho$ اعتمادًا صحيحًا جبريًا وصحيحَ المعاملات على $\rho$ من كثيرات الحدود الصحيحة المعاملات هو أن تكون محصلة الأجزاء ذات أعلى درجة من هذه الكثيرات الحدود وحدةً من $[\Omega]$ --- $R_0=\eps$.}

وتنتج من هذه المبرهنة المساعدة، كما في \S\ 12 تمامًا، لكل كثيرة حدود $F(x)$ في منظومة منتظمة صحيحة المعاملات، الصيغة (3) من \S\ 12، حيث تكون الآن $\Gamma_1(G_i)\rcdots\Gamma_k(G_i)$ كثيرات حدود صحيحة المعاملات في $G_i$. ثم إن تطبيق مبرهنة هيلبرت الثانية (Math. Ann. 36)، بعد توسيعها إلى كثيرات الحدود ذات المعاملات الصحيحة الجبرية، مقترنًا بتعريف أصغر منطقة تكاملية صحيحة المعاملات حاوية، يعطي، كما في \S\ 12، وجود القاعدة التكاملية الصحيحة المعاملات؛ أي:

\emph{المبرهنة التاسعة$'$: لكل منظومة منتظمة صحيحة المعاملات قاعدة تكاملية صحيحة المعاملات.}

ولنشر، مثالًا على منطقة صحيحة نسبيًا ذات معاملات صحيحة من النوع الأول --- قياسًا على المثال 1 في \S\ 13 ---، إلى مجموع $\Gdom_{nn}$ لكثيرات الحدود الصحيحة المعاملات في حقل من النوع اللاغرانجي. وينتج كون $\Gdom_{nn}$ من النوع الأول وذات معاملات صحيحة من الهوية الآتية:
\[
  x_k^n-\sigma_1(x)x_k^{n-1}+\cdots\pm\sigma_n(x)=0
  \qquad(k=1,2,
\ldots,n);
\]
ولما كانت محصلة الدوال المتناظرة الأولية تساوي $\pm1$، كانت $\Gdom_{nn}$ أيضًا منظومة منتظمة صحيحة المعاملات.

ويعطي المبرهنةُ المساعدة في \S\ 14 مثالًا آخر، هو مجموع الدوال المتناظرة كثيرة الحدود الصحيحة المعاملات في $n$ من صفوف المقادير.

إرلنغن، مايو 1914.

\clearpage

\providecommand{\Hgrp}{\mathfrak{H}}
\providecommand{\GaloisRes}{\Phi}
\providecommand{\pol}[2]{\left(#1\frac{\partial}{\partial #2}\right)}
\providecommand{\Deltaop}{\Delta\!\left(\frac{\partial}{\partial x},\frac{\partial}{\partial y},\frac{\partial}{\partial z}\right)}
\providecommand{\Omegaop}{\Omega}
\providecommand{\NoetherSrcNote}[2]{\textsuperscript{#1}\begingroup\renewcommand{\thefootnote}{}\footnotetext{\textsuperscript{#1}~#2}\endgroup}


\clearpage
\setcounter{footnote}{0}
\section*{٧. مبرهنة تناهي لامتغيرات الزمر المنتهية}
\begin{center}
Math. Ann. ٧٧ (١٩١٦)، ص ٨٩--٩٢
\end{center}

\begin{center}
{\Large مبرهنة تناهي لامتغيرات الزمر المنتهية.}\par
\vspace{1.2em}
بقلم\par
\vspace{0.8em}
\srcspaced{إيمي نوتر} في إرلنغن.
\end{center}

سنقدم فيما يأتي برهانًا أوليًا تمامًا لتناهي لامتغيرات الزمر المنتهية --- لا يقوم إلا على نظرية الدوال المتناظرة --- ويعطي في الوقت نفسه تعيينًا فعليًا للمنظومة الكاملة؛ في حين أن البرهان المعتاد، المستند إلى مبرهنة هيلبرت في قاعدة المودول (Ann. 36)، ليس إلا برهان وجود.\NoetherSrcNote{*)}{قارن مثلًا Weber, \emph{Lehrbuch der Algebra} (الطبعة الثانية)، المجلد ٢، \S\ 57.}

لتتألف الزمرة المنتهية $\Hgrp$ من $h$ تحويلات خطية (محدد كل منها غير منعدم) $A_1,\ldots,A_h$، وليكن $A_k$ ممثلًا للتحويل الخطي
\[
  x_1^{(k)}=\sum_{\nu=1}^{n}a_{1\nu}^{(k)}x_\nu,\ldots,
  x_n^{(k)}=\sum_{\nu=1}^{n}a_{n\nu}^{(k)}x_\nu
\]
أو، اختصارًا، $(x^{(k)})=A_k(x)$. وهكذا تحول الزمرة $\Hgrp$ الصف $(x)$، الذي عناصره $x_1,\ldots,x_n$، إلى الصفوف $(x^{(k)})$، التي عناصرها $x_1^{(k)},\ldots,x_n^{(k)}$. ولما كان لا بد أن يكون التحويل المحايد بين $A_1,\ldots,A_h$، كانت الصفوف $(x^{(k)})$ تشتمل على الصف $(x)$ نفسه. ونقصد بلامتغير مطلق صحيح نسبيًا للزمرة دالة صحيحة نسبية في $x_1,\ldots,x_n$، أي دالة كثيرة حدود، تبقى مطابقةً بلا تغير عند تطبيق $A_1,\ldots,A_h$؛ ومن ثم يصدق لكل لامتغير من هذا القبيل $f(x)$:
\begin{equation}
  f(x)=f(x^{(1)})=\cdots=f(x^{(h)})=
  \frac1h\sum_{k=1}^{h}f(x^{(k)}).
\tag{1}
\end{equation}

\subsection*{١.}
تعبر الصيغة (1) عن أن $f(x)$ دالة متناظرة صحيحة نسبية، أي كثيرة حدود، في صفوف المقادير $(x^{(1)}),\ldots,(x^{(h)})$؛ بل نحن هنا في أبسط حالة، تسمى عادةً \emph{أحادية الصف}، لأن كل حد $f(x^{(k)})$ لا يشتمل إلا على الصف الواحد $(x^{(k)})$. وبحسب المبرهنة المعروفة في الدوال المتناظرة في صفوف المقادير\NoetherSrcNote{*)}{قارن الملاحظة على ٢.} يمكن إذن تمثيل $f(x)$ تمثيلًا كثير الحدود بالدوال المتناظرة الأولية لهذه الصفوف، أي بمعاملات $G_{\alpha\alpha_1\ldots\alpha_n}(x)$ في «كثيرة حدود غالوا المُحَلِّلة»:
\[
\begin{aligned}
  \Phi(z,u)
  &=\prod_{k=1}^{h}\bigl(z+u_1x_1^{(k)}+\cdots+u_nx_n^{(k)}\bigr) \\
  &=z^h+
    \sum G_{\alpha\alpha_1\ldots\alpha_n}(x)
    z^{\alpha}u_1^{\alpha_1}\cdots u_n^{\alpha_n},
\end{aligned}
\qquad
\left(\begin{array}{c}
\alpha+\alpha_1+\cdots+\alpha_n=h\\
\alpha\ne h
\end{array}\right).
\]
وهنا تكون $G_{\alpha\alpha_1\ldots\alpha_n}(x)$ لامتغيرات درجتها $\alpha_1+\cdots+\alpha_n$ في $x$. وبذلك ثبت ما يأتي:

\emph{تؤلف معاملات $G_{\alpha\alpha_1\ldots\alpha_n}(x)$ في كثيرة حدود غالوا المُحَلِّلة منظومة كاملة من لامتغيرات الزمرة، بحيث يمكن تمثيل كل لامتغير تمثيلًا كثير الحدود بهذه اللامتغيرات المنتهية العدد.}

\subsection*{٢.}
يمكن أيضًا، انطلاقًا من (1)، إجراء النظر الآتي، وهو أكثر أولية،\NoetherSrcNote{**)}{وهو في جوهره برهان مبرهنة الدوال المتناظرة في صفوف المقادير للحالة أحادية الصف المذكورة أعلاه.} ويؤدي في الوقت نفسه إلى منظومة كاملة ثانية. لنضع
\[
  f(x)=a+b x_1^{\mu_1}\cdots x_n^{\mu_n}
       +\cdots+c x_1^{\nu_1}\cdots x_n^{\nu_n},
\]
حيث تدل $a,b,\ldots,c$ على ثوابت؛ فعندئذ تعطينا (1):
\[
\begin{aligned}
 h\cdot f(x)
 &=h\cdot a+b\sum_{k=1}^{h}(x_1^{(k)})^{\mu_1}\cdots(x_n^{(k)})^{\mu_n}
     +\cdots
     +c\sum_{k=1}^{h}(x_1^{(k)})^{\nu_1}\cdots(x_n^{(k)})^{\nu_n} \\
 &=h\cdot a+b\cdot J_{\mu_1\ldots\mu_n}+\cdots+c\cdot J_{\nu_1\ldots\nu_n}.
\end{aligned}
\]
ومن ثم يمكن تركيب كل لامتغير تركيبًا خطيًا، من غير مقامات في هذه المولدات، من اللامتغيرات الخاصة
\[
  J_{\mu_1\ldots\mu_n}
  =\sum_{k=1}^{h}(x_1^{(k)})^{\mu_1}\cdots(x_n^{(k)})^{\mu_n}
\]
ولذلك يكفي إثبات المطلوب لهذه اللامتغيرات الخاصة. غير أن $J_{\mu_1\ldots\mu_n}$ هو، مع إغفال عامل عددي، معامل $u_1^{\mu_1}\cdots u_n^{\mu_n}$ في التعبير
\[
  S_\mu=\sum_{k=1}^{h}
  \bigl(u_1x_1^{(k)}+\cdots+u_nx_n^{(k)}\bigr)^{\mu},
\]
حيث وضعنا $\mu=\mu_1+\cdots+\mu_n$، وهذا التعبير هو مجموع القوى من الرتبة $\mu$ للأشكال الخطية $h$ الآتية:
\[
  \xi_1=u_1x_1^{(1)}+\cdots+u_nx_n^{(1)},\ldots,
  \xi_h=u_1x_1^{(h)}+\cdots+u_nx_n^{(h)}
\]
ومن المعلوم أن مجاميع القوى غير المنتهية العدد $S_\mu$ تراكيب كثيرة حدود في
\[
  S_1,S_2,\ldots,S_h,
\]
التي تُعطى معاملاتها باللامتغيرات
\[
  J_{\mu_1\ldots\mu_n}\qquad(\mu_1+\cdots+\mu_n\le h)
\]
وبذلك نحصل على منظومة كاملة ثانية:

\emph{تعطى منظومة كاملة من لامتغيرات الزمرة بجميع اللامتغيرات $J_{\mu_1\ldots\mu_n}$ التي تحقق $\mu_1+\cdots+\mu_n\le h$، حيث يدل $h$ على رتبة الزمرة.}

وتظهر الصلة بالفقرة ١ بملاحظة أن مجاميع القوى $S_1,\ldots,S_h$ يمكن أن تُستبدل بها الدوال المتناظرة الأولية في $\xi_1,\ldots,\xi_h$، وهو ما يؤدي إلى المنظومة الكاملة من $G_{\alpha\alpha_1\ldots\alpha_n}(x)$ المعطاة هناك. وتبين النتيجتان أن جميع اللامتغيرات يمكن التعبير عنها على صورة كثيرات حدود في اللامتغيرات التي لا تتجاوز درجتها في $x$ رتبة الزمرة $h$.

\subsection*{٣.}
يمكن استخلاص نتائج للتمثيل النسبي مما ثبت للتو. فكل لامتغير مطلق نسبي يمكن --- كما يتبين مباشرةً بتوسيع البسط والمقام بضربهما في مترافقات المقام بالنسبة إلى الزمرة --- تمثيله خارج قسمة لامتغيرين مطلقين كثيرَي الحدود، لا يلزم أن يكونا أوليين فيما بينهما. وينتج من ذلك أن كل لامتغير مطلق نسبي يمكن التعبير عنه دالةً نسبيةً في معاملات $G_{\alpha\alpha_1\ldots\alpha_n}(x)$ لكثيرة حدود غالوا المُحَلِّلة $\Phi(z,u)$. ويمكن أيضًا إثبات هذه القضية بسهولة بطرائق شتى من غير المرور بمبرهنة التناهي؛ وقد وردت مصوغةً من قبل في Weber II, \S\ 58.\NoetherSrcNote{*)}{غير أن البرهان المعطى هناك غير سليم؛ إذ لم يُثبت إلا أن معاملات الدالة $\Psi(t)$ في الصيغة (7) لامتغيرات، ولم يُثبت أن هذه اللامتغيرات يمكن التعبير عنها دوالَّ نسبيةً في معاملات $\Phi(t)$. وتُتلافى هذه الثغرة إذا استُعملت، لتمثيل $x_i^{(k)}$ بدلالة $\xi_k$، طريقة معروفة قائمة على التفاضل بدل صيغة لاغرانج للاستيفاء. وهذه هي العلاقة الناتجة من تفاضل الهوية $\Phi(-\xi_k,u)=0$ بالنسبة إلى جميع $u$:
\[
 \left[\frac{\partial\Phi}{\partial u_i}
      -x_i^{(k)}\frac{\partial\Phi}{\partial z}\right]_{z=-\xi_k}=0.
\]
وعليه ينبغي أن تحل، بدل الصيغتين (7) و(8) عند Weber، الصيغة الآتية لكل دالة نسبية $\omega(x)$:
\[
 \omega(x_1^{(k)},\ldots,x_n^{(k)})=
 \left.
 \omega\!\left(
   \frac{\partial\Phi/\partial u_1}{\partial\Phi/\partial z},\ldots,
   \frac{\partial\Phi/\partial u_n}{\partial\Phi/\partial z}
 \right)\right|_{z=-\xi_k},
\]
وهي لا تشتمل بعد ذلك إلا على معاملات $\Phi(z,u)$؛ وجمعها على جميع $k$ يعطي، حين تكون $\omega(x)$ لامتغيرًا، التمثيل المطلوب.}

\subsection*{٤.}
وأخيرًا نذكر أن النتائج المستخرجة هنا واردة ضمنًا في بحثي «Körper und Systeme rationaler Funktionen»؛\NoetherSrcNote{*)}{Math. Ann. 76, S. 161 (1915).} فالمنظومة الكاملة المعطاة في الفقرة ١ واردة في المبرهنتين VI وVII، وبرهان التمثيل النسبي المستقل عن مبرهنة التناهي هذه وارد في المبرهنة III.\NoetherSrcNote{**)}{أما اللامتغيرات النسبية، ففي المبرهنتين VIII وIX برهان تناهٍ يقوم هو أيضًا على أساس غير الأساس المعتاد، ولكنه ليس إلا برهان وجود؛ وسبب ذلك أن اللامتغيرات النسبية لا تؤلف حقلًا.} غير أن مسار البرهان والنتيجة يتبسطان كثيرًا في الحالة الخاصة الحاضرة بالقياس إلى النظرية العامة. وقد خطرت لي فكرة تطبيق البحوث العامة على لامتغيرات الزمر المنتهية من خلال أحاديث مع السيد إ. فيشر؛ وإليه يرجع أيضًا، من حيث المضمون، البرهان المعطى في الفقرة ٢ --- مع أن المنظومة الكاملة التي تظهر هناك لا تكون، في الحالة العامة، قابلة للتعريف نسبيًا --- والملاحظة على الفقرة ٣.

إرلنغن، مايو ١٩١٥.

\clearpage
\section*{٨. في تمثيل لامتغيرات منظومة مؤلّفة من عدد اعتباطي من الأشكال الأساسية بكثيرات الحدود}
\begin{center}
Math. Ann. ٧٧ (١٩١٦)، ص ٩٣--١٠٢
\end{center}

\begin{center}
{\Large في تمثيل لامتغيرات منظومة مؤلّفة من عدد اعتباطي\\
من الأشكال الأساسية بكثيرات الحدود.}\par
\vspace{1.2em}
بقلم\par
\vspace{0.8em}
\srcspaced{إيمي نوتر} في إرلنغن.
\end{center}

في ختام بحثه في تذكار شفارتس «\srcLTR{Über die Invarianten eines Systems von beliebig vielen Grundformen}»، يطرح د. هيلبرت الحدس القائل إنه يمكن \emph{تمثيل هذه اللامتغيرات تمثيلًا صحيحًا نسبيًا، أي كثير الحدود، باللامتغيرات المنتهية العدد للمنظومة $(J,PJ)$}، حيث يدل $J$ على المنظومة الكاملة للامتغيرات الخاصة بالأشكال المستقلة خطيًا من بين الأشكال الأساسية، وتنشأ اللامتغيرات $PJ$ من $J$ بعمليات الاستقطاب.

أبيّن فيما يأتي في I أن هذا الحدس صحيح بالفعل، وذلك بوصفه نتيجة بسيطة لمبرهنة الاختزال التي تقضي \emph{بأن كل شكل في عدد اعتباطي من الصفوف، في كل منها $n$ متغيرات، يمكن تمثيله مجموعًا لأشكال قطبية لأشكال لا تحتوي إلا على $n$ صفوف ثابتة، وتكون بدورها مشتقة من الشكل الأصلي بعمليات الاستقطاب}. ومبرهنة الاختزال هذه حالة خاصة صاغها أولًا ف. مرتنس من التعميم الذي أعطاه أ. كابلي وف. مرتنس وج. ديرويتس لتوسعة كليبش--غوردان المتسلسلة إلى الأشكال ذات عدد اعتباطي من الصفوف، في كل منها $n$ متغيرات؛ وهو تعميم ثبت بتحويل صوري لعمليات التفاضل.\NoetherSrcNote{*)}{F. Mertens: «\srcLTR{Über eine Formel der Determinantentheorie}» (الصيغة ٥)، \srcLTR{Sitzb. d. Ak. d. Wiss. Wien}، المجلد ٩١، القسم II (١٨٨٥)؛ وبصورة أوسع، مع تطبيقات، في: «\srcLTR{Über ganze Funktionen von $m$ Systemen von je $n$ Unbestimmten}»، \srcLTR{Monatsh. f. Math. u. Phys.}، السنة ٤ (١٨٩٣). ولبعض براهين كابلي، منذ ١٨٨٢، قارن مثلًا \srcLTR{Math. Ann.} ٣٧، أو المحاضرات المنسوخة «\srcLTR{Lezioni sulla teoria delle forme algebriche}» (١٩٠٢). J. Deruyts: \emph{\srcLTR{Essai d'une theorie générale des formes algébriques}}، \srcLTR{Mém. d. l. Soc. Roy. des Sciences de Liège} (2) 17 (1892).}

وأعطي أيضًا في II برهانًا أكثر مفهومية لمبرهنة الاختزال، إذ أصوغ المسألة بوصفها مسألة تكافؤ عائلات خطية من الأشكال. وقد أخذت هذه النظرة، وكذلك جزءًا جوهريًا من الاعتبارات المستعملة، من بحث «\srcLTR{Über algebraische Modulsysteme}»\NoetherSrcNote{*)}{E. Fischer: \srcLTR{Über algebraische Modulsysteme und lineare homogene partielle Differentialgleichungen mit konstanten Koeffizienten}، \S\ 1--4، \srcLTR{J. f. M.} 140 (1911).} ومن مبرهنات غير منشورة لإ. فيشر متصلة به، وقد تفضل فأذن لي باستعمالها.

وأبيّن أخيرًا في III كيف تفضي الاعتبارات المناظرة أيضًا إلى أعم التوسعات المتسلسلة.

\subsection*{I.}
ليكن $\theta$ شكلًا متجانسًا في كل صف من الصفوف $N>n$
\[
  A_1,A_2,\ldots,A_N
\]
بحيث يجوز، على العموم، أن يتألف الصف $A_k$ من العناصر $n$
\[
  A_k^{(1)},A_k^{(2)},\ldots,A_k^{(n)}
\]
يمكن أن تكون معاملات $\theta$ متغيرات غير معيّنة أو مقادير من حقل معيّن. وليرمز $P$، فضلًا عن ذلك، إلى دالة صحيحة نسبية، أي كثيرة الحدود، --- \emph{ذات معاملات عددية نسبية} --- في عمليات الاستقطاب:
\[
  P_{hk}=\left(A_h\frac{\partial}{\partial A_k}\right)
  =A_h^{(1)}\frac{\partial}{\partial A_k^{(1)}}+
   A_h^{(2)}\frac{\partial}{\partial A_k^{(2)}}+
   \cdots+
   A_h^{(n)}\frac{\partial}{\partial A_k^{(n)}},
\]
ويُفهم من الجداء $P_{hk}P_{ij}\theta$ تطبيق $P_{hk}$ على $P_{ij}\theta$.

وعندئذ تُعطى \emph{مبرهنة الاختزال} المذكورة بالهوية
\begin{equation}
  \theta=\sum PZ,
\tag{1}
\end{equation}
حيث لا تعود الأشكال $Z$ تحتوي إلا على الصفوف
\[
  A_1,A_2,\ldots,A_n
\]
وتكون مشتقة من $\theta$ بعمليات الاستقطاب $P$.

لننظر الآن في منظومة الأشكال الأساسية $(F')$، المؤلفة من $N$ من الأشكال المتساوية الدرجة وذات العدد نفسه من المتغيرات، وهي
\[
  F_1,F_2,\ldots,F_N,
\]
والتي تُتخذ معاملاتها، على الترتيب، صفوفًا $N$
\[
  A_1,A_2,\ldots,A_N
\]
ولتدل $(J)$ على المنظومة الكاملة --- المؤلفة من عدد منته من اللامتغيرات --- للأشكال $F_1,\ldots,F_n$، بحيث يمكن التعبير عن كل لامتغير لـ$F_1,\ldots,F_n$ تعبيرًا صحيحًا نسبيًا، أي كثير الحدود، باللامتغيرات $(J)$. \emph{وعندئذ يكون حدس هيلبرت المراد إثباته هو أن للامتغيرات المتزامنة $S$ للأشكال $F_1,\ldots,F_N$ منظومة كاملة تؤلفها اللامتغيرات المنتهية العدد $(J,PJ)$، حيث يدل $PJ$ على جميع اللامتغيرات التي يمكن اشتقاقها بعمليات الاستقطاب $P$ من $J$.}

والواقع أن كل لامتغير متزامن $S$ لمنظومة $(F')$، ذي معاملات عددية نسبية، يكون شكلًا متجانسًا في كل صف من الصفوف $N$، $A_1,\ldots,A_N$ --- وذا معاملات عددية نسبية ---، ولذلك تنطبق عليه الهوية (1). ولما كان تطبيق عمليات الاستقطاب $P$ على $S$ يولّد، كما هو معلوم، لامتغيرات من جديد، كانت الأشكال $Z$ أيضًا لامتغيرات، وبالتحديد لامتغيرات متزامنة للأشكال $F_1,
\ldots,F_n$؛ لأنها لا تعود تحتوي إلا على الصفوف $A_1,
\ldots,A_n$. ومن ثم فهي، بحسب الفرض، دوال صحيحة نسبية، أي كثيرة الحدود، في اللامتغيرات $(J)$. وهكذا تتحول الهوية (1) إلى
\[
  S=\sum PY,
\]
حيث تدل $Y$ على جداءات من قوى اللامتغيرات $(J)$. غير أن التنفيذ الفعلي للعمليات $P$ على $Y$ يتألف، بحسب تعريف $P$، من التنفيذ المتعاقب المتكرر عددًا منتهيًا من المرات لعمليات الاستقطاب البسيطة $P_{hk}$، وهي تفضي، بحسب قاعدة تفاضل الجداءات، إلى مجاميع جداءات من اللامتغيرات $(J)$ و$(PJ)$. ويصدق الأمر نفسه على تطبيق $P_{hk}$ على جداء من $(J)$ و$(PJ)$؛ ومن ثم لا يعطي $PY$ أيضًا إلا تراكيب صحيحة نسبية، أي كثيرة الحدود، في $(J,PJ)$. \emph{وبذلك ثبت إمكان تمثيل جميع اللامتغيرات $S$ تمثيلًا صحيحًا نسبيًا، أي كثير الحدود، بواسطة $(J,PJ)$.}

\subsection*{II.}
نمهّد لبرهان مبرهنة الاختزال ببعض الكلام على العائلات الخطية من الأشكال. ونقصد هنا بـ\emph{عائلة خطية من الأشكال في $K$} منظومةً من أشكال متساوية الدرجة، مغلقةً بالمعنى الآتي: متى انتمى إليها $\theta_1$ و$\theta_2$، انتمى إليها أيضًا الشكل $c_1\theta_1+c_2\theta_2$، حيث إن $c_1,c_2$ مقداران من الحقل المعين $K$، على أن يحتوي $K$ حقل معاملات أشكال $\theta$. ويوجد بين هذه الأشكال دائمًا عدد منته --- وليكن $\rho$ --- من الأشكال المستقلة خطيًا على $K$، في حين تخضع كل $\rho+1$ من الأشكال لعلاقة خطية ذات معاملات من $K$؛ ويقال إن للعائلة \emph{الرتبة $\rho$} بالنسبة إلى $K$.\NoetherSrcNote{*)}{نحصل على عائلات خطية من الأشكال أعم --- قد لا تكون رتبتها منتهية --- إذا أسقطنا قيد تساوي الدرجة، أو أسقطنا أيضًا قيد احتواء $K$ حقل معاملات أشكال $\theta$.} وتصح \emph{الحقيقة} السهلة الآتية: إذا كانت $\mathcal L$ و$\mathcal T$ عائلتين خطيتين في $K$ لهما الرتبة نفسها $\rho$، وكانت $\mathcal T$ عائلة جزئية من $\mathcal L$، فإن $\mathcal L$ و$\mathcal T$ متطابقتان (متكافئتان).

بعد هذه التوطئة ننتقل إلى مبرهنة الاختزال. والآن تتحول المتطابقة (في I)
\[
  \theta=\sum PZ
\]
إلى متطابقة مماثلة لـ$P\theta$ إذا طبّقنا عملية الاستقطاب نفسها $P$ على الطرفين؛ ومن ثم تعطي مبرهنة الاختزال في آن واحد تمثيلًا \emph{لجميع} الأشكال القطبية لـ$\theta$ بمجاميع من الأشكال القطبية الخاصة $PZ$. ولما كانت هذه الأشكال القطبية الخاصة داخلةً يقينًا في مجموع جميع الأشكال القطبية، فإن \emph{مبرهنة الاختزال تقرر تكافؤ هاتين العائلتين الخطيتين من الأشكال}؛ وبوجه خاص تقرر أيضًا تكافؤ العائلتين الجزئيتين اللتين تكون أشكالهما في كل صف \emph{على حدة} من الدرجة نفسها التي لـ$\theta$. ولإثبات هذا التكافؤ ندخل أيضًا عائلة وسيطة تحددها الأشكال $Z$. وللتبسيط نعطي \emph{البرهان أولًا في حالة ثلاثة صفوف ثنائية $(N=3,n=2)$}، ثم نلمح بإيجاز إلى البرهان العام الموازي له.

لتكن إذن $\theta(x,y,z)$ ذات الدرجات $\alpha,\beta,p$، على الترتيب، في الصفوف
\[
  x_1x_2;\quad y_1y_2;\quad z_1z_2;
\]
ولتكن معاملات $\theta$ في البدء غير معينات $a$ (أو أعدادًا نسبية أيضًا). وعندئذ ننظر في العائلات الخطية الثلاث الآتية من الأشكال:

١. العائلة $\mathcal L$، المؤلفة من جميع الأشكال $f(x,y,z)$ التي نشأت من $\theta$ بعمليات استقطاب $P$، والتي لها في كل صف من الصفوف $x,y,z$ الدرجة نفسها التي لـ$\theta$؛ وتحتوي $\mathcal L$ أيضًا، بوجه خاص، على $\theta$. وإذا عُدّت $f(x,y,z)$ أشكالًا في $x,y,z$ وفي غير المعينات $a$، كان حقل المعاملات هو الحقل $R$ للأعداد النسبية؛ ولا بد أيضًا من أن يؤخذ لـ$K$ الحقل $R$، لأنه لا يصبح، إلا عندما يكون $c_1,c_2$ عددين نسبيين، $c_1P_1+c_2P_2$ بدوره عمليةً $P$. ولتكن لـ$\mathcal L$ الرتبة $\rho$ بالنسبة إلى $R$.

٢. العائلة $\mathcal S$، المؤلفة من جميع الأشكال $g(\lambda;x,y)$ المعرفة بالعلاقة
\[
  g(\lambda;x,y)=f(x,y,\lambda_1x+\lambda_2y),
\]
أو، إذا عُبّر عنها بعمليات الاستقطاب:
\[
\begin{aligned}
  g(\lambda;x,y)
  &=\frac1{p!}\left[(\lambda_1x+\lambda_2y)\frac{\partial}{\partial z}\right]^p f(x,y,z)\srcfnmark{*)} \\
  &=g_0(xy)\lambda_1^p+g_1(x,y)\lambda_1^{p-1}\lambda_2+\cdots+g_p(xy)\lambda_2^p,
\end{aligned}
\]
حيث:
\[
  i!(p-i)!\,g_i(xy)=
  \left(y\frac{\partial}{\partial z}\right)^i
  \left(x\frac{\partial}{\partial z}\right)^{p-i}f(xyz).
\]
\srcfntext{*)}{المقصود، كما في I، هو:
\[
 \left(t\frac{\partial}{\partial x}\right)=t_1\frac{\partial}{\partial x_1}+t_2\frac{\partial}{\partial x_2},
\]
ومن ثم، بوجه خاص:
\[
 \left[(\lambda_1x+\lambda_2y)\frac{\partial}{\partial z}\right]
 = (\lambda_1x_1+\lambda_2y_1)\frac{\partial}{\partial z_1}
 + (\lambda_1x_2+\lambda_2y_2)\frac{\partial}{\partial z_2},
\]
\[
 \left[z\left(\frac{\partial}{\partial\lambda_1}\frac{\partial}{\partial x}
 +\frac{\partial}{\partial\lambda_2}\frac{\partial}{\partial y}\right)\right]
 =z_1\left(\frac{\partial}{\partial\lambda_1}\frac{\partial}{\partial x_1}
 +\frac{\partial}{\partial\lambda_2}\frac{\partial}{\partial y_1}\right)
 +z_2\left(\frac{\partial}{\partial\lambda_1}\frac{\partial}{\partial x_2}
 +\frac{\partial}{\partial\lambda_2}\frac{\partial}{\partial y_2}\right).
\]}
ومن ثم تعطي $g_i(xy)$ الأشكال $Z$ الواردة في المتطابقة (١). والأشكال $g$ ذات معاملات عددية نسبية في $x,y,\lambda$ وفي غير المعينات $a$؛ ولتكن لـ$\mathcal S$ الرتبة $\sigma$ بالنسبة إلى $R$.

٣. العائلة $\mathcal T$، الناتجة من $\mathcal S$ بعملية الاستقطاب العكسية، كما نتجت $\mathcal S$ من $\mathcal L$؛ وتُعرّف أشكال $h$ من $\mathcal T$ بالعلاقة
\[
\begin{aligned}
 h(x,y,z)
 &=\frac1{p!}\left[z\left(\frac{\partial}{\partial\lambda_1}\frac{\partial}{\partial x}
 +\frac{\partial}{\partial\lambda_2}\frac{\partial}{\partial y}\right)\right]^p g(\lambda;xy)\\
 &=h_0(xyz)+h_1(xyz)+\cdots+h_p(xyz),
\end{aligned}
\]
حيث، بحسب ٢:
\[
  h_i(xyz)=
  \left(z\frac{\partial}{\partial x}\right)^{p-i}
  \left(z\frac{\partial}{\partial y}\right)^i g_i(xy).
\]
ومن ثم يكون كل $h_i$ شكلًا مفردًا، ويكون $h$ بالتالي مجموعًا من هذه الأشكال؛ فالأول شكل من نوع $PZ$، والثاني مجموع من أشكال هذا النوع؛ ولتكن لـ$\mathcal T$ الرتبة $\tau$ بالنسبة إلى $R$.

يتبين من بناء أشكال $h(x,y,z)$ في $\mathcal T$ أنها نشأت بعمليات استقطاب $P$ من $\theta$، وأن لها في كل صف من الصفوف $x,y,z$، على حدة، الدرجة نفسها التي لـ$\theta$؛ ومن ثم \emph{تشكل العائلة $\mathcal T$ عائلةً جزئية من $\mathcal L$}. وبحسب \emph{الحقيقة} المذكورة أعلاه، لا يبقى لإثبات التكافؤ إلا إثبات تساوي الرتبتين. وفي هذا الإثبات لا تعود البنية الخاصة للأشكال $f(xyz)$ --- أي كونها أشكالًا قطبية لشكل واحد بعينه $\theta$ --- مستعملةً.

\emph{أ) لمقارنة رتبتي $\mathcal L$ و$\mathcal S$ نستعمل اللمّة المساعدة (أ): يلزم من $f(x,y,\lambda_1x+\lambda_2y)=0$ أن يكون $f(x,y,z)=0$.} وينتج ذلك مباشرةً من العلاقة المتطابقة القائمة بين ثلاثة صفوف ثنائية
\[
  (xy)z+(yz)x+(zx)y=0,
  \qquad \text{حيث }(xy)=(x_1y_2-x_2y_1),\ldots,
\]
التي يترتب عليها:
\[
  f[x,y,(zy)x+(xz)y]=(xy)^p f(x,y,z).
\]
فمن $f(x,y,\lambda_1x+\lambda_2y)=0$ (بوصفها متطابقةً في $\lambda$) يلزم إذن، كما يتضح من التعويض $\lambda_1=(zy)$ و$\lambda_2=(xz)$، وبما أن $(xy)\ne0$، أن يكون $f(x,y,z)=0$.

وبحسب هذه اللمّة \emph{يوجد تقابل واحد لواحد بين أشكال $\mathcal S$ وأشكال $\mathcal L$}. فمن
\[
  f_1(x,y,\lambda_1x+\lambda_2y)=g(\lambda;xy),\qquad
  f_2(x,y,\lambda_1x+\lambda_2y)=g(\lambda;xy)
\]
يلزم بالضرورة:
\[
  f_1(xyz)=f_2(xyz).
\]
وبموجب هذا التقابل، \emph{تبقى كل علاقة بين أشكال $\mathcal S$ قائمةً بين الأشكال المناظرة لها واحدًا لواحد في $\mathcal L$ (والعكس بالعكس)}. ذلك أن
\[
  c_1g^{(1)}(\lambda;xy)+c_2g^{(2)}(\lambda;xy)+\cdots+c_rg^{(r)}(\lambda;xy)=0
\]
يستلزم باللمّة المساعدة بالضرورة:
\[
  c_1f^{(1)}(x,y,z)+c_2f^{(2)}(x,y,z)+\cdots+c_rf^{(r)}(x,y,z)=0.
\]
\emph{وبذلك ثبت تساوي رتبتي $\mathcal L$ و$\mathcal S$؛ $\rho=\sigma$.}

\emph{ب) ولمقارنة رتبتي $\mathcal S$ و$\mathcal T$ نستعمل، على نحو مماثل، اللمّة المساعدة (ب): يلزم من $h(x,y,z)=0$ أن يكون $g(\lambda;xy)=0$.} وللبرهان نلاحظ أن $h(x,y,z)=0$، بحسب ٣، يكافئ انعدام جداءات القوى المنفردة
\[
 \left(\frac{\partial}{\partial\lambda_1}\frac{\partial}{\partial x_1}
 +\frac{\partial}{\partial\lambda_2}\frac{\partial}{\partial y_1}\right)^{p_1}
 \left(\frac{\partial}{\partial\lambda_1}\frac{\partial}{\partial x_2}
 +\frac{\partial}{\partial\lambda_2}\frac{\partial}{\partial y_2}\right)^{p_2}g(\lambda;x,y),
 \qquad p_1+p_2=p,
\]
ومن ثم يلزم، بمزيد من الاشتقاق، انعدام جداءات القوى
\[
\begin{aligned}
&\left(\frac{\partial}{\partial x_1}\right)^{\alpha_1}
 \left(\frac{\partial}{\partial x_2}\right)^{\alpha_2}
 \left(\frac{\partial}{\partial y_1}\right)^{\beta_1}
 \left(\frac{\partial}{\partial y_2}\right)^{\beta_2} \\
&\quad\cdot
 \left(\frac{\partial}{\partial\lambda_1}\frac{\partial}{\partial x_1}
 +\frac{\partial}{\partial\lambda_2}\frac{\partial}{\partial y_1}\right)^{p_1}
 \left(\frac{\partial}{\partial\lambda_1}\frac{\partial}{\partial x_2}
 +\frac{\partial}{\partial\lambda_2}\frac{\partial}{\partial y_2}\right)^{p_2}
 g(\lambda;xy)
\end{aligned}
\]
أيضًا. ومن ثم ينعدم كل تركيب خطي من جداءات القوى هذه، وبالتالي ينعدم كل شكل
\[
  f\left(\frac{\partial}{\partial x},\frac{\partial}{\partial y},
  \frac{\partial}{\partial\lambda_1}\frac{\partial}{\partial x}
  +\frac{\partial}{\partial\lambda_2}\frac{\partial}{\partial y}\right)g(\lambda,xy).
\]
فإذا اختير $f$ على وجه الخصوص بحيث يكون
\[
  f(x,y,\lambda_1x+\lambda_2y)=g(\lambda;xy),
\]
نتج أيضًا:
\[
  g\left(\frac{\partial}{\partial\lambda};
    \frac{\partial}{\partial x},
    \frac{\partial}{\partial y}\right)g(\lambda,x,y)=0;
\]
أو، إذا وضعنا
\[
  g(\lambda,xy)=\sum G_i\lambda_1^{\varkappa_1}\lambda_2^{\varkappa_2}
  x_1^{\varkappa_3}x_2^{\varkappa_4}y_1^{\varkappa_5}y_2^{\varkappa_6},
\]
حيث تكون $G_i$ أشكالًا خطيةً في غير المعينات $a$ بمعاملات من الأعداد النسبية، نتج
\[
  \sum \varkappa_1!\varkappa_2!\cdots\varkappa_6!\,G_i^2=0
\]
ومن ثم $g(\lambda;xy)=0$.

وينتج من اللمّة المساعدة (ب)، تمامًا كما في (أ)، \emph{تساوي رتبتي $\mathcal S$ و$\mathcal T$؛ $\sigma=\tau$.}

ومن ثم نحصل من (أ) و(ب) على: $\rho=\tau$؛ وبذلك، لما كانت $\mathcal T$ عائلةً جزئية من $\mathcal L$، ثبت تكافؤ هاتين العائلتين الخطيتين. ولذلك يمكن تمثيل كل شكل من $\mathcal L$، ولا سيما $\theta$، بتركيب خطي بمعاملات من الأعداد النسبية من $\rho$ شكلًا مستقلًا خطيًا من الأشكال $h(x,y,z)$:
\[
  \theta=\sum c_i h^{(i)}(x,y,z)=\sum PZ.
\]
غير أن هذه المتطابقة، المستنبطة لمعاملات $a$ غير المعينة للشكل $\theta$، تبقى صحيحةً إذا استُبدلت بغير المعينات $a$ مقادير من أي حقل؛ \emph{وبذلك تكون مبرهنة الاختزال قد بُرهنت بوجه عام في حالة ثلاثة صفوف ثنائية.}

يسير البرهان العام الآن سيرًا موازيًا تمامًا، وذلك لـ$N$ صفوف في كل منها $n$ متغيرات. ولتكن $\theta(A)$ ذات الدرجة $\alpha_i$ في الصف $A_i$. وننظر من جديد في العائلات الخطية الثلاث الآتية من الأشكال:

١. العائلة $\mathcal L$، المؤلفة من جميع الأشكال $f(A)$ الناشئة من $\theta(A)$ بعمليات الاستقطاب $P$، والتي تكون في كل صف منفرد $A_i$ من الدرجة نفسها التي لـ$\theta(A)$؛

٢. العائلة $\mathcal S$، المؤلفة من جميع الأشكال $g(\lambda;A_1\cdots A_n)$ الناشئة من $f(A)$ بالتعويضات
\[
\begin{aligned}
 A_{n+1}&=\lambda_{n+1}^{(1)}A_1+\lambda_{n+1}^{(2)}A_2+
        \cdots+\lambda_{n+1}^{(n)}A_n,\\
 &\vdotswithin{=}\\
 A_N&=\lambda_N^{(1)}A_1+\lambda_N^{(2)}A_2+
        \cdots+\lambda_N^{(n)}A_n
\end{aligned}
\]
؛ وعندئذ تُعطى، بمعاملات جداءات القوى المنفردة لـ$\lambda$ في $g(\lambda;A_1\cdots A_n)$، أشكال $Z$ الواردة في الهوية (1)؛

٣. العائلة $\mathcal T$، المؤلفة من جميع الأشكال $h(A)$ المعرّفة بـ
\[
\begin{aligned}
 h(A)=&\left[A_{n+1}\left(
 \frac{\partial}{\partial\lambda_{n+1}^{(1)}}\frac{\partial}{\partial A_1}
 +\cdots+
 \frac{\partial}{\partial\lambda_{n+1}^{(n)}}\frac{\partial}{\partial A_n}\right)\right]^{\alpha_{n+1}} \\
&\cdots
\left[A_N\left(
 \frac{\partial}{\partial\lambda_N^{(1)}}\frac{\partial}{\partial A_1}
 +\cdots+
 \frac{\partial}{\partial\lambda_N^{(n)}}\frac{\partial}{\partial A_n}\right)\right]^{\alpha_N}
 g(\lambda;A_1\cdots A_n);
\end{aligned}
\]
ويكون $h(A)$ مجموعًا من أشكال $PZ$.

والآن تبقى اللمّة المساعدة (أ) صحيحةً بسبب العلاقة المتطابقة القائمة بين كل مجموعة من $n+1$ صفوف؛ وكذلك تصح اللمّة المساعدة (ب)، إذ لم يُعتدّ أصلًا بعدد المتغيرات في برهانها. ومن ثم فالعائلتان $\mathcal L$ و$\mathcal T$ متساويتا الرتبة، وهما — لما كانت $\mathcal T$ عائلةً جزئية من $\mathcal L$ — متطابقتان. ومن ثم تصح الهوية
\[
  \theta=\sum PZ
\]
للمعاملات غير المعينة، ومن ثم تبقى صحيحةً أيضًا إذا أُخذت المعاملات مقادير من أي حقل؛ \emph{وبذلك تكون مبرهنة الاختزال قد بُرهنت في الحالة العامة.}

\subsection*{III.}
نبين أيضًا بإيجاز كيف نحصل، باعتبارات مماثلة، على \emph{التوسيع العام بحسب الصفوف} (عندما $N\le n$). ولتكن $Z$ صيغةً متجانسةً في كل صف على حدة من الصفوف $n$، وهي $A_1,
\ldots,A_n$، يتألف كل منها من $n$ مقادير؛ ولتكن $\Delta$ محددة هذه الصفوف. ونثبت أولًا الهوية المناظرة لـ(1):
\begin{equation}
  Z\equiv\sum PH \pmod{\Delta},
\tag{2}
\end{equation}
حيث لا تحتوي الأشكال $H$ بعد ذلك إلا على الصفوف $A_1,\ldots,A_{n-1}$، وهي مشتقة من $Z$ بعمليات $P$.

يقوم البرهان على تحويل العلاقات المستنتجة في II إلى تطابقات بترديد $\Delta$. وللتبسيط، لنأخذ \emph{ثلاثة صفوف ثلاثية} $x,y,z$ أساسًا؛ ولنفترض أن المعاملات غير معينة.

ولتكن العائلات الخطية الثلاث $\mathcal L,\mathcal S,\mathcal T$ معرّفةً كما في II في حالة الصفوف الثنائية؛ ولتكن رتبها $\rho,\sigma,\tau$، بينما ترمز $\rho^*$ و$\tau^*$ إلى رتبتي $\mathcal L$ و$\mathcal T$ بترديد $\Delta$ (أي توجد أشكال عددها $\rho^*$، لا $\rho^*+1$، من $\mathcal L$، مستقلة خطيًا بترديد $\Delta$). فعندئذ تصح، كما في II، اللمّة المساعدة الآتية:

\emph{اللمّة المساعدة (أ): من $f(x,y,\lambda_1x+\lambda_2y)=0$ يلزم بالضرورة $f(xyz)\equiv0\pmod{\Delta}$ (والعكس).} ذلك أنه، بحكم العلاقة القائمة بين أربعة صفوف ثلاثية،
\[
  \Delta_1x-\Delta_2y+\Delta_3z=\Delta t,
  \qquad \text{حيث }\Delta_1=(yzt),\ldots,
\]
تقوم دائمًا بين ثلاثة صفوف ثلاثية علاقة خطية بترديد $\Delta$:
\[
  \Delta_1x-\Delta_2y+\Delta_3z\equiv0\pmod{\Delta},
\]
ومن ثم يكون لدينا، كما في II:
\[
  f(x,y,-\Delta_1x+\Delta_2y)\equiv\Delta_3^p f(x,y,z)\pmod{\Delta}.
\]
ومن $f(x,y,\lambda_1x+\lambda_2y)=0$ (انعدامًا مطابقًا في $\lambda$) يترتب بالضرورة — لما كانت كثيرة الحدود $\Delta$ غير قابلة للاختزال، وكانت $\Delta_3$ لا تقبل القسمة على $\Delta$ — أن $f(x,y,z)\equiv0\pmod{\Delta}$. والعكس واضح لأن $\Delta(x,y,\lambda_1x+\lambda_2y)=0$. ومن هنا نستنتج كما في II: $\rho^*=\sigma$.

وتبقى اللمّة المساعدة (ب) صحيحةً ببرهانها، ومن ثم يكون لدينا: $\sigma=\tau$.

ولإثبات $\tau=\tau^*$ نستعمل اللمّة المساعدة الآتية:

\emph{اللمّة المساعدة (ج): من $h(x,y,z)\equiv0\pmod{\Delta}$ يلزم بالضرورة $h(x,y,z)=0$.}
ذلك أن $h(x,y,z)$، بوصفه شكلًا قطبيًا، ينعدم عند تطبيق عملية $\Omega$ المعروفة؛ ويُعبّر عن هذا بالمعادلة التفاضلية
\[
  \Delta\!\left(\frac{\partial}{\partial x},\frac{\partial}{\partial y},\frac{\partial}{\partial z}\right)h(xyz)
  -\left(\frac{\partial}{\partial x}\frac{\partial}{\partial y}\frac{\partial}{\partial z}\right)h(x,y,z)=0.\srcfnmark{*)}
\]
\srcfntext{*)}{لدينا بحسب 3.:
\[
 c\cdot\Omega(h)=-\left(
\frac{\partial}{\partial \xi}\frac{\partial}{\partial \eta}
\left(\frac{\partial}{\partial\lambda_1}\frac{\partial}{\partial \xi}
+\frac{\partial}{\partial\lambda_2}\frac{\partial}{\partial \eta}\right)
\right)
\left(x\frac{\partial}{\partial \xi}\right)^{\alpha-1}
\left(y\frac{\partial}{\partial \eta}\right)^{\beta-1}
\left[
z\left(\frac{\partial}{\partial\lambda_1}\frac{\partial}{\partial \xi}
+\frac{\partial}{\partial\lambda_2}\frac{\partial}{\partial \eta}\right)
\right]^{p-1}\cdot g(\lambda;\xi,\eta)
\]
ومن ثم تنعدم $\Omega(h)$ بسبب انعدام عامل المحددة.}
وإذا كان الآن $h(x,y,z)=\Delta(x,y,z)\cdot k(x,y,z)$، فإن مزيدًا من الاشتقاق يبيّن أن المقدار الآتي ينعدم أيضًا:
\[
  k\!\left(\frac{\partial}{\partial x},\frac{\partial}{\partial y},\frac{\partial}{\partial z}\right)
  \Delta\!\left(\frac{\partial}{\partial x},\frac{\partial}{\partial y},\frac{\partial}{\partial z}\right)
  =
  h\!\left(\frac{\partial}{\partial x},\frac{\partial}{\partial y},\frac{\partial}{\partial z}\right)h(x,y,z),
\]
ومن ثم ينعدم $h(x,y,z)$ أيضًا. وهكذا نحصل على $\rho^*=\sigma=\tau=\tau^*$، وبما أن $\mathcal T$ عائلة جزئية من $\mathcal L$، \emph{تكون الهوية (2) قد ثبتت}؛ ويؤدي استدلال مماثل إلى البرهان في حالة $n$ من المتغيرات.

وبتطبيق الهوية (2) على العامل الذي يضرب $\Delta$، نحصل على التوسعة:
\[
  Z=\varphi_0+\Delta\varphi_1+\Delta^2\varphi_2+\cdots+\Delta^\mu\varphi_\mu,
\]
حيث ينعدم كلٌّ من $\varphi_i$ عند تطبيق عملية $\Omega$. ومن ثم يعطي تكرار العملية $i$ مرة --- إذ لدينا عمومًا: $\Omega(\Delta\cdot\psi)=c_1\psi+c_2\Delta\cdot\Omega(\psi)$ --- ما يأتي:
\[
  c\cdot \Omega^i(Z)\equiv\varphi_i\pmod{\Delta};
\]
ومن جهة أخرى، لدينا وفقًا لـ(2):
\[
  c\cdot\Omega^i(Z)\equiv\sum PH_i\pmod{\Delta},
\]
حيث تُشتق $H_i$ من الشكل $\Omega^i(Z)$ بالطريقة نفسها التي تُشتق بها $H$ من $Z$. وبحسب اللمّة المساعدة (ج) (بعد توسيعها إلى $n$ من المتغيرات)، نحصل إذن على: $\varphi_i=\sum PH_i$، ومن ثم:
\begin{equation}
  Z=\sum PH+\Delta\cdot\sum PH_1+\Delta^2\cdot\sum PH_2+
  \cdots+\Delta^\mu\cdot\sum PH_\mu
\tag{3}
\end{equation}
\emph{بوصفها أعمَّ توسعة بحسب الصفوف لشكل ذي $n$ صفوف، في كلٍّ منها $n$ من المتغيرات، بدلالة الأشكال القطبية لأشكال لا تضم إلا $n-1$ صفًا.}

وتُستخرج منها، كما هو معلوم، التوسعة بحسب الصفوف لأشكال ذات $N$ صفًا من $n$ من المتغيرات $(N<n)$ بتعويض بسيط. نضع (في حالة $N=3$ و$n>3$):
\[
  u=\xi_1x+\xi_2y+\xi_3z,
  \quad v=\eta_1x+\eta_2y+\eta_3z,
  \quad w=\zeta_1x+\zeta_2y+\zeta_3z
\]
ونوسّع $H(u,v,w)=Z(\xi,\eta,\zeta)$ وفقًا لـ(3). وعندئذ، لما كانت $\xi,\eta,\zeta$ لا تظهر إلا في التركيبات $u,v,w$، يكون لدينا:
\[
  \left(\eta\frac{\partial}{\partial \xi}\right)=
  \left(v\frac{\partial}{\partial u}\right)\cdots,
  \qquad
  \Omega=\sum_{ikl}\left(\frac{\partial}{\partial u}
  \frac{\partial}{\partial v}
  \frac{\partial}{\partial w}\right)_{ikl}(xyz)_{ikl}
  =\nabla_{uvw}(xyz).
\]
وبموجب التخصيص $\xi_1=\eta_2=\zeta_3=1$؛ $\xi_2=\xi_3=\cdots=\zeta_2=0$، يتحول $H(u,v,w)$ إلى $H(x,y,z)$، ويتحول $\nabla_{uvw}(xyz)$ إلى مضاعف عددي من $\nabla_{xyz}(xyz)=\nabla_{xyz}$، لأن تطبيق $(y\partial/\partial x)$ إلخ على المحددات $(xyz)_{ikl}$ يُعدمها. ولما كان الأمر المناظر يصح أيضًا عند وجود $N$ صفًا، تنتج من (3) التوسعة:
\begin{equation}
  H=\sum P\Phi+\sum P\Phi_1+\cdots+\sum P\Phi_\mu,
\tag{4}
\end{equation}
حيث تنشأ $\Phi_i$ من $\nabla^i_{A_1\ldots A_N}H$ بعمليات $P$، ولا تعود تحتوي إلا الصفوف $A_1,
\ldots,A_{N-1}$، بصرف النظر عن تراكيب المحددات الواقعة في $\nabla^i$ --- وهي، كما ذُكر، لا تدخل في الاعتبار عند الاستقطاب.

ولذلك، إذا استُبدلت هذه المحددات في $\nabla^i$ بمتغيرات غير معينة $p$، أمكن توسيع الأشكال الناتجة من جديد وفقًا لـ(4)؛ فنحصل أخيرًا على التوسعة
\begin{equation}
  H=\left[\sum P\psi(p)\right]_{p=\text{محددات }A_1\cdots A_N},
\tag{5}
\end{equation}
حيث تُشتق $\psi$ من الشكل الأصلي بعمليات $P$ و$\nabla_{A_1\ldots A_N}(p)$ و$\nabla_{A_1\ldots A_{N-1}}(p)$ إلخ. وبهذه التوسعات وتركيباتها --- بأن نطبق مثلًا (3) على الأشكال $Z$ الواقعة في (1)، ثم نطبق (4) و(5) --- \emph{تُستنفد جميع التوسعات المعروفة لأشكال ذات عدد اعتباطي من الصفوف، في كل منها $n$ من المتغيرات المتوافقة التدرج}.

إرلنغن، ٥ يناير ١٩١٥.



\clearpage
\setcounter{footnote}{0}

\section*{9. أعمّ المجالات المؤلَّفة من أعداد متسامية صحيحة}
\begin{center}
{\Large\bfseries 9. أعمّ المجالات المؤلَّفة من أعداد متسامية صحيحة}\par
\vspace{0.75em}
بقلم إيمي نوتر في إرلنغن.\par
\vspace{1.2em}
\srcLTR{Math. Ann. 77 (1916)}، ص. ١٠٣--١٢٨
\end{center}

بمقتضى مبرهنة الترتيب الحسن، أنشأ إ. زيرميلو مجالًا من
\glqq{}الأعداد المتسامية الصحيحة\grqq{}، أي مجالًا عدديًا
$\mathfrak G$ تتوافر فيه الخواص \emph{المجرّدة} الآتية:
\begin{enumerate}
\item[I.] مجموع عددين من $\mathfrak G$ وفرقُهما وحاصلُ ضربهما عددٌ من
$\mathfrak G$ أيضًا.
\item[II.] كل عدد حقيقي أو عقدي هو خارج قسمة عددين من $\mathfrak G$.
\item[III.] كل عدد صحيح نسبي (أو عدد صحيح جبري، بحسب الحال) هو عنصر من $\mathfrak G$.
\item[IV.] لا ينتمي إلى $\mathfrak G$ أي عدد نسبي غير صحيح (أو عدد جبري غير
صحيح، بحسب الحال).\srcfn{*)}{\srcLTR{E. Zermelo, Über ganze transzendente Zahlen,
Math. Ann. 75}، ص. ٤٣٤ (١٩١٤).}
\end{enumerate}

تقوم هذه البنية على أن مبرهنة الترتيب الحسن تُظهر وجود
\emph{أساس جبري} لجميع الأعداد؛ أي \emph{منظومة} $H$ \emph{من الأعداد}
$\eta$، \emph{لا تقوم بينها علاقات جبرية، ولكن يمكن التعبير عن جميع الأعداد
الأخرى جبريًا بدلالتها}. وبسبب استقلال أعداد الأساس $\eta$ جبريًا، يمكن
اعتبارها \emph{متغيرات غير معيّنة}، وبذلك يتحول المجال المطلوب إلى
\emph{منطقة تكاملية من الدوال النسبية والجبرية في المتغيرات غير المعيّنة}
$\eta$،\srcfn{**)}{ولهذا السبب تسير اعتبارات هذه المقالة جزئيًا بموازاة
اعتبارات المقالة السابقة للمؤلفة: «حقول ومنظومات الدوال النسبية»
(\srcLTR{Körper und Systeme rationaler Funktionen, Math. Ann. 76, S. 161
(1915)}).} حيث تنتمي معاملاتها إلى الحقل $K$ المؤلف من جميع الأعداد الجبرية.
ويتميّز مجال زيرميلو الخاص $\mathfrak G_\eta$ عندئذ بالخواص
\emph{الأساسية} الآتية:

يحتوي $\mathfrak G_\eta$ على:
\begin{enumerate}
\item[1)] جميع أعداد الأساس $\eta$،
\item[2)] جميع كثيرات الحدود في $\eta$ ذات المعاملات الصحيحة\srcfn{***)}{يُقصد
بـ\glqq{}صحيحة المعاملات\grqq{} في جميع ما يأتي أن المعاملات تنتمي إلى
المنطقة التكاملية $[K]$ لجميع الأعداد الصحيحة الجبرية. ولتعريف
$\mathfrak G_\eta$ في البنود 1) إلى 5) يكفي أيضًا الاقتصار على معاملات من
الأعداد الصحيحة النسبية؛ غير أن الخواص الأساسية للمجالات الأعم تصير عندئذ
أقل بساطة.}،
\item[3)] جميع الدوال الصحيحة جبريًا في $\eta$ ذات المعاملات الصحيحة.
\end{enumerate}
يستبعد $\mathfrak G_\eta$ ما يأتي:
\begin{enumerate}
\item[4)] جميع الدوال النسبية الكسرية في $\eta$ ذات المعاملات الصحيحة،
\item[5)] جميع الدوال الجبرية الكسرية في $\eta$ ذات المعاملات
الصحيحة.\srcfn{*)}{زيرميلو، المرجع السابق، § ٣.}
\end{enumerate}

يتبين الآن بسهولة (قارن الأمثلة في § ٢ و٣) أن ثمة، بالقياس إلى مجال زيرميلو
$\mathfrak G_\eta$، مجالات $\mathfrak G$ تختلف عنه
\emph{اختلافًا جوهريًا} --- أي مجالات $\mathfrak G$ لا تتوافر فيها، عند
\emph{أي ترتيب حسن}، جميع الخواص الأساسية 1) إلى 5)، ومن ثم لا يمكن
توليدها ببناء زيرميلو عند أي ترتيب حسن؛ وبعبارة أخرى: ثمة مجالات
$\mathfrak G$ لا تقبل تطبيقًا تقابليًا وتشاكليًا على $\mathfrak G_\eta$.
وهكذا تنشأ مسألة \emph{أعم المجالات المؤلفة من أعداد متسامية صحيحة}،
وهي المسألة التي سيجاب عنها جوابًا كاملًا فيما يأتي.

ولهذا ينبغي أولًا تعيين أي الخواص الأساسية يترتب على خواص التعريف المجرد،
وأيها يترتب على البناء الخاص. ويتبين (§ ١) أنه يوجد لكل مجال معطى
$\mathfrak G$ ترتيب حسن تتحقق معه الخاصتان الأساسيتان 1) و2)؛ ولذلك
\emph{يمكن إنشاء كل مجال $\mathfrak G$ بواسطة أساس جبري لجميع الأعداد}.
أما الخواص الأساسية الباقية فلا يلزم أن تتحقق واحدة منها (§ ٢ و٣). وينتج
من ذلك، بوجه خاص، أن خاصية مجردة أخرى لمجال زيرميلو لا تتوافر في جميع
المجالات $\mathfrak G$، وهي: \emph{الانغلاق تحت الاعتماد الصحيح جبريًا}.
ويمكن صوغها كما يأتي:
\begin{enumerate}
\item[V.] كل عدد يعتمد على عدد منتهٍ من أعداد $\mathfrak G$ اعتمادًا صحيحًا
جبريًا وصحيحَ المعاملات ينتمي إلى $\mathfrak G$.
\end{enumerate}

ولْنسمِّ المجالات الخاصة التي تتوافر فيها V إلى جانب I--IV مجالات
$\mathfrak H$ مؤلفة من \emph{أعداد متسامية صحيحة جبريًا}. وتنتج للمجالات
$\mathfrak H$ الأعم النتائج البسيطة الآتية (§ ٤):
\begin{enumerate}
\item[a)] \emph{تقاطع جميع المجالات القابلة للإنشاء بواسطة الأساس نفسه
$H$ ضمن فئة المجالات $\mathfrak H$، هو مجموعة جميع الدوال الصحيحة جبريًا في الأساس --- أي مجال
زيرميلو $\mathfrak G_\eta$، وهو نفسه مجال $\mathfrak H$.} (وبذلك نحصل في
الوقت نفسه على توصيف مجرد لمجال زيرميلو.)
\item[b)] \emph{ينشأ كل مجال $\mathfrak H$ عن طريق الإلحاق الصحيح
جبريًا\srcfn{**)}{تفسير العبارة في § ٤.} لمنظومة كيفما كانت
$\mathfrak S(\eta)$ بالمجال $\mathfrak G_\eta$، على أن تستوفي
$\mathfrak S(\eta)$ شرطًا واحدًا فقط: ألا يدخل بالإلحاق أي عدد جبري كسري
في المجال.}
\end{enumerate}

أما النتائج المتعلقة بأعم المجالات $\mathfrak G$ عند اتخاذ أساس جبري منطلقًا
فهي أقل بساطة. ففي هذه الحال لا يكون تقاطع جميع المجالات $\mathfrak G$ نفسه
مجالًا $\mathfrak G$؛ ولذلك يصعب أيضًا الإحاطة ببناء المجالات الأعم (§ ٥).

وللحصول على نتائج تماثل نتائج المجالات $\mathfrak H$، يُستعاض عن الأساس
الجبري \emph{بقاعدة نسبية} $\Theta$ --- أي \emph{بمنظومة} $\Theta$
\emph{من الأعداد} $\vartheta$، \emph{تتيح التعبير عن جميع الأعداد نسبيًا،
في حين لا يمكن، عند ترتيب حسن واحد على الأقل، التعبير نسبيًا عن أي عدد من
أعداد القاعدة بدلالة عدد منتهٍ من سابقيه}. ولا بد أيضًا من فرض الشرط الجانبي
التالي على هذه القاعدة النسبية $\Theta$: ألا تحتوي المنطقة التكاملية
المؤلفة من جميع كثيرات الحدود في $\vartheta$ ذات المعاملات الصحيحة، أي
$[\Theta]$، على أي عدد جبري كسري. (ويُبيَّن في § ٧ إنشاء قاعدة كهذه مع الشرط
الجانبي.) وعندئذ تنتج مرة أخرى للمجالات $\mathfrak G$ الأعم --- التي يمكن
أيضًا تسميتها مجالات مؤلفة من \emph{أعداد متسامية صحيحة نسبيًا} --- النتائج
البسيطة الآتية (§ ٦):
\begin{enumerate}
\item[a)] \emph{تقاطع جميع المجالات القابلة للإنشاء بواسطة القاعدة النسبية
نفسها $\Theta$ ضمن فئة المجالات $\mathfrak G$، هو مجموعة جميع الدوال الصحيحة نسبيًا في
القاعدة --- أي المجال $[\Theta]$، وهو نفسه مجال $\mathfrak G$.}
\item[b)] \emph{ينشأ كل مجال $\mathfrak G$ عن طريق الإلحاق الصحيح نسبيًا
لمنظومة كيفما كانت $\mathfrak S(\vartheta)$ بالمجال $[\Theta]$، على أن
تستوفي $\mathfrak S(\vartheta)$ شرطًا واحدًا فقط: ألا يدخل بالإلحاق أي عدد
جبري كسري في المجال.}
\end{enumerate}

ولكي يكون التشابه مع المجالات $\mathfrak H$ تامًا، كان ينبغي ألا تُدرج
الأعداد الجبرية في خاصيتي التعريف III وIV للمجالات $\mathfrak G$. ومع هذا
التعديل للخاصيتين III وIV يمكن نقل إنشاء العناصر الصحيحة بواسطة قاعدة نسبية
إلى \emph{أي} حقل كيفما كان (§ ١٠)، في حين يفترض بناء زيرميلو الانغلاق الجبري
للحقل.

\emph{إذا جرى $H$ على كل أساس جبري، وجرى $\Theta$، على الترتيب، على كل
قاعدة نسبية تستوفي الشرط الجانبي، أعطت النتيجتان a) وb) مجموع جميع المجالات
$\mathfrak H$ و$\mathfrak G$.} وإذا جُمعت جميع المجالات المتشاكلة في فئة
واحدة، أمكن انتقاء مجموعة جزئية من هذا المجموع تمثل مجموع جميع الفئات ---
التي لا يقل عددها، وفق § ٩، عن عدد لا نهائي قابل للعد --- (§ ٨).



\clearpage
\section*{تعليقات نقدية على مقالة سنة 1911 (الصفحات الأصلية 118--154)}
\addcontentsline{toc}{section}{تعليقات نقدية على مقالة سنة 1911 (الصفحات الأصلية 118--154)}
\begin{enumerate}
\item الصفحة الأصلية 118، إحالة Study. يذكر المطبوع الأصلي المجلد 1. نُشرت مقالة Eduard Study، \emph{\srcLTR{Zur Differentialgeometrie der analytischen Curven}}، في \emph{\srcLTR{Transactions of the American Mathematical Society}}، المجلد 10، العدد 1 (1909)، الصفحات 1--49؛ المعرّف الرقمي: \url{https://doi.org/10.1090/S0002-9947-1909-1500824-2}.
\item الصفحة الأصلية 129، الحجة التالية لـ(16). لا توجد في المطبوع الأصلي نقاط استمرار بين $a^{(i+1)}$ و$a^{(\rho)}$. وتستند القراءة $a^{(i+1)}\cdots a^{(\rho)}$ إلى المتطابقة (16) في الصفحة نفسها، وتدل على متتالية العوامل المتعاقبة فيها.
\item الصفحة الأصلية 151، بعد (70). تعدّ الأعداد $\alpha_\rho$ الطيات الأساسية من كل نوع، ويجوز أن تكون صفرًا. ومن ثم فالشرط هو $\alpha_\rho\in\mathbb Z_{\ge0}$. وتُدخل حاشية الصفحة الأصلية 152 القيم الصفرية صراحة في الأعداد المناظرة. وتُضبط هنا عبارة الأصل، أعداد صحيحة موجبة، بما يشمل الصفر. وفي الصفحة الأصلية 153 توصف النقطة الشبكية المناظرة أيضًا بأن إحداثياتها صحيحة موجبة. ولأنها تمثل الأعداد نفسها، فالشرط هنا أيضًا هو $\alpha_\rho\ge0$؛ وتقابل الصيغة الابتدائية من دون طي القيمة $\alpha=0$. وهذا تحديد للأعداد، لا للاتجاهات الموجبة المذكورة لاحقًا في الفقرة.
\end{enumerate}
\smallskip
\noindent{\small تعليقات نقدية وتصحيحات منتجة بالذكاء الاصطناعي: OpenAI Codex، GPT-6 Astra، مستوى الاستدلال Ultra.}

\clearpage
\section*{ملاحظة تحريرية على P06}
\addcontentsline{toc}{section}{ملاحظة تحريرية على P06}
الصفحة المطبوعة 174، بعد (6). يتطلب القول بعدم انعدام البسط مطابقًا الشرط $f^*\ne0$. للدالة الصفرية المسموح بها يمكن اختيار $F^*=0$ و$H^*=1$؛ فتكون جميع البسوط المترافقة صفرًا، وتبقى المقامات غير صفرية. يبقى التكميم على جميع الدوال والشرط $H^*(\xi)\ne0$ في (7)؛ ولا تعطي الدالة الصفرية هناك سوى $0=0$.
\smallskip
\noindent{\small تعليقات نقدية وتصحيحات منتجة بالذكاء الاصطناعي: OpenAI Codex، GPT-6 Astra، مستوى الاستدلال Ultra.}
\end{document}
